% ============================================================
%  R. Nielsen, "Den speculative Methodes Anvendelse paa den hellige Historie."
%  Paa Dansk udgivet af B. C. Bøggild, Cand. theol.
%  Kjøbenhavn: Forlagt af H. C. Klein, Vimmelskaftet Nr. 23;
%  Trykt i det Seidelinske Officin, hos Louis Klein, 1842.
%  TRANSCRIPTION (Danish) — from the Royal Danish Library scan.
%
%  WHAT THIS BOOK IS. The Danish version, made by B. C. Bøggild, of Nielsen's
%        1840 Latin doctoral dissertation »De speculativa historiae sacrae
%        tractandae methodo«. It is his earliest book, and the fullest statement
%        of the speculative programme he spent the 1850s and 1860s dismantling
%        from within. Bøggild's preface places it against H. L. Martensen's
%        »Den menneskelige Selvbevidstheds Autonomie« and refers the reader to
%        the review in »Fædrelandet« Nr. 329 (1840). Bøggild signs himself at
%        »Kjøbenhavn, den 24de Januar 1842« — with the j, not the Kiøbenhavn
%        of the title page.
%
%  Scan: ~/bibliotek/Nielsen, Rasmus/speculative-methodes.pdf
%        (KB, 175 PDF pp., 1-bit JBIG2, 4048x6920 px at a nominal 300 ppi,
%        14 MB. Because the stored bitmap IS the working resolution, page
%        images are extracted with `pdfimages -png` rather than rasterised;
%        see cache.sh.)
%  TYPE: Fraktur (Danish body). Latin, French and English words and all figures
%        are set in antiqua against it -> \textit{}. The German quotations in
%        the notes are FRAKTUR, not antiqua, but Latin words inside them are
%        antiqua. Letterspacing (Sperrsatz) -> \emph{}; image-verify by zoom.
%
%  OFFSET (VERIFIED, UNIFORM): PDF = printed + 9.
%        Body printed pp. 1--162 = PDF 10--171. No leaf was scanned twice and
%        the offset never changes. Verified four ways:
%          (a) Fraktur OCR of every page 11--171 searched for the printed folio
%              in its first six lines: 123 of 161 pages show exactly PDF - 9,
%              and NOT ONE shows a different number (the other 38 lost the
%              folio to scan-edge noise).
%          (b) a token-overlap sweep over every adjacent page pair found no
%              pair above Jaccard 0.30, so no leaf was scanned twice.
%          (c) five structural landmarks taken from the printed Indhold land on
%              the predicted PDF page: § 4 (p. 18) on PDF 27, § 14 (p. 63) on
%              PDF 72, § 16 (p. 82) on PDF 91, § 20 (p. 119) on PDF 128, and
%              Slutning (p. 161) on PDF 170.
%          (d) the endpoints read off the image: PDF 10 = printed 1,
%              PDF 76 = printed 67, PDF 171 = printed 162.
%        pagemap.py is the single source of truth — do not hard-code +9.
%
%  EXTENT: body printed pp. 1--162 (PDF 10--171).
%        Front matter: title page PDF 6; the LAST page of B. C. Bøggild's
%        translator's preface PDF 7; Indhold PDF 8; Rettelser PDF 9.
%        PDF 2--5 and 172--175 are binding, a shelfmark leaf and blanks.
%
%  *** A LEAF OF THE PREFACE IS MISSING FROM THIS SCAN. ***
%        PDF 7 begins mid-sentence — »os imøde i dette Skrift, vistnok ville
%        bringe Læseren til at oversee hiin Mangel.« — and ends with the
%        signature »Kiøbenhavn, den 24de Januar 1842. B. C. Bøggild.« So at
%        least one preceding page of the preface is not in the file.
%        Recto/verso parity settles what kind of gap it is: even PDF pages are
%        rectos throughout (their text block runs to the left paper edge) and
%        odd ones versos, and that parity is unbroken from PDF 6 to PDF 171.
%        A single missing page would have inverted it. So a WHOLE LEAF — two
%        pages — is absent between the title page and PDF 7. PDF 2--5 were
%        checked and are genuinely blank, so the leaf is not merely misfiled.
%        The KB colour PDF (kb.dk/e-mat/dod/110208016515-color.pdf) may be a
%        different digitisation and has NOT been checked. Until it is, the
%        preface is transcribed as the fragment it is, and marked as such.
%
%  STRUCTURE (from the Indhold, PDF 8):
%        Første Deel. Forholdet imellem Philosophien og Historien, betragtet i
%          Almindelighed.
%            § 1  Forholdet imellem det Aposterioriske og det Aprioriske ... 1
%            § 2  Den subjective Tro og den speculative Logik .............. 5
%            § 3  Den vulgaire Empirie .................................... 11
%            § 4  Overgang til den speculative Empirie .................... 18
%            § 5  Forholdet imellem den vulgaire og speculative Empirie ... 25
%            § 6  Videnskabens Begreb .................................... 30
%            § 7  Den speculative Methode ................................ 36
%                 [BUT the printed head at p. 36 reads "Den videnskabelige
%                  Methode." Both readings are recorded at the site; the
%                  printed head is what is set. Do not normalise them.]
%            § 8  Begriben og Anskuen .................................... 42
%            § 9  Vor Viden er relativ ................................... 47
%            § 10 Forholdet imellem Dogmatiken og den hellige Historie .... 51
%        Anden Deel. Den hellige Historie.
%            § 11 Katholicismen .......................................... 53
%            § 12 Katholicismens Grundmodsigelse ......................... 56
%            § 13 Overgang fra Katholicismen til Protestantismen .......... 60
%            § 14 Troen og den hellige Skrift ............................ 63
%            § 15 Symbolerne. Luthersk Scholasticisme .................... 71
%            § 16 Fortsættelse. Overgang til Kriticismen ................. 82
%            § 17 Den subjective Rationalisme. Den hellige Histories
%                 Vanhelligelse .......................................... 87
%            § 18 Tro og Viden i den kritiske Theologie .................. 96
%            § 19 Undersøgelse af den kritiske Theologie ................ 104
%            § 20 Fortsættelse. Overgang til Panlogismen ................ 119
%            § 21 Panlogismens methaphysiske Prøvelse ................... 135
%            § 22 Tilbagevenden til Kirken. Den speculative Theologie ... 143
%            § 23 Miraklet og den hellige Historie ...................... 150
%        Slutning ...................................................... 161
%        The §§ run CONTINUOUSLY 1--23 across both Deele; they do not restart.
%        The list above is read from the printed Indhold ON THE IMAGE (the
%        contents page's OCR layer is unreliable — it gave "BG." for 36, "12"
%        for 42 and "113" for 143). Each argument is still to be confirmed
%        against its own printed § head as the batches reach it. Where
%        the Indhold and the printed head disagree, RECORD BOTH READINGS — the
%        disagreement is a datum, not something to normalise away (see
%        CLAUDE.md, "Editorial stance").
%
%  DISPLAY FORMS (settled on the image, PDF 10 and PDF 72):
%        Deel head : the Deel line ("Første Deel.") in large bold Fraktur,
%                    centred, WITH A PERIOD; below it the Deel argument in
%                    bold Fraktur, LETTERSPACED, over two lines.
%        § head    : "§ 14." centred in bold ANTIQUA; below it the § argument
%                    in Fraktur, LETTERSPACED and not bold.
%        folio     : centred bold antiqua at the head of the page. There are NO
%                    running heads in the original. Do not transcribe folios.
%
%  FOOTNOTES: this book has many, and they are the hard part of it.
%        They are marked 1) 2) 3) and RESTART AT 1) ON EVERY PRINTED PAGE.
%        LaTeX cannot do that by itself and [perpage] is wrong (it resets on
%        the TYPESET page, not the printed one), so \thefootnote is set to
%        \arabic{footnote}) and every printed page carrying a note gets an
%        explicit \setcounter{footnote}{0} after its page marker. check.py
%        tests for a missing reset, because that failure prints a plausible
%        wrong number instead of crashing.
%        The notes are set several points smaller than the body and are mostly
%        German (Fraktur) and Latin (antiqua): Luther's Werke, Melanchthon,
%        Bretschneider, D. F. Strauß, H. N. Clausen, Schleiermacher. Neither
%        OCR witness reads them well — tesseract's Fraktur model garbles the
%        antiqua and the ABBYY layer garbles the German — so every note is read
%        off the image.
%
%  ERRATA — the page is headed "Trykfeil," / "der bedes rettede, inden Bogen
%        læses." (PDF 9). TWELVE entries, read off the image at 2--5x. APPLY
%        THEM INLINE, each flagged % ERRATA at its site:
%
%          p.  4  l.17 fra oven        naar den        -> naar det
%          p.  4  l.26 fra oven        de maa          -> de maae
%          p.  7  l.11 fra oven        Subfaction      -> Subtraction
%          p. 10  l. 3 fra neden       almindeligt,    -> almindelige
%          p. 12  l. 1 fra neden       gaae ind        -> slaae ind
%          p. 15  l. 4 fra neden       under et nyt    -> under Skikkelse af et nyt
%          p. 41  l. 1 fra oven        næsten          -> næste
%          p. 42  l.13 fra oven        fast            -> faste
%          p. 48  l. 6 fra oven        Virken          -> Viden
%          p. 71  l.12 fra oven        Scholastisme    -> Scholasticisme
%          p. 89  l. 5 fra oven        κεκαυτερισμένυος -> κεκαυτερισμένους
%          p.132  Note, l. 1           das Mystische   -> das Mythische
%
%        Notes on the list itself:
%          * only the first entry spells out "læs:"; the rest use the
%            abbreviation "l.".
%          * the p. 71 entry sets BOTH readings letterspaced, so the word is
%            emphasised at that site in the body too.
%          * the p. 89 correction is a transposition only — the stem is
%            κεκαυτερισ-, with epsilon not eta and no alpha.
%          * every entry but the last is in the body; p. 132's is in a note.
%
%  GREEK: polytonic Greek occurs inline and in the notes. Type it as raw
%        Unicode; textalpha renders it. NORMALISE THE FOUNT'S VARIANT
%        LETTERFORMS — script theta and cursive kappa are typed as plain θ and
%        κ, per the repo's standing convention (see
%        texts/nielsen/religionsphilosophie/transcription.tex). This is also a
%        hard build requirement: U+03D1 ϑ and U+03F0 ϰ are FATAL errors under
%        textalpha/greek-fontenc in text mode. Same for ϕ ϖ ϱ ϐ ϵ ϲ -> φ π ρ β ε σ.
%        Accents and breathings are NOT fount variants — reproduce them exactly.
%
%  PAGE-JOINT WORD BREAKS: a word broken across a printed page boundary is
%        joined with `\-%` at the end of the last line of the earlier page —
%        a discretionary hyphen plus a comment that eats the newline — so that
%        the word reads whole and the break point is recorded by the
%        `% --- p. N ---` marker alone. A BARE hyphen there is a defect: it
%        typesets as "Ufuld- kommenhed", hyphen and space inside the word.
%        Nine such joints were repaired after pp. 1--108 were spliced.
%
%  OCR NOTE: the Fraktur model has a private ligature alphabet on this scan,
%        established by a census over all 170 pages and corrected in ocr.sh:
%        ſ -> s (8730), œ -> æ (552), ff fi fl ſt -> ff fi fl st, < -> ch (53,
%        "dur<" = durch), > -> ck (42, "Stü>e" = Stücke), and a bare backslash
%        for ſ before a consonant ("\ulde" = skulde). The commonest residual
%        error is k read as f (iffe = ikke, 203 times).
%
%  PRINTER'S DEFECTS -- transcribed AS PRINTED, each flagged at its site.
%        This is a diplomatic transcription: the compositor is not corrected.
%        PDF 7 (preface)  "Anmendelse" for "Anmeldelse"
%        p.  1  "Apostioriske"
%        p.  2  "absolnt" (turned u)
%        p.  3  "Phænomerne"; "Tiugen" (turned n)
%        p. 11  the § 3 head reads "Deu vulgaire Empirie." with a turned n,
%               against the Indhold's "Den vulgaire Empirie."
%        p. 35  "stundom saavidt; at den faaer" — semicolon where the sense
%               wants a comma, read at 3.5x
%        Also recorded, as setting rather than defect: the § 2 head is printed
%        "§. 2." with a point after the section sign, where §§ 1 and 3 are
%        printed "§ 1." and "§ 3."
%
%  FOOTNOTE MARKS: the standing form is 1) 2) 3) restarting per printed page,
%        but p. 12's single note is marked *) instead. It is set as an ordinary
%        \footnote with the printed mark recorded in a % comment at the site.
%
%  ONE HEBREW WORD: p. 33 sets pointed אֱלֹהִים. No Hebrew font is loaded here, so
%        it follows the precedent of texts/treschow/menneskelige-natur p. 115:
%        transliterated, with the substitution and the cjhebrew route documented
%        at the site. OPEN ITEM — see RESUME-NOTES.
%
%  GREEK, one further normalisation beyond the fount variants below: this fount
%        puts the breathing or circumflex of an -ου- diphthong over the OMICRON.
%        Unicode cannot encode a perispomeni on omicron, so the placement is
%        normalised to the standard one (οὐ, ποιοῦσι, ἐμοῦ) and logged at the
%        site. Printed accents that ARE encodable are kept exactly as printed
%        and noted where they depart from the grammar: Αληθῆ without a
%        breathing, σύ with an acute for a grave, ἀν without its grave, Θεος
%        without an accent.
%
%  QUOTATION MARKS: *** THIS BOOK DOES NOT USE THE DANISH LOW-HIGH PAIR „…“
%        ANYWHERE. *** It sets a RAISED mark, U+201D ”, at BOTH ends of every
%        quotation — in Bøggild's preface and throughout the body alike, the
%        opening mark at cap height, never on the baseline.
%        This was NOT assumed: it was found independently by four transcribers
%        working from the image at 2--6x — the front-matter pass and the passes
%        for pp. 1--12, 13--24 and 25--36 — and pp. 1--36 contain 64 raised
%        marks and not one „ or “. It is reproduced as printed.
%        The consequence for verification is that there is no signed balance to
%        take. check.py therefore runs a PARITY test on the single character
%        instead, which is a weaker instrument: it cannot tell an unclosed
%        quotation from an unopened one, and two defects on one page cancel
%        invisibly. So every quotation-mark defect must ALSO be logged by hand,
%        both at its site and in check.py's QUOTE_DEFECTS table.
%
%  QUOTE BALANCE: (none logged yet — see check.py)
%
%  STATUS: front matter transcribed (title page, the surviving page of the
%        preface, the Indhold and the Trykfeil), and body pp. 1--36 of 162.
%        Five errata applied so far (p. 4 twice, pp. 7, 10, 12).
%        Production scaffolding for this book (pagemap.py, cache.sh, crop.py,
%        ocr.sh, spacing.py, check.py, splice.py, BATCH-AGENT.md,
%        RESUME-NOTES.md) lives in this directory on the editor's machine but
%        is deliberately NOT tracked in git — see the root .gitignore for why.
%        Everything needed to verify this edition against the KB scan is in
%        this header.
% ============================================================
\documentclass[12pt]{book}
\usepackage[utf8]{inputenc}
\usepackage[T1]{fontenc}
\usepackage[danish]{babel}
\usepackage[protrusion=true,expansion=true]{microtype}
\usepackage[margin=1.4in]{geometry}
\usepackage{setspace}
\usepackage{amsmath}
\usepackage{graphicx}    % for \reflectbox (the old Danish »ɔ:« id-est mark)
\DeclareUnicodeCharacter{0254}{\reflectbox{c}}  % ɔ (U+0254) = reversed c
\usepackage{libertinus}
\usepackage{libertinust1math}
\usepackage{textalpha}   % polytonic Greek typed directly
\usepackage{fancyhdr}
% plainpages=false + pdfpagelabels: the title page, the roman front matter and
% the arabic body would otherwise all produce a PDF destination called
% "page.1", and pdfTeX drops the duplicates with a warning on every build.
% The PDF metadata is set with \hypersetup AFTER \begin{document}, not as a
% package option: as an option the "æ" reaches hyperref while fontenc still has
% it as the T1 command \T1\ae, which hyperref cannot put in a PDF string.
% hyperfootnotes=false: this book restarts its footnote counter on every printed
% page, so the same \thefootnote value recurs constantly and hyperref's footnote
% anchors collide — "destination with the same identifier (Hfootnote.N) has been
% already used", once per collision, on every build. The anchors buy nothing here
% (a printed edition's notes are read where they sit, not followed as links), so
% they are switched off rather than made unique.
\usepackage[colorlinks=true,linkcolor=black,urlcolor=black,
  unicode=true,plainpages=false,pdfpagelabels,hyperfootnotes=false]{hyperref}
\setlength{\headheight}{15pt}
\pagestyle{fancy}
\fancyhf{}
\fancyhead[LE,RO]{\thepage}
\fancyhead[RE]{\textit{Den speculative Methodes Anvendelse}}
\fancyhead[LO]{\textit{\leftmark}}
\setlength{\parindent}{1.5em}
\setlength{\parskip}{0pt}
\setstretch{1.3}

% Footnote marks as the 1842 printing sets them: 1) 2) 3), RESTARTING on every
% printed page. \setcounter{footnote}{0} after the page marker of every page
% that carries a note is what makes the restart happen; check.py enforces it.
% Not footmisc[perpage], which resets on the typeset page instead.
\renewcommand{\thefootnote}{\arabic{footnote})}

% ------------------------------------------------------------
%  STRUCTURAL MACROS — use these, do not hand-roll centre blocks.
%  They exist so that every batch renders the book's three recurring display
%  forms identically, and so that check.py can count them.
% ------------------------------------------------------------

% A Part opening, as printed at the head of an ordinary text page (NOT a
% display page of its own).
%   #1 = the Deel line as printed, WITH its period ("Første Deel.")
%   #2 = the Deel argument, which IS letterspaced in the printing
\newcommand{\deel}[2]{%
  \par\vspace{2.2\baselineskip}
  \begin{center}
    {\Large\bfseries #1}\\[0.7\baselineskip]
    {\large\bfseries\emph{#2}}
  \end{center}
  \vspace{0.9\baselineskip}
  \markboth{#1}{#1}%
  % \phantomsection gives each head its own anchor. Without it hyperref makes
  % the bookmark point at the enclosing \chapter's anchor, and every build
  % warns "The anchor of a bookmark and its parent's must not be the same".
  \phantomsection
  \addcontentsline{toc}{chapter}{\texorpdfstring{#1\quad #2}{#1 #2}}%
  \par\nopagebreak}

% A numbered section. The § number is bold ANTIQUA in the printing; the
% argument below it is Fraktur and letterspaced.
%   #1 = the number alone ("14")   #2 = the argument ("Troen og den hellige Skrift.")
\newcommand{\parag}[2]{%
  \par\vspace{1.6\baselineskip}
  \begin{center}
    {\bfseries §~#1.}\\[0.45\baselineskip]
    {\emph{#2}}
  \end{center}
  \vspace{0.5\baselineskip}
  \markright{§ #1. #2}%
  \phantomsection
  \addcontentsline{toc}{section}{\texorpdfstring{§~#1.\quad #2}{§ #1. #2}}%
  \par\nopagebreak}

% An unnumbered top-level division — used for "Slutning".
\newcommand{\division}[1]{%
  \par\vspace{2.0\baselineskip}
  \begin{center}{\large\bfseries #1}\end{center}
  \vspace{0.7\baselineskip}
  \markboth{#1}{#1}%
  \phantomsection
  \addcontentsline{toc}{chapter}{#1}%
  \par\nopagebreak}

% The old Danish "that is" sign ɔ: — never a plain "o:", never the raw glyph.
\newcommand{\dsi}{\reflectbox{c}:}

\title{\textbf{Den speculative Methodes Anvendelse\\ paa den hellige Historie}}
\author{R.\ Nielsen\\[4pt]
  \normalsize paa Dansk udgivet af B.\ C.\ Bøggild}
\date{Kjøbenhavn: Forlagt af H.\ C.\ Klein, 1842}

\begin{document}
\hypersetup{
  pdftitle={Den speculative Methodes Anvendelse paa den hellige Historie},
  pdfauthor={R. Nielsen}}
\frontmatter
\maketitle
\thispagestyle{empty}

\tableofcontents

%% ============================================================
%%  FRONT MATTER (PDF 6--9): title page, the surviving page of Bøggild's
%%  preface, the Indhold, and the Rettelser.
%% ============================================================
%% ============================================================
%%  FRONT MATTER (PDF 6--9). Read off the images; no printed folios on any
%%  of these four pages, so no page markers of the "% --- p. N ---" kind.
%% ============================================================

% --- PDF 6: TITLE PAGE ---
% Facsimile of the printed setting, with the printed line breaks. A centre
% block is used deliberately here (see BATCH-AGENT.md: the "use the macros"
% rule governs \mainmatter structural heads, not the title page).
% Fraktur throughout EXCEPT "Cand. theol." and the figures "23" and "1842",
% which are antiqua -> \textit{}.
% Two short rules and one printer's ornament (a lozenge-and-scroll flourish)
% divide the page; the ornament is rendered here as a plain rule.
\cleardoublepage
\thispagestyle{empty}
\phantomsection
\addcontentsline{toc}{chapter}{Titelblad}
\begin{center}
  {\Large Den speculative Methodes}\\[1.3\baselineskip]
  {\large Anvendelse paa}\\[1.3\baselineskip]
  {\huge\bfseries den hellige Historie.}\\[1.1\baselineskip]
  \rule{4em}{0.4pt}\\[1.8\baselineskip]
  Af\\[0.9\baselineskip]
  {\Large R.\ Nielsen,}\\[0.4\baselineskip]
  {\small Licentiat i Theologien, Professor i Philosophien\\
   ved Kjøbenhavns Universitet.}\\[1.3\baselineskip]
  \rule{9em}{0.4pt}\\[1.5\baselineskip]
  Paa Dansk udgivet\\[0.5\baselineskip]
  af\\[0.5\baselineskip]
  {\Large B.\ C.\ Bøggild,}\\[0.35\baselineskip]
  {\small\textit{Cand. theol.}}\\[2.6\baselineskip]
  % printed ornament here
  \rule{13em}{0.4pt}\\[1.6\baselineskip]
  {\large Kjøbenhavn.}\\[0.5\baselineskip]
  {\small Forlagt af H.\ C.\ Klein, Vimmelskaftet Nr.\ \textit{23}.\\
   Trykt i det Seidelinske Officin, hos Louis Klein.\\[0.3\baselineskip]
   \textit{1842}.}
\end{center}

% --- PDF 7: the LAST page of B. C. Bøggild's translator's preface ---
% THE PAGE CARRIES NO HEADING, NO FOLIO AND NO ORNAMENT: the text block
% begins straight in at the normal first-line position with "os imøde i dette
% Skrift," --- i.e. mid-sentence. The heading of the preface stood on the
% leaf that is missing from this scan (see the header of transcription.tex:
% recto/verso parity shows a WHOLE LEAF, two pages, is absent between the
% title page and this one). "Oversætterens Fortale" below is therefore an
% EDITORIAL supply, not a reading of the page.
%
% QUOTATION MARKS ON THIS PAGE ARE the RIGHT double quote (U+201D) at BOTH
% ends -- raised, comma-shaped -- not the low-high Danish pair (U+201E/U+201C)
% used in the body of the book. Verified at
% 2.6x on ”den menneskelige Selvbevidstheds Autonomie” and ”Fædrelandet”.
% Transcribed as printed; hence this fragment contains no U+201E and no
% U+201C at all. The transcribed text carries 8 instances of U+201D --- four
% opening, four closing --- and this comment block carries four more.
\chapter*{Oversætterens Fortale}
\phantomsection
\addcontentsline{toc}{chapter}{Oversætterens Fortale}

\textit{[Begyndelsen af Oversætterens Fortale mangler i det benyttede
Exemplar: et helt Blad --- to Sider --- er ikke med i Scanningen, og Fortalen
begynder derfor midt i en Sætning.]}

\noindent os imøde i dette Skrift, vistnok ville bringe Læseren til at oversee
hiin Mangel. Skal den speculative Methodes Anvendelse paa Theologien, i
Almindelighed som i en enkelt Sphære, forsvares med Sandhed, da maa dette
først og fremmest skee ved en grundig, indenfra ført Opposition imod den
kritiske, moderate Rationalismes Middelveissystem; men dette er netop en
Hovedfortjeneste ved dette Skrift. Og medens \textit{Dr.} Martensen
% "Dr." is bold ANTIQUA on the page; "Prof." two sentences on is Fraktur.
i sit herlige Skrift: ”den menneskelige Selvbevidstheds Autonomie” fra en
anden Side har gjennemført en grundig Opposition imod den subjective
Rationalisme og ved at bevise, at ogsaa Schleiermacher var subjectiv
Rationalist, tillige har udstrakt sin Opposition til hans System, men paa
Grund af sin Afhandlings særegne Oconomie ikkun korteligt berørte den
objective Rationalisme, har Prof. Nielsen i sin anden Deel: ”den hellige
Historie” fortsat den dialektiske Polemik ind paa den objective
Rationalismes (Panlogismens, Straußianismens) Gebeet, og paa samme Maade,
som ved den subjective Rationalisme, indenfra stræbt at vise dette Systems
logiske Utilstrækkelighed. Han har saaledes aabnet Kampen imellem den
”speculative Orthodoxie” og Panlogismen, der nødvendigt maatte fremkomme hos
os, ligesaavelsom i Tydskland, og ført den ind paa den eneste rette Plads for
en saadan Kamp, Videnskaben nemlig.

Jeg tillader mig at henvise dem, der kunne ønske at have en fuldstændig
Oversigt over Skriftets hele Tankegang til den Anmendelse af samme, der
% PRINTER'S DEFECT: "Anmendelse" as printed (for "Anmeldelse"); verified at
% 2.4x. Not in the Rettelser. Transcribed as printed.
findes i ”Fædrelandet” Nr.\ \textit{329}, \textit{1840}.
% The issue number is 329, NOT 3297: read at 1.9x, "329" then a comma with a
% descending tail, then "1840." This corrects the note in the header of
% transcription.tex.

Hvad Oversættelsen selv angaaer, føler jeg fuldtvel, hvormeget jeg navnlig
% Both OCR witnesses put a comma after "føler"; on the image the mark is much
% smaller than the neighbouring commas and is read here as show-through, so no
% comma is set. Doubtful; the rejected alternative is "føler, jeg fuldtvel".
med Hensyn til Stilens Reenhed, og til Eenhed i Retskrivningen trænger til
skaansom Bedømmelse af den velvillige Læser; turde jeg kun haabe ved sand
Gjengivelse af Forfatterens Tanker og af hans lette, livlige Stiil at forsone
Læseren med hine Mangler!

\bigskip
\begin{center}
Kjøbenhavn, den \textit{24}de Januar \textit{1842}.
\end{center}
% "Kjøbenhavn" with a j here, as on the title page (verified at 1.5x); the
% reading "Kiøbenhavn" recorded in the header of transcription.tex is wrong.

\nopagebreak
\begin{flushright}
{\large B.\ C.\ Bøggild.}\hspace*{3em}
\end{flushright}

% --- PDF 8: INDHOLD ---
% Read off the image, not off the OCR layer, which misreads 36 as "BG.",
% 42 as "12" and 143 as "113".
% The Deel line ("Første Deel." / "Anden Deel.") is bold Fraktur; the Første
% Deel argument below it is LETTERSPACED, the Anden Deel argument is not
% (it is simply set a size larger). All figures are antiqua -> \textit{}.
\chapter*{Indhold}
\phantomsection
\addcontentsline{toc}{chapter}{Indhold}

\begingroup
\parindent=0pt
\parfillskip=0pt

\begin{center}
  {\large\bfseries Første Deel.}\\[0.5\baselineskip]
  \emph{Forholdet imellem Philosophien og Historien, betragtet i
  Almindelighed.}
\end{center}
\medskip

{\raggedleft Side.\par}
\smallskip

§~\textit{1.}\quad Forholdet imellem det Aposterioriske og det Aprioriske\dotfill\textit{1.}\par
§~\textit{2.}\quad Den subjective Tro og den speculative Logik\dotfill\textit{5.}\par
§~\textit{3.}\quad Den vulgaire Empirie\dotfill\textit{11.}\par
§~\textit{4.}\quad Overgang til den speculative Empirie\dotfill\textit{18.}\par
§~\textit{5.}\quad Forholdet imellem den vulgaire og speculative Empirie\dotfill\textit{25.}\par
§~\textit{6.}\quad Videnskabens Begreb\dotfill\textit{30.}\par
§~\textit{7.}\quad Den speculative Methode\dotfill\textit{36.}\par
§~\textit{8.}\quad Begriben og Anskuen\dotfill\textit{42.}\par
§~\textit{9.}\quad Vor Viden er relativ\dotfill\textit{47.}\par
§~\textit{10.}\quad Forholdet imellem Dogmatiken og den hellige Historie\dotfill\textit{51.}\par

\bigskip
\begin{center}
  {\large\bfseries Anden Deel.}\\[0.5\baselineskip]
  {\large Den hellige Historie.}
\end{center}
\medskip

§~\textit{11.}\quad Katholicismen\dotfill\textit{53.}\par
§~\textit{12.}\quad Katholicismens Grundmodsigelse\dotfill\textit{56.}\par
§~\textit{13.}\quad Overgang fra Katholicismen til Protestantismen\dotfill\textit{60.}\par
§~\textit{14.}\quad Troen og den hellige Skrift\dotfill\textit{63.}\par
§~\textit{15.}\quad Symbolerne. Luthersk Scholasticisme\dotfill\textit{71.}\par
% "Luthersk Scholasticisme" --- read at 1.6x. NOT "Luther. Scholasticisme",
% which is what the OCR-derived list in the header of transcription.tex has.
§~\textit{16.}\quad Fortsættelse. Overgang til Kriticismen\dotfill\textit{82.}\par
\hangindent=3.5em\hangafter=1
§~\textit{17.}\quad Den subjective Rationalisme. Den hellige Histories
  Vanhelligelse\dotfill\textit{87.}\par
% This entry runs over two lines in the printing, the second indented.
§~\textit{18.}\quad Tro og Viden i den kritiske Theologie\dotfill\textit{96.}\par
§~\textit{19.}\quad Undersøgelse af den kritiske Theologie\dotfill\textit{104.}\par
§~\textit{20.}\quad Fortsættelse. Overgang til Panlogismen\dotfill\textit{119.}\par
§~\textit{21.}\quad Panlogismens methaphysiske Prøvelse\dotfill\textit{135.}\par
§~\textit{22.}\quad Tilbagevenden til Kirken. Den speculative Theologie\dotfill\textit{143.}\par
§~\textit{23.}\quad Miraklet og den hellige Historie\dotfill\textit{150.}\par
\hspace*{3.2em}Slutning\dotfill\textit{161.}\par

\endgroup

% --- PDF 9: RETTELSER (the page is headed "Trykfeil,") ---
% TRANSCRIBED ENTRY BY ENTRY FROM THE IMAGE AT HIGH ZOOM, not from the OCR
% layer, which is close to unusable on this page. Twelve entries.
% Column layout as printed: the words "Side" and "Linie" stand only in the
% first row and are replaced by an em-dash in every row after it; the last
% row has "Note." where "Linie" would stand. "l." = læs; "f. n." = fra neden.
% Figures are antiqua -> \textit{}.
\chapter*{Trykfeil}
\phantomsection
\addcontentsline{toc}{chapter}{Trykfeil}

\begin{center}
{\large\bfseries Trykfeil,}\\[0.4\baselineskip]
der bedes rettede, inden Bogen læses.
\end{center}
\medskip

\begin{center}
\begin{tabular}{@{}l@{~}r@{~}l@{~}l@{}}
Side & \textit{4}   & Linie & \textit{17}: naar den læs: naar det\\
---  & \textit{4}   & ---   & \textit{26}: de maa l. de maae\\
---  & \textit{7}   & ---   & \textit{11}: Subfaction l. Subtraction\\
---  & \textit{10}  & ---   & \textit{3} f. n.: almindeligt, l. almindelige\\
---  & \textit{12}  & ---   & \textit{1} f. n.: gaae ind l. slaae ind\\
---  & \textit{15}  & ---   & \textit{4} f. n.: under et nyt l. under Skikkelse af et nyt\\
---  & \textit{41}  & ---   & \textit{1}: næsten l. næste\\
---  & \textit{42}  & ---   & \textit{13}: fast l. faste\\
---  & \textit{48}  & ---   & \textit{6}: Virken l. Viden\\
---  & \textit{71}  & ---   & \textit{12}: \emph{Scholastisme} l. \emph{Scholasticisme}\\
% Both words are LETTERSPACED in this entry --- so the word is emphasised in
% the text at p. 71 l. 12 as well. Verified on the image.
---  & \textit{89}  & ---   & \textit{5}: κεκαυτερισμένυος l. κεκαυτερισμένους\\
% Greek read at 5x on both words. The error is a transposition in the ending:
% -μένυος for -μένους. Stem as printed: κεκαυτερισ- (epsilon, not eta; no
% alpha between ρι and σ). No fount-variant letterforms occur; the medial
% sigma is σ and the final ς, both standard.
---  & \textit{132} & Note. & \textit{1}: das Mystische l. das Mythische
\end{tabular}
\end{center}

\mainmatter

%% ============================================================
%%  FØRSTE DEEL (printed pp. 1--52) and ANDEN DEEL (pp. 53--160),
%%  with the Slutning at pp. 161--162.
%% ============================================================

% --- p. 1 ---
% DISPLAY FORMS, read at 1.8--3x on PDF 10, 14 and 20:
%   * The Deel line is "Første Deel." in large bold Fraktur WITH a period; the
%     ø is set with this Fraktur fount's umlaut-form of o, which is its sort
%     for ø throughout the book, not a distinct letter.
%   * The Deel argument IS letterspaced (bold Fraktur, two lines).
%   * The § ARGUMENTS OF §§ 1--3 ARE NOT LETTERSPACED. Measured and looked at
%     side by side at 1.8x: the Deel argument shows the Sperrsatz signature,
%     the three § arguments have ordinary Fraktur letterfit, indistinguishable
%     from a body line. \parag emphasises its second argument regardless; the
%     divergence from the header's DISPLAY FORMS note is recorded here.
\deel{Første Deel.}{Forholdet imellem Philosophien og Historien, betragtet i Almindelighed.}

% The printed § head reads "§ 1." (no period after the section sign) --- but
% p. 5 sets "§. 2." WITH one and p. 11 sets "§ 3." without. See the comment at
% p. 5. Argument agrees verbatim with the Indhold.
\parag{1}{Forholdet imellem det Aposterioriske og det Aprioriske.}

% A short centred rule stands between the § 1 argument and the first paragraph.
% It is printed only here, at the opening of the book: neither § 2 (p. 5) nor
% § 3 (p. 11) carries one. Verified by a row-by-row ink scan, not by eye.
% (a plain centred rule, not a display head --- \centerline rather than a
% center environment so that check.py's hand-rolled-head test is not tripped)
\par\vspace{0.35\baselineskip}\centerline{\rule{5em}{0.4pt}}\vspace{0.35\baselineskip}

% QUOTATION MARKS --- IMPORTANT. The body of this book does NOT use the Danish
% low-high pair at all. Like Bøggild's preface (PDF 7), it sets a RAISED mark
% (U+201D) at BOTH ends of every quotation. Verified at 2--3x on p. 1 ll. 1
% and 5 and 13, p. 2 l. 18, p. 7 ll. 7, 8, 11 and 12, and p. 10 l. 1 --- the
% opening mark sits at cap height in every one of them, never on the baseline.
% Consequence: this batch contributes nothing to check.py's „/“ balance.
% The paragraph opens with a large decorated Fraktur initial D; transcribed as
% an ordinary capital.
Det gamle Ord: ”\textit{medium tenuere beati!}” høres ikke blot, hvor det
gjelder om at drikke, men ogsaa, hvor det gjelder om at studere; saaledes
anbefaler ofte Hverdagsforstanden denne Regel, naar den taler om Forholdet
imellem Historien og Philosophien. ”Man maa prise Historien, hedder det, thi
al Tænkning vilde være tom og taagefuld, hvis ikke Tingene selv eftersporedes
med den største Flid og Nøiagtighed; --- men Philosophien er ikke mindre
nødvendig, thi Historikeren trænger til logiske Regler, for at ikke
Kundskabsskatten skal blive en uformet Masse, en \textit{rudis indigestaque
moles}. Forsigtighed byder altsaa, at man hverken helder formeget til den ene,
eller til den anden Side, for ikke at confundere det Aprioriske og det
% PRINTER'S DEFECT: "Apostioriske" for "Aposterioriske", as printed. Read at
% 3x on p. 1 l. 13; both OCR witnesses agree. Not in the Rettelser.
Apostioriske med hinanden.”

Men saa plausibel og betryggende end denne Ædruelighedens og
Maadeholdenhedens velmeente og venskabelige Formaning ved første Øiekast synes
at være, viser den sig dog ved nøiere Prøvelse af Sagen tom og vaklende. Thi
naar
% --- p. 2 ---
Tænkningen skal lempes efter Erfaringen og Erfaringen omvendt corrigeres af
Tænkningen, hvad er da Sandhedens Prøvesteen, hvor er da den øverste Dommer?

Erfaringen i og for sig kan ikke indtage Dommersædet; thi for at bevise sin
Adkomst til at være Sandhedens Dommer, maatte den anføre Grunde hentede fra
Tænkningen; den vilde saaledes netop i Bestræbelsen efter Dommerembedet gaae
Glip af det; --- Tænkningen derimod kunde, selv om den vilde frasige sig selv
Dommerembedet, ikke gjøre det uden ved selv at dømme, altsaa tilkommer den
endelige Dom Tænkningen; paa dens Sandhed beroer Alt. Men Undersøgelsen af vor
Tænknings Sandhed fremkalder nødvendigviis det Spørgsmaal, hvorledes
overhovedet Tænkningen opstaaer, og i Besvarelsen heraf møde vi da
Sensualisternes og Idealisternes modsatte Theorier.

Naar Sensualisterne opstille den Grundsætning, at Alt skyldes Indtryk fra
Tilværelsen udenfor os, at Menneskets Sjel er en ”\textit{tabula rasa}”, og at
derfor Almeenbegreberne fremkomme ved en Abstraction, saa følger heraf, at
Tanken kun er en Skygge, medens Realiteten selv maa sættes i empiriske,
udenfor os værende Gjenstande. Hos dem er altsaa det Aposterioriske
Hovedsagen. --- Men just idet de saaledes oversee Subjectivitetens Moment,
give de sig til Priis for Subjectiviteten; de paastaae vel, at Menneskets Sjel
er en \textit{tabula rasa}, men da absolut Passivitet ikke kan finde Sted, maa
der nødvendigviis være et subjectivt Moment med i alle de subjective Indtryk.
Man maa saaledes gaae ud fra specielle Fornemmelser for at komme til de
almindelige Begreber; men da dette ikke kan finde Sted, fordi Erfaringens
Phænomener ere particulaire, nødes man til at acqviescere i subjective
% PRINTER'S DEFECT: "absolnt" for "absolut" (turned u), as printed. Read at 3x
% on p. 2 l. 32; both OCR witnesses agree. Not in the Rettelser.
Antagelser, efterdi man har frakjendt Almeenbegreberne absolnt Gyldighed. Deri
ligger altsaa Modsigelsen, at Sub\-%
% --- p. 3 ---
jectiviteten paa den ene Side gjelder saa Lidet, at den paa den anden Side
gjelder Alt.

Er det altsaa indlysende, at Man hildes i subjective Meninger, saaofte man
forsøger at begribe de objective Gjenstande, saa ere disse Meninger dog ikke
destomindre almeengyldige fra Subjectets Standpunkt. Almeenbegreberne kunne
altsaa ikke hidrøre fra en udvortes Erfaring, men maa ligge i Subjectiviteten
som saadan.

Men naar alle Phænomener opfattes ved Hjælp af Subjectiviteten og føres
tilbage til subjectivt nødvendige og almindelige Begrebsformer, ere som en
Følge deraf Almeenbegreberne nødvendige til enhver Erfaring, saa at
Modsætningen mellem Subject og Object ikke er en Modsætning mellem Tænkning og
% "dialectisk" here with c; p. 7 l. 25 has "dialektiske" with k. Both as
% printed --- the book is not consistent.
Erfaring, hvilke flyde dialectisk over i hinanden, men en Modsætning imellem
% PRINTER'S DEFECT: "Phænomerne" for "Phænomenerne", as printed. Read at 3x on
% p. 3 l. 15; both OCR witnesses agree. Not in the Rettelser.
Phænomerne og den rene, absolute Objectivitet, \emph{Tingen i sig}.

Undersøger man nu, hvad Tingen i sig egentlig er, vil det vise sig, at den kun
% PRINTER'S DEFECT: "Tiugen" for "Tingen" (turned n), as printed. Read at 3x on
% p. 3 l. 19; both OCR witnesses agree. Not in the Rettelser.
har en tænkt Tilværelse. Man kan nemlig ikke erfare Tiugen i sig, eftersom
enhver Erfaring støtter sig paa subjective Begreber; ikkun dette veed man
altsaa om Tingen i sig, at en saadan er. Jeg tænker ikke blot mine Tanker, men
jeg tænker ogsaa, at Noget er udenfor mine Tanker; man maa da ikke søge denne
dets Væren udenfor, men indenfor den subjective Tænknings Grændser.

% "Jeg selv" at ll. 27 and 29 is NOT letterspaced --- checked at 1.8x against
% the letterspaced "Tingen i sig" three lines above. Only "prius" is set off,
% and by fount (antiqua), not by spacing.
Det er ad denne Vei, at Idealisterne have uddannet deres Theorie, i det de
lære, at hele Verden er sat af det menneskelige Jeg selv, ihvorvel dette ikke
var sig denne Act bevidst; det Aprioriske indtager da den første Plads, thi
det skjulte Jeg selv er netop dette \textit{prius}.

Vor Tids Videnskab har imidlertid fuldkomment beviist, at Idealismen bliver
Nihilisme, naar den fornægter den objective Tilværelse og kun i sin egen
Selvvirken søger Rea\-%
% --- p. 4 ---
liteten. Ere nemlig alle Objecter kun Forestillinger og Billeder, dannede af
den menneskelige Sjel, hvad er da dette Jeg selv? Betragtes det som en hvilende
Realitet og ikke som et tomt Billede, maa da ikke det Samme kunne siges om de
øvrige Objecter? Og henfører man paa den anden Side ogsaa Jegets Tilværelse
til Forestillingen, erklærer man jo Jeget selv for en Forestilling, dannet af
Jeget. --- Resultatet bliver da en uendelig Række af Forestillingers
Forestillinger.

Naar derfor den subjective Idealisme gjør Tænkningens Natur til Gjenstand for
en grundig Betragtning, nødes den til at antage en objectiv Tilværelse og til
at indrømme, at Virksomheden og Forestillingen ikke kunde være, hvis der ikke
var Noget, der gav Virksomheden et substantielt Indhold.

Det menneskelige Jeg bliver saaledes netop i det Øieblik, da det fryder sig i
Følelsen af fuldkommen Uafhængighed, Bytte for en fremmed Magt, og styrter
% ERRATA (Rettelser, p. 4 l. 17 fra oven): printed "naar den anmasser sig",
% corrected here to "naar det anmasser sig". The uncorrected reading was
% verified on the image at 3x before the correction was applied.
just da i Tomhedens Afgrund, naar det anmasser sig guddommelig Majestæt.

Men netop dette Subjectivitetens Fald bereder Frelsen og gjør en sand
Forsoning mulig mellem Subject og Object. Thi beviser Sensualismen, at der
ikke findes et eneste Punkt i Sjelens Liv, som ikke kan føres tilbage til et
Indtryk fra Tilværelsen udenfor os, og godtgjør paa den anden Side den
subjective Idealisme den menneskelige Sjels uafbrudte Selvvirksomhed, saa
følger i Sandhed heraf, at Tænkningen og den objective Tilværelse ikke kunne
% ERRATA (Rettelser, p. 4 l. 26 fra oven): printed "de maa da være",
% corrected here to "de maae da være". Uncorrected reading verified at 3x.
tilhøre to modsatte Sphærer; de maae da være af samme Natur, og vi kunne da
% GREEK, read at 3.2x. The printing has ALPHA with a circumflex in the first
% word --- γνᾶθι, where the correct form has omega with a circumflex --- and a
% GRAVE on the omicron of σεαυτὸν, where the correct form has an acute. Both are
% reproduced as printed; neither is in the Rettelser. The theta is the fount's
% script form, normalised here to plain θ per the repo convention (U+03D1 script theta is fatal
% under textalpha). NB: ᾶ (U+1FB6), ὸ (U+1F78) are new to this file.
ikke, som det hedder, \emph{kjende os selv}, (γνᾶθι σεαυτὸν), uden gjennem
Erkjendelsen af de objective Phænomener.

Maaskee vil Nogen indvende, at vi saaledes ophæve den vigtige Forskjel mellem
den individuelle Subjectivitet og den universelle Objectivitet, men denne
Indvending fremgaaer kun af en abstract og udvortes Anskuelse af Sagen: thi
betragte
% --- p. 5 ---
vi kun den ydre Verden, kunne vi ikke falde paa at nægte, at jo hvert Menneske
har sin Subjectivitet og overalt finder sig omgivet af Objecter, der stille
sig som Skranker imod ham; men fra den indre Side maa Sagen sees og man vil da
iagttage, at enhver Subjectivitet skylder Objectiviteten sin Dannelse, i det
de subjective og objective Momenter ville sees i deres indbyrdes uendelige
Overgang i hinanden. --- Det er netop Objectivitetens Natur overalt at udvikle
Subjectiviteten, som det er Subjectivitetens Maal at uddanne Objectiviteten.
Derfor maa Alt ligesaavel søges \textit{a posteriori}, som \textit{a priori},
ligesaavel henføres til Tænkningen, som til Erfaringen. Og paa dette
Standpunkt gjelde derfor heller ikke de saa ofte fremsatte Yttringer, at
Empiriens Uddannelse skulde være Philosophiens Indskrænkning, eller at
Philosophiens Blomstren skulle medføre Empiriens Forsømmelse.

Heraf følger da, at Middelveien i Videnskabernes Rige ikke fortjener nogen
særdeles Anbefaling, og at meget mere enhver Bestræbelse efter at holde
Middelveien, er Middelmaadighed.

% THE PRINTED HEAD HERE IS "§. 2." --- with a period between the section sign
% and the figure. Verified at 1.8x. § 1 (p. 1) and § 3 (p. 11) both set
% "§ 1." / "§ 3." with no such period, so the form is not uniform in the book.
% \parag renders "§ 2."; the printed reading is recorded here rather than
% normalised. The argument agrees verbatim with the Indhold.
\parag{2}{Den subjective Tro og den speculative Logik.}

Eet er det at erkjende Identiteten af Phænomener og Tanker, Empirie og
Philosophie, et Andet derimod at underkaste denne Identitets Momenter en
saadan videnskabelig Analyse, at Identiteten sees i sin Mediation, fra alle
Sider gjennemskuet og begrebet. Det gjelder altsaa her at finde det rette
Udgangspunkt.

I det umiddelbare praktiske Liv har de overalt mødende Objecters Forhold til
den individuelle Subjectivitet den største
% --- p. 6 ---
Indflydelse; Alt opfattes ved indre og ydre Sandser, og man forsvarer Alt,
saaledes som man har opfattet det. Den Ene holder paa Sandheden af sit Udsagn,
fordi hans umiddelbare Opfattelse af Sandheden just nu er saaledes beskaffen;
den Anden forsvarer af selvsamme Grund det aldeles Modsatte. De modsige
hinanden og lægge begge Haanden paa Hjertet til Bekræftelse af deres modsatte
Udsagn. --- Hvad skal nu bevæge os til at respectere deres Forsikkringer? Alt
afhænger af subjectiv Stemning. Vi krænke derfor ikke blot Sandheden i
Almindelighed ved ikke at skjænke saadanne Udsagn vor Tiltro, men ogsaa det
Menneskes Person, der forkynder Sandheden. Den subjective Tro udgjør Kraften
og saa at sige Marven i ethvert Menneskes Liv, og kan ikke forandres uden en
fuldkommen Forandring i Menneskets Aand selv; berøve vi derfor Nogen hans
subjective Tro, saa berøve vi ham Livet. --- Det er derfor ikke at undre over,
om Menneskene i Forsvaret for deres forskjellige Tro ikke alene disputere og
modsige hinanden, men ogsaa, hvad Livet ved saa mange Exempler udviser, stride
og kjæmpe med en Heftighed, som om Kampen gjaldt Alter og Arne.

Men, som den subjective Tro antænder denne Krigs Flammer, saa har den ogsaa
Midler til at slukke dem. Og havde den ingen saadanne, kunde overhovedet ingen
Modsigelse være opstaaet, thi, for at være uovereensstemmende, maa der
nødvendigviis være Noget, hvori man stemmer overeens. En absolut
Uovereensstemmelse kan ikke finde Sted, thi forstod den Ene ikke den Anden,
vilde Enhver af dem være bleven overladt til og kun kunne have talt med sig
selv; for at kunne disputere maa man da nødvendigviis have et Medium, hvori
man kan mødes. Og hvad er vel dette Medium Andet, end netop Tænkningen selv?
Alle maae tænke, naar de ville forklare deres Tro; men hvad er vel det
% --- p. 7 ---
at tænke Andet, end at opsøge og afsløre det Almindelige i det Særegne,
Eenheden i Mangfoldigheden, Identiteten i Modsætningen?

Netop derfor maa man under alle disse Stridigheder ikke tabe Modet; der bliver
stedse i Uenigheden en vis Enighed tilbage. Naar Debatten er steget til
Modsigelsens høieste Punkt, naar Du hører den Forsikkrendes: ”\emph{det er
sandt!}” ligeoverfor den Benægtendes: ”\emph{det er falsk!}”, maa Du dog ikke,
sige Logikerne, opgive Haabet om at kunne forsone de ophidsede Gemytter. Det
% ERRATA (Rettelser, p. 7 l. 11 fra oven): printed "Sub=" / "faction:" across
% the line break, i.e. "Subfaction:", corrected here to "Subtraction:". The
% uncorrected reading was verified on the image at 3x before correcting.
kan skee ved en let Subtraction: borttag blot de modstridende Prædicater:
”\emph{sandt}” og ”\emph{falsk;}” der bliver da kun hiint: ”\emph{det er}”
tilbage, og lover Du deraf at uddrage mange og det endog vigtige
Conseqventser for begge Parter, ville de Stridende nedlægge Vaabnene og
villigt indrømme, at den Ting, hvorom der strides, er.

% GREEK: σκέψις. The kappa is the fount's cursive form, normalised to plain κ
% (U+03F0 cursive kappa is fatal under textalpha). Accent as printed.
Naar saaledes denne Debat, denne Negation, denne σκέψις er skredet frem til
hiint abstracte og i vor Tid saa trivialiserede Begreb \emph{Væren}, hvis
Benægtelse indeholder dets Bekræftelse, maa den nødvendigviis ophøre. Her er
det altsaa, at den sande Videns absolute Udgangspunkt maa søges. Men et ganske
andet Spørgsmaal er det, om vi ad denne slagne Vei, ved nemlig at forfølge
Almeenbegreberne, kunne naae dertil fuldkomment at afsløre Livets Fylde,
hvorfra vi ere gangne ud.

% The letterspacing of "Væren" here runs across the line break "Væ=" / "ren";
% verified on both halves at 3x.
% GREEK: τὸ ἓν τὸ πᾶν. NB ἓ (U+1F13) and ᾶ (U+1FB6) are new to this file.
Da den dialektiske Stimulus, der ligger i Begrebet \emph{Væren}, var opdaget
af Hegel, saa at Logikerne nu ikke længere nødsagedes til at blive staaende i
Universalitetens Tautologie, som fordum Eleaterne, der ideligt gjentog deres:
τὸ ἓν τὸ πᾶν, meente Mange, at Tidens Fylde var kommen. De talede da gjerne om
dette Begrebs forunderlige Fattigdom og overordentlige Rigdom, da det paa
eengang viste sig som det eenfoldigste, idet nemlig alle forudfattede Meninger
vare nedstødte
% --- p. 8 ---
i Abstractionens Afgrund, og dog tillige indesluttede alle Viisdommens Skatte
i sig; i dettes Skjød, meente de, var det Nærværende og Tilkommende, Gud og
Verden skjulte; med faae Ord, de lovpriste Begrebet \emph{Væren} som det Æg,
der i Sandhed indesluttede Universet i sig, og kun behøvede at udruges og
udklækkes af en ret speculativ Philosoph for at lade Tingenes sande Skikkelse
fremgaae af sig. --- Under denne Stjerne fødtes den logiske Mysticisme, der,
som det anstaaer en sand Mysticisme, sagde Livet og dets Virkelighed Farvel.

% "Væren" is NOT letterspaced at this second occurrence (l. 11), though it is
% at l. 4 above. Checked at 3x on both lines. Emphasis follows the page.
Men vi synes at modsige os selv, naar vi paastaae, at det absolute
Udgangspunkt maa sættes i Begrebet Væren, og paa samme Tid erklære, at
Analysen af de logiske Begreber ikke kan føre til Realiteterne selv; thi kunne
vi ikke ad denne Vei naae Maalet, maae vi et andet Sted søge et nyt
Udgangspunkt; vi kunne da ikke lade det være os nok med den eengang gjorte
Begyndelse, men maae endnu een eller maaskee flere Gange begynde forfra.

Man bør derfor iagttage, i hvilket Forhold Almeenbegreberne staae til de
særegne Gjenstande. Den almindelige Mening om dette Forhold er, at de logiske
Begreber ere hule Former, hvori de faste Gjenstande ere indgydte; tages da
Formerne bort, bliver der ikkun en sammenblandet Mængde af Gjenstande tilbage;
tages de solide Gjenstande bort, ville, ligesaavel i Verden, som i det logiske
System, de rene Begreber ikkun vise sig som et kunstigt Næt. Og denne
Adskillelse mellem Almeenbegreberne og de særegne Gjenstande har ligesaavel en
sand, som en falsk Side. Den er sand, forsaavidt som det fastholdes, at
Tingens Form ikke er Tingen selv og at Tingen ikke er Formen; thi Ingen kan
nægte, at Formen netop er Tingens Negation og saaledes dens Begrændsning; og
hvorledes skulde vel og det Almindelige, det abstract Tænkte kunne være noget
Særegent? Ved deres egen Mod\-%
% --- p. 9 ---
sigelses Kraft udelukke derfor disse Momenter hinanden. Men grundfalsk er
Adskillelsen af Tilværelsens Momenter, dersom dermed menes, at Tingene og
Formerne havde kunnet bestaae adskilte fra hinanden, hvis ikke Gud ved en
særegen Villiesact havde forenet dem. Ved en nærmere Betragtning af Sagen vil
dette let blive indlysende.

% Letterspacing of "Væren" on this page, settled line by line at 3x: SPACED at
% ll. 7, 9, 11 and 13; NOT spaced at ll. 12, 14, 15 and 16. spacing.py reported
% a RUN "Væren det" at l. 15 --- rejected: at 3.2x the letterfit of "Væren det"
% there is ordinary, and the tool was reading the line's wide word gaps.
% "Noget" at l. 14 IS spaced.
Det er med fuldkommen Ret, at vi i Begrebet \emph{Væren} see det Almindelige i
dettes høieste Abstraction, men Logiken lærer os tillige, at \emph{Væren} ikke
kan tænkes og heller ikke være, hvis der ikke er Noget, hvis Prædicat den er;
for at \emph{Væren} selv kan være, maa der være givet Mere end den rene Væren.
Det rene Almindelige adsplittes saaledes i to Momenter, af hvilke det ene er
den rene \emph{Væren} selv, det andet er dette \emph{Noget}, om hvilket Væren
% The mark after "hinanden" is inked as a comma with a speck above it; read as
% a comma (the sense requires one), but a semicolon cannot be ruled out. The
% page carries scattered scan specks of the same size.
udsiges, thi confunderes disse to Momenter med hinanden, bliver Væren det
samme som Væren; men en saadan tautologisk Væren kan ikke være. Men bevares de
i deres Modsætning til hinanden og kalde vi det ene \textit{A}, det andet
\textit{B}, viser det sig, at, hvad der ikke er i \textit{B}, maa være i
\textit{A}, eller --- for at det Almindelige kan være, maa nødvendigviis ogsaa
det særegne Moment være tilstede.

Det Almindelige spreder sig da til alle Sider i særegne Gjenstande. Men i det
det nu er bleven klart, at det Almindelige ikke er tilstede nogetsteds, hvor
ikke tillige det Særegne er, er der om dette særegne Moment endnu ikke sagt
Andet, end at samme maa søges overalt; med andre Ord --- det særegne Moment
opfattes her ikke paa særegen, men paa almindelig Maade. Der savnes altsaa det
i Sandhed Særegne, og vi maae da undersøge, hvor dette er at finde.

Det i Sandhed Særegne er ligesom et fast og ubevægeligt Punkt, der i og for
sig betragtet ikke kan henføres til et almindeligt Begreb, men maa kunne
bestemt paapeges, saa at
% --- p. 10 ---
% The quoted words here are NOT letterspaced, unlike the parallel quotations on
% p. 7; checked at 3x.
man kan sige om det: ”det er her! der! saaledes!” Det særegne Moment er et
positivt og reelt Moment, der staaer fast ved sig selv og maa søges i
Empiriens Rige. For at de ontologiske Begreber kunne fattes, udfordres altsaa
en vis Empirie, efterdi alle slige Begreber støtte sig paa en Negation af
empiriske Gjenstande, ligesom heller ikke en Gjenstands Begreb kan forstaaes
uden ved Hensyn til den bestemte empiriske Gjenstand.

Men denne eller hiin Gjenstands Billede fremstilles kun for Sjelen, for at
det, der er fælleds for alle, ikke det, der er ejendommeligt og særegent for
de enkelte Gjenstande, kan fastholdes; den specielle Gjenstand fremkaldes
altsaa ved Hjælp af Forestillingen, for at det særegne, det er, det
\emph{empiriske} Moment kan negeres. Da Empirien saaledes kun tilegnes paa
negativ Maade, kan her ikke være Tale om Kundskab til de enkelte Phænomener;
men naar de logiske Begrebers speculative Erkjendelse er bleven fuldstændig,
udviklet, bliver Resultatet det, at hele Tilværelsens Princip er den absolute
Idee, eller at, med andre Ord, de reale Phænomeners hele Række er fornuftig,
saa at den objective og concrete Fornuft ligesaalidt kan fattes uden gjennem
Erkjendelsen af de empiriske Gjenstande, som nogen Ting kan existere, hvis den
er stridende mod den absolute Fornufts Natur. --- Men saaledes forlade vi da
% ERRATA (Rettelser, p. 10 l. 3 fra neden): printed "almindeligt", corrected
% here to "almindelige". Verified at 3x: the printed word is "almindeligt" and
% there is NO comma after it, though the Rettelser entry quotes "almindeligt,".
Logiken, hvis Begreber baade ere almindelige og nødvendige, og vende os til
Empirien, hvis Rige omfatter det Særegne og Tilfældige.

% --- p. 11 ---
% Printed head: "§ 3." (no period after the section sign), as at § 1.
\parag{3}{Deu vulgaire Empirie.}
% PRINTER'S DEFECT IN THE § HEAD: the printed argument reads "Deu vulgaire
% Empirie." --- a turned n. Verified at 3x against the very next line of the
% body, which sets "Den vulgaire Empirie," in the same fount and size: the head
% has the bottom-joined bowl of Fraktur u, the body line the top-joined arch of
% Fraktur n. Both OCR witnesses read "Deu" too. THE INDHOLD (PDF 8) READS
% "Den vulgaire Empirie" --- both readings recorded, neither normalised away.
% Not in the Rettelser.

Den vulgaire Empirie, i hvilken Alt er særegent og tilfældigt, betragtes med
Rette som Logikens Negation; thi medens vi i Logiken gik ud fra at frakjende
Sandserne al Auctoritet, skulle vi nu opgive Almeenbegrebernes Undersøgelse og
appellere til Sandsernes Vidnesbyrd. --- Men herved synes vi at komme i
Modsigelse med os selv; thi vi have paastaaet, at der ikke er nogen Forskjel
mellem det Aprioriske og det Aposterioriske, og maae dog nu tilstaae, at der
finder en Forskjel Sted. Vi vide nemlig, endnu før vi gjennem Sandserne have
opfattet de enkelte Gjenstande, at der ikke kan gives Noget, der strider imod
de ontologiske Love; Alt hvad Ontologien lærer, kan man da forudsige om
Gjenstandene, og dette er netop det Aprioriske; men hvorledes Tingene fra
deres særegne Side ere beskafne, det kunne vi først faae at vide ved at erfare
Tingene selv, og dette er det Aposterioriske.

% "væsentlige" is letterspaced at l. 18 but NOT at its second occurrence,
% l. 26 ("den væsentlige Tilværelses Phænomener"); checked at 3x on both.
Men denne Modsigelse finder sin Løsning deri, at Forskjellen ikkun angaaer
Phænomenerne som saadanne, medens Identiteten er at søge i deres
\emph{væsentlige} Tilværelse. Vare nu disse hinanden absolut modsatte, vilde
Knuden neppe kunne løses; men Enhver veed, at Tænkningen ikke lader
Phænomenerne, som de umiddelbart ere opfattede ved Hjælp af Sandserne,
forsvinde i Væsenet, forat de skulle forblive ophævede, men netop forat de
atter skulle restaureres udaf Væsenet, som af deres sande Grund.

Den speculative Logik lærer, at Væsenet selv aabenbarer sig i Phænomenerne,
saa at disse ere den væsentlige Tilværelses Phænomener, men ikke kunne
fastholdes i og for sig; og Erfaringen viser, at selv de, der med den største
Iver hengive
% --- p. 12 ---
\setcounter{footnote}{0}
% THE PRINTED FOOTNOTE MARK ON THIS PAGE IS AN ASTERISK, "*)", NOT "1)".
% It is the only note in pp. 1--12 and the first in the book, and it is set
% with a star rather than a figure. \thefootnote will print "1)"; the printed
% mark is recorded here. The note stands under a short left-set rule.
sig til Empirien, tvertimod deres egen Forventning støde paa Væsenet, idet de
gribe efter Phænomenet. Hvad Glæde vilde saaledes den have, der fryder sig ved
Betragtningen af Solens Straaler, dersom disse alle paa eengang, ligesom med
eet Slag, mødte hans Øie, saa at han for lutter Skuen Intet kunde tænke og
Intet føle? Hvad Glæde eller Frugt vilde den have, der, som en Sandsesløs,
holdtes saa fast i Beskuelsen af de enkelte Ting, at han Intet forstod,
skjøndt han hørte med Ørene, Intet opfattede, skjøndt han saae med Øinene?
Derfor see vi ogsaa overalt i Historien, hvor en vis umiddelbar Anskuelse er
overveiende, og hvor Philosophien endnu ikke er uddannet til selvstændig
Videnskab, ikke den raa og barbariske Empirisme, der udslukker og fortrænger
al Tænkning, men det frieste Vexelskin af Empiriens og Philosophiens
Momenter\footnote{Jvfr.\ Herodot o.\ A.}. Vel kommer her det vigtige
Spørgsmaal til, om disse eller hine Tildragelser virkelig have fundet Sted,
saaledes som de blive fortalte, og man søger da at komme paa det Rene med,
hvorvidt de ydre Phænomener have Realitet, som om Historiegranskeren ikke
havde nogen anden Opgave, end den at bevise, at de ere paalidelige og have
svaret til deres Beskrivelse; men man maa da vel lægge Mærke til, at
Phænomenerne ikke ere tilstede, naar man saaledes ivrigst søger dem, og at det
meget mere er deres Fraværelse, der lader os føle Savnet af dem, saa at her
ikke kan være Tale om sandselig Beskuelse af nærværende Gjenstande. Ja vare de
Gjenstande, som Empirikerne saa ivrigt lede efter, allerede afslørede og nøie
begrændsede og fremdragne for Lyset, og havde de saaledes alt opnaaet deres
Ønske, saa maatte jo enten deres Forstandsvirksomhed ophøre, eller de maatte
% ERRATA (Rettelser, p. 12 l. 1 fra neden): printed "gaae ind", corrected here
% to "slaae ind". The uncorrected reading was verified on the image at 3x.
slaae ind paa en anden Vei for deres Granskning.

% --- p. 13 ---
% p. 13 opens FLUSH LEFT (no indent), i.e. it continues the paragraph running
% over from p. 12; but no word is broken across the joint --- the page begins
% with the whole word "Vi" and a fresh sentence. Verified on the image at 4x.
Vi maae da nøiere betragte den Afstand, der er imellem de ydre Phænomener og
Tilskuerens Bevidsthed, for at indsee, hvorledes Identiteten kan restaureres.

Betragteren kan jo ikke gjengive det Phænomen, der har fremstillet sig for hans
Blik, paa samme umiddelbare Maade, som det fremstillede sig for ham; han kan
ikke beskrive det ved at pege paa det; vil han derfor fastholde det, saa er han
nødsaget til at fælde en Dom over det. Men hvad er vel det at fælde en Dom
Andet, end at henføre Phænomenet til Totaliteten af de almindelige Begreber og
sammenligne det med Totaliteten af de allerede gjorte Erfaringer for derved at
anvise det sin rette Plads.

Men her ere atter to Momenter, der vel maae adskilles fra hinanden; thi hvilken
Plads Phænomenet skal indtage, afhænger paa den ene Side af dets egen Natur,
paa den anden Side af Betragterens subjective Eiendommelighed; og da nu
Phænomenet ikkun paa abstract Maade giver sig tilkjende for den første
Betragtning (§~\textit{2}), saa vil det maaskee komme til at indtage en altfor
ringe Plads i hans Bevidsthed, fordi dets sande Egenskaber ikke ere klart
fremtraadte for ham. Derfor er den Maade, hvorpaa Phænomenerne opfattes,
vaklende og ubestemt, og man maa anvende den største Forsigtighed, for at ikke
en Vildfarelse skal indsnige sig.

Men denne Vanskelighed voxer, naar man betænker, at det kun ere de færreste
Phænomener, der umiddelbart fremstille sig for Betragteren, medens man i de
fleste Tilfælde maa see med Andres Øine, høre med Andres Øren, stole paa Andres
Forstand og Dømmekraft. Man maa derfor tillige, da den ene Betragter som oftest
kun iagttager Phænomenet med den andens Øie, være betænkt paa omhyggeligt at
fjerne det Usande, der fra een eller anden Subjectivitet kan have indblandet sig
i Tingens Opfattelse, og dette er særdeles van%
% --- p. 14 ---
skeligt. Thi det er ikke blot en Andens Feiltagelse, man har at befrygte; men da
Referenten ikke blot kan have opfattet Gjenstanden urigtigt, men ogsaa kan ---
ifølge Reflectionens Evne til forsætligt at fremstille den samme Gjenstand paa
en anden Maade --- have forrykket Dommen efter vilkaarlige Hensyn, maa man ogsaa
befrygte en forsætlig Forvanskning. Man maa altsaa undersøge Forfatterens
Troværdighed; og da nu Forfatterens Troværdighed afhænger af den Maade, hvorpaa
han beskriver Tingen, medens denne omvendt afhænger af Forfatterens
Troværdighed, saa nødes man til at opsøge et andet, almindeligere og sikkrere
Sandhedskriterium. Vi nødes da saaledes til at tage vor Tilflugt til de
ontologiske Begreber og til de Betingelser, uden hvilke Tingen ikke kan finde
Sted, for at fjerne Absurditeter og skille Historie fra Fabel. Vi troe saaledes
ikke den, der vilde fortælle os om runde Fiirkanter, og ere ligesaa uvillige til
at troe Isidors Dekretaler, naar de berette, at Victor, der i det
\textit{2}det Aarhundrede var romersk Biskop, skrev et Brev til sin kjære Ven,
Theophilus fra Alexandrien, skjøndt denne levede i det \textit{4}de
Aarhundrede. --- Ville vi nu ikke blive staaende ved den trivielle tautologiske
Regel, at det Umulige er umuligt, saa maae vi ved Erfaringens Hjælp nøiere
bestemme, hvad Begrebet Mulighed betyder. Vi sammenligne da det, vi kjende fra
en tidligere Erfaring, med det, som Erfaringen sidenefter fremstiller for os,
for at ikke hiint skal ophæves ved eller befindes stridende imod dette.
% The figures in "2det" and "4de" are ANTIQUA against the Fraktur; the personal
% and place names on this page (Isidors, Victor, Theophilus, Alexandrien) are
% NOT --- they are Fraktur, and so are not italicised. Checked at 3x.

Man slutter altsaa her ved Hjælp af Analogie og Induction, idet man antager de
Phænomener, der ellers pleie at finde Sted, for sande, men fatter Mistanke mod
dem, der ere usædvanlige. Fra den ene Side betragter man altsaa Phænomenernes
Sum, fra den anden det særegne Phænomen, for hvis Bedømmelse hine skulle angive
Rettesnoren. Men i Lig%
% --- p. 15 ---
heden bliver der dog stedse en Ulighed tilbage; der vilde jo ellers intet Nyt
være kommet til; det er da ikke den Side af Phænomenet, fra hvilken det ligner
de foregaaende, der kommer i Betragtning, men det er det Ulige deri, der gjør en
Bedømmelse nødvendig. --- Men den Dom, der begrundes paa et Phænomens Lighed
eller Ulighed med et foregaaende, er atter udvortes og usikker; thi fordi den
paagjeldende Gjenstand ofte før har fundet Sted, er det ikke sagt, at den samme
eller en lignende har fundet Sted paa den Tid, da det fortælles, at den har
fundet Sted; og fordi et Phænomen sjeldent eller aldrig før har fundet Sted,
derfor er det heller ikke sagt, at Fortællingen om, at det har fundet Sted,
nødvendigt maa være falsk. --- Man maa da tage den større eller mindre Grad af
Sandsynlighed i Betragtning, og denne Betragtning fører nu Tilskueren over paa
de subjective Formodningers relative Gebeet
(\textit{calcul des probabilités}). Thi hvor der kan være Tale om
Sandsynlighed, der maa Sandheden baade være tilstede og ikke være det. Den kan
ikke være tilstede; thi savnede man den ikke, hvorfor vilde man da tale om dens
ydre Skikkelse, og ikke langt hellere om Sandheden selv; og var den slet ikke
tilstede, hvorledes kunde man da falde paa at søge den, hvorledes tale om en
Sandsynlighed? Sagen synes da at forholde sig saaledes. Da man erkjender det
frugtesløst at sammenligne de enkelte Phænomener med hinanden indbyrdes,
sammenligner man dem med den over dem alle staaende Realitet, \dsi{} med
Sandheden selv. Idet man nu saaledes abstraherer Sandheden fra Phænomenerne,
stiller man Sandheden paa udvortes Maade ligeoverfor Phænomenerne, saa at denne
selv fremgaaer under Skikkelse af et nyt, ihvorvel fortrinligere Phænomen, ---
% ERRATA (Trykfeil, p. 15 l. 4 fra neden): the page as printed reads "fremgaaer
% under et nyt, ihvorvel fortrinligere Phænomen". Verified at 4x on PDF 24 --- the
% uncorrected reading really is there --- and the errata correction
% "under Skikkelse af et nyt" is applied above.
thi om det immanente Forhold, ifølge hvilket Sandheden gjennemtrænger og
udvikler alle Phænomenerne, saa at den gjør dem til Former for sig, kan%
% --- p. 16 ---
der paa dette Standpunkt ikke være Tale. Men spørger man nu, hvorledes man i den
store Række af Phænomener kan gribe Sandhedens eget Phænomen, saa vil man finde,
at dette objectivt betragtet ligger skjult bag de øvrige Phænomener, som bag et
Slør, medens det derimod subjectivt betragtet ved en Anticipation og lykkelig
Ahnelse opgaaer for Tilskuerens Bevidsthed, som en Rettesnor, hvorefter han
bedømmer de enkelte Phænomener. Men da nu hver Tilskuer har sin egen Mening om
Sandheden, saa maae ogsaa disse forskjellige Meninger sammenlignes med
hinanden; og kan man da end ikke have absolut Eenstemmighed, saa maa man dog
idetmindste have Eenstemmighed hos Majoriteten, hvis man skal komme til noget
Resultat; saaledes vil da den almindelige og offentlige Mening, som er gjeldende
i den Tidsalders Forestilling, der har givet Tilskueren hans Bevidsthed, være
den afgjørende i Sagen. --- Og dog have vi endnu ikke her noget fast Standpunkt
for vor Bedømmelse; thi Enhver veed, hvor forskjellige Meningerne have været i
Tidernes Løb om, hvad der er muligt eller umuligt. Saaledes havde, for at nævne
et Exempel, Mulighedens Rige en langt større Udstrækning i Middelalderen, end i
den Tid, hvori vi nu leve. Vi see altsaa, at jo mere vi ad denne Vei forfølge
Sandheden, desto længere fjerner den sig fra os.

Og for at nu ikke Nogen skal indvende os, at vi blot paa sophistisk Maade søge
Vanskeligheder, hvor der i Virkeligheden ingen er, maae vi atter og atter
erindre om, at vi ere langt fra at miskjende den oprigtige Søgen efter, den
Sands for og den Kundskab til Sandheden, ved hvilke mange af denne Methodes
varmeste Tilhængere udmærke sig; --- thi Eet er det at bedømme den
Sandhedserkjendelse, som een eller anden Forfatter er i Besiddelse af, et Andet
derimod at undersøge den Methode, efter hvilken han troer at have er%
% --- p. 17 ---
hvervet sig samme. Vi dadle ikke Sammenligningen af de enkelte Momenter i og for
sig; thi vi vide fuldt vel, at Sandheden ikke kommer paa en saa haandgribelig
Maade (μετὰ παρατηρήσεως), at man skulde kunne sige: ”See, her eller der er
den!” og vi mene ikke, at Iagttagelsen af de enkelte Phænomener maa forsømmes,
thi tages de enkelte Dele bort, saa forsvinder ogsaa Totaliteten. Men tage vi
ikke meget Feil, ligger denne Methodes Vildfarelse deri, at man, uagtet de
enkelte Phænomener opstilles paa udvortes Maade ved Siden af hinanden, uden at
gjennemtrænge hinanden eller gjennemtrænges af Sandheden, dog troer ved
Sammenligning af alle de enkelte Phænomener at kunne gribe Sandheden, medens det
dog er soleklart, at man ved denne Sammenligning ikke kan komme til noget andet
Resultat, end det, overalt at see, ikke Sandheden selv, men kun Skinnet,
Alterationen af den. --- Det gaaer den, der ad denne Vei søger Sandheden, som
det gaaer den, der i et fremmed Land opsøger en Ven, træder ind i Herberget,
hører, at han for kort Tid siden var der, men netop i dette Øieblik er gaaet
ombord, og nu, naar han er kommen ned i Havnen, erfarer, at han alt er ude paa
Havet. Thi søger man Sandheden i Subjectiviteten, saa henviser Kriticismen til
Objectiviteten; ønsker man da at see Sandhedens objective, almindelige Idee, saa
analyserer den de enkelte Phænomener, og, bliver nu Talen atter om dette eller
hiint Phænomen, saa appellerer den til Totaliteten.
% Greek read at 6x: μετὰ παρατηρήσεως (Luke 17:20). Grave on the alpha of μετὰ,
% acute on the eta of παρατηρήσεως; medial σ, final ς. No fount-variant
% letterforms occur on this page --- the rho is the ordinary ρ.

Heraf følger da, at Kritikens Regler enten ere saa almindelige, at de ere
aldeles indholdsløse, eller at de, naar de bestemme Noget om de speciellere
Forhold, ere af den Natur, at man kan sige om dem, hvad der som oftest gjelder
om den tydske Grammatik i dens Bestemmelser om Substantivernes Kjøn:
”\textit{quot regulae, tot exceptiones!}”
% The quotation marks are the Fraktur fount's, the Latin inside them antiqua.
% (The centred antiqua "2" at the foot of this page is the sheet signature, not
% a folio; not transcribed.)
% --- p. 18 ---
\parag{4}{Overgang til den speculative Empirie.}
% The printed § head agrees WORD FOR WORD with the Indhold reading
% ("Overgang til den speculative Empirie.", p. 18) --- nothing to record as a
% divergence. "§ 4." is centred bold antiqua, the argument centred letterspaced
% Fraktur; read at 3x on PDF 27.

De ydre Phænomener maae altsaa ikke længere betragtes i og for sig; hvad der
bliver tilovers af dem, er ikkun de subjective Meninger om dem, og i disses Sum
maa da Sandheden sættes. --- Om jeg saaledes end ikke veed, hvorvidt jeg har
bedømt et eller andet ydre Phænomen rigtigt eller ei, veed jeg dog idetmindste,
at jeg har dømt paa denne eller hiin Maade; og er det end ikke sikkert, om en
eller anden Begivenhed har fundet Sted eller ikke, er dog saameget vist, at den
er bleven fortalt og derfor ogsaa maa være bleven tænkt af Nogen; og kan jeg end
tvivle om, at Mirakler have fundet Sted i Oldtiden, kan jeg dog ikke tvivle om,
at de Gamle have troet paa Mirakler. Det er saaledes ikke Historikerens Sag at
undersøge, hvorvidt de enkelte Ting have fundet Sted paa den bestemte udvortes
Maade, hvorpaa de berettes; men, som ingen Gjerning kan gjøre sig selv, men
enhver Gjerning forudsætter En, der gjør den, saaledes kan heller ingen
Tildragelse fastholdes i og for sig, men maa kunne tjene til at give os en sand
Forestilling om den Persons Standpunkt, hvem den tillægges. Herved blive altsaa
de ydre Phænomener aldeles negerede. --- Det aandelige Liv fremtræder i
sandselige Phænomener; Phænomenerne udtrykke Sjelens væsentlige Meninger, men de
væsentlige Meninger ere Tanker; det bør derfor være Historikerens Opgave at
paavise Aandens Skikkelse til de forskjellige Tider, Folkenes Opdragelse,
% DOUBTFUL READING: "Folkenes" is broken over the line as "Folke=" / "nes", and
% the three letters "nes" at the head of the line are badly under-inked on this
% copy. Read at 6x: three glyphs, n-e-s. Rejected alternative: "Folkets"
% (only two letters would follow the break, so it does not fit the space).
afspeilet i Menneskeslægtens Bedrifter som i et blankt Speil. --- Betragter man
Sagen fra denne Side, kommer man ind paa en, saa at sige, historisk Idealismes
Standpunkt, hvilken idetmindste har Ret deri, at Tænkningen og Fornuften søge
sig selv i Historien. Det er imidlertid saa%
% --- p. 19 ---
\setcounter{footnote}{0}
langt fra, at denne Idealisme frakjender den vulgaire Empirie al Auctoritet, at
den meget mere bekræfter samme, ikkun med den Forskjel, at Historiens tvende
Momenter, det Reale og det Ideale, ikke opfattes, som om deres Forening skulde
foregaae paa en umiddelbar Maade, idet Beretningen om de historiske Data
kryddres med tilfældige subjective Bemærkninger, og der fra de fromme Følelsers
dunkle Kilde udvælder en aandelig Strøm for at vande den jordiske Tilværelses
golde Marker. Medens Idealismen forkaster Fastholdelsen af de ydre Phænomener,
som saadanne, og griber efter det indre Væsen, erkjender den tillige, at det kun
er igjennem Phænomenerne, at Adgangen staaer aaben til den aandelige Tilværelse.
Det vilde jo ogsaa være umuligt at danne sig et Billede af de enkelte Menneskers
eller et heelt Folks Aand, dersom man kun kjendte de nøgne Begivenheder og ikke
tillige kjendte de handlende Personers Navne, eller dersom, for at nævne et
Exempel, Cæsars Bedrifter tillagdes Pompeius; umuligt at kjende og forstaae
Aandens Virksomhed, dersom man ikke havde nøie Kundskab om, hvilke ydre Skranker
den havde fjernet, hvilke raae Naturmasser den havde bearbeidet, hvilke
objective Sphærer den havde formet\footnote{\textit{Avec le monde a commencé une
guerre, qui doit finir avec le monde, et pas avant; celle de l'homme contre la
nature, de l'esprit contre la matière, de la liberté contre la fatalité.
L'histoire n'est pas autre chose, que le récit de cette indéterminable lutte.
Michelet: Introduction a l'histoire universelle.}}. Det er for Exempel af den
største Vigtighed for den rette Opfattelse af Jødefolkets Aand, om vi antage, at
de mange Mirakler, der tillægges Jehova, virkeligt fandt Sted i den ydre Natur,
eller om vi ikkun antage dem for at være Producter af fromme Individers
subjective Forestillinger.
% The ONLY footnote in pp. 13--24. Printed mark "1)" --- the plain next number.
% The whole note is French set in ANTIQUA against the Fraktur body, so the whole
% of it is \textit{}. Read off the image at 4x, both text layers being useless
% here. Neither "Michelet" nor any other word in the note is letterspaced: the
% apparent extra weight of "Michelet" at low magnification is inking, not a
% second fount. The unaccented "a" in "a commencé" and in "Introduction a
% l'histoire" is as printed. (The "2*" below the note is the sheet signature.)
% --- p. 20 ---
Den sande Historie fordrer altsaa vel, at man skal søge den objective Fornuft i
de ydre Phænomener, men tillige, at Phænomenernes Masse fuldkomment skal
forbrændes af Tankens Ild, for at opstaae som en Phoenix af sin Aske. Men det
Arbeide, som denne grundige Negation og derpaa følgende Restauration af
Phænomenerne medfører, forekommer de fleste Historikere vanskeligt og
betænkeligt. De undslaae sig vel ikke for at tænke; de ere meget mere villige
til at slaae deres Betragtningers Islæt imellem de berettede Facta's Rendegarn;
thi denne Maade at tænke paa synes at være den sikkreste; skulde det nemlig nu
ogsaa hænde sig, at deres subjective Betragtninger befandtes usikkre og
inconseqvente, have de dog stedse de berettede Facta tilbage; til dem kunne de
da i ethvert Tilfælde holde sig og sige: ”Factum selv er her, og det er
idetmindste utvivlsomt, at det har tildraget sig paa den Tid eller det Sted.”
Ønsker man derimod at forlade de empiriske Phænomener og at henvende
Opmærksomheden paa den objective Fornuft i dem, som saadan, vende de os Ryggen
af Frygt for, hvad de kalde, philosophiske forudfattede Meninger og vilkaarlig
apriorisk Construction.
% "Facta's", "Facta" and "Factum" are all FRAKTUR on this page, not antiqua ---
% checked at 4x --- so they take no \textit{}. The quotation is not letterspaced.

Dersom det nu virkeligt forholdt sig saa, at vi, efterat have negeret alle
Phænomenerne, vendte os til dette eller hiint System for vilkaarligt og paa
udvortes Maade at anvende visse forudfastsatte Formler paa dem, vilde det i
Virkeligheden være frugtesløst for Philosophen at anvende Tid og Flid paa
Historien. Men det forholder sig ikke saaledes med den philosophiske Negation.
--- Naar man i daglig Tale nævner Begrebet Intet eller Negation, tænker man
gjerne paa en, saa at sige, tom Nullitet; men det er let at indsee, at der ere
ligesaamange Negationer, som der ere Modificationer af de positive Ting, da jo
Negationen er Tingens egen Begrændsning eller Bestemmelse. Vilde for Exempel
Nogen blot sige:%
% --- p. 21 ---
\emph{”jeg har Intet,”} saa vilde man ikke ret forstaae ham; sagde han derimod:
\emph{”før havde jeg Rigdom, nu har jeg Intet!”}, saa vilde man let begribe,
baade hvad dette Intet stod i Modsætning til, og, af hvad Natur det selv var, og
forstaae, at Manden var bleven fattig. --- Det daglige Liv anbefaler ogsaa selv
denne Maade at slutte paa. Naar vi saaledes bedømme vore fortrolige Venners
Handlinger, blive vi jo ikke staaende ved deres Handlinger, som blot ydre
Phænomener, men see tillige stedse hen til den Sindsstemning, de Aarsager og
Grunde, hvoraf de fremgik. Enhver Handling har sin Grund, men Grunden er
Phænomenets Negation; Phænomenet angiver altsaa det Negative af sig selv. ---
Derfor er det ogsaa netop med den speculative Negations Sandhed, at alle
Beviserne for Guds Tilværelse staae og falde; thi ethvert af dem gaaer netop ud
paa at negere et almindeligt Phænomen, saaledes som Verden umiddelbart
fremstiller sig for os, og fører os igjennem dettes Negation til Guds Væsen,
hvoraf da det negerede Phænomen paany restaureres.

Betragte vi saaledes først det kosmologiske Beviis, møde vi i dette først
Iagttagelsen af \emph{et empirisk Phænomen:} \emph{”See, Verden er!”}; men
\emph{Negationen} træder strax til, idet man nemlig slutter fra Verdens
Tilfældighed: \emph{”Betragtes Verden i og for sig, saa er den Intet, har i sig
ingen Substants”}, og først nu følger \emph{Slutningen:} \emph{”Den absolute
Substants er Gud,”} hvorpaa \emph{Phænomenet atter restaureres:} \emph{”Fra Guds
Substants erholder Verden sin Substants.”}
% Letterspacing read at 3x on PDF 30. The spaced runs are exactly: the quoted
% "jeg har Intet," and "før havde jeg Rigdom, nu har jeg Intet!" at the head of
% the page; then, in the kosmologiske-Beviis paragraph, "et empirisk Phænomen:"
% (the two-letter "et" IS spaced --- confirmed at 6x, e and t stand apart),
% "See, Verden er!", "Negationen", the whole "Betragtes ... ingen Substants",
% "Slutningen:" (broken "Slutnin=" / "gen:", both halves spaced), "Den absolute
% Substants er Gud,", "Phænomenet atter restaureres:", and "Fra Guds Substants
% erholder Verden sin Substants." Everything between those runs is plain.

Heraf sees tillige, at Væsenet selv bestemmes efter det Phænomen, der er
negeret; thi medens vi i det kosmologiske Beviis negerede Verdens Substants og
igjennem denne Negation kom til Guds Substants, vilde vi, hvis vi havde negeret
et andet Phænomen, f. Ex. den menneskelige Tanke, hvil%
% --- p. 22 ---
ket skal foregaae i det ontologiske Beviis, have faaet en anden Form for Guds
Tilværelse. Ogsaa her kan man tydeligt see, hvor farligt det er at blive
staaende paa Middelveien, Noget, vi alt eengang have søgt at bevise.

\emph{Pointet} i det ontologiske Beviis kan angives saaledes: \emph{”Jeg tænker,
altsaa er Gud til.”} Det Phænomen, der umiddelbart er tilstede, er altsaa den
menneskelige Tanke og Selvbevidsthed. \emph{Den absolute Negation lyder} da
saaledes: \emph{”Den menneskelige Tanke er Intet af sig selv”}; gjennemføres
denne Negation, fremgaaer følgende Slutning: \emph{”Den absolute Tanke \dsi{}
Guds Selvbevidsthed er”}, og nu restaureres Phænomenet: \emph{”Er Guds Tanke og
den absolute guddommelige Selvbevidsthed, da kan ogsaa den menneskelige Tankes
Tilværelse deraf forklares.”}
% Letterspacing read at 3x on PDF 31. "Pointet" alone in its clause; then each
% quoted formula. "Den absolute Negation lyder" is spaced across the line break
% ("ly=" / "der") and stops there --- "da saaledes:" is plain, and so are
% "gaaer følgende Slutning:" and "og nu restaureres Phænomenet:", which are NOT
% spaced even though their counterparts on p. 21 are. The \dsi sign stands
% inside the spaced run, as printed.

Men da nu Kant tog sig for at forkaste det ontologiske Beviis, førte hans
Adskillelse mellem en reel og en ideel Tilværelse, hvorpaa hans hele
Beviisførelse beroede, ham netop ind paa Middelveiens slette Standpunkt. ”At
Mennesket er og kan tænke,” sagde han, ”er et Factum; fastholdes dette, er det
fuldkommen vist, at Mennesket formaaer, saa ofte det vil, ligesaavel at lade
være at tænke Gud, som at tænke ham.” Han fastholder altsaa dette empiriske
Phænomen, lader der ikke komme til nogen sand Negation og faaer saaledes let
Beviset til at vakle, efterdi den guddommelige Tanke ikke paa udvortes Maade kan
stilles ved Siden af den menneskelige Tanke. Havde han derimod opkastet sig det
Spørgsmaal, hvorfra ogsaa denne Menneskets Evne til at tænke var opstaaet, vilde
hans Idealisme have udsondret de raae, ufordøiede Phænomener af den men%
% --- p. 23 ---
neskelige Bevidsthed, og han vilde da maaskee være kommen dertil ret at
anerkjende og udvikle det ontologiske Beviis.

Ligesom den slette Middelvei har havt den skadeligste Indflydelse paa
Theologien, saaledes har den ogsaa havt det og vil fremdeles vedblive at have
det paa Historien, saa længe indtil man faaer gjennemført en sand Negation af
Phænomenerne. Thi det forholder sig med de historiske Ideer, ganske som med de
theologiske Ideer, vi nyligt omtalte. Man bør allerførst bestræbe sig efter
gjennem Negationen af de berettede Facta at trænge ind til Principet selv. Vil
f. Ex. Nogen begribe Katholicismens Aand, maa han først saavidt muligt betragte
Summen af de derhenhørende Phænomener, for gjennem deres Negation at komme til
det katholske Princip; thi gjør han ikke dette, da maa han uvilkaarligt gribe et
Princip af Luften, og da vil dette bestandigt støde an mod de historiske
Phænomener. Ligesom altsaa i det kosmologiske Beviis Verdens Phænomen først var
givet i sin Umiddelbarhed og Tilfældighed, saaledes maa ogsaa Katholicismens
almindelige Phænomen først iagttages i sin Tilfældighed, før man kan trænge ind
til dets Grund. Og har man saaledes tilegnet sig Principet, maa man atter vel
erindre, at dette ikke maa fastholdes i sin Abstraction, men at det ogsaa i sin
specielle Udvikling kun kan tilegnes paa rette Maade gjennem Negationen af de
specielle Phænomener, saa at Tanken aldrig forlader det historiske Object. ---
Der ere da tvende Vildfarelser, for hvilke man især maa vogte sig, deels at man
ikke af et saadant almindeligt Princip udvikler reent abstracte Begreber med
Forsømmelse af de specielle Phænomener, deels at man ikke paa udvortes Maade
overfører det almindelige Princip paa alle Enkelthederne. Der er nemlig Intet i
Historiens Studium, der mere røber den slette Middelvei, end det, at anvende de
almindelige Ka%
% --- p. 24 ---
tegorier: Frihed, Despotisme, Hierarchie eller andre lignende, som Indklædning
for de enkelte Begivenheder, saa at man ikke udvikler disse Begreber selv, men
kun udfylder dem med udvortes Enkeltheder. Skal f. Ex. Hierarchiets Begreb
historisk udvikles, kan man af en abstract Definition ikke vide, at det i
Aarhundredernes Løb forandres, og, spørges der om dets enkelte Forandringer, da
er det ikke nok at opregne Irenæus's, Cyprians og Gregors enkelte Handlinger
eller Yttringer. --- Der ere da tre Momenter at adskille i en saadan Udvikling:
det første er Hierarchiets almindelige Begreb; det andet de enkelte Handlinger
og Begivenheder, der ere Udtryk af den hierarchiske Idee; det tredie Identiteten
af begge Momenter. Og hvad der menes med denne Identitet, er ikke vanskeligt at
forstaae; thi spørges der, hvorfor der skal skee en Anvendelse af det Almindelige
paa det Enkelte, kan der, om vi ikke tage Feil, ikke svares Andet, end, at det
almindelige Begreb, Ideen, udvikler sig paa særegen Maade i de enkelte
Phænomener, at saaledes Hierarchiets Begreb fremtræder paa en anden Maade hos
Gregor den Syvende, end hos Irenæus. Men heraf følger da aabenbart, at Gregors
Foretagender kunne betragtes fra en dobbelt Side; fra den ene Side betragtede
ere de særegne, saa at vi ikke kunne gjenfinde de samme noget andet Sted i
Historien, og, som saadanne, udgjøre de det faste og urokkelige Moment, vi før
omtalte; men paa den anden Side aabenbare de paa speciel Maade Hierarchiets Idee
og ere saaledes overalt gjennemtrængte af det almindelige Moment. De individuelle
Momenter have saaledes modtaget det Almindeliges Præg, medens omvendt den
almindelige Idee har erholdt sine særegne Modificationer. --- Men heraf
fremgaaer da det tredie Moment, der hverken kan bestemmes som en enkelt
Erfaringsgjenstand ved et: ”saa eller saa!” og heller ikke%
% p. 24 is a single unbroken paragraph --- no indent anywhere on the page, and
% the names Irenæus's, Cyprians, Gregors are Fraktur, not antiqua (checked 4x).
%
% ============================================================
% QUOTATION MARKS IN THIS RANGE --- A CORRECTION TO THE HEADER.
% The header says the BODY uses the Danish low-high pair „…“ and that only the
% preface on PDF 7 sets a RAISED mark at both ends. That is NOT true of pp.
% 13--24. Every quotation on these pages --- pp. 17, 20, 21, 22, 24, eleven
% pairs in all --- sets the RAISED mark (U+201D) at BOTH ends, exactly as the
% preface does. Checked at 4--6x at each site: the opening mark sits at
% ascender height, never on the baseline. So this batch is transcribed with
% ”…” throughout and contributes „=0, “=0 to check.py's balance, which is
% therefore still 0 and still correct. Whether pp. 1--12 do the same is for
% that batch to settle; the claim in the header should not be relied on until
% it has been.
% ============================================================

% --- p. 25 ---
% p. 25 opens mid-paragraph: p. 24 ends "...og heller ikke" (complete word, no
% broken syllable at the joint --- verified on PDF 33 at the foot of the page).
% NOTE TO THE SPLICER: this first line must NOT be preceded by a blank line, or
% the paragraph carried over from p. 24 is broken in two.
kan siges at være et tomt almindeligt Begreb, anvendeligt paa
hvilkensomhelst Maade, men maa kaldes Begrebets historiske
Begrændsning. Og en saadan Empirie, der fuldkomment svarer til og er
overeensstemmende med den tilgrundliggende Fornuft, kaldes med Rette
speculativ.

% Printed head agrees with the Indhold reading (p. 25).
\parag{5}{Forholdet imellem den vulgaire og speculative Empirie.}

Forskjellen mellem den vulgaire og speculative Empirie er
iøinefaldende; hiin iagttager Phænomenerne udvortesfra, denne
indenfra, og, medens hiin sammendynger en Mængde tilfældige
Phænomener, skildrer denne derimod Momenternes organiske Forbindelse
paa en saadan Maade, at intet Phænomen kommer i Betragtning, uden
forsaavidt det erkjendes at være fremgaaet af Ideen selv. Derfor kan
Omfanget af de enkelte Phænomener være langt større i den vulgaire
Empirie, end i den speculative; thi i den vulgaire Empirie spørge vi
kun, om Begivenhederne have fundet Sted, ikke, \emph{hvorfor} de have
% "hvorfor" letterspaced: measured intra-letter gaps 22--28 px against 4--13 px
% for the rest of the line. The only emphasis on this page.
fundet Sted, og have derfor kun den Opgave nøie at iagttage og
bestemme saa mange enkelte Phænomener som muligt. Kritikernes Værker
frembyde mange Exempler paa det rige Udbytte af en saadan omhyggelig
% A short stray horizontal mark follows "Udbytte" on the page (a broken sort or
% an ink speck at mid-height, not a quotation mark: no quotation is open here,
% and the ABBYY layer reads it as an apostrophe, the Fraktur layer as a quote).
% Not transcribed.
og ufortrøden Undersøgelse. Philosopherne forfølge derimod stedse
Fornuftens constitutive Princip og oversee derfor, som tilfældige og
ubetydelige, mange Phænomener i Historien, der af Kritikerne behandles
med største Nøiagtighed. Imidlertid maa man ikke troe, at denne
Ligegyldighed er en nødvendig Følge af Sagens Natur, som om denne
maatte medføre en vis speculativ Hovmod; thi objectivt betragtet kan
Alt tilegnes og begribes, --- den er meget mere en Følge af den
menneskelige Tænknings Ufuld\-%
% --- p. 26 ---
kommenhed, i det denne vel kan adskille den objective Fornufts
almindelige Lineamenter fra hinanden, men ikke de utallige
Nuanceringer, i hvilke den fremtræder, saa at Philosophien netop
opfordrer til Beskedenhed; --- og det er saa langtfra, at Speculationen
skulde svække Interessen for de empiriske Phænomener, at den tvertimod
forøger den. Een enkelt Handling kan kun udtrykke eet Moment af den
almindelige Idee; skulle derfor de øvrige Momenter træde til, maae
ogsaa de øvrige Handlinger tages i Betragtning. For at udfinde det
objective og historiske Begreb, maa man altsaa have et Apparat af
mange Phænomener; jo flere Phænomener vi derfor optage af den vulgaire
Empirie, desto bedre vil det lykkes os at udsondre og udvikle de
væsentlige Phænomener.

Vil man derfor paa vulgair Maade beskrive Gregor den Syvendes
Hierarchie, ved nemlig at anføre hans enkelte Yttringer og
Foretagender, behøver man kun hist og her at indflette sine
Definitioner og forresten udpille og i en beqvem Orden fremstille
saadanne Steder, der ere skikkede til at anskueliggjøre det; (og naar
da et klækkeligt Antal af saadanne Steder, sammendyngede allestedsfra,
prange i Definitionernes Række, ville de tillige være et glimrende
Beviis paa den store Lærdom, hvoraf Forfatteren er i Besiddelse). Vil
man derimod begribe denne Paves Hierarchie i dets hele Sandhed, kan
man ikke lade sig nøie med at have en Sum af udvalgte Steder; man vil
da finde, at man ikke paa nogen anden Maade kan træffe Begrebets
finere Traade end ved omhyggeligt at oplyse alle de Grunde, der have
givet det dets ejendommelige Charakteer. Sandhedens Begreb er saaledes
i begge disse Arter af Empirie aldeles forskjelligt. I den vulgaire
Empirie er det sandt, som i Virkeligheden har fundet Sted, uden Hensyn
til, om det er noget Betydeligt eller noget Ubetydeligt; i den
speculative derimod afhænger Sandheden
% --- p. 27 ---
af det Forhold, hvori den bestemte Gjenstand staaer til det
constitutive Princip, og der kan være foregaaet Meget, hvori den
speculative Empirie kun finder liden Sandhed. --- Medens saaledes
Isidors Dekretaler absolut forkastes af Kritikeren paa Grund af deres
Usandhed, have de dog, dersom Phænomenet betragtes fra dets anden
Side, en dybere Sandhed; og vi have alt bemærket, at disse tvende
Betragtningsmaader gjensidigt maae veies mod hinanden. --- En lignende
Forskjel finder Sted, hvor Talen er om historiske Beviser; thi vel
kunne vi ofte dokumentere den omhandlede Sag ved bestemte Udsagn eller
Begivenheder, men vi paastaae ogsaa ofte, om vi end ikke have saadanne
at holde os til, at Sagens Sandhed fremgaaer af alle de enkelte
Gjenstandes almindelige Forhold. Det er netop herved, at Forskjellen
fremkommer mellem de saakaldte juridiske Beviser og dem, der kun
støtte sig paa en moralsk Overbeviisning; og som vi henføre hine til
den vulgaire Empirie, saaledes maae vi, naar vi ikke ville blive
staaende paa Sandsynlighedens vaklende Grund (\textit{cfr.~§.~3})
% The whole parenthesis "(cfr. §. 3)" is set in ANTIQUA against the Fraktur ---
% brackets, abbreviation, section sign and figure alike. Verified at 1500 px.
henføre disse til den speculative.

Hvis nu den vulgaire Empirie selv kunde løse den Opgave, den har for
sig, vilde Overgangen til den speculative Empirie være let og sikker;
men det er let at indsee, at den vulgaire Empiries Sandhed ikke kan
bestemmes efter den speculative Empiries Rettesnor alene; skal man
fuldkomment kunne begribe, hvorledes disse to hverken kunne være
umiddelbart overeensstemmende eller absolut uovereensstemmende, maa
det vise sig, at Modsætningen imellem dem har en dybere Modsætning til
sin Grund, nemlig en Modsætning imellem den vulgaire og speculative
Logik.

Den udvortes Opfattelsesmaade af Tilværelsens Gjenstande er deri
forskjellig fra den, der trænger ind til Gjenstandenes Indre, at hiin
sammenstiller Gjenstandene paa en
% --- p. 28 ---
saadan Maade, at de vel i enkelte Punkter berøre hinanden, men
indenfor de Grændser, der omgive dem, forblive adskilte, medens denne
overalt lader dem falde sammen og gjennemtrænge hinanden. Heraf følger
da, at den vulgaire Logik gaaer ud fra Phænomenernes indbyrdes
Forskjellighed og kun tillægger deres Eenhed en formel Betydning,
medens den speculative fører os fra Eenheden ud i Forskjellighederne.
I den vulgaire Logik, hvor de abstracte Begreber fastholdes for sig,
de reelle Gjenstande for sig, er der Intet, man mere maa vogte sig for
end Modsigelse; thi naar Tingene gjensidigt berøre hinanden og det dog
er Tingens Natur at holde alt Fremmed udenfor sig, hvorledes kan da
Tingen \textit{A} paa eengang være det Samme og dog ikke det Samme som
Tingen \textit{B}. Naar derimod Betragtningen føres indenfra, sees
% The algebraic letters A and B are ANTIQUA against the Fraktur.
Tingene ikke i en abstract Adskillelse fra hinanden; Momenterne gaae
uendeligt over i, svare uendeligt til hinanden, og saaledes ophæver
den absolute Modsigelse sig selv.

Det bliver nu let at forklare sig de mange Stridigheder imellem
Hverdagsforstanden og den speculative Fornuft. Thi Hverdagsforstanden
oprøres ved at see Regler, som den, at \emph{der ikke kan blive Noget
til af Intet}, eller at \emph{Tre ikke kunne være Eet og Eet ikke
Tre}, Regler, den selv har erkjendt for sikkre og urokkelige, traadte
% TWO letterspaced runs, and the words between them are NOT spaced. Measured
% per word (intra-letter gaps in px): "der ikke kan blive Noget til af Intet"
% 19--30; "eller"=12,12,12 and "at"=10 --- ordinary; "Tre ikke kunne være Eet og
% Eet ikke Tre" 18--30; "Regler, den selv" 5--13 --- ordinary.
under Fødder, om man saa tør sige, af den speculative Logik. Det gaaer
imidlertid med disse to Arter af Logik, ligesom med de to Arter af
Empirie, at de ikke ligefrem ophæve eller kuldkaste hinanden; alle
disse Modsætninger opstaae deraf, at man har et forskjelligt
Udgangspunkt for Betragtningen, saa at der finder en mærkelig
\textit{concordia discors} Sted imellem Empirien og Philosophien. Naar
nemlig Philosopherne omtale den absolute Modsigelse og lære den deri
liggende Identitet, gjøre de det ikke i Hverdagsforstandens Navn; de
indrømme gjerne, at
% --- p. 29 ---
alle Hverdagsforstandens Regler maae antages af den, der staaer paa
Hverdagsforstandens Standpunkt; de erkjende fuldtvel, at Middelveien
fortjener Anbefaling, hvis man vil holde sig til de almindelige
Meninger om Philosophien og Historien; men de ville have Standpunktet
selv forandret; de ville ikke forandre den vulgaire Logiks Love, men
den vulgaire Logik selv. Empirikerne ville paa den anden Side ikke
forandre deres Standpunkt, fordi de ikke kunne lade sig overbevise om,
at deres Methode virkeligt er skjæv og eensidig. --- Det er da en
Selvfølge, at den vulgaire og speculative Empirie ikke paa udvortes
Maade kunne stilles ved Siden af hinanden for ligesom ved et
subjectivt Kunststykke at forsones med hinanden, da den Vaklen imellem
begge, der derved vilde opstaae, vilde fremkalde en, begge
undergravende, Skepticismus, men at Overgangen fra den ene til den
anden paa utallige Maader skeer og maa skee ved Tingens egen objective
Magt. Thi det Ydre maa stedse have det Indre i sig, --- og det beroer
ikke paa et subjectivt Indfald, om den speculative Logik skal have
Fortrinet for den vulgaire, eller ikke; det ligger meget mere i den
vulgaire Logiks egen Natur, at den, naar den blot gjennemføres og ikke
bliver staaende paa Halvveien, maa føre til den speculative Logik. Det
er netop dette Forhold, der gjør, at den speculative Empirie erkjender
en væsentlig Identitet imellem sig og den vulgaire Empirie. --- Heraf
er det ogsaa indlysende, med hvilken Frihed det speculative Begreb kan
uddanne sig. Vel afskrækkes Mange, naar man anbefaler dem et Begreb
eller et System, fordi de forkaste enhver med indre Nødvendighed
strengt afmaalt Begrebsbygning, som en utaalelig og generende
Jernharnisk, hvori deres Tanker indpines; men denne Frygt er ikke
begrundet i Sagens sande Natur, men er et Foster af deres subjective
Indbildningskraft.
% --- p. 30 ---
Thi vel fordrer Systemet, at de nødvendige Begreber skulle fastholdes,
men Aarsagen hertil er ingen anden end den, at disse i Livet selv
overalt gjøre deres Magt gjældende. Det fordrer ligesaafuldt, at
Begrebet skal oplyses og paavises gjennem alle dets særegne
Modificationer, og bevarer saaledes i fuldeste Forstand Tanken i dens
Frihed. Baade Systemet og Begrebet ere i Sandhed i Besiddelse af den
frieste Bevægelighed; thi hvad vil vel Systemet Andet end udvikle alle
Begrebets Momenter, forsaavidt de gjensidigt optage hinanden i sig; og
hvad Andet gaaer vel Begrebet ud paa end netop Udviklingen af den
prægnante Idees Fylde? Af Ideens indre Fylde fremgaae Gjenstandens
ydre Modificationer, og, betragtes disse udvortesfra, ere de vel
utallige og tilfældige, men betragtes de indenfra, sees de fremgaaede
af Ideens egen i sig selv begrændsede Magt og ere saaledes begrændsede
midt i deres Ubegrændsethed. Begrebets Nødvendighed er saaledes ikke i
nogen Modsætning til Livets Frihed.

Heller ikke oversees det subjective Frihedsmoment, fordi vi ikke ville
overlade den sande Udvikling af Livets Momenter til den subjective
Vilkaarlighed. Thi da den speculative Erfaring paa eengang uendeligt
svarer til og tillige er uendeligt afvigende fra Hverdagserfaringen,
udfordres der jo netop et genialt Blik, en Evne til umiddelbart at
søge Ideen, hos den, der rigtigt skal kunne løse de indviklede
Momenters Knude.

% Printed head agrees with the Indhold reading (p. 30).
\parag{6}{Videnskabens Begreb.}

Hvad vi hidtil have søgt at vise, er med faae Ord det, at der ere
trende Stadier, som vor Erkjendelse har at gjennemløbe. Det første
byder os at forlade Hverdagserfaringen med dens Umiddelbarhed og
Tilfældighed og at gaae over til
% --- p. 31 ---
den speculative Logik; det andet negerer Logiken og fører os ind paa
den vide Mark, som den omhyggelige Prøvelse af Hverdagserfaringen
aabner for os; det tredie fremstiller den objective Fornuft
anskueliggjort i den speculative Erfarings Speil.

Vilde man nu paa udvortes Maade stille disse tre, hinanden negerende
Stadier ved Siden af hinanden, da behøvede man kun paa tre Steder og
tre Gange at søge Negationen og Restaurationen. Men hvad der saaledes
i Almindelighed er henført til tre Stadier, gjentager sig atter
uendeligt indenfor ethvert af disse; hvad der nemlig gjælder om den
speculative Logik, at den overhovedet seer bort fra Erfaringen, det
Samme gjenfinde vi atter indenfor denne Logiks Grændser, thi den
udvikler sig netop gjennem en uafbrudt Negation af empiriske
Momenter; og hvad der er sagt om Hverdagserfaringen, at den staaer i
Modsætning til den logiske Fornuft, det Samme gjelder ogsaa indenfor
% "gjelder" here, but "gjælder" eleven lines above and twice on p. 35 --- the
% compositor's own variation, transcribed as printed at each site.
denne Erfarings Grændser, dens Momenter gaae aldrig over i hinanden og
i dem alle see vi Fornuftens Alteration, ikke Fornuften selv. Det
almindelige Resultat, hvortil vor Erkjendelse igjennem alle disse
Stadier kommer, er da det, at see Modsætningerne i deres uendelige
Overgang i hinanden.

Men heraf fremgaaer Videnskabens Begreb, thi det er netop
Videnskabens Opgave at føre til en absolut Identitet af de hinanden
absolut modsatte Momenter. Vil man derfor have en almindelig og formel
Typus for den absolute Videnskab, kunne vi finde en saadan i vor egen
Selvbevidsthed; thi da det menneskelige Jeg bevarer Identiteten med
sig selv, idet det skiller sig fra sig selv, kunne vi med Rette sige,
at det veed sig selv. Da nu ingen Selvbevidsthed kan være til, uden at
medieres igjennem den reelle Tilværelse, \dsi{} gjennem det Negative af
Bevidstheden, saa er vel denne Begrændsning, hvis vi
% --- p. 32 ---
see hen til de specielle Gjenstande, en slet og ret Negation af
Bevidstheden og altsaa en Ophævelse af al Viden; men see vi derimod
hen til den reale Negation i Almindelighed, kunne vi med Sandhed sige,
at vor Viden udvides, jo mere Bevidstheden begrændses.

Thi kan Selvbevidstheden ikke være uden igjennem et Andet, der er den
modsat, maa den ogsaa paa ethvert Punkt være gjennemtrængt af dette
fremmede Moment, og, saa ofte den griber sig selv, tillige gribe
dette; der er derfor i Selvbevidstheden paa eengang en Adskillelse
imellem dens eget, det formelle Moment, og det fremmede, det reelle
Moment, og tillige en Identitet af begge; den veed derfor sit Andet
paa samme Tid, den veed sig selv. Og heraf forstaaes det ogsaa let i
Almindelighed, at man paa eengang kan vide Gud og Naturen uden
nogensomhelst Sammenblanden af Skabningen og Skaberen. De, der derfor
fordreie Philosophernes bekjendte Sætning, at kun det Ligeartede kan
fatte og erkjende det Ligeartede, (\textit{simile modo simile capere
et cognoscere}), glemme desværre, at Selvbevidsthedens absolute
% The Latin tag is ANTIQUA, not letterspaced Fraktur: spacing.py flags it as a
% RUN only because the antiqua fount is wider than the Fraktur around it.
Identitet med det Object, der skal erkjendes, ene og alene resulterer
af den absolute Modsætning imellem dem.

Det følger derfor af Videnskabens egen Natur, at enhver Art af
Erkjendelse, der stiller Tænkningen i et udvortes Forhold til de
Gjenstande, der skulle erkjendes, maa bortvises fra dens Helligdom.
Hvad der paa nogen Maade henhører til Hverdagserfaringen og
Hverdagslogiken, er for umiddelbart og ufordøiet til at optages i
Videnskabens Rige.

Men da Hverdagserfaringen, hvad vi alt forhen have bemærket, ikke kan
undværes og dog staaer udenfor Videnskabens Grændser, maa man vel
skjelne imellem de Elementer, der forberede Videnskaben og
Videnskaben selv; og vigtige ere de Følger, der flyde af denne
Adskillelse.
% --- p. 33 ---
Vi geraade da for det Første ikke oftere paa den Afvei, under Navn af
Videnskab at falbyde Rudimenta blandede med trivielle Begreber, i den
% "Rudimenta" is set in FRAKTUR here, not antiqua --- so no \textit{}.
% Verified at 1500 px; contrast "Deus" nine lines below, which is antiqua.
Tanke, at vi havde tilegnet os Videnskaben, en Afvei, der vilde føre
os længere og længere bort fra al sand Videnskab. At foredrage
Historiens Videnskab er saaledes ikke at fortælle Historier; kun det,
vi have begrebet og gjennemskuet i Historien, bør vi gjengive; alt
Øvrigt maae vi lægge tilside som Apparater, vi paany skulle undersøge
for at komme paa det Rene med os selv, om de kunne indtage en Plads i
Videnskaben, eller ikke. Men saavidt har den Uskik grebet om sig at
fremsætte kritiske Elementer under Navn af positiv Videnskab og at
svare den, der søger et alle Enkeltheder gjennemtrængende Begreb:
”det er noget Historisk; her skal ikke philosopheres!”, at
% QUOTATION MARKS. The two quotations on this page, and the one on p. 34, are
% set with a RAISED mark (U+201D) at BOTH ends --- exactly as the preface on
% PDF 7 sets its four quotations, and NOT with the low-high pair „…“. The same
% form stands at the foot of p. 24. Reproduced as printed; pp. 25--36 therefore
% contribute nothing to check.py's „/“ balance.
Dogmatikere, der saagodt som ikke lære os Andet om Gud, end at der er
et høieste Væsen til, der paa Hebraisk hedder \emph{Elohim}, paa
% *** HEBREW IN THE PRINTING --- SUBSTITUTION, NOT A TRANSCRIPTION. ***
% The page sets the pointed Hebrew word itself: aleph-lamed-he-yod-final mem
% with hataf-segol, holam and hiriq (U+05D0 U+05B1 U+05DC U+05B9 U+05D4 U+05D9
% U+05B4 U+05DD), read at 4x. This edition is built with pdfLaTeX + Libertinus
% and has NO Hebrew font, so the word is rendered by its name, following the
% precedent set at texts/treschow/menneskelige-natur/transcription.tex p. 115.
% If cjhebrew is installed the glyphs should be restored: add
%   \usepackage{cjhebrew}
% to the preamble and set "\cjRL{'/E lo hi:m}" in place of "\emph{Elohim}".
Græsk \textit{Θεος}, paa Latin \textit{Deus} og paa Tydsk Gott
% Θεος is printed WITHOUT its accent (standard Θεός). As printed.
% "Gott" is FRAKTUR, "Deus" antiqua --- the fount follows the language.
o.~s.~f., og at disse Benævnelser ere blevne forklarede paa forskjellig
Maade, vinde et stort Navn som Dyrkere af den ”positive Videnskab”. Ja
naar den speculative Dogmatik udelukker Alt, hvad den religieuse
Selvbevidsthed ikke kan begribe, som noget Fremmed, fremføre de endog
i den hellige Skrifts og Kirkens Navn heftige Klager imod den, som om
speculative Begreber ikke krævede vedholdende Studium af den hellige
Skrifts og Kirkens Lære. Og naar de see Dogmerne i uafbrudt
Tankefølge fremgaae af hinanden, beskylde de Dogmatiken for at være
udledet af et eller andet philosophisk System, medens de selv mene
bedst at sikkre sig Overgangen fra den ene Deel af Dogmatiken til den
anden ved at anføre Parallelsteder af Skriften og herved røbe, at de
ikke ere frie for en vis Overtro, da man lige saalidt holder dogmatisk
Vildfarelse borte fra sig ved at stille
% --- p. 34 ---
Skriftsteder i Spidsen, som man fordriver Djævelen ved at slaae Kors
for sig.

Det bliver dernæst ved hiin Adskillelse indlysende, hvor megen
Anstrængelse og Flid der udfordres for at komme til den sande
Videnskab. Naar Empirikerne og Kritikerne lykønske hinanden, som de
Eneste, der vove at give sig i Kast med Videnskabernes store Masser,
medens Philosopherne ene og alene give sig af med deres egne tomme
Speculationer, vise de i Sandhed, at de have altfor store Tanker om
sig selv og altfor ringe Forestillinger om Tænkningens Virksomhed. Men
de Samme, der fælde saa haard en Dom over den speculative Tænkning,
sætte snart Videnskaben selv altfor høit, snart altfor lavt. De have
altfor store Forestillinger om den, naar de paastaae, at Menneskene
ved deres Tænkning ikke kunne tilfulde og med absolut Sandhed begribe
Noget, og derfor sætte det hele Udbytte af deres Studium i Forøgelsen
og Anordningen af Hverdagserfaringens Stof; altfor ringe Forestillinger
om den, naar de mene med ringe Uleilighed og paa egen Haand selv at
kunne medbringe den fornødne Tænkning, naar blot de positive
Gjenstande ere tilstede. Havde de ret betænkt, hvormegen Flid der
hører til paa rette Maade at forfølge et absolut Begreb gjennem alle
dets Modificationer, vilde de vistnok tilstaae, at Jehovas Ord til
Adam: ”I dit Ansigts Sved skal Du æde dit Brød!” ogsaa i denne
Henseende kan finde Anvendelse paa Philosophen.

Skjøndt det saaledes ikke er Kritikerne og Empirikerne alene, hvem
Æren for Anstrængelse og Flid tilkommer, skydes dog i en anden
Henseende det Besværlige ved Arbeidet over paa dem; thi man maa vel
skjelne mellem den Tanke, der tænkes for dens egen Skyld, og den
forberedende Undersøgelse, der har sit Maal udenfor sig. Hiin tilhører
Speculationen,
% --- p. 35 ---
denne Hverdagsforstanden og Kritiken. Thi det ligger i de kritiske
Arbeiders Natur aldrig at være tilfredsstillende i og for sig; og
dette gjælder ikke blot om de Undersøgelser, der gaae ud paa at
efterforske et Skrifts Authentie, uden Hensyn til, om det er et
helligt Skrift, eller ei, men det gjælder ogsaa om den Maade, hvorpaa
de behandle Ideerne selv. Naar Kriticismen udvikler dogmatiske
Begreber, f.~Ex. Skabelsen, Inspirationen \textit{ell.~a.~l.}, naaer
% "f. Ex." is FRAKTUR on this page; "ell. a. l." is ANTIQUA. Both checked.
den vel stundom saavidt; at den faaer Problemet rigtigt opstillet, men
% PRINTER'S DEFECT, transcribed as printed: "saavidt;" with a SEMICOLON where
% the sense requires a comma ("naaer den vel stundom saavidt, at den faaer
% Problemet rigtigt opstillet"). Read at 3.5x: the mark is unmistakably a dot
% over a comma. The Fraktur OCR reads ";" here, the ABBYY layer "-".
Problemet selv løser den aldrig.

Verdens Skabelse, hedder det, maa hverken tænkes saaledes, at vi
sammenblande Verden med Gud, som Pantheisterne gjøre, eller saaledes,
at Verden antages at have sat sig selv, som Atheisterne lære; thi
saadanne Tanker ere ugudelige og fordærvelige i deres Følger; den maa
tænkes som en fri Handling af Gud. Men denne Frihed driste de sig ikke
selv til at tænke; de fordre altsaa Noget, som de ikke selv gjøre
(\textit{λέγουσι γὰρ καὶ οὐ ποιοῦσι}).
% GREEK, read at 3x. Fount variants normalised per the standing convention:
% cursive kappa -> κ (καὶ), cursive rho -> ρ (γὰρ).
% ACCENT PLACEMENT: this Greek fount consistently sets the breathing or
% circumflex of an ou diphthong over the OMICRON, not the upsilon --- so the
% page reads (in effect) "o-with-psili + upsilon" for οὐ and
% "omicron-with-perispomeni + upsilon" for ποιοῦσι, and the same again in the
% Plutarch quotation on p. 36. Omicron
% cannot carry a perispomeni in Unicode at all, and the displacement is a habit
% of the fount rather than a distinct reading, so standard placement is set
% (οὐ, ποιοῦσι) and the printed placement recorded here.

Paa samme Maade behandles Inspirationen: den maa hverken tænkes
saaledes, at Gud har overskredet Lovene for den menneskelige Aands
Natur, eller saaledes, at Menneskets Aand slet ikke antages at være
berørt af Guds Aand, men saaledes at Guds Aand har oplyst Menneskets
Aand. De fordre altsaa med tydelige Ord, at Inspirationen skal tænkes
paa denne Maade, og dog opfordre vi dem forgjeves til at tænke den. Vi
undre os imidlertid ikke saa meget over, at de ikke selv ville besvare
Spørgsmaalet, --- de kunne jo svare, at de ikke have antaget det for
deres Opgave at løse Problemet, men kun at opstille det; men idet de
kritiske Undersøgelser have indtaget den sande Videnskabs Plads, viser
deres høieste Modsigelse sig deri, at de, uagtet de opfordre os til at
tænke, dog selv erklære, at disse Spørgsmaal om, hvorledes
% --- p. 36 ---
Gud skabte Verden eller hvorledes han inspirerede Propheterne og
Apostlene, ere unyttige og umulige at besvare.

Hvad vi nu her have søgt at udvikle, maa ikke fremstilles saaledes,
som om vi havde meent, at Empirikernes og Kritikernes Undersøgelser
vare at foragte eller dog at betragte med Ligegyldighed; thi af det
Forhold, der efter vor foregaaende Udvikling bestaaer imellem den
vulgaire og speculative Erfaring, følger netop, at de ere aldeles
nødvendige; kun maae de ikke forvexles med den egentlige Videnskab.
% NOT emphasised, despite appearances. The ABBYY layer spaces out "de ere
% aldeles nødvendige" and the line does look opened up, but that is justification:
% measured word gaps 96--132 px with intra-letter gaps of only 4--12 px, against
% 19--30 px in the genuinely letterspaced runs on p. 28. No \emph{} here.
Vilde Philosophien derfor forekaste Kritiken det Ubetydelige i dens
Arbeider og prale af selv at være i fuldkommen Besiddelse af
Sandheden, kunde denne med største Ret svare den med de samme Ord, som
Themistokles hos Plutarch lader Festdagen svare den efterfølgende Dag:
\textit{Αληθῆ λέγεις. ἀλλ' ἐμοῦ μὴ γενομένης σύ οὐκ ἀν ἦσθα.}
% GREEK, read at 2.2--2.8x. Fount variants normalised: script theta -> θ
% (Αληθῆ, ἦσθα), cursive kappa -> κ (οὐκ). Accents ARE reproduced as printed,
% and the print differs from the received text at three points, each verified:
%   Αληθῆ  -- capital alpha carries NO breathing (received: Ἀληθῆ)
%   σύ     -- acute, not grave (received: σὺ)
%   ἀν     -- smooth breathing only, no grave (received: ἂν)
% "ἐμοῦ" is normalised: the fount sets the circumflex over the omicron of the
% ou diphthong, as with ποιοῦσι on p. 35; see the note there.
% "Themistokles" and "Plutarch" are FRAKTUR, so no \textit{}.

% *** THE PRINTED HEAD AND THE INDHOLD DISAGREE HERE. ***
% The printed § head on p. 36 reads "Den videnskabelige Methode."
% The Indhold (PDF 8) reads  "Den speculative Methode. ... 36".
% Both readings verified on the image; the printed head is set, per the
% standing rule that the page governs, and the Indhold reading is recorded
% here rather than normalised away.
\parag{7}{Den videnskabelige Methode.}

I sin Bestræbelse efter at føre til en absolut Identitet af Tanken og
det tænkte Object har Videnskaben stedse det Maal for Øie, at udvikle
denne gjennem alle Objectets særegne Modificationer. Kan derfor
Selvbevidstheden ikke være uden igjennem et Andet, og det endog et
reelt Andet, bliver det Opgaven at tilveiebringe en fuldkommen
Forsoning af dette Positive, der træder frem i den hele Mængde af
positive Elementer, og Selvbevidstheden, saa at hiint paa ethvert
Punkt kan svare til denne. Enhver Videnskab afhænger af denne
Forsoning, men Forsoningen afhænger atter af den
\emph{videnskabelige Methode}.
% The letterspacing covers only "videnskabelige Methode." --- the whole of the
% last line of the page. Measured: intra-letter gaps 19--37 px there, but 9--16
% px in "afhænger atter af den" on the line above, so "den" is NOT spaced.

% --- p. 37 ---
Den Modsigelse, hvorfra den sande Methode nødvendigt maa gaae ud, viser sig
deri, at, medens den formelle Viden og Tilværelsens reelle Fylde synes
indbyrdes at udelukke hinanden, stemme de dog, som vi forhen have viist, deri
overeens, at de begge \emph{ere}. Og det er i dette Identitetspunkt, at den
% p. 37 opens a NEW paragraph, indented on the image; p. 36 ended with a full
% sentence ("... af den \emph{videnskabelige Methode}."), so no word is broken
% across the joint. Verified on the image, not taken from the brief.
% "ere" IS letterspaced --- measured against the neighbouring words on the same
% line, which are tight; spacing.py misses it because the word is short.
hele Modsætning og enhver Uovereensstemmelse imellem dem skal finde sin
Ophævelse. Vi vide altsaa, at det er den objective Væren, der constituerer den
subjective, at den formelle Væren fremgaaer af den reelle, vor Tænknings Væren
af Substantsens Væren.

Men have vi nu ophævet denne Modsætniug imellem det formelle og reelle Moment,
% PRINTER'S DEFECT, transcribed as printed: "Modsætniug" with a TURNED n, read
% at 4x. The same page sets "Modsætning" correctly eight lines above, so this
% is the compositor and not a spelling. Both OCR witnesses read "Modsætniug".
viser der sig atter en ny Modsigelse. Thi den fuldkomne Overeensstemmelse med
den objective Realitet opnaaer Selvbevidstheden kun ved en Accomodation, i det
den gjør Afkald paa ethvert andet Prædicat, end det simple, at den \emph{er},
% "Afkald", not "Affald": read at 5x, the third letter is the Fraktur k. This
% is the scan's commonest trap (k read as f), and tesseract falls into it here.
uagtet den dog har mange flere Momenter i sig end dette; og see vi hen til det
Objective, da finde vi, at ogsaa den reelle Tilværelse foretager sig samme
\textit{συγχατάβασις}: skjøndt den strømmer over i uendelig Fylde, lægger ogsaa
% GREEK, read at 8x. *** AS PRINTED: the fourth letter is CHI, not kappa. ***
% The received word is συγκατάβασις. The glyph here is two crossing strokes
% whose left arm descends well below the baseline --- a chi; this fount's kappa
% (cf. οὐκ, p. 36) has no descender and does not cross. So this is a
% compositor's error, transcribed as printed per the editorial stance, NOT a
% fount variant to be normalised. The acute over the second alpha is as printed.
den alle sine Prædicater tilside undtagen det ene, at den \emph{er}. Derfor
indtræder, saa ofte vi ville acqviescere i Begrebet Væren, den Modsigelse, at
% "acqviescere" with QV here; p. 41 sets the same word "acquiescere" with QU.
% Both read at 5x. The two spellings stand as printed.
vi om et Object, der har et langt rigere Indhold end den blotte Væren, ikke
udsige Andet, end dette Ene, at det \emph{er}. Og vi kunne ikke undgaae denne
Modsigelse, fordi Selvbevidstheden, der er det Medium, igjennem hvilket vi
komme til dette Udsagn, stedse vil gjøre sin Ret gjældende. Hvad der nu
saaledes gjælder om Begyndelsen, det Samme gjælder ogsaa om ethvert
mellemliggende Trin og om Slutningen selv. Enhver concret Bestemmelse eller
Betegnelse, man er nødt til at tillægge Selvbevidstheden, maa ogsaa
nødvendigviis udsiges om den reelle Tilværelse. Er det derfor ikke
tilstrækkeligt at betegne Forholdet
% --- p. 38 ---
imellem Jeg og Mig, for saaledes at benævne den formelle Videns Subject og
Object, ved Forholdet mellem Væsenet og Phænomenet, Aarsagen og Virkningen,
% "Phænomenet" has ordinary n's --- checked at 5x. tesseract's "Phæuomeuct" is
% noise, not a turned-n defect.
kunne vi ogsaa være visse paa, at denne Bestemmelse ikke engang tilfulde
udtrykker det reelle Objects Natur. Heri ligger netop den dialektiske Stimulus,
som ikke lader Tænkeren finde Hvile, før han er kommen til Selvbevidsthedens
udviklede Begreb. Men har man endelig naaet dette Udviklingens høieste Punkt,
vil det ogsaa vise sig, at der ikke blot er gjort et subjectivt Fremskridt, men
at den selvsamme Bevægelse ogsaa er foregaaet i den reale Væren. Som vi derfor
i Begyndelsen af Udviklingen sige, at den objective og reale Væren constituerer
den subjective og formale Væren, saaledes forstaae og begribe vi nu, naar vi
komme til Systemets Slutning, at det er en objectiv og reel Selvbevidsthed, der
constituerer den subjective og formelle Selvbevidsthed.

Og dog ophører Modsigelsen endnu ikke her; thi den Identitet imellem den
formelle og reelle Selvbevidsthed, vi saaledes have godtgjort, indeholder
tillige en væsentlig Forskjel. --- Medens den absolut reelle Selvbevidsthed har
hele Realitetens Fylde i sig, viser det sig jo, at den formelle
Selvbevidstheds Identitet med den kun fandtes i det reent Formelle. Den
Realitet, vi før for Identitetens Skyld opgav, maae vi nu paa ny restaurere.
Men saaledes driver den dialektiske Stimulus os igjennem hele den speculative
Theologie.

At saaledes Gud tænkes som den absolute Frihed, medfører ingenlunde, at man maa
opgive den dialektiske Nødvendighed; thi hvor der er Frihed, \emph{maae} der
ogsaa være Yttringer af Friheden. Men paa den anden Side ophæves Guds absolute
Frihed ligesaalidt ved denne Begrebets Nødvendighed; der er nemlig følgende
Cirkel i Beviisførelsen: giennem Negationen af Verden kom vi til den Frihed,
% "giennem" with an I, not "gjennem": read at 6x --- the letter after the g is
% a dotted i sitting on the baseline, where this fount's j carries a dot AND
% descends (cf. "gjør", p. 37, and "gjensidige", pp. 42 and 43, all measured).
% Both OCR witnesses also read "gien". The book's usual form is "gjennem"; the
% only other "gi-" spelling in this batch is "giensidige" on p. 47. As printed.
der maa
% --- p. 39 ---
have givet Verden sin Oprindelse; vi maae som en Følge deraf, hvis vi ville
udtømme denne Friheds Begreb, atter fra Friheden vende tilbage til den Verden,
den har skabt; saalænge vi derfor beskjæftige os med Udviklingen af den absolute
Frihed, betragtet i og for sig, ville vi bestandigt føle denne dialektiske
Nødvendighed til videre Fremadskriden. Den samme dialektiske Nødvendighed, der
gjør, at \emph{Grunden} for at være Grund maa have en \emph{Følge}, den samme
% "Grunden" and "Følge" letterspaced, "Grund" between them NOT --- read at 4x.
gjør altsaa ogsaa, at Friheden for at være absolut maa have en Verden. --- Og
see vi hen til de specielle Momenter i Historien, gjenfinde vi ogsaa i dem det
selvsamme Forhold. Vi kunne f.~Ex. ikke begribe alle de enkelte Data, der
% "f. Ex.", "Data" and "Facta" are all FRAKTUR on this page --- checked at 5x,
% so no \textit{}. Only the section signs and the figures are antiqua.
skildre Forholdet mellem Tydskland og Italien, imellem den politiske og
kirkelige Magt paa Pornokratiets Tid, uden ved gjennem Totaliteten
(§~\textit{4}) af de berettede Facta paa en saadan Maade at trænge ind til
Principet i dets simpleste og almindeligste Form, at den selvsamme Modsigelse,
vi ovenfor omtalte, ogsaa her vil vende tilbage, dersom vi nu ville blive
staaende ved det første, abstracte Udgangspunkt (smlgn.\ §.~\textit{10}).
% The two section references on this page are set DIFFERENTLY: "(§ 4)" without
% a point after the sign, "(smlgn. §. 10)" with one. Both read at 5x and both
% reproduced. Cf.\ the § 2 head on p. 5, printed "§. 2." where §§ 1 and 3 are
% printed "§ 1." and "§ 3.".
Saaledes er ogsaa selv Menneskets Vilkaarlighed, som saa ofte antages at mangle
enhver Fornuftnødvendighed og at være aldeles tilfældig, bunden til objective
Love uden dog derved at ophøre at være Vilkaarlighed; og var det ikke saa,
kunde vi hverken antage de Indvendinger for sande, der gjøres mod
Inditerministerne, eller rigtigt forstaae Forhærdelsens Natur og Væsen. Og
Enhver, der har lagt Mærke til, at Menneskets Modstand imod Sandhedens Object
ogsaa selv bestemmes og modificeres ved dette Objects Natur, vil indrømme, at
Messalinas og Aggrippinas Grusomhed var af en anden Beskaffenhed end
Fredegundes og Brunahildis, og at atter Theodoras og Marjucc'ias Grusomhed
% "Marjucc'ias" --- the apostrophe is printed; read at 5x.
havde en særegen Charakteer. Saaledes var ogsaa den Forhærdelse, Katholi\-%
% --- p. 40 ---
kerne udviste, da de forkjettrede Luther, af en anden Beskaffenhed end den,
Jøderne udviste, da de forfulgte Christus, fordi Sandheden her viser sig i
forandret Skikkelse.

Som vi i den sandselige Verden med begge Øine betragte det samme Punkt, og
begge derfor saa at sige kun udgjøre eet Øie, saaledes tillægges paa den anden
Side vel Sjelen kun eet Øie, men dette betragter i Virkeligheden to Punkter
% NOT emphasised, despite spacing.py's flag on "betragter" (ratio 2.59): the
% line is set TIGHT, not open --- "betragter i Virkeligheden to Punkter" has no
% word space to speak of, which is what inflates the ratio. Checked at 4x.
paa eengang og kan derfor opfattes som to Øine. Thi saasnart vi begynde at
tænke, støde vi stedse paa den os umiddelbart imødekommende Realitet, hvis
Negation Tænkningen vil finde; derfor maa den Tænkendes Sjel stedse med det ene
Øie fastholde den positive Fylde, medens den med det andet søger at finde det
simpleste og reneste Moment i dennes Tilværelse. Afstanden mellem
Udgangspunktet og den Fylde, der er resulteret af dette, udgjør altsaa den
Grundmodsigelse, af hvilken alle de midt imellem begge liggende Modsigelser
fremgaae, saa at enhver Videnskab i denne Betydning kan kaldes hypothetisk.

Forholder nu dette sig saa, kan man ogsaa med Rette paastaae, at Methodens
Hovedopgave er den, fuldstændigt at afsløre Grundmodsigelsen i alle
Tilværelsens enkelte Sphærer. Thi Forskjellen imellem disse enkelte
Modsigelser beroer paa Forskjellen mellem de positive Realiteter, fra hvis
Negation man gaaer ud.

Det første Positive, hvorfra vi gaae ud, er jo vor egen Selvbevidsthed, der
støtter sig paa den objective Realitet, og da dette er det alleralmindeligste
Grundlag, saa følger heraf, at dets Grundmodsigelse maa udstrække sig overalt,
og at vi ikke kunne antage denne for at være knyttet til nogen Hypothese, da
Hypothesen netop er Tænkningens egen nødvendige Betingelse. Denne Modsigelse
kunne vi da kalde den \emph{ontologiske}.
% The letterspacing runs over the line break: "onto-" ends the penultimate line
% and "logiske." the last. Both halves measured; \emph{} covers the whole word.
% --- p. 41 ---
Det næste positive Grundlag fremkommer, naar man gaaer ud fra at negere
% ERRATA, applied: the page prints "næsten" (verified on the image at 6x, and a
% contemporary hand has struck the final n through in ink). The Rettelser leaf
% gives  p. 41 l. 1 fra oven: næsten l. næste.  Corrected here to "næste".
Realiteten selv, hvis Ophav den absolute Selvbevidsthed er, og da navnlig den
Realitet, vi nærmest have for Øie, vor Verden nemlig; ogsaa heraf fremgaaer en
Grundmodsigelse, men dog saaledes, at det Hypothetiske her træder langt
stærkere frem; thi vel er det vist, at den absolute Selvbevidsthed kan
acquiescere i sig selv, men hvilken Realitet der skal fremgaae af den, det
% "acquiescere" with QU on this page, against "acqviescere" with QV on p. 37.
% Both read at 5x; both stand as printed.
synes at være hypothetisk og at forudsætte Friheden som sin nødvendige
Betingelse, da vor Verden vel ikke er den eneste, i hvilken det guddommelige
Liv kan aabenbare sig. Betragte vi derfor i Dogmatiken Theologien (Faderens
Rige), Christologien (Sønnens Rige) og Pneumatologien (Aandens Rige), kunne vi
bemærke, at det, vi i det første Rige skulle negere, er Realiteten i dens rene
Almindelighed, medens vi derimod i det andet og tredie Rige mere specielt fæste
Blikket paa vor Verdens Realitet, saa at i disse Modsigelsen og Nødvendigheden
sættes inden visse bestemte Grændser; thi det er en bestemt Verden, vor Verden
nemlig, i hvilken Adam syndede og Christus incarnerede sig, der her bliver
negeret. Vi kunne altsaa kalde denne Grundmodsigelse den \emph{dogmatiske}.

Men som den almindelige Realitet gaaer over til at blive denne Verdens
Realitet, saaledes udtræder ogsaa denne Verdens Realitet i speciellere
Realiteter; herved komme vi til den tredie Skikkelse af Grundmodsigelsen, i
hvilken Hypothesen vel endnu langt mere gjør sin Magt gjældende, men som
desuagtet ogsaa indenfor sine Grændser medfører den dialektiske Nødvendighed;
og dette er den \emph{historiske}.

Vi troe da saaledes at have beviist, at enhver Methode fører igjennem
Modsigelsen, at enhver Videnskab udfordrer, at dens Momenter gjensidigt ophæve
hinanden, og at al Dialektik er negativ.
% --- p. 42 ---
% The printed head and the Indhold AGREE here: both read "Begriben og Anskuen."
% Checked against the Indhold list in the header. (§ 7 on p. 36 is so far the
% only place where the two diverge.)
\parag{8}{Begriben og Anskuen.}

Den objective Fylde, vi ved Videnskabens Hjælp søge grundigt at optage i os,
maa betragtes fra følgende tvende Sider. Paa den ene Side maa man foretage en
Analyse af de enkelte Momenter for at faae dem at see i deres rette gjensidige
Belysning; men paa den anden Side fremstiller ogsaa Totaliteten af alle
Momenterne sig for vort Øie og fordrer saaledes tillige en umiddelbar Anskuen.
Ethvert Begreb er saaledes bestemt paa en doppelt Maade; thi deels maa det
indenfor sine egne Grændser have en særegen Eiendommelighed, hvorved det
adskiller sig fra de andre Begreber og holder dem uden for sig, --- og denne
udgjør det faste ubevægelige Individualitetsmoment, vi før have omtalt
% ERRATA, applied: the page prints "det fast ube-vægelige" (line 13 from the
% top, counting the two lines of the § head, verified on the image at 5x). The
% Rettelser leaf gives  p. 42 l. 13 fra oven: fast l. faste.  Corrected here.
(§~\textit{2}), --- men paa den anden Side maa det, for at kunne fastsættes som
et Begreb, nødvendigviis tillige kunne opløses i andre Begreber. At opfatte
Forholdet mellem de enkelte Begreber ved paa eengang at sammenfatte flere
Momenter kan kaldes at \emph{begribe}, medens derimod det Særegne, der tilhører
ethvert af dem, maa henføres til \emph{Anskuelsen}; Videnskaben optager derfor
begge Momenter i sig, baade Begriben og Anskuen. Vilde man ene og alene
bestemme et Begreb ved at angive dets Forhold til andre beslægtede Begreber,
saa vilde derved \textit{eo ipso} enhver Bestemmelse ophøre; thi kan man kun
% "eo ipso" is ANTIQUA, and so are the letters A, B, C below --- all checked at
% 5x against the surrounding Fraktur.
bestemme, hvad \textit{A} er, ved at henvise til \textit{B}, og atter hvad
\textit{B} er, ved at henvise til \textit{C}, forsvinder Bestemmelsen i det
Uendelige, og idet Skygge følger paa Skygge, slukkes ethvert Erkjendelsens Lys.
Og fremstilles paa den anden Side Gjenstanden som et hvilende, dødt Object uden
indre Bevægelighed, faaer man vel en forstenet Positivitet, men heller ikke her
nogen sand Erkjendelse. Ogsaa Identiteten af disse to modsatte
% --- p. 43 ---
Sider er formaliter tilstede i den formelle Viden, men er realiter og absolut
% "formaliter" and "realiter" are FRAKTUR here, not antiqua --- read at 6x, so
% no \textit{}. The ABBYY layer's "formali'tcr" invites the opposite call.
tilstede i Gud.

Som Tænkeren er, saaledes ere ogsaa hans Tanker. Fra denne Sætning maa man gaae
ud, dersom man vil bestemme den guddommelige Tænknings Natur. Den absolut
reelle Selvbevidsthed veed sig tillige som den absolute Positivitet. Den
concentrerer sig saa at sige i sin Villie, i det den bekræfter sig selv og
derved har alle sit Livs Elementer i fuldkommen Besiddelse. I det den gjør sig
selv til Object for sig, betragter den ikke en tom Skygge, men den fyldigste
Realitet, og derfor faae ogsaa alle de Tanker, der mediere denne dens
Selvbetragtning, Deel i den samme Realitet; thi hvis Gud ikke selv udtaler sit
absolute: ”Her er Jeg!” (\emph{Hinneni}), vilde heller ikke de Ting, han
% *** HEBREW IN THE PRINTING --- SUBSTITUTION, NOT A TRANSCRIPTION. ***
% The page sets the pointed Hebrew word itself, read at 30x from the native
% bitmap: he + hiriq, nun + dagesh + tsere, nun + hiriq, yod --- הִנֵּנִי,
% "hinneni", "Here am I" (Gen. 22:1 etc.), glossing the Danish ”Her er Jeg!”
% just before it. Code points U+05D4 U+05B4 U+05E0 U+05BC U+05B5 U+05E0 U+05B4
% U+05D9. This edition is built with pdfLaTeX + Libertinus and has NO Hebrew
% font, so the word is rendered by its name, exactly as the second Hebrew word
% in this book is at p. 33 (\emph{Elohim}), which follows
% texts/treschow/menneskelige-natur p. 115. If cjhebrew is installed the glyphs
% should be restored: add \usepackage{cjhebrew} to the preamble and set
% "\cjRL{hi n.E ni y}" in place of "\emph{Hinneni}" --- that input string is
% reconstructed from the pointing and has NOT been compile-tested. This is the
% SECOND Hebrew word in the book; see the OPEN ITEM in RESUME-NOTES.
sætter, komme til nogen sand Individualitet. Men idet han saaledes sætter sig
selv og sætter Alt, anskuer han tillige Alt. Han anskuer, i det han sætter, og
aabenbarer sig saaledes som den absolute Substants; han sætter, idet han
anskuer, fordi han er den absolute Frihed: ”\emph{Sandheden er, fordi Gud vil
den.}” --- Dersom nu dette Forhold indskrænkede sig til Guddomslivet alene,
% QUOTATION MARKS: both marks RAISED (U+201D), as everywhere in the body ---
% verified at 12x on this page. Contrast the footnote on p. 48 below, which is
% the first place in the book to set the low-high pair.
vilde hele Existentsen adsplitte sig i Monader, der, som Leibniz siger, ikke
havde Vinduer, igjennem hvilke de kunde see ind til hverandre. Men af Guds egen
individuelle Tilværelse følger det, at ogsaa hans Tanker gaae uendeligt over i
og gjennemtrænge hinanden. Da det er den samme Gud, den samme Selvbevidsthed,
der tænker Alt, maae ogsaa hans forskjellige Tanker staae i indbyrdes
Sammenhæng med hinanden. Gud anskuer nemlig ikke blot de enkelte Ting
stykkeviis, thi hans Viden er simultan, men han udleder tillige den ene af den
anden, betragter tillige det gjensidige Forhold mellem dem alle; og forholdt
det sig ikke saa, vilde han hverken kunne klart adskille eller
% --- p. 44 ---
fuldkomment gjennemtrænge Tiugene med sin Viden; --- uden Abstraction, ingen
% PRINTER'S DEFECT, transcribed as printed: "Tiugene" with a TURNED n, read at
% 5x. Cf. "Tiugen" on p. 3, already logged in the header --- the same sort.
Distinction.

Vel udfordrer den guddommelige Tænknings absolute Idealitet, at ethvert af hans
enkelte Objecter har Totaliteten af alle hans øvrige Tanker i sig; men skal
dette ikke være en blot tautologisk Tænkning, er det paa den anden Side ligesaa
nødvendigt, at Ideernes Fylde maa være tilstede i de forskjellige Objecter paa
forskjellig Maade. Gud kan da ikke tænke Objectet \textit{A} uden tillige
derfra at abstrahere Objectet \textit{B}; --- Differentspunktet er
Abstractionens Udgangspunkt. --- Derfor maae de, der udhæve Guds Viden som den,
der er simultan og paa eengang anskuer Alt, ikke glemme, at de guddommelige
Tanker udgjøre et System, og at Gud i Anskuelsen af de enkelte Ting
gjennemløber den hele lange Tankevei fra Begyndelsespunktet, det rene Intet,
til den enkelte concrete Ting, saa at hans umiddelbare Viden tillige fra
Evighed af har været og stedse vil vedblive at være en medieret.
% "været" and "medieret" both end in the Fraktur t, not r --- settled at 6x by
% comparing them with the t of "at" on the same line. Both OCR witnesses give
% "-er" for both words.

Vil Nogen indvende, at vi herved overføre upassende Anthropomorphismer paa Gud,
og anvende Athanasii bekjendte Ord: ”\textit{o ludicram doctrinam aedificantem
simul et demolientem}” paa os, da maae vi bede ham vel erindre, at
% The Latin quotation is ANTIQUA against the Fraktur; the two quotation marks
% around it are RAISED, read at 6x.
Abstractionen og Mediationen ikke paa udvortes Maade adskilles fra den
simultane Viden, men at meget mere Begriben og Anskuen falde sammen. Idet Gud
begriber sig selv og hans Tanker ved et Nødvendighedens Baand hænge sammen med
hverandre, har han den absolute Sandhed: ”\emph{Gud vil, fordi det er
Sandhed.}” Den nøiere Udvikling heraf maa forøvrigt overlades den speculative
Theologie; vi have kun berørt disse Punkter for at fremstille den menneskelige
Videns Urtypus.
% --- p. 45 ---
Den menneskelige Selvbevidsthed, der er et Afbillede af den guddommelige,
indeholder paa formel Maade de ovennævnte Momenter i sig. \emph{Jeg} gaaer
nemlig ideligt over i \emph{Mig} og \emph{Mig} i \emph{Jeg}; den ene Tanke
% Only these four words are letterspaced. "Jeg" and "Mig" further down the page
% ("thi naar Jeg betragter Mig", "hvad Jeg er") are NOT --- checked at 4x.
optager stedse den anden, opløser sig stedse i den anden; der maa saaledes
finde en Begriben Sted; --- men Anskuelsen er ikke mindre tilstede; thi naar
Jeg betragter Mig, maa hiint stedse see dette som noget Positivt; var dette
nemlig ikke Tilfældet, kunde man heller ikke fatte, hvad Jeg er, og der kunde
da slet ingen Selvbevidsthed være til. Men at det Positive og Anskuelige er et
nødvendigt Moment, dersom Selvbevidstheden skal have Realitetens Fylde, er let
at indsee. Thi skal man kunne komme til en Begriben, maa man ogsaa kunne anskue
den reelle Fylde, og en saadan Anskuelse er ikke alene nødvendig i
Almindelighed, men kræves ogsaa i de enkelte Begreber, for at hvert enkelt
Begreb derved kan erholde sin sande Intensititet. --- At dette forholder sig
% PRINTER'S DEFECT, transcribed as printed: "Intensititet" for "Intensitet",
% the syllable -ti- set twice ("Jntensi-" at the foot of one line, "titet." at
% the head of the next). Read at 4x; both OCR witnesses read it the same way.
saa, viser ogsaa selve Dialektiken; thi naar man vil udvikle et Begreb af de
Begreber, der gaae foran det, antyder og beskriver man det vel først med andre
Udtryk, men gaaer derfra over til dets egen sande og nødvendige Betegnelse. Men
Tilføielsen af en saadan Kategorie kan ikke have nogen anden Grund end den, at
denne medfører Anskuelsen af et vist positivt og særegent Moment, og at denne
er nødvendig, for at Begrebet i Virkeligheden kan erholde Individualitet. ---
Man kan derfor ogsaa i Ontologien selv paavise Anskuelsens Betydning. Synes der
saaledes ikke at ligge en langt kraftigere Stimulus til dialektisk
Fremadskriden i den Lære, at den rene Væren gaaer over til det rene Intet, end
i den, at Verdenssjelens Idee nødvendigt maa gaae over til den abstracte
Personligheds Idee? Dette kan neppe benegtes; men Grunden dertil er ikke
saameget hiin Læres større Begribelighed, som langt snarere dens
% --- p. 46 ---
\setcounter{footnote}{0}
% THE PRINTED FOOTNOTE MARK ON THIS PAGE IS AN ASTERISK, "*)", NOT "1)" ---
% read at 5x. It is the second such in the book (cf. p. 12) and the only note
% in pp. 37--48 that is not a plain figure. \thefootnote will print "1)".
% THE NOTE RUNS OVER ONTO p. 47: the block below the rule on p. 47 carries no
% mark of its own and completes the second quotation ("saa maa Mennesket frit
% underkaste sig Gud, ... troe for at erkjende.” Samme, P. 15."). It is set
% here as one \footnote, so p. 47 gets NO \setcounter and no note of its own.
større Anskuelighed. Thi først den, der umiddelbart har erfaret
Personlighedens Betydning, kan erkjende, at Universet ikke kan fastholdes eller
begribes uden ved at sees udgaaet fra en Person. Ja Nødvendigheden af en saadan
forudgiven Anskuelse vil stedse paa ny gjentage sig, efterdi Guds absolute
Frihed ikke blot har skabt Realiteten i Almindelighed, men ogsaa denne bestemte
Verden og alle de bestemte Phænomener, der møde os i denne. Man maa da have sat
sig i Besiddelse af den hele Kirkelære, før man kan udvikle, hvorledes denne
Frihed er beskaffen, og først efter at have udviklet Friheden i Almindelighed,
kan man gaae over til de enkelte Phænomener i Historien. Heraf indsees da, hvad
det vil sige, at det er nødvendigt at troe, en Fordring, der saa ofte gjøres,
naar der tales om den sande Videnskab. Jeg fatter Guds Personlighed ved at
gjennemføre Negationen af min egen Personlighed; havde jeg nu ikke selv en reel
Personlighed, kunde jeg heller ikke have noget høiere Personlighedsbegreb; men
skal den være reel, maa den ogsaa være religieus; --- at kjende Gud er altsaa
at bekjende Gud \dsi at overgive sig til Gud. Som derfor objectivt betragtet
Alt gaaer op i en ikke blot tænkende, men ogsaa villende Tilværelse, saaledes
ligger ogsaa subjectivt betragtet det at kjende Gud ikke blot i Erkjendelsen af
ham, men ogsaa i Villiens Beredvillighed til at underkaste sig
ham.\footnote{”Da man altsaa i Samvittighedens Lys kommer til sand Erkjendelse
af et \emph{Ophav} for al menneskelig Selvbevidsthed, kunne vi nøiere bestemme
den som det Lys, i hvilket Mennesket aabenbares som \emph{Guds Skabning}; i
denne Erkjendelse ligger Spiren til al Religion og Gudsdyrkelse.”
\par H.\ Martensen, Selvbevidsthedens Autonomie, oversat af Petersen,
P.\ \textit{10}.
\par ”--- --- --- da desuden Friheden er Princip for enhver Erkjendelse, idet
det ene Fornuftvæsen kun gjennem Frihedens Mediation kan meddele sig til det
andet og leve og røre sig i det: saa maa Mennesket frit underkaste sig Gud, for
at den guddommelige Idee reelt skal kunne trænge ind i hans Aand og inderlig
forenes med hans Bevidsthed; Mennesket maa, med andre Ord, \emph{troe for at
erkjende}.” Samme, P.\ \textit{15}.}
% The note is FRAKTUR throughout, so nothing in it is \textit{} except the two
% page figures, which are antiqua. Letterspacing inside the note, all read at
% 4x off the image: "Ophav"; "Guds Skabning"; "troe for at erkjende" (that last
% on the p. 47 run-over). Both quotations take the RAISED mark at both ends.
% "P." is the Fraktur P (= Pagina), not antiqua.
% --- p. 47 ---
% NO \setcounter here: the block under the rule on this page is the tail of
% p. 46's note, carries no mark of its own, and is set inside that \footnote.
% This page begins WITHOUT an indent, so the paragraph runs on from p. 46.
Det er netop i dette Punkt, det ontologiske Beviis for Guds Tilværelse falder
sammen med den religieuse Selvbevidstheds Vidnesbyrd; som hiint taler til vor
% NOT emphasised, despite spacing.py's RUN on "taler til": the letters are
% tight and only the word spaces are open (justification). Checked at 4x.
Begriben, saaledes taler dette til vor Anskuelse. Man maa da have religieus Tro
og Erfaring for ret at kunne udvikle den christelige Religions Dogmer; --- at
man maa have tilegnet sig Kirkelæren, have vi ovenfor antydet --- thi det er
gjennem disse, at Sjelen fyldes med indre Sandhed; det er derfor ogsaa først
igjennem dem, vi blive istand til at gjennemtrænge de dialektiske Momenter. Men
denne Tro er netop Genialiteten i dens høiste Skikkelse \dsi den religieuse
Genialitet.

% The printed head and the Indhold AGREE here: both read "Vor Viden er relativ."
\parag{9}{Vor Viden er relativ.}

Af det giensidige Forhold imellem Begriben og Anskuen følger imidlertid, at vor
% "giensidige" with an I on this page, against "gjensidige" with a J on pp. 42
% and 43 --- all three read at 6x, and the j of this fount both carries a dot
% and descends, while the i does not. As printed; cf. "giennem" on p. 38.
Viden ikke kan blive fuldkommen, før Gudslivet i dets hele Fylde er blevet
aabenbaret, eller, med andre Ord, før Salighedens Rige er kommet. Thi vort
Begrebs Klarhed er jo stedse vunden gjennem et større eller mindre Tab af den
reelle Fylde, og jo mere vi henvende vor Opmærksomhed paa concrete og særegne
Gjenstande, desto mere drages vi ind i Frihedens hypothetiske Sphære. Ja dette
har endog sin naturlige Grund i det gjensidige Afhængighedsforhold, hvori
Videnskabens Discipliner staae til hverandre, hvilket medfører, at ingen af dem
kan blive fuldkom\-%
% --- p. 48 ---
\setcounter{footnote}{0}
ment afsluttet, fordi den enes Fuldstændighed beroer paa den andens. Saalænge de
tvende forskjellige Betragtningsmaader af Gjenstandene, den vulgaire gjennem
Sandserne og den speculative gjennem Tænkningen, kun optage hinanden, idet de
uendeligt forudsætte hinanden, det er, saalænge vi færdes i dette Liv, saalænge
vedbliver der i vor Viden at herske et relativt Mørke
% ERRATA, applied: the page prints "i vor Virken at herske" (line 6 from the
% top, verified on the image). The Rettelser leaf gives
%   p. 48 l. 6 fra oven: Virken l. Viden.  Corrected here to "Viden".
(\textit{βλέπομεν γὰρ ἀρτι δι' ἐσόπτρου ἐν αἰνίγματι}). Dette maa ikke blot
% GREEK (1 Cor. 13:12), read at 6x. Fount variants normalised per the standing
% convention: the cursive rho -> ρ (γὰρ, ἐσόπτρου).
% ACCENT AS PRINTED: "ἀρτι" carries its smooth breathing but NO acute (received
% ἄρτι). The mark over the alpha is the same curl as the psili over the epsilon
% of ἐν and ἐσόπτρου, and quite unlike the straight grave over the alpha of
% γὰρ on the same line. Recorded rather than corrected --- cf. "Αληθῆ" without
% its breathing and "ἀν" without its grave on p. 36.
tilskrives en Mangel i vor Bevidsthed; vor Videns Klarhed formørkes meget mere
ved den sandselige Verdens egen Uklarhed. Vi kunne ikke erkjende Noget, før det
fremtræder i det objective Lys\footnote{Enhver Videnskab er en Viden om Noget,
som Videnskaben ikke frembringer, men, som der er nafhængigt af den.
\emph{Mynster}: „Om Begrebet Dogmatik.“ §~\textit{4}. Smlgn.\ \emph{Sibbern}:
„Om Erkjendelse og Grandskning.“ §~\textit{1}. \textit{Scimus, quia
accepimus}: Al Erkjenden er en Finden, Modtagen og Skuen.}. Vi kunne derfor
% *** THE ONLY LOW-HIGH QUOTATION MARKS FOUND IN THIS BOOK SO FAR. ***
% This note, and this note alone, sets the Danish pair „…“ where the body and
% every other quotation in pp. 1--48 set the raised mark ” at BOTH ends. It was
% settled by measurement, not by eye, on the native bitmap at 12--30x:
%   * the two OPENING marks (before "Om Begrebet" and before "Om Erkjendelse")
%     have their feet exactly ON the baseline --- their ink ends at the same y
%     as the bottom of the neighbouring "O" and of the preceding full stop ---
%     and they are top-heavy (upper/lower ink ratio 1.3--1.5), i.e. ball at top:
%     that is „.
%   * the two CLOSING marks hang from the x-height, 18 px clear of the baseline,
%     and are bottom-heavy (ratio 0.84--0.89), i.e. ball at bottom: that is “.
%   * the same measurement on the raised marks of p. 46's note gives ratios of
%     1.9--2.6 with the ink sitting at cap height --- the book's usual ”.
% So the book's rule ("” at both ends") holds everywhere else and is broken
% here. Reproduced as printed. „ and “ balance within the note (2 and 2), so
% check.py's ” parity test is unaffected.
% PRINTER'S DEFECT in the same note, transcribed as printed: "nafhængigt" with
% a TURNED u, for "uafhængigt". Read at 4x; both OCR witnesses read "naf".
% "Mynster" and "Sibbern" are letterspaced; the Latin tag "Scimus, quia
% accepimus" is antiqua. The rest of the note is Fraktur.
hverken i Naturen eller i Historien og ligesaalidt i Religionen selv vente at
komme til absolut Viden, da det netop er denne Verdens (\textit{τοῦ αἰῶνος
τούτου}) ejendommelige Charakteer, at alle Livets Elementer i den adsplittes i
% GREEK (2 Cor. 4:4), read at 6x. Normalised placement: this fount sets the
% perispomeni and the acute of an -ου- diphthong over the OMICRON, so the page
% reads in effect tau + omicron-with-perispomeni + upsilon and "τόυτου". Unicode cannot carry a perispomeni on
% omicron at all, so the standard placement is set (τοῦ, τούτου) and the
% printed placement recorded here --- the same habit as οὐ / ποιοῦσι on p. 35
% and ἐμοῦ on p. 36. The circumflex on αἰῶνος is over the omega, as it should
% be, and the psili over the iota of αἰ- is as printed.
løsrevne Brudstykker. Og hvorledes skulde vi ogsaa kunne gribe Sandheden, før
den selv har naaet en adæqvat Form, hvori den fuldkomment kan aabenbare sig, da
det jo netop er Sandheden selv, der uddanner Selvbevidstheden i os?

Det er netop denne Indskrænkning i vor Viden, der betegnes i den vulgaire
Yttring, at der vel ere adskillige Ting, som Menneskene kunne begribe, men at
de ikke kunne begribe Alt, ikke det Høieste.

Men forsøger man nu at henføre alle Ting til to Klasser, de begribelige og de
ubegribelige, saa at de skulle staae afsondrede hver paa sin Side, vil man i
Sandhed komme til det Resultat, at de enten alle ere begribelige, eller at slet

% --- p. 49 ---
\setcounter{footnote}{0}
% JOINT WITH p. 48, verified on the image at 0.42x: p. 48 ends "eller at slet"
% and p. 49 opens "ingen af dem er det." --- "slet", not "let" (the Fraktur OCR
% reads "let"; the ABBYY layer and the image both give "slet").
ingen af dem er det. Thi see vi først hen til de Ting, vi kalde begribelige, da
begrændses og bestemmes de jo netop ved deres Forhold til dem, vi kalde
ubegribelige. Og gaae vi ud fra de ubegribelige Ting, maae de jo nødvendigviis
træde ud i de Ting, der ere begribelige, og blive saaledes selv begribelige. Vi
kunne af Dogmatiken selv anføre Exempler herpaa: netop ved at urgere, at den
menneskelige Natur har bevaret sine Kræfter uskadte, mene Pelagianerne at
forklare alle det religieuse Livs Phænomener, medens tvertimod Augustinianerne
for at tolke Syndens og Gjenfødelsens Natur henvise til Adams Fald som
ødelæggende for hele Menneskeslægten, samt til Guds Naade og Retfærdighed, med
det Tilføiende, at Gud ikke havde forudbestemt Adams Fald; ja Calvinisterne gaae
endnu videre, idet de lære, at Gud har forudbestemt Alt. De ville saaledes Alle
\emph{begribe} det religieuse Livs Phænomener, idet de søge at forklare dem og
% "begribe" IS letterspaced (checked at 0.50x: clear intra-word gaps against the
% tight setting of the same word two lines below). "Ingen" in "kommer Ingen af
% dem" is NOT --- spacing.py flagged it, the image refutes it.
forene dem med hinanden; men i Bestræbelsen efter at naae dette Maal komme de
Alle, Hver paa sit Sted, til at indsmugle et ubegribeligt Moment, og netop derfor
kommer Ingen af dem til en sand Begriben. Vel antager man i Almindelighed, at
Pelagianernes Sætninger ere lettere at begribe, end Augustinianernes og
Calvinisternes; men da det Forhold mellem den guddommelige og menneskelige
Frihed, som Pelagianerne lære, aabenbart staaer i Modsætning til Begrebet om Guds
Almagt og Naade, og saaledes alle menneskelige Handlinger ere Yttringer af en
ubegribelig menneskelig Frihed, forraade jo alle disse Handlinger overalt et
ubegribeligt Moment og ere derfor selv ubegribelige\footnote{At
\emph{Augustinus} var Infralapsarier, giver \emph{Neander} (Allgem. Gesch. d.
Rel. und Kirche. \textit{2} Band. \textit{3}te Abtheil, \textit{p.\ 876})
Anledning til følgende Bemærkning: „Aber es war dies eine schöne Inconseqvenz,
welche aus dem Siege seines religiössittlichen Gefühls über seine
dialektisch=speculative Richtung herrührte.}. --- Paa den anden Side kan Calvin,
% *** QUOTATION MARKS IN THE FOOTNOTES ARE NOT THE BODY'S. *** The header records
% that the BODY sets a raised ” at both ends. That holds on every body quotation
% in this range (pp. 54, 56, 57 --- all four marks checked at 4x). But the note
% on this page opens its German quotation with a LOW „ (U+201E): read at 4x, the
% two marks sit squarely on the baseline, exactly as the comma of "Bemærkning,"
% does, and are unmistakably distinct from the cap-height marks of ”Og Ordet
% blev Kjød on p. 54. The same low opening recurs three times in the note on
% p. 50. Where those notes close, the mark is the raised ” again --- so the notes
% run „…”, a mixed pair. NOT universal, though: the note on p. 60 opens its Latin
% quotation with a RAISED ”. Each is transcribed as printed.
% PRINTER'S DEFECT: this quotation is never closed. The note ends "…Richtung
% herrührte." and nothing follows (checked at 2x on the line end); the „ has no
% partner. Transcribed as printed. Contributes „=1, ”=0 to the parity count.
% "Neander" and "Augustinus" are both letterspaced in the note; "Infralapsarier"
% is not. Read at 1.15x. "Abtheil," carries a comma, not a period.
da han
% --- p. 50 ---
\setcounter{footnote}{0}
har givet Præmisserne, ikke undlade til Gunst for den menneskelige Erkjendelse at
fuldende sin Slutning; han refererer derfor Alt til et \textit{decretum
absolutum} hos Gud, der overgaaer al menneskelig Forstand, og kommer saaledes i
sin Bestræbelse efter at begribe Forholdet til et Resultat, der er aldeles
ubegribeligt.

Vi finde da saaledes, at det ikke afhænger af Sagens egen Natur, hvor dette den
menneskelige Erkjendelses Modermærke, om vi saa tør sige, er at finde, eller,
hvorledes det skal skjules, men at det meget mere har en anden Grund, hvad enten
vi nu ville kalde denne subjectiv Vilkaarlighed eller religieus
Tact\footnote{Derfor kunne ogsaa de Indvendinger, der endnu i vor Tid ere
fremførte imod Calvins Conseqventser i denne Henseende, vel fortjene at høres.
Saaledes siger \emph{Guerike}: Kirchengesch. P. \textit{1020}, Not.
\textit{216}: „Die heilsam inconseqvente Vorstellungsweise der Concordienformel
im Verhälltnisse zu der selbst gotteslästerlich conseqventen Calvins ist also
kurz diese: Seligkeit und Hoffnung ist nach beiden nur ein Werk der göttlichen
Gnade, Verdamniß eine Folge der eignen Schuld. Während aber nun die
Concordienformel demüthig bei diesem letzteren Satze stehen bleibt . . . . .
argumentirt Calvin weiter etc.” Men naar Guerike fandt for godt at tilføie:
„Gott also hat selbst die Sünde der Menschen gewollt. Diese letzte,
nothwendige, ob auch noch so bemäntelte Schlußfolge . . . . . synes han ved at
tilføie denne nye Slutning selv at have syndet imod den fromme Ærefrygt for det
Hellige, som han anpriste hos Concordieformlen, og jeg tvivler ikke paa, at
Calvin, hvis han endnu var ilive, vilde svare ham med de samme Ord, hvormed han
besvarede sin Tids irreligieuse Stræben efter at begribe de guddommelige
Mysterier: \textit{Arripiunt subito, fatcor, profani homines in prædestinationis
materia, quod vel carpant, vel cavillentur, vel allatrent, vel subsannent. At si
nos absterret eorum procacitas, celanda crunt præcipua quæquc fidei dogmata.
(Instit. christ. rcl. De prædest cap. XIV.).}”}. Det er altsaa aabenbart, at
% NOTE ON THIS NOTE --- it is the heaviest in the range and runs on to p. 51.
% Setting: "Guerike" is letterspaced at its first occurrence and NOT at its
% second ("Men naar Guerike fandt"); read at 0.62x. The German is Fraktur, so no
% \textit{}; the Latin and every figure are antiqua, so they take it.
% PRINTER'S DEFECTS, all transcribed as printed:
%   * "Verhälltnisse" with a doubled l, for "Verhältnisse".
%   * "Verdamniß" with a single m, for "Verdammniß".
%   * the „Gott also… quotation is never closed --- the ellipsis is followed
%     straight by the Danish "synes han ved at tilføie". Read at 0.62x.
%   * the Calvin Latin is never OPENED: the line begins flush left with
%     "Arripiunt", with no mark before it (checked at 2.5x, 150 px of clear
%     paper), yet it CLOSES with a raised ” after "XIV.)." on p. 51.
%   * FOUR e/c SUBSTITUTIONS IN THE LATIN, all transcribed as printed and all
%     verified at 8x: "fatcor" for fateor, "crunt" for erunt, "quæquc" for
%     quæque, "rcl." for rel. (Calvin, Instit. III.xxi: "Arripiunt subito,
%     fateor, profani homines … celanda erunt praecipua quaeque fidei dogmata".)
%     THE ALTERNATIVE READING WAS TESTED AND REJECTED. The tempting explanation
%     is that the 1-bit JBIG2 scan dropped four thin 'e' crossbars, in which case
%     the correct Latin should be set. Two things rule it out. (a) At 8x the
%     glyphs are crisp, well-formed c's with clean upper and lower terminals ---
%     not ragged, half-inked e's --- while genuine e's on the very same lines
%     ("fidei", "De", "celanda", "absterret", "eorum") carry unmistakable
%     crossbars. (b) A connected-component census of the two footnote blocks
%     (pp. 50--51, 312 components) found 304 DISTINCT bitmaps: no two glyph
%     instances are pixel-identical, so this file is not symbol-substituted
%     JBIG2 and its glyphs are what is on the film. Four e-for-c mis-picks in one
%     stint of unfamiliar Latin, in the smallest fount on the page, is an
%     ordinary compositor's failure. Worth a look at the KB colour PDF.
%   * the Latin ends "XIV.)." --- period, paren, period --- then the raised ”.
%     "prædest" carries no point.
% Quote parity contributed by this note: „=2, ”=2. Two of the three quotations
% are defective, and they cancel in a parity test, which is why they are logged.
man ikke paa udvortes Maade kan opstille hiin Adskillelse.
% --- p. 51 ---
% NO \setcounter HERE, AND NONE IS MISSING: p. 51 carries no note of its own.
% The block below its rule is the tail of p. 50's note ("subsannent. At si nos
% absterret…"), and the body of p. 51 has no reference mark anywhere on it.
Maa nu saaledes Mysteriet gjennemtrænge alle Momenter, naar det først er tilstede
paa eet Punkt, kunne vi ogsaa omvendt slutte, at vor Viden maa kunne udstrække
sig til alle Gjenstande, saasnart den blot er tilstede paa eet Punkt; thi maatte
den bukke under for Mysteriet, kunde den overhovedet slet ikke være til, da jo
Mysteriet overalt gjør sin Magt gjældende; med andre Ord: vor Videns
Indskrænkning er ikke metaphysisk, men empirisk. De Love, hvortil vor Erkjendelse
er bunden, gjælde i al Evighed, og dette stemmer ogsaa med Sandhedens Natur; thi
om Sandheden gjælder det Samme, der gjælder om Gud, at den ikke er langt borte
fra Nogen af os (\textit{Ἡμῖν δὲ ἀπεκάλυψε ὁ θεὸς διὰ τοῦ πνεύματος αὐτοῦ}
\textit{1} Cor. \textit{2, 10}).
% GREEK, read at 4x. Fount variants normalised per the standing convention:
% script theta -> θ (θεὸς), cursive kappa -> κ (ἀπεκάλυψε).
% ACCENT PLACEMENT, normalised and logged as the header requires: this fount sets
% the perispomeni of an -ου- diphthong over the OMICRON, so the page reads (in
% effect) "omicron-with-perispomeni + upsilon" in both τοῦ and αὐτοῦ. Unicode
% cannot encode that, so standard placement is set. The SAME habit was recorded
% on pp. 35--36.
% Accents that ARE encodable are kept exactly as printed, and two depart from the
% received text of 1 Cor. 2, 10: the print has ἀπεκάλυψε without the final ν, and
% Ἡμῖν carries its rough breathing over the capital H (verified at 4x --- the
% mark is the c-shaped dasia, not the psili that stands over the α of
% ἀπεκάλυψε).
% "Cor." is FRAKTUR; only the figures are antiqua.

Derfor skal den menneskelige Bevidsthed ikkun urokkeligt holde fast ved det
guddommelige og reelle Object og omhyggeligt stræbe efter at gjennemtrænge dettes
Momenter; den vil da i sit Indre finde, hvad den søger (\textit{quod petis hic
est}), og den absolute Erkjendelses Grændser ville saaledes indenfra mere og mere
udvide sig for den.

% *** THE PRINTED HEAD AND THE INDHOLD DISAGREE HERE --- the second such case in
% this book, after § 7 on p. 36. ***
% The printed § head on p. 51 reads "Forholdet mellem Dogmatiken og den hellige
%   Historie." (read at 0.38x; the word is plainly "mellem", four letters after
%   the initial m of the second word, no leading i).
% The Indhold (PDF 8) reads  "Forholdet imellem Dogmatiken og den hellige
%   Historie. ... 51."
% The printed head is set, per the standing rule that the page governs; the
% Indhold reading is recorded here rather than normalised away. The body of the
% section immediately below also uses "mellem", which supports the head.
% The head is printed "§ 10." --- no point after the section sign.
\parag{10}{Forholdet mellem Dogmatiken og den hellige Historie.}

For fuldstændigere at udvikle vor Videns relative Begrændsning maae vi nøiere
betragte Forholdet mellem den speculative Dogmatik og den speculative Historie.
--- Vor Viden er absolut, deels fordi Bevidstheden er sig sit Indhold fuldkom\-%
% --- p. 52 ---
ment klar, deels fordi de reelle Momenter, den omfatter, ere aldeles universelle
og almeengyldige. Heraf følger, at Dogmatiken er den reelleste af alle
Videnskaber. Ontologiens Begreber ere nemlig aldeles abstracte, og den Viden, som
Historien meddeler, er altfor hypothetisk. Derfor kommer hiint positive Moment,
der, som vi forhen have bemærket, stedse er nødvendigt for at fremkalde
Dialektikens negative Stimulus, paa Historiens Gebeet ikke til nogen sand
anskuelig Klarhed, men bevirker ofte, at ogsaa den dialektiske Braad sløves,
fordi det selv er omgivet af sensuel Dunkelhed. Dogmatiken derimod tager vel
Hensyn til den positive Realitet, men ikke til dennes mere specielle
Modificationer, og bevirker netop derved, at Modsigelsen stedse træder frem, saa
at Begrebet bevares. Men denne Ejendommelighed ved Dogmatiken fremkalder en vis
Modsætning. Thi udslettes de individuelle Phænomener, forsvinder ogsaa den
almindeligere Realitet selv; man maa da tye til Hverdagserfaringen, for ved
dennes Hjælp allestedsfra at vinde den manglende Fylde, --- men da enhver saadan
Erfaring er en Negation og Begrændsning af Videnskaben, maa Hverdagserfaringen
forvandle sig til en speculativ Erfaring, \dsi{} Historien maa blive Videnskab.
% p. 52 carries no note and no letterspacing anywhere --- spacing.py returns
% nothing for it, and a page-wide look at 0.38x confirms it.

De to forskjellige Arter af Historien udfordre hver sin Tro; den i den vulgaire
Historie nødvendige Tro medieres gjennem denne Verdens Aand, der ikkun betragter
den trivielle Virkelighed; den videnskabelige Histories Tro derimod medieres
gjennem den guddommelige Aand, der sætter den ideelle Virkelighed. Om nogle
Begivenheder ere begge Arter af Tro enige; om andre eller rettere om de fleste
kunne de derimod ikke komme overeens.

Vi have her i Forbigaaende berørt disse Forhold, skjøndt de senere ville blive
Gjenstand for en nøiere Betragtning, for at gjøre det indlysende, at Dogmatiken
er en nødvendig Forud\-%
% --- p. 53 ---
sætning for enhver historisk Viden. Man kan nemlig ikke begribe Menneskeslægtens
Bedrifter uden at have opfattet Frihedens Idee i dens Sandhed; men for at forstaae
Menneskets Frihed er det nødvendigt at have forstaaet Guds Frihed; man kan derfor
først begribe Historien, naar man har tilegnet sig Religionens og Aabenbaringens
Idee. Det er da saalangt fra, at den hellige Historie krænkes, naar vi gjennem
Philosophien søge at begribe den, at den meget mere profaneres, saa ofte
Phænomenerne fremstilles paa udvortes Maade, løsrevne fra Gudsideen. Da saaledes
den vulgaire Historie stedse er profan, den speculative derimod hellig, maa den
hellige Historie være speculativ.

% A SWELLED RULE stands between the last line of the Første Deel and the Anden
% Deel head --- thick through the middle and tapering to a point at each end, not
% the plain hairline of p. 1. Measured, not eyeballed: it runs x 1167--2354 on a
% text block of roughly x 300--3150, i.e. 0.42 of the measure, centred, and its
% ink profile rises from 28 px of ink in its first row to 811 in its ninth and
% back to 10, which is the taper. \rule cannot taper, so a plain centred rule of
% the same proportion is set and the printed form recorded here. \centerline
% rather than a center environment, per the p. 1 precedent, so that check.py's
% hand-rolled-head test is not tripped. REPORTED to the caller: if a \swelledrule
% macro is wanted, this is the first site.
\par\vspace{0.6\baselineskip}\centerline{\rule{0.42\textwidth}{0.6pt}}\vspace{0.6\baselineskip}

% THE ANDEN DEEL OPENS HERE, and in exactly the display form of the Første Deel
% on p. 1: the Deel line in large bold Fraktur WITH A PERIOD --- "Anden Deel." ---
% and below it the argument in bold letterspaced Fraktur, "Den hellige Historie.",
% also with a period. It differs from p. 1 only in that the argument is short
% enough to sit on ONE line. Not a colon in either place. spacing.py returns both
% lines as RUNs, and the letterspacing of the argument is plain at 0.38x.
% The argument agrees verbatim with the Indhold.
\deel{Anden Deel.}{Den hellige Historie.}

% The head is printed "§ 11." --- no point after the section sign --- and the §§
% do NOT restart in the Anden Deel: § 11 follows § 10. Argument agrees with the
% Indhold.
\parag{11}{Katholicismen.}

% The paragraph opens with a large decorated Fraktur initial D, as § 1 does on
% p. 1; transcribed as an ordinary capital, following that precedent.
Den hellige Historie kræver, at den usynlige guddommelige Natur fremstiller sig i
synlige Phænomener; den ophæver altsaa Modsætningen mellem de skabte Tings
sandselige og den guddommelige Tilværelses aandelige Sphære, idet de aandelige
Gjenstande træde sandseligt frem, for at de sandselige Gjenstande kunne forklares
i Aandens Lys; den ophæver ligeledes Antinomien mellem Skaberens og Skabningens
Frihed, idet Gud paalægger sig Endelighedens Baand, for at give den skabte Frihed
Leilighed til at vælge og handle. Medens den saaledes medfører en umiddelbar
Identitet af det over\-%
% --- p. 54 ---
naturlige og naturlige Rige, maa den dog ligesaavist lade en Modsætning mellem
begge komme tilsyne. Thi skal Guds Historie være hellig, maa Gud æres som den
Hellige; og skal Gud æres som den Hellige, maa han ophøies over den skabte Natur
i uendelig Majestæt. Men denne Selvmodsigelse fjerner den hellige Historie
% "fjerner", read at 4x: the second letter carries both a dot above and a
% descender, so it is the Fraktur j. Both OCR witnesses give "sterner", which is
% not a word; the alternative "stemmer" was considered and rejected on the image.
derved, at den henviser til en Modsætning mellem forskjellige Tider og Steder.
Paa nogle Steder og til nogle Tider stemmer nemlig Verdenslivet, saa at sige,
ikke med Gudslivet og viser sig løsrevet fra dette, men paa andre Steder og til
andre Tider fremtræde de i umiddelbar Forsoning. I denne Anledning opstaaer ingen
følelig Splid mellem den vulgaire og speculative Erfaring; thi i Sandseverdenen
skue vi overalt med Troens Øine de overnaturlige Væsener, og saa ofte
Verdenslovene afbrydes af en høiere Tingenes Orden, saasnart Mysterier lade sig
tilsyne, giver Hverdagsbevidstheden sig villig fangen uden al philosophisk
Tumult; men neppe er den naturlige Orden vendt tilbage, før den atter indtræder i
sine gamle Rettigheder. Hvor der gives en Aabenbaring, der indtræder stedse dette
gjensidige Forhold mellem begge Sphærer; i den absolute, den christelige
Religion, fremtræder det derfor i høieste Maal.
% ON ø IN THIS FOUNT, since it bears on every page from here on. The Danish ø of
% this Fraktur is a slashed o whose stroke closes the counter, so at 3--4x it
% reads very like the Fraktur e --- which is why tesseract renders it now as "e"
% (heiere, gjere, fer) and now as "o" (Sporgsmaal, horer). It is ø, and is set
% as ø throughout, exactly as pp. 1--36 already do. Checked at 3.2x on
% "løsrevet" and "høiere" on this page.

Da Apostlene saae Christus, om hvis Fremtræden Evangelisten paa høitidelig Maade
vidner: ”Og Ordet blev Kjød og tog Bolig iblandt os, og vi saae hans Herlighed,
en Herlighed, som den Eenbaarnes hos Faderen, fuld af Naade og Sandhed”,
% BOTH marks raised (U+201D), read at 4x: the opening mark before "Og" sits at
% cap height, ball uppermost and tail sloping down to the left --- the same
% 9-shaped sort as the closing mark after "Sandhed". This is the book's body
% convention, and it is NOT what the footnotes on pp. 49--50 do.
betragtede de i Sandhed hverken en blot sandselig Skikkelse eller en abstract
aandelig Idee, men de skuede i ham den guddommelige og menneskelige Natur i
fuldkommen Identitet, saa at hans Historie med fuld Ret kunde kaldes hellig og
guddommelig. Og da han var gaaet bort til Faderen og havde udsendt Aanden,
gjennemtrængte Gudslivet paa en saadan Maade Apostlene og alle de Troende, var
% --- p. 55 ---
\setcounter{footnote}{0}
Evnen til at gjøre Mirakler i saa fuldt Maal tilstede, aabenbarede Himlens Engle
sig saa klart for Alle, kundgjorde de guddommelige Raadslutninger sig saa
tydeligt i Modstandernes (\textit{τῶν θεομάχων}) Rænker, at vi ogsaa med største
% GREEK: script theta normalised to θ per the standing convention. The
% perispomeni of τῶν sits over the ω, where it belongs --- the omicron habit
% logged at p. 51 applies only to the -ου- diphthong.
Ret kalde den apostoliske Historie en hellig Historie. --- Selv efter Apostlenes
Død, da denne Verdens Efterstræbelser forøgedes og Christendommens uforsonlige
Fjende \emph{udtænkte nye Rænker for under det christelige Navns eget Skjul at
vildlede de Letsindige}\footnote{Smlgn. \textit{Cyprian: de unitate ecclesiae,}}%
% LETTERSPACING, settled at 1.1x on the line: "Christendommens uforsonlige
% Fjende" is set tight and "udtænkte nye Rænker" opens up in the middle of the
% same line, so the emphasis begins at "udtænkte", not at the head of the clause.
% The reference mark stands after "Letsindige" and before the comma.
% PRINTER'S DEFECT in the note: it ends with a COMMA --- "de unitate ecclesiae,"
% --- where a period is wanted. Read at 1.8x; the mark has a clear descending
% tail. Transcribed as printed.
% "Smlgn." is Fraktur; "Cyprian" is a heavier antiqua than the title that follows
% it, but both are antiqua against the Fraktur and both take \textit{}.
, syntes Aabenbaringens Kilde ikke at kunne udtømmes; jo mere Menneskene
sammenblandede Sandheden og Løgnen, desto større var Trangen til, i hellige
Personer, i ubedragelige Læresætninger, i udvortes Kriterier paa de sande
Christnes Menighed, kort paa enhver Maade, at have Gud synligt nærværende.
Saaledes uddannede Kirken i denne Verdens Trængsler og Mørke, \emph{underviist af
den hellige Aand, der hver Dag meddeelte den den hele Sandhed}, sine fortrinlige
Institutioner. Præsterne fremtraadte som Guds Repræsentanter, en synlig Styrer
grundfæstede den synlige Eenhed, forat ikke Christi Bruds hulde \emph{Aasyn
skulde faae nogen Plet eller Rynke.} I Sacramenternes mystiske Nimbus seer man
% THE LETTERSPACING BEGINS AT "Aasyn", NOT AT "Christi". Checked line by line at
% 1.1x: "fæstede den synlige Eenhed, forat ikke Christi Bruds hulde" is set tight
% throughout, and the whole of the next line up to "Rynke." is opened out. The
% emphasis therefore starts mid-phrase, and the period belongs inside it.
Christus selv, dagligt voxer Miraklernes Tal, og i det den guddommelige Natur
stedse paany nedstiger i Kjødets Skikkelse, aabner Himlen sig for de Troende, og
alle Hverdagserfaringens Skranker overskrides: man hører Englenes Myriader love
den Almægtiges Majestæt, og omringet af seierrige Martyrers Hærskare tilsmiler
den fromme \emph{Gudsmoder} de bedende Pilegrime, medens Paradisets Indvaanere
juble og Helvedes Uhyrer bæve.
% --- p. 56 ---
\setcounter{footnote}{0}
Saaledes holder Kirken paa eengang Jorden og Himmelen, det Nærværende og det
Tilkommende, det Naturlige og det Overnaturlige i sin Favn, saa at Kirkens
Historie i Ordets fuldeste Forstand er Katholikerne en hellig Historie.

% Head printed "§ 12."; argument agrees verbatim with the Indhold.
\parag{12}{Katholicismens Grundmodsigelse.}

Men denne Katholicismens Herlighed indeholder en Grundmodsigelse. Det er stedse
Væsenets Natur at skabe sit Phænomen; da saaledes den hellige Aand nødvendigt maa
skabe Former og Traditioner, siger man med Rette: ”\emph{Hvor Aanden er, der maa
ogsaa Kirken være}” (\textit{Ubi spiritus, ibi ecclesia}).

Da nu den hellige Aands Institutioner ere ufeilbare, da Traditionens enkelte
Sætninger ere Ord, dicterede af Gud og derfor uforanderlige, er Sandhedens Object
absolut tilstede, saa at den usynlige og synlige Kirke fuldkomment falde sammen.
Vil derfor den menneskelige Aand gribe Sandheden, behøver den kun ganske at kaste
% "Vil derfor den menneskelige Aand" is NOT letterspaced; spacing.py flags it,
% and 1.1x on the line shows ordinary tight setting under a loose justification.
sig i Kirkens Arme; dens hele Kundskab bestaaer i den Erkjendelse, at Kirken har
Sandheden, og Alt, hvad der er sandt, vil den finde udtalt \emph{i den hellige
Moders} Forskrifter og Læresætninger\footnote{Baur: Gegens. d. Kath. und Protest.
\textit{p.\ 501---504.}\\
Sollte jener freiere Begriff der Tradition seine Gültigkeit haben, so dürfte das
eigene Urtheil des Einzelnen nicht schlechthin ausgeschlossen werden, es müsste
jedem überlassen bleiben, die Unterscheidung, die zwischen Form und Inhalt,
zwischen dem Wesentlichen und Unwesentlichen, dem Allgemeinen und Besondern
gemacht werden soll, auch wieder auf alles dasjenige anzuwenden, was die Obern
der Kirche als allgemeine Lehre der Kirche geltend machen. \textit{p.\ 505.}}.
% The emphasis covers "i den hellige Moders" only --- "Forskrifter og
% Læresætninger" is set tight on the same line. Read at 1.1x.
% The note's German block carries NO quotation marks at either end: it simply
% begins on a fresh indented line after the Baur reference. Verified at 1.35x.
% "müsste" is set with long-s + s, not ß. "p. 501---504." and "p. 505." both
% read at 1.35x (the Fraktur OCR gives 509 for the second and the ABBYY layer
% gives 501 for the first; the image gives 504 and 505).
% --- p. 57 ---
\setcounter{footnote}{0}
Men herved glemmer man at foretage den Negation og Restauration af
Phænomenerne, vi ovenfor have omtalt (§ 4); der indfinder sig da en vis kirkelig
Sensualismus, idet den hellige Aands uforanderlige Former paatrykkes det
religieuse Gemyt som en \textit{tabula rasa}; dette bliver ikke et frit Organ for
% NO quotation mark before "dette" --- both OCR witnesses put one there, and 1.2x
% on the line shows only a speck of scan dirt after the semicolon.
Sandheden, men en Form for den hellige Aands Former, betragter derfor
Phænomenerne udvortesfra igjennem en vulgair Empirie og acqviescerer i tomme,
subjective Meninger\footnote{Den (\dsi{} den catholske Kirke) udleder ikke sin
Auctoritet af sin Christelighed, men sin Christelighed af sin Auctoritet.\\
\textit{Dr.} H. N. Clausen, Cathl. og Protest \textit{p.\ 10. 11.}}; saaledes
% The note's second line stands on its own, indented, as an attribution. "Dr." is
% antiqua; "H. N. Clausen, Cathl. og Protest" is Fraktur; the figures are
% antiqua. "Cathl." and the missing point after "Protest" are as printed.
kommer da den Sætning til at gjælde: ”\emph{Hvor Kirken er, der er Aanden}”
(\textit{ubi ecclesia, ibi spiritus}). Skulle nu de kirkelige Institutioner
vedligeholdes, maa den hellige Aand have visse umiddelbare Organer, \dsi{} Kirken
maa have Præster. Da disse hverken nedstige fra Himlen eller fødes som saadanne,
maae de udvælges af Lægfolkenes Midte; derfor er Overgangen fra Lægfolkenes
Klasse til Præstestanden af gjennemgribende Betydning. Da alle Lægfolk ere slagne
med lige Blindhed, saa at der, naar Spørgsmaal opstaae om Forbedringer eller
Rettelser i kirkelige Læresætninger, hverken tages Hensyn til Forskjel i
Sjelsevner eller i Dannelse, kan der i denne Henseende ikke blive Tale om
forskjellige Grader af Oplysning. Før Ordinationen er Aanden kun umiddelbart
tilstede; først ved den lovmæssige Ordination træder den hellige Aand til og
paatrykker den Indviede Ufeilbarhedens Præg (\textit{character indelebilis}). Det
er saaledes vel Aanden, der gjør ham til Præst, men der er intet andet
Kjendemærke paa Aandens Præsents, end Ordinationen, \dsi{} Ingen bliver Præst,
fordi han har Aanden, men
% --- p. 58 ---
\setcounter{footnote}{0}
Enhver, der er Præst, har Aanden\footnote{\textit{Quodsi rasuram, unctionem et
longam tunicam tantum possunt ostendere pro suo sacerdotio, permittimus illis
gloriari. Lutherus de instituendis ministris.} Smlgn. \textit{Dr.} F. C. Baur.
Der Gegensatz des Chatol. und Protest, \textit{P.\ 475.}}. Det er heller ikke
% PRINTER'S DEFECT: "Chatol." for "Cathol." --- note 3 on this same page spells
% it correctly, so the two stand side by side. Transcribed as printed.
% "Protest," carries a comma before the page reference, where notes 1 and 3 on
% p. 59 use a period. As printed. "P." here is Fraktur; the figures are antiqua.
gjennem Præstens Lære, at den ham iboende Aands Sandhed giver sig tilkjende; thi
uagtet hver enkelt Præst er i Besiddelse af Aanden, vedblive de dog Alle at være
svage Dødelige, der kunne synde og fare vild. Man kan da ikke godtgjøre den
enkelte Præsts Ufeilbarhed uden ved Kirkens Vidnesbyrd; Præsten er derfor Kirken
samme Lydighed skyldig, som Lægmanden Præsten. Men gjennem hvis Mund kan Kirken
vel tale, uden igjennem Præsternes? Derfor maa hver enkelt Præst være alle de
Øvriges forenede Mening underkastet: \emph{Majoriteten er den afgjørende.} Skulde
nemlig den Ene overbevise den Anden om Sandheden ved Sandhedens egen Magt, maatte
man jo negere Formens Ufeilbarhed, som saadan, men havde en saadan Negation og
fri Undersøgelse først paa eet Punkt gjort sig gjældende, vilde den trænge
igjennem til alle Kirkens Dele, de mindre med de større, og da vilde
Katholicismens egen Idee forsvinde. Enhver Undersøgelse støtter sig meget mere
paa ydre Traditioner, og ved en saadan træder nu Enhvers umiddelbare Mening til,
saa at det altid er noget Tilfældigt, hvilken Mening der vil faae Pluraliteten
for sig\footnote{Denn ich glaube weder dem Pabst noch den Concilien allein nicht,
weil es am Tag und offenbar ist, daß sie oft geirrt haben, und ihnen selbst sind
widerwärtig geweßt. Luther.}. Det bidrager kun lidet til Bortfjernelsen af denne
% This note too carries NO quotation marks. "geweßt" is as printed --- the sort
% is the same ß that stands in "daß" three words earlier, not the ſt ligature of
% "ist" and "selbst" on the same line (read at 1.1x). Luther's own text has
% "gewest".
Vanskelighed, at Katholicismen har indstiftet forskjellige Grader i
Præstestandens og den kirkelige Jurisdictions Hierarchie\footnote{Smlgn. Clausen
Cathol. og Protest. \textit{P.\ 91} f.}; Muligheden af Sandhedens For\-%
% --- p. 59 ---
\setcounter{footnote}{0}
vanskning forgrenes kun herved endmere. Hvad hjælper det nemlig, at Presbyternes
Meninger rette sig efter Biskoppernes, dersom Biskopperne modsige hverandre
indbyrdes. Det bliver da nødvendigt at have en synlig Styrer, til hvem man kan
tye hen for at høre den hellige Aands levende Røst. Vi ere herved komme til de
bekjendte Stridigheder mellem Episcopalister og Curialister\footnote{Samme: Cath.
og Prot. Det curialistiske System \textit{P.\ 14---26.} Det episkopale
Kirkesystem \textit{P.\ 43---52.}}. Men at Episcopalismen er det System, der
stræber at holde den slette Middelvei i den katholske Kirke, er bekjendt.

Katholicismens Herlighed, som dens Grundmodsigelse, culminerer i Pavens Person;
paa den ene Side er Paven et svagt og syndigt Menneske, paa den anden Side er han
Guds Repræsentant, i hvem den levende Aands Fylde tager Bolig; men Hvo skal nu
afgjøre, naar han taler som Pave, dreven af den hellige Aand, og naar han handler
efter egen subjectiv Vilkaarlighed? Concilierne, som skulde have det aandelige
Livs hele Fylde i sig, ere jo underlagte Paven. Derfor har ogsaa den
protestantiske Videnskab forlængesiden beviist\footnote{Clausen \textit{l.\ l.\
P.\ 58.}\\
Curialismen danner i Kirken et theocratisk Monarchie, saa uindskrænket, at det i
Virkeligheden maa gaae over til Despotie, thi det bekjæmper ikke alene ydre
Forhold og Forbindelser, men ogsaa ethvert enkelt Menneskes aandelige
Personlighed, enhver Yttring af selvstændig Tænkning og Villie; og til Gjengjeld
fratager Kirken sine i Troen lydige Borgere al Samvittigheds Ansvar, og skjænker
dem uforstyrret Tryghed og Ro; men det er Dødens Ro paa slagne Fienders Grav.},
% Note 1 spells "curialistiske" with c and "episkopale" with k, in the same
% sentence; as printed. Note 2's Danish quotation from Clausen carries no
% quotation marks either --- the pattern holds for every prose note in this
% range except p. 60's Latin. "P. 43---52." read at 1.05x (the Fraktur OCR gives
% 13---52).
at Katholicismens objective Idee --- netop der, hvor den har naaet
Objectivitetens høieste Spidse --- er gaaet over til og forsvundet i en
vilkaarlig Subjectivitet, hvis hele Ufeilbarhed bestaaer deri, at den erklærer
sig for ufeilbar. Derfor kunne
% --- p. 60 ---
\setcounter{footnote}{0}
alle katholske Institutioner, uagtet de ere af en saadan Beskaffenhed, at de
kunne være Former for den hellige Aand, tillige, naar de tilegnes paa udvortes
Maade, fastholdes i og for sig.

Kirken fremstiller saaledes ofte tomme Former, og det Falske i Katholicismen
ligger netop deri, at den, der er i Besiddelse af Formerne, ogsaa antages at
besidde Aanden. Under et saadant Forhold er det ikke at undre over, at de ydre
Phænomener ofte i den Grad have været forladte af Aanden og Sandheden, at ikke
engang Ærbarhedens Skin er bleven bevaret, og at Katholicismen ofte har maattet
rødme over Historiens Vidnesbyrd, skjøndt det dog ere disse, hvortil den har
slaaet sin Lid\footnote{\textit{Baronius} ”\textit{significavit suo Jac.
Sirmondo, nihil se æque reformidasse, quam pervenire ad hoc tempus (\dsi{} magni
Schismatis Occidentis) de quo, quid statuendum, non esset libere
pronuntiaturus}” \textit{l.\ l.\ p.\ 49.}}.
% *** THIS NOTE OPENS WITH THE RAISED ”, NOT THE LOW „ OF THE NOTES ON
% pp. 49--50. *** Read at 1.15x: the mark before "significavit" stands at cap
% height, matching its partner after "pronuntiaturus". So the low opening is not
% a rule for the notes of this book but a variation within them, and both are
% reproduced as printed. This note is balanced: ”=2.
% The Danish id-est sign stands INSIDE the Latin, "(\dsi{} magni Schismatis
% Occidentis)" --- Bøggild's gloss on "hoc tempus". Whole note is antiqua.

% *** THE PRINTED HEAD AND THE INDHOLD DISAGREE HERE TOO --- the third case, and
% the second in this batch. ***
% The printed § head on p. 60 reads "Overgangen fra Katholicismen til
%   Protestantismen." (read at 0.38x; the definite form, with -en).
% The Indhold (PDF 8) reads  "Overgang fra Katholicismen til
%   Protestantismen. ... 60."
% The printed head is set; the Indhold reading is recorded, not normalised.
% The head is printed "§ 13." --- no point after the section sign.
\parag{13}{Overgangen fra Katholicismen til Protestantismen.}

I Forsøget paa at fremstille Sandhedens absolute Object for Betragteren,
bestræber Katholicismen sig stedse for at bringe den menneskelige Selvbevidsthed
til at bortvende Blikket fra sig selv til ydre Gjenstande. Men heraf følger, at
Gjenstandene i Guds Rige ved at fixeres i concrete Billeder blive de jordiske
Gjenstande sideordnede, saa at det Hellige profaneres, medens paa den anden Side
de himmelske Væseners Nedstigen i jordiske Skikkelser medfører, at timelige
Gjenstande tillægges guddommelig Hellighed, fordi de indeslutte det Guddommelige
i sig, saa at det i sig Vanhellige bliver Gjenstand for Til-

% --- p. 61 ---
% ELISIONS. Throughout pp. 61--72 the notes set elisions as three, four or six
% spaced points. They are uniformly rendered \ldots here; the printed count is
% not reproduced.
\setcounter{footnote}{0}
bedelse\footnote{Samlgn. Hegel, Geschichte der Philosophie. \textit{3}ter
Band, P.~\textit{253} f.}. Da denne barbariske Confusion af modsatte
Elementer havde lagt sig klart for Dagen i Bestyrelsens Theokratie, i
Videnskabens Scholasticisme, i den Betydning, der tillagdes de gode
Gjerninger, begyndte Selvbevidstheden, der alt længe havde været forviist fra
sit eget Indre, ad forskjellige Veie at vende tilbage til sig
selv\footnote{Hegel \textit{l.~l.} P.~\textit{267}. Die andere Seite ist, daß
das Ewige, was an und für sich wahr ist, auch erkannt, aufgefasst werde durch
das reine Herz selbst; der eigene Geist mach sich für sich das Ewige zu
eigen. Das ist der lutherische Glaube ohne anderes Beiwesen (die Werke, wie
man es nannte).}.
% PRINTER'S DEFECT in note 2, transcribed as printed: "mach" for "macht"
% (Hegel, Gesch. d. Philosophie: "der eigene Geist macht sich für sich das
% Ewige zu eigen"). Read at 6x --- the word ends in "ch", with no t.
% "l. l." (loco laudato) is ANTIQUA; the "P." beside it is FRAKTUR. Both
% settled at 6x on the same line, so only "l. l." takes \textit{}.
% "aufgefasst" is set with double s + t, not ß --- checked at 6x.
Undersøger man nu, hvorledes det var muligt, at en saa stor Spaltning kunde
opstaae i Kirken, vil man finde, at det eneste Middel til grundigt at
reformere Kirken, var at negere Phænomenerne for at trænge igjennem til
Væsenet, at drage en skarp Adskillelse mellem Ideen og Formen, mellem den
usynlige og synlige Kirke, og som en Følge deraf at træde i aaben Kamp mod
Hierarchiet. Thi hvor Bevidstheden ikke kan Andet, end blive Træl under en
fremmed Auctoritet, kan der aldrig fremkomme nogen sand Eenhed af Subject og
Object; man seer derfor i en saadan Positivisme en bestandig Vaklen mellem
Lydighed og Oprør. Lydigheden er blind, saalænge Subjectet kan bindes til at
søge Substantsen udenfor sig; der bliver Oprør, naar denne Substants ikke
sees at udgjøre Subjectets egen Natur, men som en fremmed Magt ikkun indjager
Frygt; derfor er Katholikernes \emph{Tro} et \emph{Trældoms}forhold og
indeholder altid en vis Mistro.
% The letterspacing covers "Tro" and "Trældoms" only --- "forhold", set solid,
% closes the compound. Read at 3x.

Skal det Sande skues i Sandhedens Lys, maa den objective Fornuft
nødvendigviis gjennemtrænge den subjective; vi kunne ellers ikke være ligesaa
visse og faste i Besiddelsen af Tingen selv, som vi ere det i Besiddelsen af
vor
% --- p. 62 ---
\setcounter{footnote}{0}
egen Selvbevidsthed. Først hvor dette Forhold er umiddelbart tilstede, er der
en sand Tro. --- Det er ved en saadan Tro saa langt fra, at det subjective
Liv hæmmes af de objective Baand, at Subjectet meget mere med fuld Tillid
overgiver sig til Objectet, i Erkjendelsen af, at det netop er af det
Objective, det drager sin sande Næring. Dette er den protestantiske Tro, der
bryder Positivismens Lænker og begrunder en virkelig Frihed. Hvilket skuffende
Phænomen skulde kunne vildlede Bevidstheden, hvilken udvortes Magt kunne
ryste Sjelens Styrke, naar man har fundet dette Frihedens Palladium? Derfor
stod den Guds Mand, hvem vi med taknemmelige Hjerter prise som
Protestantismens Fyrste, urokkelig fast; omgjordet med Sandhedens Belte
spredte han Traditionernes og Fablernes Taage; beskyttet af Troens Skjold
udslukkede han det kirkelige Tyrannies brændende Pile\footnote{Es thut dir
wehe im Herzen mein fröhlicher großer Muth. Ich bin aber, und will, ob Gott
will, auch bleiben, gegen dir und Ecken, Pabst und eurem Hauffen, auch dem
Teufel, mit Gottes Hülfe, in einem beständigen, hochmüthigen, unerschrocknem
Geist, und euch trotzen und verachten \ldots\ Luthers Schutz=Schriften v.
An.~\textit{1524} wied. Hier. Emsern u. Paris. Theologos etc.}: med Aandens
dragne Sværd, Guds Ord, angreb han Overtroens vanhellige
Afguder\footnote{Hörst du es, Pabst, nicht der Allerheiligste, sondern der
Allersündigste, daß Gott deinen Stuhl vom Himmel auf's schierste zerstöre und
in Abgrund der Hölle senke. Smlgn. Pfizer L.~L. P.~\textit{161}.}. Fra ham
udgik den første Krisis, --- den falske Skjønhed svandt bort, og de
menneskelige Vildfarelsers Historie blev fremdraget for Lyset.
% "L. L." here is FRAKTUR, unlike the antiqua "l. l." on p. 61 --- compared
% with the Fraktur H of "Hölle" on the same line at 6x. So no \textit{}.
% "Schutz=Schriften" and "Allersündigste" are both broken over a line with the
% Fraktur double hyphen; only the first is a real compound rule in the setting.
% --- p. 63 ---
\setcounter{footnote}{0}
% The printed § head agrees with the Indhold here: "Troen og den hellige
% Skrift." Verified on the image at 3x.
\parag{14}{Troen og den hellige Skrift.}

Og dog ophævede denne første Krisis ikke Baandet mellem det naturlige og det
overnaturlige Rige; vel ombyttede Troen de ydre Elementer med Aandens indre
Momenter, saa at den brogede Billedpragt forsvandt; men ikkedestomindre
vedbleve endnu de gode Engle med deres hulde Aasyn at skue ned til
Menneskene gjennem Himmelens Skyer\footnote{Luther: Daß man also Cherubim
verstehe für Engel, die da erscheinen, nicht mit einem runtzlichten, noch
traurigen Angesichte, sondern mit einem fröhlichen Geberde, und vollem
feistem Angesichte \ldots\ Ausleg. des \textit{3}ten Capit. des \textit{1}sten
Buch Mose.}, vedblev Tjævelen, \emph{dette Vrængebillede af
Gud}\footnote{”\textit{Simia dei}” en Allusion til Luthers Ord om Djævelen:
”der Affe Gottes.”\hfill Overs.\ Anm.}, under utallige Forvandlinger at friste
Menneskenes Børn\footnote{Luther: Er (der Teufel) ficht mich selbst offtmahls
so gewaltig an, und überfället mich so hefftig mit schweren und traurigen
Gedanken, daß ich meines Herrn Christi gar vergesse \ldots\ ist er dennoch
heimlich bei uns, schleicht Tag und Nacht um uns her \ldots\ Ausleg. d.
Epistel a.~d. Galater Cap.~\textit{3}.}.
% *** PRINTER'S DEFECT, transcribed as printed: "Tjævelen" with a Fraktur T.
% Settled at 8x by direct comparison of four initials on this one page: the
% body's "Tjævelen" matches the T of "Troen" (body l. 3) exactly, while note 2
% of the same page spells "Djævelen" with the D of "Den katholske" (body
% l. 15). So the compositor set T for D in the body only. Both OCR witnesses
% also read T. NOT corrected --- it is not in the Rettelser list.
% "runtzlichten" (note 1) is the tz ligature, not ß: the glyph carries a
% descender, which the ß of "Daß" on the line above does not. Read at 6x.
% QUOTATION MARKS in note 2: RAISED (U+201D) at BOTH ends of both quotations,
% as everywhere else in this book. Verified at 8x. Four marks on this page.
% "Simia dei" is antiqua; "der Affe Gottes." is Fraktur.
Heller ikke gik den Frihed, hvormed man negerede Phænomenernes Sandhed,
saavidt, at der ikke skulde være bleven nogen ufeilbar Skikkelse af Sandheden
tilbage; --- jo mere den bestaaende Kirkes Fordærvelse var til Forargelse,
desto hellere søgte man hen til dens Forbillede: jo dristigere man var i
Forkastelsen af Traditionernes Guddommelighed, desto fastere sluttede man sig
til Evangeliet i dets oprindelige Skikkelse. Den katholske Kirke havde paa
særegen Maade ophøiet den apostoliske Tidsalder og anerkjendt dens Skrifter
som kanoniske; da nu den religieuse Bevidsthed i disse Bøger havde anskuet
Sandheden i uforfalsket Skikkelse, greb man med Begjerlighed \emph{den hellige
Skrift} som Troens sande Grundpille.
% --- p. 64 ---
\setcounter{footnote}{0}

Dog er dette ikke at forstaae saaledes, som om man ved at urgere den hellige
Skrift kun vilde paa positivistisk Maade opstille en forstenet Norm, saa at
Katholicismen, der alt længe havde gjort sig til af sin Læres
Overeensstemmelse med den apostoliske, endnu længe uforstyrret kunde have
fortsat denne Selvroes, hvis ikke nogle Theologer tilfældigviis havde stødt
paa den hellige Skrift, læst den ret omhyggeligt og anstillet en
Sammenligning. Dette var umuligt; thi havde de ikke allerede næret stor Tvivl
om mange af Katholicismens Sætninger, vilde de jo have udviklet Skriften
efter Kirkens angivne Rettesnor. Sandheden er overalt, og til enhver Tid
tilstede; derfor aabenbarede Aanden, der aldrig forlader sin Kirke, sig
strax, saasnart man havde bortkastet det falske Slør, der havde tilhyllet
den. --- Man maa derfor vel erindre, at naar Protestanternes Forkjæmpere
ideligt appellerede til Skriften som Sandhedens høieste Vidnesbyrd, kunde de
ene og alene gjøre dette, fordi Sandheden selv var Skriftens Vidne, saa at de
ærede Aanden i samme og ikke dens Bogstaver. --- Da Reformatorerne havde
bemærket Traditionernes Uovereensstemmelse med Skriften\footnote{Was Christum
nicht lehret, das ist noch nicht apostolisch, wenn es gleich S.~Petrus oder
S.~Paulus lehrete. Wiederum was Christum prediget, das wäre apostolisch, wenns
gleich Judas, Hannas, Pilatus und Herodes thät. Luther: Vorrede auf die
Epistel St.~Jacobi. Smlgn. hermed: Ein Stück der Vorrede auf das
N.~Testament: Welches die rechten und edelsten Bücher des Neuen Testaments
sind. An.~\textit{1521}.} og fundet, at de mere havde forvansket, end udviklet
den christelige Idee, laa det i Sagens Natur, at de paa eengang maatte
tillægge Skriften \emph{Tydelighed} og \emph{Fuldstændighed}
(\textit{perspicuitas et sufficientia})\footnote{Da erdenken sie eine neue
Lüge, finden Degen und Spies und dergleichen Narrenwerk, und sprechen: Die
Schrift sey so finster, daß wir sie nicht mögen verstehen ohne der heiligen
Väter Auslegung \ldots\ Wenn sie denn nun einen Spruch der Väter wieder mich
aufbringen, so läuten sie alle Glocken, schlagen alle Drummeln, und schreien
feindlich, sie haben gewonnen. Luther: Von dem bleyern Degen Bocks Emsers.

Og jo mere Alt blev henført til Troens Substantialitet og ikke til
Gjerningernes Tilfældighed, desto lettere var det at vise Skriftens
Fuldstændighed. Derfor sagde han, da Katholikerne henviste til Joh.
\textit{21, 25}: Es ist nicht alles geschrieben, was wir thun sollen. Danck
habt ihr guten Gesellen, ihr wisset der Schrifft Auslegung wohl zu geben.
S.~S.}; men man maa vel mærke, at de, opdragne som de
% THE PRINTING SPLITS THIS NOTE. Note 2 begins in the footnote block of p. 64
% and runs on through the whole upper half of the footnote block of p. 65,
% where its second paragraph ("Og jo mere Alt...") is set indented and in
% Danish. It is given entire here, at its mark, and p. 65's own note block
% below therefore begins at the printed "1)".
% --- p. 65 ---
\setcounter{footnote}{0}
vare, i den katholske Kirke og den af Hjertet hengivne, vare saa
gjennemtrængte af Kirkelæren i dens Heelhed, at de, hvad ogsaa deres
Fortolkning tilstrækkeligt beviser, ikke saameget øste Troens Grunddogmer af
Skriften, som erkjendte dem nedlagte i den; og uagtet de i Modsætning til de
Allegorier, Katholikerne anvendte for saaledes at kunne udtyde ethvert Sted
til deres Fordeel, indskjærpede den bogstavelige Betydning og en
Parallelisering af Skriftstederne, vilde de dog ikke destomindre have Skriften
forklaret i Totalitetens Lys, saa at det gamle Testament ikke kunde udvikles
uden ved Hjælp af det ny Testament\footnote{Die heilige Schrifft dringet
vielmehr auf den Sohn, als auf den Vater: denn die ganze Schrifft ist um des
Sohnes willen geschrieben: darum sind auch im alten Testament mehr Sprüche
oder Zeugnisse vom Sohne als vom Vater. --- Luth. Schrifft. \textit{III.~4.}
Denn da steckt es, da bleibt es, wer diesen Mann, der da heist Jesus Christus
Gottes Sohn, den wir Christen predigen, nicht recht und rein hat, noch haben
will, der lasse die Bibel zu frieden, das rathe ich, er stöst sich gewisslich
und wird, je mehr er studiret, je blinder und toller.

Wenn es nun sollte Wünschens und Wählens gelten, entweder, daß ich
S.~Augustini und der lieben Väter, das ist der Apostel Verstand in der
Schrifft sollte haben, mit dem Mangel, daß S.~Augustinus zuweilen nicht die
rechten Buchstaben oder Worte im Ebräischen hat, wie die Jüden spotten
\ldots\ ich liesse die Jüden mit ihrem Verstande und Buchstaben zum Teufel
fahren, und führe mit S.~Augustini Verstande ohne ihre Buchstaben zum Himmel.
Auslegung der letzten Worte Davids \textit{2}~Sam. \textit{23, 1---7}. Schr.
\textit{IV. 303}.}. Og ikke var dette at undre over, thi trods Ti-
% "øste" (from "øse", to draw or scoop) --- read at 4x: the initial is ø, not
% e, and the second letter is a long s, not f. Both OCR witnesses fail here
% ("este", "oste").
% "stöst" and "gewisslich" are set with s, not ß --- checked at 4x.
% THIS NOTE ALSO RUNS ON: note 1) begins in p. 65's footnote block and its
% closing citation stands at the head of p. 66's footnote block. Given entire
% here; p. 66's block below therefore begins at the printed "1)".
% --- p. 66 ---
\setcounter{footnote}{0}
dernes Forskjelligheder havde den hellige Aand været den samme, og den kunde i
Sandhed have havt Christus i Tanken, skjøndt den henvendte sine Ord til
Jøderne\footnote{Saaledes kan Luther selv ikke afholde sig fra, tvertimod sit
eget Princip, undertiden at tilføie en allegorisk Betydning, f.~Ex. Genes
\textit{c.~9} om Noah: Darum willst du es ohne Gefahr deuten, so deute es auf
den Herrn Christum. Denn wie \emph{Noah den Weinberg pflanzet und des Weins
trincket,} davon \emph{truncken wird} und einschläfft, \emph{und bloß in der
Hütten lieget,} und wird von diesem verlachtet, aber von andern
\emph{zugedecket:} also ist es auch Christo ergangen.}.

Her er det altsaa, at den Negation og Restauration af Phænomenerne, vi
ovenfor omtalte, tager sin Begyndelse. Thi naar Luther foretager sig at læse
Skriften og at bedømme Traditionernes større eller mindre Sandhed, er
Sandhedens Aand tilstede i hans Bevidsthed, saa at han ikke blot dømmer efter
et subjectivt Skjøn. Men hvorledes fik han denne Aand? Ikke ved nogen særegen
Inspiration, som \emph{Enthusiasterne}\footnote{Saaledes skriver han til
Amsdorff om de nye Propheter: Durch die Zwickauer Propheten lasst Euch nicht
so schnell bewegen! Ihr habt ja Zeugnisse der Schrifft, die Euch sicher
machen, daß Ihr nicht sündiget, so Ihr sie hinausschiebet und zuerst die
Geister prüfet, ob sie aus Gott sind \ldots\ Mir wahrlich ist das dem ersten
Anblick nach gar verdächtig, daß sie sich der Gespräche mit der göttlichen
Majestät rühmen..}; men gjennem det Totalindtryk, han modtog under det
vedholdende Studium af Traditionerne og Skriften. Der var saaledes foregaaet
en sand Forsoning af det objective og subjective Moment; at han gjorde sin
subjective Frihed gjældende, sees deraf, at ingen udenfra given Mening
fremkaldte en blind Underkastelse hos ham; men den objective Realitet var
ligesaavel bevaret, thi han erkjendte, at man kun gjennem Til-
% LETTERSPACING IN NOTE 1, read off the image at 4x. The compositor spaced the
% Noah allegory in four separate runs, leaving three stretches inside it solid:
% "Noah den Weinberg pflanzet und des Weins trincket," / "truncken wird" /
% "und bloß in der Hütten lieget," / "zugedecket:" are spaced, and "davon",
% "und einschläfft," and "und wird von diesem verlachtet, aber von andern" are
% not. Measured, not eyeballed: intra-letter gaps 21--34 px in the spaced runs
% against 5--13 px in the solid ones. Reproduced exactly.
% The footnote block of p. 66 opens with the last line and a half of p. 65's
% note 1), which is set there, at its mark; the printed "1)" of this page is
% the Noah note above.
% "rühmen.." is transcribed with TWO points, as printed: at 8x both marks sit
% on the baseline and are the same size. Probably a foul sort, not an ellipsis.
% --- p. 67 ---
\setcounter{footnote}{0}
egnelsen af de objective Phænomener kan komme til Sandheden\footnote{So Gott
sein Heiliges Evangelium hat aus lassen gehen, handelt er mit uns auf
zweierley Weise. Ein äusserlich; das andere mahl innerlich. Aüsserlich handelt
er mit uns durch mündliche Wort des Evangelii und durch die leiblichen
Zeichen, als da ist, Tauffe und Sacrament. Innerlich handelt er mit uns, durch
den H.~Geist und Glauben \ldots\ also, daß er's beschlossen hat, keinem
Menschen die innerlichen Stücke zu geben, ohne durch die äusserlichen Stücke.
Wieder die himmlischen Propheten. Schrifft. \textit{XIX., 186}.}. Paa dette
Standpunkt er da ogsaa en sand Overgang mulig fra Aandens rene Almindelighed,
til dens enkelte concrete Bestemmelser, thi naar Sandhedens almindelige Idee
ved sit Lys har negeret Traditionerne og Skriftstederne i deres Enkelthed, vil
den altid paany stræbe at restaurere de negerede Former, og naar Luther
urgerede, at den hellige Skrift skulde være Norm for Udviklingen af Fædrenes
Meninger, Troen paa Christus en nødvendig Betingelse for Udviklingen af de
enkelte Skriftsteder, var det ikke Christusideen løsrevet fra dennes enkelte
Bestemmelser, han urgerede; denne omhyggelige Undersøger af Skriftens
allermindste Yttringer\footnote{Ach es ist Dollmetschen ja nicht eines
Jeglichen Kunst, wie die tollen Heiligen meinen; es gehöret dazu eine recht
fromm, treu, fleissig, furchtsam, christlich, gelehrtes, erfahren, geübt Herz.
Samlgn Pfizer S.~S. P.~\textit{843---852}.}, gjorde aldrig Christus til
Gjenstand for sin Tanke uden tillige at see hen til alle de enkelte
Skriftsteder, der skildre Christus for os\footnote{Also weiset er uns
allenthalben in die Schrifft, warum das? Darum daß du kannst bewahren den
Christlichen Glauben. Es ist mir wohl selbst begegnet, wenn ich das Wort habe
fahren lassen, daß ich Gott, Christum und alle mit einander verlohren habe.
\emph{Es ist auch kein leichterer Weg alle Artickel des Glauben zu verlieren,
denn ausser der Schrifft daran zu gedencken.} \textit{III.} Predigt über das
Evangelium am Ostermont. Luc. \textit{24, 13, 25}, zu Coburg geh.
\textit{1530}. Schrifft. \textit{XV. 631}.}.
% "Aüsserlich" is transcribed as printed: at 5x the two dots stand over the u,
% not over the A. (The first occurrence, three words earlier, is the ordinary
% "äusserlich".)
% "Luc. 24, 13, 25," --- read at 6x. The reference is presumably to the Emmaus
% pericope of Easter Monday, Luc. 24, 13--35, but the third figure on the page
% is a 5, and there is a comma, not a dash, after the 13. As printed.
% Note 3's letterspaced run covers "Es ist auch ... zu gedencken." exactly.
% THIS NOTE RUNS ON: note 3) begins in p. 67's block and its last line and a
% half stand at the head of p. 68's block. Given entire here.
% --- p. 68 ---
\setcounter{footnote}{0}

I det Momenterne saaledes flyde uendeligt over i hinanden, see vi den
objective Aand være reelt tilstede i Luther, see den gjøre hans Aand til et
Organ for sig, i det han som en \textit{ἄνθρωπος πνευματικὸς} prøver Alt
\emph{igjennem Troen}\footnote{Ich will mich hören lassen, und wie St.~Petrus
lehret, meiner Lehre Ursach und Grund beweisen für alle Welt, und sie
ungerichtet haben von jedermann, auch von allen Engeln. Denn sintemal ich ihr
gewiß bin, will ich durch sie, euer und auch der Engel (wie St.~Pauius
spricht) Richter seyn, daß wer meine Lehre nicht annimmt, daß der nicht möge
selig werden.}.
% GREEK, read at 5x. Fount variants normalised per the standing convention:
% script theta -> θ (ἄνθρωπος), cursive kappa -> κ (πνευματικὸς). The
% breathing and both accents are reproduced as printed.
% PRINTER'S DEFECT in the note, as printed: "St. Pauius" for "St. Paulus"
% --- at 6x the letters are u-i-u-s. Both OCR witnesses read it the same way.
% The letterspacing covers "igjennem Troen" only; "prøver Alt" before it is
% set solid (measured at 5x).
Det laa derfor i Sagens Natur, at Forkyndere af Ordet, der optraadte, drevne
af den urokkelige Overbeviisning om dettes Sandhed, maatte erklære deres
Modstandere for enten at være bedragne eller Bedragere, om de end derved
syntes at staae i aabenbar Modsigelse til sig selv, da de selv vare svage
Dødelige og netop prædikede imod Ufeilbarhed hos nogen Dødelig. Denne
Modsigelse ophævedes imidlertid derved, at de skarpt skjelnede imellem Læren,
forsaavidt den var af dem selv, og den objectivt begrundede Sandhed, der
havde sit umiddelbare og oprindelige Udtryk i den hellige
Skrift\footnote{Nicht also du Narr, höre und laß dir sagen. Zum Ersten bitte
ich, man wollt meines Namens schweigen und sich nicht Lutherisch sondern
Christen heissen. Was ist Luther? Ist doch die Lehre nicht mein. Luthers tr.
Ermahn. a.~all. Christen Schrifft. \textit{XVIII.} P.~\textit{289}.}. Af
Hensyn til Tilhørernes Tarv maatte de gjøre denne Adskillelse, der var
ligesaalangt fra at svække deres egen Overbeviisnings Kraft, som den var
langtfra at være fremgaaet af en hyklet Beskedenhed; thi selv om deres Ord
vare de fuldkomneste Udtryk af Sandheden, kunde Tilhørerne dog ikke gjøre
det, de havde hørt, til deres sande Eiendom, uden ved selv at gjøre det til
Gjenstand for en Undersøgelse; heller ikke rokkedes herved Talerens
Overbeviisning
% --- p. 69 ---
\setcounter{footnote}{0}
om sin Læres Sandhed, da han jo netop haabede ved en grundig og aaben
Undersøgelse af Sagen ogsaa at kunne overbevise Andre om Sandheden af det,
hvortil han selv ad Overbeviisningens Vei var kommen, og at kunne tilbageføre
den Meningsforskjel, der muligt kunde blive tilovers, til en blot subjectiv
Ejendommelighed hos de stridende Parter; man søgte altsaa ikke Enigheden
deri, at Alle skulde have de selvsamme Meninger, men deri, at Sandhedens
Fylde skulde tilfredsstille Alles religieuse Trang. Sandheden er altsaa kun
for den prøvende Bevidsthed\footnote{Für das erste (neml. werden die Geister
geprüfet) durch ein innerlich Urtheil, da ein jeder Christ durch den Heiligen
Geist und Gottes Gnade, für sich und sein Gewissen, also erleuchtet ist, daß
er aufs allergewisseste schliessen und urtheilen kann von allen Lehren.

Und diese Gewißheit gehört zum Glauben, und ist von nöthen einem jeden
Christen, ob er gleich nicht ein Prediger oder in öffentlichem Amte ist. Luth.
w. Crasm. Schr.: \textit{XIX}.}, nøies ikke med døde Sætninger, men attraaer
levende Subjecter; den fremtræder derfor i særegen Skikkelse i hver Enkelts
Bevidsthed; der ere altsaa ligesaamange Sandheder, som der ere
Subjectiviteter. Men for nu ikke selv at forsvinde ved subjectiv
Particularisme i Adsplittelsen af sine Elementer, samler den sine Bekjendere
og opfordrer dem til gjensidig Meddelelse og Discussion; saaledes fremkaldes
Religionssamtalen. Derfor kan Reformatorernes Tidsalder snarere kaldes
Colloqviernes, end Conciliernes Tidsalder. Concilier egne sig for Hierarcher,
der, skjøndt de prætendere en præstelig Ufeilbarhed, behandle Sandheden paa
Lægmands Viis, ved paa udvortes Maade at sammenligne Meninger og optælle
Stemmer; Colloqvier derimod egne sig for Lægfolk, der aflægge Ufeilbarhedens
laante Nimbus og i den objective Sandheds Navn sætte Tro ligeoverfor Tro,
Subjectivitet ligeoverfor Subjectivitet, for gjennem denne Ildprøve at komme
% PRINTER'S DEFECT in the note, as printed: "Crasm." for "Erasm." (Luther
% wider Erasmus). Settled at 8x: the initial is the C of "Colloqviernes" and
% "Conciliernes" on this same page, and is plainly not the Fraktur E of
% "Enigheden" seven lines above it. NOT in the Rettelser list; not corrected.
% --- p. 70 ---
\setcounter{footnote}{0}
til Vished om, hvo der ere de sande Præster\footnote{Det var derfor
naturligt, at Paven ikke hastede med for Protestanternes Skyld at sammenkalde
et Concilium; thi han vilde befale, disse derimod disputere og „føre deres Sag
for et almindeligt, frit og christeligt Concilium.“}. Men i saadanne
Religionssamtaler ere tvende Momenter at iagttage; paa den ene Side har
nemlig Enhver af de Disputerende sin egen aandelige Fylde, til hvilken hans
Bevidsthed \emph{indadtil} paa universel Maade er bunden, paa den anden Side
vender han sig \emph{udadtil}, i det han vil beskrive og forsvare denne sin
Subjectivitet; Enhver bestræber sig derfor nærmere at bestemme sin
Subjectivitet ved enkelte Meninger og at befæste den ved enkelte Gruude, for
saaledes gjensidigt at rette hinanden og komme til den fuldkomne Sandhed. ---
Men skulle alle Vanskeligheder ad denne Vei fjernes, maa for det Første den
mystiske og dunkle Side, der pleier at findes i ethvert Menneskes Indre,
træde klart frem for Bevidstheden; thi først naar man er sig selv fuldkomment
klar, kan man give Andre en fuldkommen Skildring af sin Subjectivitet, ---
men dette vil neppe kunne skee, før Salighedens Rige er kommet ---; dernæst
maa den subjective Sandheds Udtalelse ikke være en blot subjectiv
Bekjendelse, men maa føres tilbage til en indre Nødvendighed, saaledes som
den sande Methode fordrer, --- men Methoden var fremmed for denne Tidsalders
Aand.
% *** QUOTATION MARKS: THIS PAGE DEPARTS FROM THE BOOK'S STANDING FORM. ***
% Both notes on p. 70 set the DANISH LOW--HIGH pair „…“ --- the opening mark on
% the baseline, the closing mark raised --- and not the raised-at-both-ends
% mark ” that pp. 1--36 and p. 63's note 2 use without exception. Read at 8x
% in both notes. Transcribed as printed: this page contributes 2 „ and 2 “,
% and no ”, to check.py's counts.
% *** PRINTER'S DEFECT, transcribed as printed: "Gruude" for "Grunde", a
% turned n. Settled at 8x by comparing the two medial letters with the n's of
% "hinanden" nine words later on the same line: the n's close at the TOP, and
% both letters of "Gruude" close at the BOTTOM, like the u of "universel" and
% "bunden" four lines above. Same defect as p. 2 "absolnt", p. 3 "Tiugen" and
% the § 3 head's "Deu". NOT in the Rettelser list; not corrected.

Hine Prædikener og Religionssamtaler gik mere ud paa at blive sig det
gjensidige Aandsslægtskab bevidst, end at komme til sand
Begrebsklarhed\footnote{Derfor bortstødte Luther paa Conciliet i Marburg den
Haand, som Zvingli med Taarer i Øinene rakte ham, efterat de, uden at blive
enige, havde afhandlet Alt ligefra Legemets metaphysiske Egenskaber til
Skriftens ubetydeligste Partikler, med de Ord: „Ihr habt einen andern Geist,
denn wir!“}. Det var Følelsen af dette Aandsslægtskab, vakt ved Mesterens
stærke Røst, der grund-
% The mark of note 2 on this page is heavily inked and its outline is broken;
% at 12x it is a figure 2, and the body's reference mark is a clean 2). Set as
% the plain next number.
% --- p. 71 ---
\setcounter{footnote}{0}
lagde den Kirke, hvortil vi bekjende os; og saa hurtigt trængte denne Følelse
igjennem, at der, uagtet den subjective Luther, om vi saa maae sige, forlod
Conciliet i Worms som en Fredløs, neppe vare \textit{10} Aar forløbne, før
den objective Luther vendte tilbage til Conciliet i Augsburg, styrket ved
Reformationens rige Velsignelse. Og ligesom Luther stedse appellerede til Guds
Ords objective Vidnesbyrd, skjøndt han holdt urokkeligt fast ved sin Tro,
saaledes anmasser heller ikke den augsburgske Troesbekjendelse sig nogen
Auctoritet ved Siden af eller over Skriften, men udleder tvertimod fra den
hele sin Gyldighed\footnote{\textit{Offerimus in hac religionis causa
nostrorum concionatorum et nostram confessionem, cujusmodi doctrinam ex
scripturis sanctis et puro verbo Dei hactenus illi in nostris terris
\ldots\ et in ec clesiis tractaverint. Præf. ad Cæs.}}.
% PRINTER'S DEFECT in the note, as printed: "in ec clesiis" --- "ecclesiis"
% broken by a space in mid-word, with no hyphen and no line break at the
% point. Read at 5x. The whole note is antiqua Latin, hence \textit{}.

% *** ERRATA APPLIED HERE, AND THE PRINTED HEAD DIFFERS FROM THE INDHOLD. ***
% The printed head reads "Symbolerne. Luthersk Scholastisme." --- verified on
% the image at 4x. The Rettelser page corrects it (p. 71, l. 12 fra oven:
% Scholastisme l. Scholasticisme), and the Indhold already reads
% "Scholasticisme", so the correction is set. Both readings are recorded here.
% The Rettelser entry sets both words letterspaced because the word stands in
% the letterspaced § argument; \parag already emphasises the whole argument,
% so no further \emph{} is needed. (The "l. 12" of the entry counts the head
% line itself as 12; the "§ 15." line above it is line 12 by a plain count.)
\parag{15}{Symbolerne. Luthersk Scholasticisme.}
% ERRATA: p. 71 l. 12 fra oven, "Scholastisme" -> "Scholasticisme".

Den ydre Organisme, den objective Aand indenfra former sig, nedbryder den
atter indenfra i Tidernes Løb. Saalænge den endnu er beskjæftiget hermed,
forbliver Organismens Grundtræk de samme, og Alt er roligt, uagtet Formen
ældes og Ungdomskraften henvisner; men har Tiden først fuldbaaret sit Foster,
da gribes den af Fødselsveer, da rystes Livselementerne og under heftige
Fødselssmerter komme Tingenes nye Skikkelse til Frembrud. Protestantismens
Aand, der ikke havde villet bære det katholske Aag, havde givet den
subjective Frihed saa slappe Tøiler\footnote{Ulrich von Huttens Ord til
Luther: Wache auf die edle Freiheit! var en Røst, der neppe havde gjenlydt
over Tydskland, før den kaldte de meest forskjellige Stridsmænd til Vaaben.},
at selv Forkjæmperne
% --- p. 72 ---
\setcounter{footnote}{0}
ofte gyste tilbage; man kan da ikke undre sig over, at de nu søgte, saavidt
muligt, at opføre en Dæmning mod den rivende Strøm. Men til at give Læren den
nødvendige Fasthed, var hiin offentlige Troesbekjendelse det bedste Middel.
Man blev sig nemlig snart den Afstand bevidst, der ifølge Sagens Natur altid
maa være imellem Skriftens objective Sandhed og den subjective religieuse
Bevidsthed, og som en Følge deraf gik den almindelige Bestræbelse ikke
saameget ud paa at bevise Skriftens ubetingede Gyldighed --- thi derom var man
alt enig ---, som paa at fortolke den rigtigt; men det eneste ubedragelige
Kriterium paa en sand Fortolkning af Skriften, er det, at Bevidstheden skal
kunne gjenfinde sig selv i den, og et saadant Kriterium afgav netop
Protestantismens Symboler.

Herved lagdes ogsaa et Baand paa den Modsætning, der i de enkelte Menigheder
var fremtraadt mellem de forskjellige Subjectiviteter. Thi da den subjective
Tro, som Protestantismen ikke kan bekjæmpe med ydre Vaaben, ved det aandelige
Slægtskabs indre Magt var bleven ledet til at give den offentlige Bekjendelse
sit Bifald og vedkjende sig den i samme fremsatte Lære i dens Heelhed, fulgte
heraf, at de Enkelte ikke længere strede med hinanden om deres subjective
Tro, men at den enkelte Subjectivitet blev bedømt efter den almindelige
Objectivitet\footnote{\textit{Omnino Philippus Melanchton privatus hanc
insignem mutationem (de confessione variata dicitur) fecit, inconsultis,
quorum communis causa agebatur. Quare statim ac ista innotuit, improbavit
illam elector Saxoniæ, Johannes Fredericus, misitque ad Melanchtonem
cancellarium Pontanum, qui eum objurgaret: \ldots\ Walch: Introductio in
libr. eccl. Luth. symb. p. 188.}

Hvor ofte han maatte forsvare sin Lære mod Indvendinger, vise følgende Ord af
ham selv: \textit{Non gigno novas opiniones, nec aliud majus scelus esse in
Ecclesia Dei sentio, quam ludere fingendic novis opinionibus, et discedere a
Prophetica et Apostolica scriptura et consensu vero Ecclesiæ Dei. \emph{Sequor
autem et amplector doctrinam Ecclesiæ Vitebergensis et conjunctarum, qui sine
ulla dubitatione consensus est Ecclesiæ catholicæ.} Melanch. Opera p.~I. p. 148
Locor. Theol. præfatio.}}. Men herved fremkom tillige lidt efter
% *** THIS NOTE CROSSES THE BATCH BOUNDARY. IT IS COMPLETE HERE. ***
% Note 1) of p. 72 begins in p. 72's footnote block and runs on into the top
% three lines of p. 73's footnote block ("et amplector doctrinam Ecclesiæ
% Vitebergensis ... Locor. Theol. præfatio."), which were read off PDF 82 and
% are included above. The batch for pp. 73--84 must NOT transcribe them again;
% p. 73's own note block begins below them, at its printed "1)".
% PRINTER'S DEFECTS in this note, both transcribed as printed and both read at
% 6x: "communis" is set with a raised, undersized u and a small n (foul
% sorts); and "fingendic" for "fingendis" --- the last letter is a c.
% TYPE: the note is antiqua Latin throughout, hence \textit{}. Within it the
% sentence "Sequor autem ... Ecclesiæ catholicæ." is set in a SLOPED antiqua
% against the upright antiqua of the rest; \emph{} inside \textit{} reverts to
% upright and so preserves that contrast, which is what it is for.

% --- p. 73 ---
\setcounter{footnote}{0}
% *** THE JOIN WITH THE PREVIOUS BATCH. *** The first three lines of p. 73's
% footnote block are the tail of p. 72's note 1 --- "et amplector doctrinam
% Ecclesiæ Vitebergensis et conjunctarum, qui sine ulla dubitatione consensus
% est Ecclesiæ catholicæ. Melanch. Opera p. I. p. 148 Locor. Theol.
% præfatio." --- and they are ALREADY in the pp. 61--72 batch. Read at 3x on
% PDF 82: the join falls immediately after "præfatio." at the end of the THIRD
% line of the block; p. 73's own note 1) begins on the fourth line. Verified
% by the setting as well as the text: the three tail lines carry no note mark
% and are set to the block's continuation indent (x=655 px), where the marks
% of notes 1) 2) 3) stand out at x=511--517 px. The first two of those three
% lines are sloped antiqua and the third upright, exactly as the pp. 61--72
% batch set them. Nothing of p. 72's note is repeated below.
lidt en vis udvortes Begrændsning; den dristige Negation af Former og
Phænomener hørte op, og da nu en Krænkelse af Troesbekjendelsens Gyldighed
tillige ansaaes for en Krænkelse af Sandheden selv, saa man holdt det for
Pligt at udelukke af Kirken eller fængsle dem, der fordreiede Sandheden til
Løgn: var ogsaa Staten strax tilrede med sine Vaaben\footnote{Man sammenligne
hermed --- for ikke at tale om Calvinisternes Kirketugt --- Begivenhederne i
de krypto=calvinistiske Stridigheder. Walch. Historische und theologische
Einleitung in die Religions=Streitigkeiten. \textit{1} Th. \textit{2} Cap.}.
--- Og dog havde Trangen til at anstille Undersøgelser og opkaste
Indvendinger en langt dybere Grund, end tom Spidsfindighed eller
Trættekjærhed. Vel laa det i den protestantiske Troes Natur, at den i
Begyndelsen af Reformationen næsten udelukkende henvendte sig til den
umiddelbare religieuse Bevidsthed, saa at den ikke blot afsondrede Dogmerne
fra de reent historiske Gjenstande\footnote{Denn wo ich je deren eines
mangeln sollte, derer Wercke oder der Predigt Christi; \emph{so wollte ich
lieber derer Wercke, denn seiner Predigt mangeln.} Denn die Wercke hülffen
mir nichts; aber \emph{seine Worte, die geben das Leben,} wie er selbst sagt
Joh.~\textit{5, 51}. Vorrede auf das N.~T. \textit{XII. 55}.}, men endog
begyndte i Dogmerne selv at adskille de moralske og praktiske Momenter fra de
objective og mysterieuse\footnote{\textit{Mysteria divinitatis rectius
adoraverimus, qvam vestigaverimus \ldots\ Proinde non est, cur multum operæ
ponamus in locis illis supremis, de Deo, de unitate, de trinitate Dei, de
mysterio creationis, de modo incarnationis \ldots\ Reliqvos vero locos,
peccati vim, legem, gratiam, qvi ignorarit, non video, qvomodo Christianum
vocem; nam ex his proprie Christus cognoscitur, siqvidem hoc est Christum
cognoscere beneficia ejus cognoscere. Melancht. Loci com. præfatio ed.
principis (ad ann. 1521 ref.)}

Hinfürder lehret er (Carlstad) uns, was Frau Hulda, die natürliche Vernunft,
zu diesen Sachen sagt: gerade als wüsten wir nicht, daß die Vernunft des
Teufels Hure ist und nichts kann denn lästern und schänden alles, was Gott
redt und thut. Wieder die himmlischen Prophet. Schrifft. \textit{XIX. p.
199} smlgn. Von Frau Hulda der klugen Vernunft Doctor Carlstads. Schrifft.
\textit{XIX. p. 208}.}; men skal denne Tro (ὑπόστασις)
% *** THIS NOTE CROSSES THE PAGE BOUNDARY THE OTHER WAY. ***
% The second paragraph of note 3 above --- the Luther "Hinfürder lehret er
% (Carlstad) uns…" --- is NOT printed on p. 73. It stands at the head of
% p. 74's footnote block, UNMARKED and set with an ordinary paragraph indent,
% above p. 74's own note 1). It is set here because it belongs to p. 73's
% note 3, and three things settle that:
%   (a) p. 73's last body line ends flush at the right margin with the closing
%       parenthesis of "(ὑπόστασις)" and carries NO fourth reference mark.
%       Measured, not eyeballed: the raw ink of that line stops at x=3177 px,
%       the same right margin as every full-measure line on the page, and the
%       next ink to the right is the scan edge at x=3790. (The ABBYY layer
%       prints a spurious "*)" there; it is not on the leaf.)
%   (b) p. 74's body carries exactly ONE reference mark, 1) after "tænke",
%       and p. 74's block carries exactly one numbered note, 1).
%   (c) the same thing happens again in this batch at pp. 77--78, where the
%       unmarked leading paragraph of p. 78's block ("Congruentiæ naturales…
%       Samme Ἔκθεσις p. 266") is unambiguously the continuation of p. 77's
%       note 3 ("…Samme p. 265 Θέσις."): same author, same work, the very next
%       page of it. So an unmarked leading paragraph in a note block is this
%       book's way of continuing the previous page's last note.
% "wüsten" for "wüßten" is as printed (long s, not ß --- the glyph has no
% descender and no ß-bowl at 8x).
% --- p. 74 ---
\setcounter{footnote}{0}
kunne finde fuldkommen Hvile i sin Substants, maa den ogsaa have den fuldeste
Tillid til den Grundvolds Fasthed, hvorpaa den støtter sig; thi har den kun
sin Grund i sig selv, da ophæver den sig selv. Derfor kan Troen ikke møde
Modsigelse uden at gaae over til at prøve sit Objects Objectivitet som
saadan, det er, til at tænke\footnote{Derfor kan heller ikke Luther afholde
sig fra, skjøndt han bekjæmper Fornuften, at forsvare sin Lære med objective
Grunde og Analogier, hentede fra de metaphysiske Forhold; derfor tager han
med bestemte Ord sin Tilflugt til Tænkningen, for ikke at forvexle
”\textit{usus passionis}” med ”\textit{factum passionis}” eller
”\textit{quod}” med ”\textit{qualiter}”: ”Gilt denken und ist genug, so will
ich auch denken, besser denn sie \ldots” smlgn. Bekenntniß vom Abendmahl
Christ. \textit{XIX. 440}. Ogsaa Melanchton viser denne Overgang fra
Subjectiviteten til Objectiviteten, i det han f.~Ex. i første Udgave af sine
\textit{loci} (\textit{1521}) forbigaaer Trinitets= og Incarnationsmysteriet
m.~fl.~a., men i de senere ikke blot tilføiede, men endogsaa søgte at
forklare dem. Smlgn. \textit{Edit. ann. 1548 absoluta (In Opp. Viteb. 1580 I.
pg 148)}.

\textit{Filius dicitur Imago et} λόγος. \textit{Est igitur imago cogitatione
Patris genita, qvod ut aliqvo modo considerari possit, a nostra mente exempla
capiamus \ldots\ Mens humana cogitando mox pingit imaginem rei cogitatæ, sed
nos non transfudimus nostram essentiam in illas imagines \ldots\ At Pater
æternus sese intuens gignit cogitationem sui, quæ est imago ipsius non
evanescens, sed subsistens communicata ipsi essentia \ldots\ Dicitur} λόγος,
\textit{qvia cogitatione generatur. Dicitur Imago, qvia cogitatio est imago
rei cogitatæ \ldots\ De Filio pg. 152.}}; ad hvilken Vei denne Overgang
finder Sted, er et andet Spørgsmaal.
% QUOTATION MARKS in this note: RAISED (U+201D) at BOTH ends of all five
% quotations, as the body sets them. Read at 5x on PDF 83. Ten marks.
% The note mixes qv and qu as printed --- "qvod", "aliqvo", "qvia" against
% "quæ est imago ipsius". Not normalised.

Den absolute Tro overgiver i ubetinget Underkastelse det hele Menneske til
Gud; den erkjender ikke blot i praktisk
% --- p. 75 ---
\setcounter{footnote}{0}
Henseende Naturens store Fordærvelse og de gode Gjerningers Afmagt, men
kalder ogsaa i theoretisk Henseende den formørkede Fornuft tilbage til
Lydighed under Troen\footnote{\textit{Hominis intellectus et ratio in rebus
spiritualibus prorsus sunt coeca, nihilque propriis viribus intelligere
possunt. Form. Concord. Cfr. Hase: Hutterus red. 3} Aufl. Pg. \textit{65
etc.}}. Heller ikke var den Strænghed, hvormed hiin Tidsalder vilde
undertrykke Fornuften, blot subjectiv Vilkaarlighed; thi ved Fornuften
forstod man Hverdagsforstanden, der paa udvortes Maade abstraherer alle sine
Begreber (smlgn. §~\textit{5}), og saaledes nødvendigviis maa træde i skarp
Modsætning til Troen, som betragter Alt indenfra. Derfor lagde hiin Tids
Theologer saamegen Vægt paa den Sætning, at en sand og saliggjørende
Gudserkjendelse kun kunde øses fra en særegen
Aabenbaring\footnote{\textit{Datur qvidem revelatio naturalis, sed prorsus
mutilata: Ut post splendidæ domus conflagrationem favillæ et adustæ
particulæ, aut post collapsum gravissimo casu ædificium rudera nonnulla
remanent, ita quoque vix tenues qvædam reliqviæ ac} ἐρείπια \textit{divinæ
imaginis et parvæ scintillæ primævæ lucis, qvæ in mente hominis ante lapsum
fulgebat, in natura jacente et misere prostrata supersunt. Qvenst. Theolog.
Pars I p. 7.} Θέσις \textit{XIII not. III.}}. --- Men hvorledes kan den
formørkede Bevidsthed være modtagelig for en sand Aabenbaring? Jo større
Adskillelsen og Modsætningen er mellem den hellige Gud og den fordærvede
Natur, desto underfuldere maa den Virksomhed være, hvorved Skillevæggen
mellem begge nedbrydes. Skal altsaa en Aabenbaring blive mulig, maa Naturen
blive en anden, maa han, der er over Naturen, vise sig synlig i den; maa der
fremtræde en umiddelbar Identitet af de hinanden modsatte Momenter, \dsi\ der
maa skee et Mirakel. Men skal Gud give sig tilkjende gjennem Mirakler, kan
han kun erkjendes i sine Phænomener paa følgende Maade. Man hører Gud tale,
man
% TYPE in note 1: the note is antiqua Latin, hence \textit{}, but "Aufl." and
% "Pg." inside it are set in FRAKTUR against it, and the figures 3, 65 and
% "etc." are antiqua. Read at 4x. The two Fraktur words are therefore put
% OUTSIDE the \textit{} so that they render upright against the italic, which
% is exactly the contrast the page prints.
% --- p. 76 ---
\setcounter{footnote}{0}
seer himmelske Skikkelser aabenbare sig; Alt, hvad der høres med Ørerne og
sees med Øinene, er Kjendsgjerninger for den fromme Troende. I Alt, hvad der
fremskinner i Phænomenerne, ere den guddommelige Sandheds Mysterier
aabenbarede. --- Men i et aabenbaret Mysteriums Begreb ligger en Modsigelse
skjult; thi forsaavidt der skeer en Aabenbaring, forsaavidt nedstiger
Mysteriet i Phænomenernes Sphære; der viser sig ved denne Nedstigen et
naturligt Moment i Mysteriet, og dette kan da i denne Henseende ikke skjelnes
fra andre Phænomener; men forsaavidt det er en Aabenbaring af Mysteriet,
indeholder det ogsaa et overnaturligt Moment og er uendeligt ophøiet over
alle andre Phænomener. Man maa altsaa vide at adskille Mysteriets Phænomen
fra alle øvrige Phænomener. Men naar man nu ikke kan anerkjende Phænomenets
Guddommelighed uden at kjende Mysteriet og ikke kan kjende Mysteriet uden at
anerkjende Phænomenets Guddommelighed, udfordres der jo en Aabenbaring for at
erkjende Aabenbaringen, er jo Erkjendelsen af Mysterierne og Miraklerne selv
et Mirakel.

Enhver speciel Gudsaabenbaring kan altsaa henføres til de to Hovedmirakler,
Inspirationen og den fra Inspirationen gjennem Kirken medierede Illumination.
--- Som hiin har givet os en fuldkommen Fremstilling af de guddommelige Ting
i den ufeilbare Skrift\footnote{\textit{\ldots\ revelationes certis
testimoniis in verbo suo proposuit, in qvo ceu foetus in alvo materna
nutrimentum attrahens ex umbilico et cotyledonibus nos qvoque sedemus
inclusi, et attrahimus agnitionem Dei et vitam ex verbo Dei, ut cum
invocemus, sicut se patefecit. Melancht. l.~c. cap.~I pg. 151 \ldots\ in
Scriptura Canonica nullum est mendacium, nulla falsitas, nullus vel minimus
error sive in rebus, sive in verbis. Qvenst. l.~l. p.~I. p. 77.}}, saaledes
oplyser denne Bevidstheden
% p. 76 carries ONE note only, and it runs the two quotations together with a
% four-dot ellipsis --- there is no second mark on the page, in the body or in
% the block. Verified at 3x.
% "pg. 151" read at 6x: the middle figure is a 5 with a flat top, not an 8;
% the Fraktur OCR's "181" is rejected.
% "Qvenst. l. l. p. I. p. 77." --- in this fount lowercase l, capital I and the
% figure 1 are the same stroke, so the third token cannot be told apart on the
% glyph alone. It is set as "I" (pars I) because p. 80's note cites the very
% same locus as "Qvenst. p. I pg. 77 Θέσις."
% --- p. 77 ---
\setcounter{footnote}{0}
til at see klart og rigtigt\footnote{\textit{\ldots\ adeo ignorans et perversa
est ratio illa, ut etiamsi ingeniosissimi et doctissimi homines in hoc mundo
Evang. de Filio Dei et promissiones divinas de æterna salute audiant, tamen
ea propriis viribus percipere et vera esse statuere neqveant. F.~C. cfr.
Hntt. red.} Pg. \textit{67.}

\textit{Fides divina, qva qvi infallibiliter credit, hoc verbum, qvod
Scriptura proponit, esse vere Dei verbum, immediate oritur a vi et potestate
S. Scripturæ intrinseca, qva Scriptura ipsa ad generandam fidem sibi debitam
efficax est. Qvenst. p.~I. pg. 98.}}. Det er da let at forstaae, hvad der
menes med den Fordring, at vi skulle tage Fornuften fangen. Alle
Aabenbaringens Phænomener maae nemlig betragtes fra en doppelt Side. Fra den
objective Side betragtet aabenbarer Gud sig i dem: han tænker i Sandhed
saaledes, som han taler; han er i sig saaledes, som han vil forestilles af
Menneskene; han er treenig, thi han har aabenbaret sig som treenig; han er
\emph{sanddru}\footnote{\textit{Causa certitudinis est revelatio Dei, qvi est
verax. Melanch. l.~l. pg. 148.}}. Men skjøndt Gud er Sandhed, og hans Natur
ikke kan modsige sig selv, fatter Menneskets Fornuft dog kun de Egenskaber i
Gud, som ligne dets egne og henhøre til vor Verden; da altsaa Gud har en
særegen Natur, erkjender Fornuften selv, at den i Henseende til denne maa
fornægte sin egen Gyldighed\footnote{\textit{Qvando objiciunt adversarii
(Socin, Armin.) religionem multa habere supra, nihil vero contra rationem,
respondeo: 1) Articuli fidei in se non sunt contra rationem, sed solum supra
rationem; per \emph{accidens} vero fit, ut sint etiam contra rationem, qvando
ratio judicium sibi de illis sumit ex suis principiis, nec seqvitur lucem
verbi, sed eosdem negat et impugnat. Qvenst. p. 45 XIX Distinct.}

\textit{Cognitio Dei naturalis non extendit sese ad essentiam divinam,
qvatenus illa relativa seu in Tribus Personis consideratur \ldots} Samme
\textit{p. 265} Θέσις.

\textit{Congruentiæ naturales et analogia creatarum rerum cum hoc mysterio
non divinam fidem, sed opinionem tantum humanam generant \ldots} Samme
Ἔκθεσις \textit{p. 266 IV. Obst}}. Men paa den anden
% *** PRINTER'S DEFECT, transcribed as printed: "Hntt." for "Hutt." in note 1
% --- a turned u. Settled at 7x: the arch of the second letter closes at the
% TOP, exactly like the n of "ignorans" and unlike the u of "qveant" three
% words earlier on the same line, whose bowl is open at the top. Note 1 of
% p. 79 spells the same abbreviation "Hutt." correctly. Both OCR witnesses
% also read Hntt/Ilntt. NOT in the Rettelser; not corrected. Same family of
% defect as p. 2 "absolnt", p. 3 "Tiugen", p. 70 "Gruude".
% EMPHASIS: "sanddru" is letterspaced in the body. Read at 6x --- s a n d d r u
% with the fount's emphasis spacing, against the ordinary setting of "treenig"
% on the same line. spacing.py missed it (a short word; a known false negative).
% "doppelt" for "dobbelt" is as printed; both OCR witnesses agree. Not corrected.
% TYPE in note 1: the note is antiqua Latin, but "Pg." between "red." and "67"
% is FRAKTUR, so it stands outside the \textit{}. Same setting as p. 75 note 1.
% "esse vere Dei verbum" in note 1's second paragraph is NOT letterspaced ---
% spacing.py flagged it, and at 8x "esse" measures no wider per letter than
% "vere" four words later. Set plain.
% "per accidens" and, on p. 79, "Hoc est", "conseqventes", "non positis" and
% "secundarii" are set in a SLOPED antiqua against the upright antiqua of the
% rest; \emph{} inside \textit{} reverts to upright and so preserves the
% contrast, following the p. 72 note.
% *** THIS NOTE CROSSES THE PAGE BOUNDARY. *** The third paragraph of note 3
% ("Congruentiæ naturales … Samme Ἔκθεσις p. 266 IV. Obst") is printed at the
% head of p. 78's footnote block, unmarked and indented, above p. 78's own
% notes 1) and 2). It is set here because it continues note 3: the paragraph
% before it cites "Samme p. 265 Θέσις" and this one "Samme p. 266", the next
% page of the same work. Same pattern as pp. 73--74.
% "Obst" is what is printed at the end of it --- O b s t, with a crossed t,
% read at 7x, where p. 78 note 1 and p. 79 spell the same abbreviation "Obs."
% Both OCR witnesses read Obst. Transcribed as printed; "Obs." is meant.
% GREEK, normalised per the header: the script theta of Θέσις and Ἔκθεσις is
% typed as plain θ/Θ. In Ἔκθεσις the breathing-and-acute is printed as a
% separate raised mark BEFORE the capital, in the older manner; it is set here
% as the precomposed Ἔ.
% --- p. 78 ---
\setcounter{footnote}{0}
Side skulle vi jo kjende Aabenbaringens Phænomener og anvende dem paa vort
eget Liv; dersom nu den menneskelige Fornuft maatte erklære dem for at
modsige hverandre, saa at vi altid maatte være i Tvivl om deres Sandhed, blev
der os jo i Virkeligheden ingen Aabenbaring til Deel. Det vil derfor være
hensigtsmæssigt i de guddommelige Phænomeners Rige at anvende Fornuften i
formel Henseende\footnote{\textit{Aliud est principia et axiomata philosophica
in Theologia adhibere illustrationis, declarationis et secundariæ probationis
gratiâ, ubi res e Scriptura definita est; et aliud, adhibere eadem decisionis
et demonstrationis causa, vel principia philosophica et ratiocinationes ex
illis factas pro principio in Theologia agnoscere, aut ex illis de rebus
fidei judicium ferre.} Samme \textit{p. 42 II Obs.} smlgn. \textit{III. IV. V
VI. Distinct.}}. --- Hvorledes lyder nu den øverste Lov, som den i sit
lovmæssige Herredømme gjenindsatte Fornuft maa fastsætte? Da de aabenbarede
Sandheder ikke kunne begribes i den Forstand, at man skulde kunne paavise den
dialektiske indbyrdes Overgang imellem dem, og den ordnende Virksomhed kun
har den Opgave at vise, at de alle kunne existere ved Siden af hinanden,
efterdi Visheden om deres Tilværelse ene hidrører fra deres
Aabenbarelse\footnote{\textit{Etsi autem omnes Angelicæ et humanæ mentes
obstupescunt admiratione hujus arcani, qvod Deus Filium genuit \ldots\ tamen
sic sentire necesse est, qvia, ut toties jam dictum est, de Deo sentiendum
est, sicut se patefecit \ldots\ qvanqvam hæc arcana non penitus perspicimus,
tamen in hac vita Deus inchoari vult hanc agnitionem, et invocationem nostram
a falsa discerni.} Samme \textit{p. 15.}}, maa man fremfor Alt fastholde den
bekjendte Sætning, at hvad der er virkeligt, ogsaa maa være muligt
(\textit{ab esse ad posse valere conseqventiam}). Denne
% "(ab esse ad posse valere conseqventiam)" is antiqua, NOT letterspaced ---
% spacing.py's RUN here, and its RUNs on pp. 79 and 83, are all antiqua Latin
% inside Fraktur, whose wider fitting the tool reads as Sperrsatz. Checked on
% the image at 4x in each case. \textit{} only.
% "Samme" and "smlgn." inside the notes are Fraktur and stand outside \textit{}.
% --- p. 79 ---
\setcounter{footnote}{0}
Lov medfører, at Logiken selv snart gjør Afkald paa, snart opfordres til
Brugen af sine egne Regler. I Treenighedsdogmet tier saaledes Logiken; thi
een Gud har aabenbaret sig i tre Personer, altsaa kan Een være i Tre; derimod
bør man høre Fornuften, naar den ifølge sine egne Regler fremsætter Læren om
Christi Legemes Allestedsnærværelse, uagtet denne ikke er bestemt fremsat i
Skriften, thi Christus er jo aabenbaret som den, der paa eengang sidder ved
Guds høire Haand og er reelt nærværende i
Nadveren\footnote{\textit{Gerhard: Qværitur, sitne corpus Christi vere et
substantialiter in coena præsens? Affirmant Eccles. nostræ, qvia Christus
dicit: \emph{Hoc est} etc. Negant adversarii ex hoc principio, qvia verum et
naturale corpus non potest simul et semel esse in pluribus locis. Addunt,
rationem renatam non posse aliter de corpore statuere, testimonium ergo ejus
audiendum esse. At inqvam, ratio renata, qvatenus talis de art. fidei statuit
ex Dei verbo, et limites ejus non egreditur. Cfr. Hutt. red. p. 71.}}. Herved
fremkommer ogsaa Inddelingen af Troens \textit{articuli fundamentales} i
\textit{primarii} og \textit{secundarii} og af dens \textit{articuli
conservativi} i \textit{antecedentes} og
\textit{consequentes}\footnote{Inddelinger, der vise, at vore Fædre ikke
undlode at tage Hensyn til Fornuftens Principer: „\textit{Articuli
\emph{conseqventes} sunt, qvi explanant fidei salvificæ conseqventia, qvibus
positis fides confirmatur, \emph{non positis} evanescit.}“ \textit{Art.
\emph{secundarii}, qvi} „\textit{licet ignorari possint illæso fidei
salutisque fundamento, pertinaciter tamen negari salvo illo non possunt,
nedum impugnari.}“ \textit{Hutt. rediv. §~17 p. 24. 25.}}. --- Da man ved
denne Behandlingsmaade ikke kunde komme til nogen Indsigt i den guddommelige
Naturs Væsen, men kun befæstede en ydre Kundskab til dens Mysterier, og disse
ikke droges frem af de Phænomener, hvori de vare fremtraadte, men meget mere
forbleve indesluttede i dem, var en Adskillelse mellem Ideen og dens ydre
Form naturligviis umulig, og man knyt-
% *** PRINTER'S DEFECT IN THE FOOTNOTE MARKS: p. 79 sets 1) TWICE. *** The
% mark of the SECOND note in the block is printed 1), not 2). Read at 12x and
% compared glyph for glyph with the 1) of note 1 on the same page (identical:
% the same flag at upper left, the same foot serif) and with the 2) of p. 83
% note 2 (quite different: a bowl and a flat base). The BODY is correct and
% sets 1) and 2). The transcription therefore prints 2) --- the plain next
% number, as \footnote gives it --- and the printed 1) is recorded here.
% QUOTATION MARKS: *** THIS PAGE DEPARTS FROM THE BOOK'S STANDING FORM. ***
% Note 2 sets the DANISH LOW--HIGH pair „…“ twice --- opening mark on the
% baseline, closing mark raised --- where the body of this book uses the
% raised ” at both ends. Read at 8x. This page contributes 2 „ and 2 “ and no
% ” to check.py's counts, and is balanced in itself.
% "consequentes" in the body is set with qu and "conseqventes" in the note
% with qv, on the same opening. Both as printed.
% --- p. 80 ---
\setcounter{footnote}{0}
tede kun det ydre Baand mellem Dogmatiken og den hellige Historie end
fastere. Man havde forkastet menneskelige Traditioner og opgivet
Catholikernes τερατολογία; man var derfor saameget nidkjærere i at værne om
det Grundlag af overnaturlige Facta, man havde ladet staae urørt af Tvivlen,
og kunde end ikke i de ubetydeligste Punkter af den hellige Historie
indrømme Muligheden af nogen Vildfarelse i den inspirerede
Skrift\footnote{\textit{\ldots\ omnia et singula sunt verissima, qvæcunqve in
illâ traduntur: sive dogmatica illa sint sive moralia, sive Historica,
Chronologica, Topographica, Onomastica, nullaque ignorantia, incogitantia aut
oblivio, nullus memoriæ lapsus Spiritus Sanctus amanuensibus in consignandis
sacris litteris tribui potest aut debet. Qvenst. p.~I pg. 77} Θέσις.

Og det var ikke blot Dogmatikere henimod Slntningen af det \textit{17}de
Aarhundrede, der fremsatte saadanne Meninger; ogsaa Historikere i den sidste
Deel af det \textit{19}de Aarhundrede have paa lignende Maade skjelnet mellem
guddommelige og menneskelige Kilder:

Ihre (\dsi\ den hellige Skrifts) Verfasser haben nicht allein nie geirret,
sondern auch nicht einmal irren können, und ihre Erzählungen sind daher nicht
allein so gewiß, als jede andere historische Wahrheit, sondern auch
untrüglich gewiß. Ch.~W.~F. Walch, krit. Nachricht v.~d. Qvellen der
Kirchengeschichte. Leipz. \textit{1770} p, \textit{23}. --- De forsøgte
derfor ikke blot efter Luthers Exempel i moralsk Henseende at gjøre
mangfoldige, frugtbringende Anvendelser af den hellige Historie paa de
Christnes Liv, men begrundede endogsaa vigtige Dogmer paa Begivenheder, der
fortælles i Skriften:

Saaledes gjendriver f.~Ex. E. Brochmand ved Sammenligning og sindrig Analyse
af nogle faae Skriftsteder (Gen. \textit{1, 26. 27}; Eph. \textit{4, 24};
Col. \textit{3, 10}) med ringe Uleilighed alle Bellarmins Indvendinger imod
Læren om det første Menneskes Fuldkommenhed og Udødelighed. \textit{Univ.
Theol. Syst. p. 124. Quæst. III. IV.}

I sit Beviis for, at der i det første Menneskes Erkjendelse var en særegen
Overeensstemmelse med Gud ”ὁ πρωτοτύπος”, siger Qvensted: \textit{\ldots\ Ex
Adami factis qvalia sint 1)} Ονοματοθεσία, \textit{sive nominum conveniens
impositio Gen. 2. 19, qvæ erat non solum \emph{Grammatica}, qvoad
nomenclaturam animalium, sed vel maxime \emph{Logica}, qvoad verissimam
definitionem.} Ja man beviste den hellige Aands Personlighed paa samme
historiske Maade; saaledes allerede Melanchton: \textit{\ldots\ impio
sophismati Ecclesia vera opponit testimonia in scripturis tradita, qvorum
primum et maxime illustre est revelatio divinitatis, facta in baptismo
Christi, ubi clare discernuntur tres personæ. Pater inqvit: Hic est Filius
meus dilectus. Est igitur alia persona Patris, alia Filii. Descendit autem
Spiritus Sanctus in specie columbæ. Jam si Spiritus esset agitatio in animis
creata, non appareret peculiari specie corporali.}

\textit{Sic et in Pentecoste apparet Spiritus Sanctus peculiari specie
corporali. Hæ patefactiones non sunt frustra factæ, immo sunt excellentia
beneficia Dei, in qvibus Deus se patefecit Ecclesiæ, et testificatus est
Spiritum Sanctum esse personam.} S.~S. Pg. \textit{137}.

Og efter hans Exempel beviser ogsaa Qvensted Triniteten „\textit{ex singulari
illa} θεοφανείᾳ \textit{facta in baptismo Christi \ldots\ ut vere dictum sit
ab antiqvis:} ”\textit{abi Ariane ad Jordanem! et videbis Trinitatem}“ Pg.
\textit{324}.}. Da Theologerne saaledes agtede Skriftens
% *** THIS NOTE RUNS OVER ONTO p. 81, WHICH HAS NO NOTE OF ITS OWN. ***
% p. 80's single note fills the whole of p. 81's footnote block as well ---
% five paragraphs on p. 80, two more on p. 81, none of them marked. p. 81's
% body carries no reference mark and p. 81 therefore gets NO
% \setcounter{footnote}{0}. Verified line by line on PDF 89 and PDF 90.
% *** PRINTER'S DEFECT, transcribed as printed: "Slntningen" for
% "Slutningen", a turned u. Settled at 7x: the letter after Sl closes at the
% TOP, matching the n of "Meninger" on the next line, where the u of
% "guddommelige" two lines down is open at the top. Both OCR witnesses read
% Sln. Same family as p. 2 "absolnt", p. 3 "Tiugen", p. 70 "Gruude", p. 77
% "Hntt.". NOT in the Rettelser; not corrected.
% *** PRINTER'S DEFECT, transcribed as printed: "Leipz. 1770 p, 23." --- a
% COMMA after p where the sentence wants a point. Read at 6x; both OCR
% witnesses read a comma. Not corrected.
% *** THE GREEK WORD Ονοματοθεσία IS BROKEN ACROSS THE PAGE BOUNDARY. *** The
% leaf sets "1) Ονο" as the last thing on p. 80's block and "ματοθεσία," as
% the first thing on p. 81's, with NO hyphen at the break. Read at 8x on both
% leaves. It is joined here. The capital omicron carries NO breathing in the
% printing (checked at 12x against the Ἔκθεσις of p. 78, whose breathing is
% plainly set); it is left bare, as printed.
% GREEK, accents as printed and normalisations noted: τερατολογία and
% πρωτοτύπος are set with the fount's CURSIVE RHO, typed here as plain ρ;
% θεοφανείᾳ is set with the SCRIPT THETA, typed as plain θ. Both per the
% header's standing rule (the variant sorts are fatal under textalpha). The accent of
% πρωτοτύπος stands over the upsilon in the printing --- πρωτοτύπος, not the
% grammatical πρωτότυπος --- and is reproduced as printed, like the σύ and
% Αληθῆ already logged in the header.
% QUOTATION MARKS on this page: the body quotation of Quenstedt's Greek,
% ”ὁ πρωτοτύπος”, is RAISED at both ends --- the book's standing form, read at
% 10x. Two marks. The marks in the p. 81 half of this note are a different
% matter; see the comment at p. 81.
% "1)" inside the Quenstedt Latin is Quenstedt's own enumeration, not a
% footnote mark.
% --- p. 81 ---
Bogstaver som det eneste og tilstrækkelige Grundlag for den positive Lære,
medens de nægtede Enthusiasterne (Schwenkfeldianerne og Weigelianerne) med
deres aandelige Indskydelser og Engleaabenbaringer al Tiltro, gik deres hele
Stræben ud paa at fastsætte Troens objective Former. De afholdt sig derfor
vel fra tomme Speculationer, thi de vilde ikke som Scholastikerne ”opstille
en centauragtig Theologie, et \textit{mixtum compositum} af guddommelige
Aabenbaringer og philosophiske Fornuftslutninger”; men overlode alligevel
Fornuften det vigtige Kald at ordne og opstille Stoffet. Man kunde nemlig i
den store Masse af Gjenstande befrygte, at den ene skulde synes at modsige og
ophæve den anden, og maatte derfor være betænkt paa efter bedste Evne at
godtgjøre Muligheden af alle Momenternes gjensidige Bestaaen.
% NO \setcounter{footnote}{0} HERE, AND THAT IS NOT AN OMISSION. p. 81 carries
% no note of its own: its whole footnote block is the tail of p. 80's single
% note, unmarked, and p. 81's body carries no reference mark. Checked on the
% image and by a line-by-line measurement of PDF 90.
% *** QUOTATION-MARK DEFECT, in the p. 81 part of p. 80's note. ***
% The last paragraph opens Quenstedt's Latin with the LOW mark
% („ex singulari illa θεοφανείᾳ …), opens the inner quotation with the RAISED
% mark (”abi Ariane ad Jordanem!), and then closes with a SINGLE high “ after
% "Trinitatem" --- which serves the inner quotation, leaving the outer one
% never closed. Read at 8x. Transcribed exactly so. Its effect on check.py's
% counts is „ +1, “ +1 (so the low--high balance is unaffected) and ” +1, so
% the RAISED PARITY GOES ODD. It needs a QUOTE_DEFECTS row (81, 0, +1, …).
% The BODY quotation on this page --- ”opstille en centauragtig Theologie …
% Fornuftslutninger” --- is the ordinary raised pair, read at 6x. Two marks.
% *** THE MARKS ON THIS PAGE THEREFORE COME TO: „=1 “=1 ”=3. ***
% "Grammatica" and "Logica" in the p. 81 half of the note are sloped antiqua
% against the upright antiqua of the rest, hence \emph{} inside \textit{}.
% --- p. 82 ---
\setcounter{footnote}{0}
Dette besværlige Arbeide maatte da Fornuften paatage sig at udføre ved sine
Bekræftelser, Benægtelser, Adskillelser, Inddelinger og Undtagelser. Og idet
den nu i den subtile Analyse af de enkelte Momenter angav τὸ ζητούμενον,
opstillede τὴν θέσιν, tilføiede τὴν ἀντίθεσιν, overtog τὴν βεβαίωσιν, vorede
Arbeidet op til en saadan Masse af Sætninger, at vi ere uvisse om, hvad vi
meest skulle beundre, --- den Udholdenhed, hvormed Fornuften rigtignok
aldeles mechanisk sammenpassede alle de enkelte Dele, eller --- den Styrke,
der satte Troen istand til at bære denne store, af de meest modsatte
Elementer opførte Bygning\footnote{Das dogmatische System des \textit{16}-und
\textit{17} Jahrhunderts kam mir vor wie einer unsrer alten deutschen Münster
mit seinen himmelstrebenden Spitzbogen und wunderlichen sinnvollen
Zierrathen \ldots\ Einen Dom wie unsre Vorfahren kann auch unsre Zeit nicht
wieder bauen \ldots\ es wird einem aber doch ganz besonders wie in einem
Gotteshause darin zu Muthe. Hase: Der theol. Facult. zu Tübingen
(\textit{Hut. rediv,})}. Dette pragtfulde Tempel i stor Stiil, som den fra
Landflygtigheden tilbagevendte Menneskeaand opførte sig, opbyggedes og
befæstedes under en bevæget Tids mangfoldige Storme; stedse i Vaaben, stedse
beredte til Kamp reiste disse Mænd deres Bygning; --- om det var Guld og Sølv
og kostelige Ædelstene, hvoraf de opførte Bygningen, eller om det var Træ og
Hø og Straa, det viser først den fremgryende Dag, den prøvende Ild.
% *** PRINTER'S DEFECT, transcribed as printed: "vorede" for "voxede". ***
% The sense is plain --- the work GREW up to such a mass of propositions ---
% but the leaf sets an r. Settled at 8x by direct comparison of three glyphs:
% the letter here has a short arm and sits on the baseline, exactly like the r
% of "tilrede" (p. 73 l. 6); the Fraktur x of "strax", four words earlier on
% that same line of p. 73, drops a clear curl BELOW the baseline, which this
% letter does not. The Fraktur OCR reads x correctly elsewhere on this scan
% ("strax", "forvexle" p. 74) and reads r here. NOT in the Rettelser; not
% corrected.
% "16-und" is as printed: a short antiqua hyphen set tight against "und",
% where the sense wants "des 16. und 17. Jahrhunderts". Read at 7x.
% "(Hut. rediv,)" is as printed: one t, and a COMMA before the closing
% parenthesis where a point is wanted. Read at 8x. p. 79 note 2 spells the
% same work "Hutt. rediv." Not corrected.
% *** THE PRINTED § HEAD AGREES WITH THE INDHOLD HERE. *** Both read
% "Fortsættelse. Overgang til Kriticismen." Verified on PDF 91 at 4x, word for
% word including both points. "§ 16." is centred bold antiqua, the argument
% below it letterspaced Fraktur, as everywhere else in this book. No
% disagreement to record, unlike §§ 7, 10, 13 and 15.
\parag{16}{Fortsættelse. Overgang til Kriticismen.}

Den absolute Overeensstemmelse mellem Væsenet og dets Phænomener, hvorpaa
denne Bygning hviler, ophæver sig selv.
% --- p. 83 ---
\setcounter{footnote}{0}
Thi da den guddommelige Natur ikke har aabenbaret sig i eet, men i flere
Phænomener, ikke er udtrykt paa een, men paa mange, hverandre modsigende
Maader, maa der i Sandhed i ethvert Udtryk for Gudsideen være Noget, mere
eller mindre uadæqvat\footnote{\textit{Quia intellectus noster finitus
infinitam et simplicissimam Dei essentiam uno conceptu adæquato adæquate
concipere nequit, ideo distinctis et inadæquatis conceptibus essentiam
divinam inadæquate repræsentantibus eandem apprehendit. Quenst.} P.
\textit{284,}}; man maa som en Følge heraf kunne tage det enkelte Ord i
forskjellig Betydning paa dets forskjellige Steder og gjøre Forskjel mellem
det, der er sagt i egentlig og i uegentlig Forstand\footnote{Herhen høre alle
Anthropomorphismer; saaledes adskiller Qvensted ”\textit{ea, quæ
imperfectionem aliquam subinferunt et Deo proprie tribui nequeunt, sed
improprie et per} ἀνθρωποπάθειαν \textit{ipsi adscribuntur, ut affectus
etorgana corporea, ut ira, odium, manus, pedes, oculi}” \textit{etc,} \ldots{}
P. \textit{296}.}. Man vender da Reglen om og slutter fra en Tings Umulighed
til dens Uvirkelighed (\textit{a non posse esse ad non esse valere
consequentiam}), --- men saaledes erholder Fornuften sin tabte Frihed
tilbage, thi det er jo den, der maa sammenligne de enkelte Udtryk. Hvor
tvende Udtryk ere hinanden saaledes modsatte, at det ene aabenbart udelukker
det andet, der bliver jo Resultatet forskjelligt, eftersom man antager det
ene eller det andet af dem for sikkert. Men hvilket, der skal blive staaende,
og hvilket, der skal bort, det kan ligesaalidt Summen af de aabenbarede
Yttringer --- thi det er netop om de enkelte Yttringer, at Sagen dreier sig
---, som Fornuften selv afgjøre, thi det er jo netop det subjective
Raisonnement paa dette Standpunkt ene og alene derom at gjøre at have et fast
og urokkeligt Udgangspunkt; men saaledes forsvinder hele den objective Læres
System og kun den subjective Tro i sin fulde Ubestemthed bliver tilbage.
% PRINTER'S DEFECTS on this page, both in the notes, both transcribed as
% printed and both read at 7x:
%   note 1 ends "Quenst. P. 284," --- a COMMA where the note wants a point.
%   note 2 sets "etorgana" for "et organa", the two words run together with no
%   space. Both OCR witnesses agree in each case. Neither is in the Rettelser.
% "P." in both notes is FRAKTUR against the antiqua Latin, so it stands
% outside the \textit{}; the figures 284 and 296 are antiqua.
% QUOTATION MARKS in note 2: RAISED (U+201D) at both ends, the book's standing
% form. Read at 8x. Two marks.
% GREEK: ἀνθρωποπάθειαν is set with the SCRIPT THETA in both of its thetas
% (the word is broken ἀνθρωποπά-θειαν across two lines of the note); typed
% here as plain θ, per the header.
% "Quia"/"Quenst."/"nequit"/"aliquam" are set with qu on this page where other
% pages use qv. As printed; not normalised.
% --- p. 84 ---
\setcounter{footnote}{0}
Man kan da indrømme, at een Gud har aabenbaret sig i tre Personer, og dog paa
samme Tid efter den ovenfor anførte Regel slutte, at, hvor der tales om tre
Personer i Gud, kan det paa Grund af Guds Eenhed kun være at tage i figurlig
Forstand\footnote{Saaledes Socinianernes Indvending: \textit{Quoniam essentia
divina tantum unica est numero, ideo nullatenus plures in eâ possunt esse
personæ. Catech. Racov. Cap. 1, p. 21.}

Ja Volkelius \textit{Libr. V. de V. R. cap. IX. fol. 401} gaaer endog
saavidt, at han tilskriver „den gamle Løgner Satan“ Læren om Triniteten; og
Joachim Stegmannus lærer, at dette Mysterium baade strider imod
\emph{Skriften} og imod \emph{Fornuften}. Smlgn. Quenst. P.
\textit{345}.}. --- Da man søgte at fastsætte Læren paa denne Maade, var det
en Selvfølge, at den subjective Tro ogsaa maatte holde over Lærens enkelte
Sætninger, ligesom den søgte at forsvare den i Skriften aabenbarede Sandhed i
Almindelighed; thi kunde man drage Sandheden af een Troessætning i Tvivl,
vilde det jo ogsaa være let paa samme Maade at svække Troen paa enhver af de
øvrige. Derfor kunde man i Religionssamtalerne ikke give efter i
nogetsomhelst Punkt; der disputeredes stedse \textit{de omni et de nullo}.
Derfra hidrørte det uheldige Udfald af de ireniske
Colloqvier\footnote{Saaledes klages f.~Ex. i heftige Udtryk over Colloquiet i
Thurn:\\
\textit{Quæ sunt gesta tribus Thoruna mensibus urbe,}\\
\textit{Hæc potuere tribus cuncta diebus agi.}\\
og\\
\textit{Quid synodus? nodus. Patrum chorus integer? æger.}\\
\textit{Conventus? ventus. Gloria? stramen. Amen.}\\
Smlgn. \textit{Dr.} K.~H. Hagenbach: Vorles. üb. W. Geschicht. d. Reform.
\textit{4} Theil, P. \textit{148---149}.

Om de syncretistiske Stridigheder, der give rige Bidrag til hiin Tidsalders
Charakteristik, smlgn. J.~G. Walch. R. Str. \textit{1} Theil, §~\textit{5},
P. \textit{216} ff.}; derfra det Frugtesløse i Syncretisternes Bestræbelser
efter at at holde Middelveien; derfra Heftigheden, Partisygen, Fana-
% *** PRINTER'S DEFECT, transcribed as printed: "efter at at holde
% Middelveien" --- "at" set twice. Read at 5x; both OCR witnesses have it.
% NOT in the Rettelser; not corrected.
% EMPHASIS in note 1: "Skriften" and "Fornuften" are LETTERSPACED --- read at
% 8x, S k r i f t e n and F o r n u f t e n, both plainly wider than the
% "Mysterium" and "strider" beside them. The ABBYY layer also spaces Skriften.
% spacing.py did not flag either (both fall in the note block, where the tool
% is unreliable).
% QUOTATION MARKS in note 1: the DANISH LOW--HIGH pair „den gamle Løgner
% Satan“ --- opening mark on the baseline, closing mark raised. Read at 8x.
% Like p. 79's note, and unlike the body, which uses the raised mark at both
% ends. Two marks, balanced.
% The verse in note 2 is set as four display lines with "og" alone between the
% couplets; the line breaks are kept with \\ inside the footnote.
% "Thurn" for Thorn is as printed, and the note itself spells the town
% "Thoruna" in the Latin two lines later. "Colloquiet"/"Colloqvier" --- qu in
% the note, qv in the body. All as printed.
% TYPE: "Dr." is antiqua; "K. H. Hagenbach: Vorles. üb. W. Geschicht. d.
% Reform." and "Theil", "P." and "ff." are Fraktur; the figures are antiqua.

% --- p. 85 ---
% JOINT WITH p. 84, verified on the image (PDF 93 foot / PDF 94 head, 3x):
% p. 84 ends "...Partisygen, Fana-" and this page opens "tismen" --- the word
% is "Fanatismen". p. 84's two notes are COMPLETE on p. 84 (they close with
% "1 Theil, § 5, P. 216 ff."); nothing runs over. p. 85 carries no footnote
% block at all, so it takes no \setcounter.
tismen og den rasende Forbittrelse hos den Tids Theologer. Det er kun ved at
udvikle sig til Extrem, at saadanne Systemer fuldendes og ophæves; derfor maa
man stedse iagttage dem under deres egen Udviklingsgang.

Troens Grundprincip er Tænkningen; derfor udfordrer Troen en umiddelbar
Identitet af det subjective og objective Moment; men i denne ligger tillige en
umiddelbar Modsætning skjult. Hiin træder først frem, denne senere. Naar
Subjectet i umiddelbar Overeensstemmelse hengiver sig til Objectet, falder det
ikke den Troende ind, at anstille en udvortes Sammenligning imellem sig selv
og Objectet for sin Tro, --- thi den Tanke at skulle anstille en omhyggelig og
smaalig Undersøgelse af Begges Egenskaber, eller at disse skulle behøve at
afpasses efter hverandre for at kunne stemme med hverandre, er unaturlig for
den første umiddelbare Tro ---; men greben af det Objects Storhed, hvorpaa han
troer, gaaer hans hele Bestræbelse ud paa at bestemme og fastsætte den
objective Troes Former. Herfra hidrørte netop den protestantiske Troes Lyst
til at opstille Systemer. --- Men naar nu Troen saaledes har nedsænket sig i
den rene Obiectivitet, forandres Forholdet. Det objective Systems absolute
% *** PRINTER'S DEFECT, transcribed as printed: "Obiectivitet" for
% "Objectivitet". Read at 3x: the letter after the O is a DOTTED i with no
% descender, against the clear descending Fraktur j of "objective" in the very
% next sentence on the same page. Not in the Rettelser; not corrected.
Gyldighed viser sig nemlig deri, at Gjenstandene gjælde, ikke fordi de
begribes, men fordi de ere factisk givne. Systemet fremtræder da som noget
Fremmed, der skal adlydes, fordi det fordrer Lydighed, \dsi{} Objectiviteten
opsluger Subjectiviteten. Men derved nedbryder ogsaa Systemet sig selv; thi
vel har man glemt dets subjective Oprindelse, men det er i Virkeligheden
Troen, der har skabt det; derfor føler den subjective Bevidsthed sig hæmmet
ved fremmede, ydre Baand, naar Systemet paa udvortes Maade vil gjøre sig
gjældende over Aanderne; den begynder at reflectere, og nu træde
Modsætningerne frem; men da vil ogsaa det objective System opløses og gaae til
Grunde,
% --- p. 86 ---
% No footnote block on this page either.
fordi det kun er paa Subjectivitetens umiddelbare Sympathie med det, om man
saa tør sige, at dets hele Magt beroer.

Da derfor Troens enkelte Sætninger vare blevne umiddelbart opstillede i denne
Theologie, gik Troen selv over til at blive en Sum af udvortes Sætninger, en
forstenet Formalisme. Man samlede med største Begjerlighed \textit{dicta
probantia} og \textit{loci communes}, ikke for indenfra i Sandhedens Navn at
retfærdiggjøre Skriftens Udsagn, men for i Skriftens Navn paa udvortes Maade
at komme i haandgribelig Besiddelse af Sandheden: man tilbad de enkelte Ord og
Tegn i Skriften. Derfor trak Aanden sig tilbage fra denne Aandens
Objectivitet; derfor formørkedes Bibelens sande Lys under denne
\emph{Bibeldyrkelse}. --- En Negation af de positive Former fulgte efter denne
Formalisme; den religieuse Bevidsthed begyndte at søge Sandheden i sit eget
Indre, og vore Scholastikeres Kræfter udtømtes i Kampe mod Pietisterne.
% EMPHASIS: "Bibeldyrkelse" is letterspaced, read at 3x (spacing.py's
% "single?" flag confirmed on the image). "dicta probantia" and "loci
% communes" are ANTIQUA against the Fraktur, hence \textit{}; the "og"
% between them is Fraktur and is not italicised.

Saaledes fremtraadte nu det andet Moment i Troen. Vel skal Subjectet af sit
ganske Hjerte hengive sig til Objectet; men --- Skade gjør klog, derfor maa
den Troende først vende Øiet indad i sin egen Bevidstheds Dybder og undersøge
sit eget Hjerte. --- Protestantismen indeholder saaledes ikke blot sin egen
Negation, men i samme Maal sin egen Restauration; som det reelle Moment i
Begyndelsen havde været det overveiende, netop fordi det umiddelbart greb
Overbeviisningen, saaledes erholdt nu lidt efter lidt det formelle Moment
større Betydning, saa at Philosophien allerede paa selvsamme Tid,
Protestantismen efter Trediveaarskrigens Slutning fik fuldkommen
Anerkjendelse og retlig Tilværelse i Staten, kunde vende sig bort fra
Theologien og lære, \emph{at man skulde tvivle om Alt}. I den concrete Troes
Navn hævdede Luther i det \textit{16}de Aarhundrede Samvittighedsfriheden; i
Frihedens Navn forsvarede i \textit{18}de og \textit{19}de Aarh. hans
% EMPHASIS: "at man skulde tvivle om Alt" is letterspaced across six words,
% read at 3x; the closing period is NOT letterspaced and is set outside the
% \emph{}. spacing.py's RUN flags "Vel fal" (= "Vel skal") and "skulde tvivle
% Alt" were checked on the image: only the second is real.
% --- p. 87 ---
% "%" below suppresses the newline: the printed page continues the previous
% page's sentence, and a bare newline here would double the interword space.
\setcounter{footnote}{0}%
sande Efterfølgere Friheden og Tænkningen og muliggjøre saaledes en
Gjenfødelse af den objective concrete Troeslære\footnote{\ldots\ selv i det
Tilfælde, at Reformatorerne havde ikke blot villet skildre Protestantismens
Aand og Tendents og derhos tillige deres egen og deres Samtidiges religieuse
Overbeviisning, men tillige i denne villet opstille en uforanderlig Troesnorm
for de følgende Slægter, da maatte man beklage Menneskehedens Lod, som er
Svaghed selv i sin største Kraft; men nogen Forpligtelse for den
protestantiske Kirke kunde heraf ikke udledes. Clausen: Cath. og Protest.
P.~\textit{331}.\par
Der Protestantismus hat, was vom Anfange in ihm lag, offenbart, die
Entwicklung der religiøsen Selbständigkeit, die durch freie Liebe und Einsicht
am kirchlichen Vaterlande festhält so daß die Treue an Luthers Geiste zum
Abfalle von seinen Buchstaben \ldots\ geworden ist. Hase: Kirchengesch.
(\textit{1834}) §~\textit{545}.}.
% DOUBTFUL READING at the head of this page: at the end of the first line,
% after "muliggjøre", there is a small faint mark on the baseline. Read at 3x
% it is lighter and smaller than the page's other commas and detached from the
% word; the two OCR witnesses disagree about it (tesseract a comma, ABBYY a
% period). Transcribed WITHOUT it; the rejected alternative is
% "muliggjøre, saaledes".
% THE NOTE. One note, its mark a superscript figure 1 followed by ")", read at
% 3x; the body's reference mark on "Troeslære" is the same. TWO PARAGRAPHS,
% the second unmarked: the Clausen extract in Danish Fraktur, then the Hase
% extract in German Fraktur. NO quotation marks anywhere in it.
% *** COMPOSITOR'S SORT, transcribed as printed: "religiøsen" for "religiösen"
% --- the o carries the Danish SLASH, not the German umlaut. Settled at 3x by
% comparing it with the ä of "Selbständigkeit" and "festhält" in the same two
% lines, both of which have proper dots. The Danish compositor reached for his
% own ø sort.
% Also as printed: "festhält so daß" with no comma before "so"; and the
% elisions are four spaced points, uniformly rendered \ldots (the convention
% adopted at p. 61).
% TYPE, recorded but not reproduced: "Kirchengesch." and "§ 545." are set in a
% heavier (bold) face than the rest of the note --- as the source citations in
% the notes generally are in this printing. Left undistinguished here, as in
% the earlier batches.

% The § 17 head, read at 3x on PDF 96: "§ 17." centred in bold antiqua, and
% below it the argument in letterspaced Fraktur over two lines,
% "Den subjective Rationalisme. Den hellige Histories Vanhelligelse."
% This AGREES word for word with the Indhold (p. 87). No divergence to record.
\parag{17}{Den subjective Rationalisme. Den hellige Histories Vanhelligelse.}

Efterat man havde forladt den Vei, hvorpaa man før søgte at opstille Læren,
kunde heller ikke Troens objective Former, idetmindste for en Tid, undgaae at
blive ringeagtede, efterdi disse jo udgjorde det forældede Systems egentlige
Grundstof; begjerlig efter Leilighed til at anstille Undersøgelser og øve sin
Tænkning lader Menneskeaanden kun den Deel af den guddommelige Sandhed staae
urørt, der umiddelbart griber det religieuse Hjerte, og vender sig derpaa til
andre Egne, der kunne frembyde et nyt og rigere Stof for dens Bearbeidelse.
--- Troens Gjenstande træde i Baggrunden, medens denne Verdens Gjenstande ere
Betragteren nærmest; de sandselige Fornemmelser og Indtryk, gjennem
% --- p. 88 ---
% No footnote block on this page.
hvilke Adgangen til Erfaringens vide Mark staaer aaben, møde først det Blik,
der søger en umiddelbar Identitet af Subject og Object. Dog vil Betragteren
ikke længe tilfredsstilles ved at vanke om i den ydre, Bevidstheden fremmede
Virkelighed; han vil snart erkjende, at det kun er empiriske Ting, der her ere
Gjenstand for hans Tanke. Tænkningen er en Negation af Tingen; derfor vender
han Blikket bort fra den sandselige Natur og søger en fornuftig Verdensorden
for i denne Objectivitet at gjenfinde sig selv. Uagtet nu Empirismen og den
abstracte Rationalisme ere hinanden aldeles modsatte, idet den ene beskjæftiger
sig med enkelte, concrete Gjenstande, den anden derimod alene med almindelige
og abstracte, række de dog gjerne hinanden Haanden, thi Begges Elementer drage
Bevidstheden, og den abstracte Fornuft trænger til Empirien for af dennes Fylde
at kunne give sine hule Begreber Indhold. Hvad nu Troeslæren angaaer, da staae
de Begge i et vaklende Forhold til den; paa den ene Side ere de Troen
beslægtede: Troens Indhold er jo ikke Andet end Summen af de aandelige
Realiteter, det Samme gjælder ogsaa om de Gjenstande, vi møde i vor
Samvittighed, og endeligt fordre Begge, at den tænkende Aand i Objectiviteten
skal kunne gjenfinde sig selv, men dette fordrer ogsaa Troens Natur; --- paa
den anden Side staae de derimod i directe Modsætning til Troen. Denne udleder
nemlig alle sine Elementer fra den hellige, over vor Natur ophøiede Gud; disse
derimod gaae ud fra Subjectet og den skabte Natur og overskride aldrig denne
vor Verdens profane Grændser.

Paa Grund af dette tvivlsomme Slægtskab træde de i et doppelt Forhold til
Kirken. Saasnart de nemlig forraade deres Principers profane Natur bortvises de
i den krænkede Religieusitets Navn, og søge de med opløftet Banner og klingende
Spil at storme Christi Borg, falde de for den mægtigere
% --- p. 89 ---
% "%" below suppresses the newline: the printed page continues the previous
% page's sentence, and a bare newline here would double the interword space.
\setcounter{footnote}{0}%
Aands Sværd; lykkes det dem derimod formummede at indsnige sig i Borgen,
begynde de en hemmelig Undermineren. Dog er disse Fjenders Trædskhed ikke
saaledes at forstaae, at de ved eet eller andet Middel skulde have underkjøbt
samvittighedsløse (κεκαυτερισμένους τὴν ἰδίαν συνείδησιν) Theologer til at
% *** ERRATA APPLIED HERE (the Greek entry). ***
% Rettelser, p. 89 l. 5 fra oven: κεκαυτερισμένυος l. κεκαυτερισμένους.
% The UNCORRECTED printed reading was verified on PDF 98 at 3--10x before the
% correction was applied: after "μέν" the page sets upsilon, then omicron, then
% final sigma --- "κεκαυτερισμένυος". The correction is a transposition only;
% the stem is κεκαυτερισ- with epsilon (not eta) and no alpha, exactly as
% printed, and the acute stays on the epsilon of -μέν-. Set here as corrected.
% GREEK FOUNT VARIANTS NORMALISED, per the header's standing rule: every kappa
% on this line is the CURSIVE kappa (the x-shaped variant, U+03F0) and is
% typed as plain kappa U+03BA.
% Accents and breathings are reproduced as printed: τὴν with a grave, ἰδίαν
% with a smooth breathing on the first iota and an acute on the second,
% συνείδησιν with the acute on the iota of -ει-. Nothing on this line needed
% the -ου- diphthong normalisation.
gjøre Religionens Sandhed mistænkt; tvertimod handle dette Slags Theologer i
Almindelighed af fuldeste Overbeviisning og med freidigt Aasyn. De Kampe, der
heraf opstaae, føres derfor heller ikke umiddelbart om, hvo der er Christi Ven
eller Fjende, men om, hvad der hører til Religionens Væsen, og hvad der kun
hører til dens Phænomener. Thi de, der blive staaende ved den ældre
Læremethode, urgere det overnaturlige Moment i de hellige Phænomener, idet de
erklære, at disse constituere de objective Troesformer og derfor ikke kunne
sættes i Klasse med de sædvanlige Begivenheder, der ere vor Dom underkastede,
ligesaalidt som de uden største Skade for Religionen selv kunne forandres eller
forflygtiges\footnote{De Stridigheder, der fandt Sted i den reformeerte Kirke i
Anledning af den hellige Skrifts kritiske Behandling, afgiver saa at sige et
Forbillede for alle lignende Debatter. Saaledes erklærede Buxdorf Ludv.
Capellus for en \textit{\emph{novator, propheta, videns}, qui somniis
indulgeret}. Capellus derimod svarede \textit{(advers. Censorem Buxdorfium
\ldots\ brevis oeconomia): Satis, mihi fuisset id ostendisse adversus Censorem
meum \ldots\ nisi illi ipsi et aliis multis altius animo insedisset et infixa
esset hæc opinio: Actum esse de Dei verbo ejusque apud nos auctoritate, atque
ex consequenti de hominum salute, (quæ non aliunde quam ex Dei verbo pendet),
si apiculi isti non statuantur esse a Mose et Prophetis, libris suis, Deo eos
docente, additi et adseripti.} Og med Rette lægger Georgius Pasor Sebastianus
Pfoechenius, der er berømt af de puristiske Stridigheder, følgende Ord i
Munden:\par
\hspace*{2em}\textit{Diligo Philologum, qui eviscerat abdita verba,}\par
\hspace*{2em}\textit{Jure novum Chrtsti sed mage \emph{fedus} amo,}\par
\hspace*{2em}\textit{Ergo \emph{novo} quisquis de \emph{federe} tenuia dicit,}\par
\hspace*{2em}\textit{En pupillam oculi vulnerat ille mei.}\par
\centerline{\rule{0.16\linewidth}{0.4pt}}\par
\hspace*{2em}\textit{Falx critica in quosvis per me grassetur. At unis}\par
\hspace*{2em}\textit{A Bibliis fas est abstinuisse manus.}}.
% *** THIS NOTE CROSSES A PAGE BOUNDARY. IT IS COMPLETE HERE. ***
% Note 1) of p. 89 begins in p. 89's footnote block and runs on into the TOP of
% p. 90's footnote block, above p. 90's own rule-and-mark: the couplet
% "Ergo novo ... ille mei.", then a broken rule, then "Falx critica ...
% abstinuisse manus." Those six lines were read off PDF 99 and are included
% above. p. 90's own note 1) begins below them.
% *** PRINTER'S DEFECTS in this note, both transcribed as printed, both read
% at 3x: "Chrtsti" for "Christi" (a t where the i belongs), and "adseripti"
% for "adscripti" (an e where the c belongs).
% Also as printed and NOT corrected: "Buxdorf" / "Buxdorfium" for Buxtorf,
% "reformeerte", and "tenuia" (for tenuis).
% TYPE. The whole quotation is antiqua Latin, hence \textit{}; within it
% "novator, propheta, videns", "novo", "federe" and "fedus" are set in a SLOPED
% antiqua against the upright antiqua of the rest, and "novum" beside "fedus"
% is upright. \emph{} inside \textit{} reverts to upright and so preserves that
% contrast, following the precedent set in p. 72's note.
% THE SEPARATOR between the second and third couplets is a PRINTED BROKEN RULE
% --- about nine short strokes, centred, roughly a sixth of the measure. It is
% not scanner damage: it was measured row by row on PDF 99 (a clean band of ink
% at y=3540 with white above and below) and it is distinct from the footnote
% rule at the left margin higher up the page. Rendered here as a plain centred
% rule; a macro for the broken form would be an improvement.

% p. 90 OPENS A NEW PARAGRAPH (its first line is indented); hence the blank
% line above the marker.
% --- p. 90 ---
\setcounter{footnote}{0}
Grundfeilen i Undersøgelsen ligger altsaa deri, at man ikke slutter fra den sig
aabenbarende Guds Sanddruhed til Sandheden af de Phænomener og Former, gjennem
hvilke Aabenbaringen skeer, men at man gjennem en Undersøgelse af Phænomenerne
troer at kunne bestemme, hvilke Momenter der ere væsentlige for Aabenbaringen.
De urgere derfor i Tillid til den menneskelige Undersøgelses Ufeilbarhed den
naturlige Side af Phænomenerne og erklære det for en \textit{petitio
principii}, hvis Nogen i Troens Navn holder fast ved de enkelte Former og
Phænomener, fordi man kun ved en Undersøgelse af dem kan komme til at bestemme
Troens Objecter. Og det er paa denne Overbeviisning, at Sandheden ikke kan
modsige Sandheden, den empiriske Sandhed ikke ophæve Religionens Sandhed, at
denne Skole støtter sin Auctoritet\footnote{\textit{Sunt tamen et pii et docti
viri, qui cum vix habeant, quod tam manifestæ veritati opponant, pergunt
nihilominus periculosas inde metuere consequentias. Quasi vero a veritate ullum
jure metuendum sit veritate periculum! Verum certe vero semper consonat fertque
opem, nunquam repugnat aut officit \ldots\ Ergo candido et christiano pectore,
et forti constantique animo dignius esse existimavi, fateri in hac causa, quod
res est, quam cum nonnullis id oppugnare, quod sine rubore ab eo, cui ingenua
est frons, negari non potest. Ludv. Capellus. Arcan. punct. Præfatio ad
lect.}\par
At Tænkningen, som Erkjendelsen selv, er algyldig, efterdi den har Troen som
saadan i sig, har Hegel tilfulde beviist i sin Phænomenologie des Geistes
(Werke P.~\textit{410}): Von dieser Seite, daß beide wesentlich dasselbe sind
und die Beziehung der reinen Einsicht durch und in demselben Elemente geschiet,
ist ihre Mittheilung eine unmittelbare, und ihr Geben und Empfangen ein
ungestörtes Ineinanderfließen \ldots\ Die Mittheilung der reinen Einsicht ist
deswegen einer ruhigen Ausdehnung oder dem Verbreiten wie eines Duftes in der
wiederstandslosen Atmosphäre zu vergleichen \ldots\ ein unsichtbarer und
unbemerkter Geist, durchschleicht sie die edlen Theile durch und durch, und hat
sich bald aller Eingeweide und Glieder der bewußtlosen Götzen gründlich
bemächtigt, und ”an einem schönen Morgen giebt sie mit dem Ellbogen dem
Kammeraden einen Schubb und Bautz! Baradautz! der Götze liegt am Boden.”}.
% *** THIS NOTE ALSO CROSSES A PAGE BOUNDARY. IT IS COMPLETE HERE. ***
% Note 1) of p. 90 has TWO paragraphs, the second unmarked: the Capellus
% extract in antiqua Latin, then the Danish-and-German paragraph on Hegel. The
% Hegel quotation runs on into the whole of p. 91's footnote block ("ihre
% Mittheilung eine unmittelbare ... am Boden.”"), read off PDF 100 and included
% above. p. 91 therefore has NO note of its own and takes NO
% \setcounter{footnote}{0}.
% QUOTATION MARKS: the German quotation inside this note uses the book's
% standing RAISED " (U+201D) at BOTH ends, read at 3x --- not the low-high
% pair (U+201E ... U+201C) that the
% notes on pp. 48 and 70 use. Contributes 2 to check.py's ” parity count.
% As printed and not corrected: "geschiet" (for geschieht), "wiederstandslosen"
% (for widerstandslosen), "Ellbogen", "Schubb", "Bautz! Baradautz!".
% The body's "den sig aabenbarende" has a speck, not a comma, after "sig";
% ABBYY reads a comma there and is wrong.
% EMPHASIS: none on this page. spacing.py's RUNs "væsentlige for" and "fast
% ved" and its "petitio"/"Auctoritet" flags were all checked on the image and
% are justification artefacts; "petitio principii" is italic because it is
% ANTIQUA, not because it is letterspaced.

% p. 91 OPENS A NEW PARAGRAPH (its first line is indented).
% --- p. 91 ---
% NO \setcounter HERE ON PURPOSE: p. 91's footnote block carries only the tail
% of p. 90's note 1) --- there is no mark and no note of its own on this page.
Under et saadant Forhold kunde der naturligviis ikke længere være Tale om at
fastsætte Troens positive Former; Aabenbaringsphænomenerne, der hidrørte fra
den graa Oldtid, og Formerne, der i Tidernes Løb alt vare fastsatte, bleve kun
ydre Gjenstande for Granskeren, og Resultaterne af dennes Undersøgelser kunde
derfor kun bestaae deri, at udfinde, hvilke Begivenheder der havde kunnet finde
Sted, hvilke der virkeligt havde fundet Sted, og hvilke der endeligt vare at
betragte som Producter af en overtroisk Phantasie. Troen var da ikke længere
det constitutive Princip i Theologien, og uagtet det, der tages til Rettesnor,
ifølge sin Natur ikke kan optage eller anerkjende Andet, end hvad der ligger
indenfor dets egne Grændser, lod man dog Naturen være Rettesnor for og Dommer
over de overnaturlige Ting. Man kan derfor ikke undre sig over, at Religionens
positive Elementer bleve færre og færre og omsider ganske forsvandt. Medens
saaledes Empirismen beviste, at de hebraiske Vocaler hverken vare opfundne af
Adam eller af Moses, ja ikke engang af Esras, at det ny Testamentes Græsk ikke
var classisk reent, at den hellige Skrifts Efterretninger baade vare
selvmodsigende og stridende mod troværdige Profanskribenters Vidnesbyrd, og at
mange af de enkelte Bøger ingenlunde vare authenti\-%
% "uagtet det, der tages til Rettesnor" --- "det," with the comma, read at 3x;
% ABBYY reads "der" and is wrong. spacing.py's RUN "clas reent" was checked on
% the image: "classisk reent" is NOT letterspaced.
% --- p. 92 ---
% The trailing "%" on the next line is load-bearing: p. 91 ends mid-word
% ("authenti\-%") and a newline after \setcounter would put a space inside
% "authentiske".
\setcounter{footnote}{0}%
ske\footnote{Smlgn. Selmers Abhandlung von freier Untersuchung des Kanon.},
godtgjorde paa den anden Side den subjective Rationalisme, (hvad der netop
viser det gjensidige Forhold mellem Empirie og Tænkning), det Urimelige i at
antage en overnaturlig Inspiration\footnote{Smlgn. Eckermann: Beweis, daß der
biblische Begriff der Offenbarung kein andrer, als Offenbarungsbegriff der
gesunden Vernunft sey.}, fremsatte Tvivl om Miraklernes
Mulighed\footnote{Die Verächter der Bibel nehmen gerade die gefährlichsten
Waffen, womit sie das Ansehen, die Ehrwürdigkeit und Glaubwürdigkeit der Bibel
bestreiten, aus der Behauptung der Vertheydiger derselben her, daß ihre
Verfasser sich einer unmittelbaren und übernatürlichen Offenbarung und
Eingebung rühmen. Die Verfasser der Bibel werden deswegen als eine Reihe von
Schwärmern angeklagt und herabgewürdigt. Samme: Handbuch der christ.
Glaubenslehre \textit{1} Theil P.~\textit{410}.\par
Aber sagt man, wenn ein Mann von erprobter Frömmigkeit, Rechtschaffenheit und
Wahrheitsliebe \ldots\ versicherte, daß Gott dies unmittelbar und übernatürlich
bewirke \ldots\ sollten wir denn das einem solchen Manne nicht glauben? Die
Vernunft antwortet: Nein! Es ist unmöglich, daß ein Mensch unmittelbaren
Wirkungen Gottes erkenne! Auch der redlichste Mensch irrt, wenn er etwa
dergleichen zu erkennen glaubt \ldots\ S.~S. P.~\textit{456}.}, og erklærede,
at Gjenfødelsen ene og alene var at tænke som en subjectiv Forandring af
Menneskets fri Villie\footnote{Saaledes omskriver Paulus (Comment. over d. N.
Test.) Joh. \textit{3, 6} τὸ γεγεννεμένον ἐκ του πνεύματος: der Mensch, welcher
durch seine Geisteskraft, Einsicht und Wollen des Guten sich selbst zeugt
\ldots\ hvilket han oplyser videre ved at henvise til Seneca, \textit{de brev.
vitæ}: ”\textit{Solemus dicere, non fuisse in nostra potestate, quos sortiremus
parentes \ldots\ nobis vero ad arbitrium (!) uasci licet.}”}.
% FOUR notes on this page, marked 1) 2) 3) 4) in sequence; note 3) has two
% paragraphs, the second unmarked. All four are complete on p. 92.
% *** PRINTER'S DEFECT, transcribed as printed: "Selmers" in note 1) for
% "Semlers" (J. S. Semler's "Abhandlung von freier Untersuchung des Canon").
% Read at 3x; both OCR witnesses agree on the l. Not in the Rettelser.
% *** PRINTER'S DEFECT, transcribed as printed: "uasci" for "nasci" in the
% Seneca quotation of note 4) --- a turned n, read at 3x, the same defect as
% p. 2 "absolnt" and p. 3 "Tiugen". The "(!)" beside it is the author's own.
% QUOTATION MARKS in note 4): the book's standing RAISED ” at BOTH ends, read
% at 3x. Contributes 2 to check.py's ” parity count.
% GREEK in note 4), read at 7x on PDF 101:
%   * every kappa is the CURSIVE variant (U+03F0) and is normalised to plain
%     kappa (U+03BA), per the header.
%   * "γεγεννεμένον" is set with EPSILON where the received text has eta
%     (γεγεννημένον). Reproduced as printed.
%   * *** LIGATURE EXPANDED: the word after ἐκ is printed as tau followed by
%     the OU-LIGATURE (U+0223) --- a single sort, with no accent of any kind over
%     it. It is expanded here to the two letters "του", exactly as the fount's
%     st- and fi-ligatures are expanded elsewhere in this edition, and the
%     circumflex of the received "tou" is NOT supplied. The raw U+0223 is
%     deliberately
%     not typed: the preamble does not handle it.
% --- p. 93 ---
% "%" below suppresses the newline: the printed page continues the previous
% page's sentence, and a bare newline here would double the interword space.
\setcounter{footnote}{0}%
udslukkedes der paa den overnaturlige Himmel. Den myndige Fornufts
Skarpsindighed opdagede, at Guds Røst og Englenes Aabenbarelser., hvorom
Skriften saa ofte beretter, vel var Noget, der kunde tiltale de Umyndige (οἱ
νήπιοι), men i Virkeligheden som oftest kun vare Følger af Naturbegivenheder,
Torden og Lynild eller Andet lignende, og Producter af en ophidset Phantasie;
saadanne Aabenbaringer bestaae da ikke deri, at de guddommelige Mysterier
aabenbare sig i den forklarede Natur, men deri, at ganske naturlige
Begivenheder indhylles i den menneskelige Phantasies Taageslør; man behøver da
kun at drage Sløret bort for et gribe den nøgne Sandhed\footnote{Fortrinlige
Exempler paa en saadan Aabenbaring findes hos G. L. Bauer (Mythologie des alten
und neuen Test. Leipzig \textit{1802}), og saa sindrige ere disse, at vi ei
kunne afholde os fra at meddele den velvillige Læser nogle af dem. Saaledes om
Moses P.~\textit{271}. \textit{1} Band: Moses stellte sich vor, ein brennender
Dornbusch sey Jehova. \emph{Kann das Geschichte seyn}? Kann der brennende Busch
würklich Gott, oder nur Gott in demselben gewesen seyn? Ist Gott ein Feuer? ---
P.~\textit{272}: Ein Blitzstrahl fuhr in einen Dornbusch, als Moses weidete am
Horeb mit seiner Heerde: Es war aber ein kalter Strahl, der nicht zündete
\ldots\ Nach Hezel wars electrisches Feuer.}. Gud har saaledes ikke med lydelig
Røst fra Tornebusken kaaret Moses til Sendebud trods hans Vægring, nei! Moses
indbildte sig at høre Jehovah i den af Lynilden antændte Tornebusk, fordi han
havde Hovedet fuldt af store Planer; og det er heller ikke Gud, der gjennem
Moses har dicteret Loven, nei, det var Moses, der i Guds Navn gjorde det, da
det er under Guds Værdighed at tale i articulerede Toner\footnote{S.~S.
P.~\textit{225} \textit{2} Band: Und daß Gott selbst vom Himmel geredet habe,
ist auch schwer zu glauben, wenn reden heissen soll articulirte Töne, Worte der
damaligen palästinischen Landessprache aussprechen; og P.~\textit{296}.
\textit{1} Band: Hätte Gott selbst gesprochen, sollten denn die Israeliten so
gar dumm, unempfindlich und hartnäckig gewesen seyn, daß sie gleich darauf sein
Verbot des Bilderdienstes vergassen? Und für die zehn Gebote, welche zwar gut
und meistens Natur= und Vernunftgebote sind, zu deren Erfindung aber kein
göttlicher Verstand gehört, sollte eine Gesetzgebung vom Himmel nöthig sein?}.
% *** PRINTER'S DEFECT, transcribed as printed: "Aabenbarelser.," --- a full
% stop AND a comma, both on the baseline, read at 3.5x. Not a speck: the two
% marks are of equal weight and are separately formed.
% *** PRINTER'S DEFECT, transcribed as printed: "for et gribe" for "for at
% gribe". Both OCR witnesses agree and it is plain on the image.
% GREEK: (οἱ νήπιοι), read at 6--10x. The eta of νήπιοι carries an acute. The
% mark over the iota of οἱ is a small comma-form whose DIRECTION could not be
% settled even at 10x; οἱ (rough breathing) is set, as the grammar requires,
% and the uncertainty is recorded here rather than silently resolved.
% EMPHASIS: "Kann das Geschichte seyn" in note 1) is letterspaced across the
% four words, read at 1.6x and unmistakable; the question mark is not
% letterspaced and is set outside the \emph{}.
% *** NOTE 2) CROSSES A PAGE BOUNDARY. IT IS COMPLETE HERE. *** Its last
% sentence ("hartnäckig gewesen seyn ... nöthig sein?") stands at the TOP of
% p. 94's footnote block, above p. 94's own mark, and was read off PDF 103.
% The Fraktur double-hyphen of "Natur= und Vernunftgebote" is kept as "=".

% p. 94 OPENS A NEW PARAGRAPH (its first line is indented).
% --- p. 94 ---
\setcounter{footnote}{0}
Ikke blot det gamle Testamentes Engleaabenbarelser, men ogsaa det nyes ere at
betragte som menneskelige Illusioner. Den sunde Fornuft lærer os jo, at da der
ingen Engle lade sig tilsyne for os, kan heller ikke Gabriel have forkyndt
Maria, at hun skulde føde Christus; og da det er vist, at der ikke kan
fremkomme Noget af Intet, og hvad der er naturligt, desuden ikke kan være til
Skam for Nogen, gjøre vi rettest i at tilstaae, at Jesus blev avlet paa
naturlig Maade af Joseph og Maria\footnote{Soll die Erscheinung nur von einer
Vision oder innern Intuition oder den eigenen Gedanken der Maria erklärt
werden, so fällt zwar der Einwurf von der persönlichen Erscheinung eines Engels
weg, aber dafür ist unaufgeklärt, wovon denn die Maria ist schwänger geworden.
Durch die bloße Einbildung hat sie doch wohl nicht schwanger werden können
\ldots\ Weit natürlicher ist's, Jesum für einen Sohn des Josephs und der Maria,
unsere Sage aber für eine später entstandenes Raisonnement über die Entstehung
des Messias zu halten S.~S. P.~\textit{192}. \textit{1} Band.\par
\textit{Neque vero fieri potest, ut religionis christianæ auctori major dignitas
aut excellentia concilietur, si vitæ ejus partes et casus non secundum
naturalem rerum humanarum viam a Deo constitntam, sed supranaturali sive
miraculosa quadam ratione evenisse dicantur. Nam si hoc esset, consequeretur,
naturales rerum in orbe gerendarum leges, a divina sapientia et omnipotentia
constitutas, quasi minorem quandam dignitatem habere, hominumque præstantiam ex
iis potius, quæ ipsis accidunt, quam ex ipsorum consiliis et factis debere
existimari. Wegscheider: Institutiones Theol. Halæ 1833. p. 442 not. a.}}. Men
hvorledes kan den, der er født paa naturlig Maadr, anvende overnaturlige
Kræfter i sin Tjeneste! --- Johannes den Døber bildte sig ligeledes ind at see
den hellige Aand, medens han dog i Virkeligheden ikke saae Andet, end en
almindelig Due, der tilfældigviis flei
% *** THIS NOTE CROSSES A PAGE BOUNDARY. IT IS COMPLETE HERE. ***
% p. 94 carries ONE note, marked 1), in two paragraphs, the second unmarked:
% the German extract, then the Latin one. The Latin runs on into the TOP of
% p. 95's footnote block ("cidunt, quam ex ipsorum consiliis ... not. a."),
% read off PDF 104 and included above. p. 95's own note 1) begins below it.
% *** PRINTER'S DEFECTS on this page, all transcribed as printed:
%   "Maadr" for "Maade" (an r for the e), read at 3x;
%   "flei" for "fløi" at the foot of the page --- the vowel is a plain e with
%     no slash, read at 4x, and the word is completed by "forbi" at the head of
%     p. 95, so the sense is "fløi forbi";
%   "constitntam" for "constitutam" in the Latin note (a turned u), read at
%     1.1x on a clean impression;
%   "schwänger" with an umlaut in the German note, against "schwanger" two
%     lines later.
% The Latin note is antiqua throughout, hence \textit{}; the German note is
% Fraktur and is not italicised. The date is printed 1833, not 1832 or 1835 as
% the two OCR witnesses variously give it.
% --- p. 95 ---
% "%" below suppresses the newline: the printed page continues the previous
% page's sentence, and a bare newline here would double the interword space.
\setcounter{footnote}{0}%
forbi\footnote{Der heilige Geist kann nicht wie eine Taube mit sichtbarer
Gestalt herabgefahren seyn, denn er ist Gottes unsichtbare überall wirksame
Kraft, wie hätte dieses sollen sichtbar werden? Bauer: \textit{l.~l.}
P.~\textit{225}. \textit{2} Band.}. Og vistnok troede Apostlene, at Christus,
forklaret paa Bjerget, havde talt med Moses og Elias; men betragter man Sagen
lidt nærmere, er det let at indsee, at Disciplene hiin Nat blot havde sovet og
drømt og derpaa, da de vare blevne vækkede, gnedet sig i Øinene og seet
Christus belyst af Morgensolen\footnote{Samme P.~\textit{213}. \textit{2}
Band.}. Man vil derfor ogsaa ved en Betragtning af Sagens sande Sammenhæng
indsee, at den hellige Historie ikke er meget forskjellig fra den profane. Thi
ogsaa Hedningene troede i den pludselige Torden at høre Guds Røst; ogsaa deres
Lovgivere udgave deres Love for at være dicterede af Gud; ja ogsaa disse toge
Varsler af Fuglenes Flugt (som Johannes den Døber af hiin Due), ligesom de
ogsaa gjerne troede paa jomfrufødte Heroer med guddommelige Fædre. Desuagtet
maa man ikke ringeagte den hellige Skrifts Historier; thi medens mange
Mennesker opvaagne af en Skindød, uden at man kan see Forsynets Hensigt
dermed, maa man love det guddommelige Forsyn, fordi det øiensynligt har styret
det saaledes, at det fuldkomneste og frommeste Menneske, der er fremstillet os
som et Mønster til Efterfølgelse, vendte tilbage til Livet for at gjenopreise
sine Disciples sjunkne Mod; og om vi end neppe kunne slutte fra hans
Opstandelse til vore Sjeles Udødelighed, kunne vi
% *** PRINTER'S DEFECT OF A KIND NOT YET MET IN THIS BOOK: p. 95's footnote
% block carries FOUR notes, marked 1) 2) 3) 4), but the page's TEXT carries
% only TWO reference marks --- 1) on "forbi" and 2) on "Morgensolen". Every
% line of the page was re-read at 2--4x hunting for the missing marks
% (including "profane.", "Guds Røst;", "hiin Due),", "Fædre.", "Skindød" and
% "Christus,", where a speck misled ABBYY): there are none. Notes 3) and 4)
% stand in the block with nothing in the text to call them.
% Transcribed as printed. They are set below with \stepcounter + \footnotetext,
% which puts them in the footnote block under their printed numbers and leaves
% no mark in the text --- exactly what the page does.
\stepcounter{footnote}\footnotetext{S.~S. P.~\textit{230}. \textit{2} Band.}%
\stepcounter{footnote}\footnotetext{S.~S. P.~\textit{192}. \textit{1} Band.}%
% --- p. 96 ---
% "%" below suppresses the newline: the printed page continues the previous
% page's sentence, and a bare newline here would double the interword space.
\setcounter{footnote}{0}%
dog nu have den Fortrøstning til Gud, at ogsaa vi, dersom vi skulde bortrives
ved en altfor tidlig Skindød, hvis ellers Gud vil, ville vende tilbage til
dette Liv for derved at fremme Erkjendelsen af en fornuftig Verdensorden; og
uagtet Christus ikke paa synlig Maade opfoer til Himmelen, men kun skiult af en
Sky ad en ubekjendt Vei steg ned af Bjerget, har han dog saaledes faaet en
langt herligere Apotheose, end Romulus og Drusilla\footnote{Derfor siger Paulus
ikke uden religieus Emphasis i Anledning af Luc. \textit{21, 51}: Welch ein
Gegensatz gegen dergleichen Apotheosen lag in Jesu ἀναλήψις, auf welche \ldots\
das alte Römergesetz würdig anzuwenden ist: \textit{Divos et eos, qui coelestes
semper habiti, \emph{colunto, et illos quos endo coelo merita locaverunt},
Herculem, Liberum, Æsculapium, Castorem, Pollucem, Quirinum, et Horatii}:\par
\hspace*{2em}\textit{”Virtus, recludens immeritis mori}\par
\hspace*{2em}\textit{Coelum, negata tentat iter via:}\par
\hspace*{2em}\textit{Coetusque vulgares et udam}\par
\hspace*{2em}\textit{\emph{Spernit} (!) humum fugiente penna.”}\par
(Comment. üb. d. N. T. \textit{3} Theil. Leip. \textit{1805}).}, og med Rette
beundre vi hans Dyd og stræbe at følge hans Exempel; kun vogte man sig for at
tilbede ham, som var han en Gud.
% *** THE PRINTED MARKS DO NOT AGREE ON THIS PAGE. *** The body's reference
% mark on "Drusilla" is a superscript figure 1 --- read at 4x, unmistakably a
% 1 --- but the single note in the block below is marked 3). Both were
% re-checked at 4x. The note is set here as an ordinary \footnote, so that it
% prints 1) and matches the body, following the precedent of p. 12's *) note;
% the printed 3) is recorded here and not reproduced.
% As printed: "skiult" with a dotted i and no descender (read at 3.5x, against
% the descending Fraktur j of "ubekjendt" on the same line), and "opfoer".
% "Luc. 21, 51" is the printing's own reference; not corrected.
% QUOTATION MARKS in the note: the book's standing RAISED ” at BOTH ends of the
% Horace quotation. Contributes 2 to check.py's ” parity count.
% TYPE in the note: the Latin is antiqua, hence \textit{}; within it "colunto,
% et illos quos endo coelo merita locaverunt" is SLOPED against the upright
% antiqua of the rest, and "Spernit" is sloped AND letterspaced. Both are set
% with \emph{} inside \textit{}, which reverts to upright and preserves the
% contrast, as in p. 72's and p. 89's notes; the letterspacing of "Spernit" is
% recorded here rather than reproduced.
% The Greek ἀναλήψις is as printed, with the smooth breathing on the alpha and
% the acute on the eta. No fount variant needed normalising in it.
% This note is set several points smaller than the other notes in this batch.

% The § 18 head, read at 3x on PDF 105: "§ 18." centred in bold antiqua, the
% argument below it in letterspaced Fraktur, "Tro og Viden i den kritiske
% Theologie." This AGREES word for word with the Indhold (p. 96). No
% divergence to record.
\parag{18}{Tro og Viden i den kritiske Theologie.}

Efterat Kriticismen saaledes tilstrækkeligt havde udviist, hvad den førte i sit
Skjold, skulde man vente, at den enten vilde være bleven udstødt af Kirken,
eller selv have kuldkastet Kirkens Bygning; og dog var ifølge Sagens Natur ingen
af Delene mulig. Vel indeholder Kriticismen en Negation af Troen, men
Spørgsmaalet er, i hvilket Punkt Negationen ligger. Idet Empirismen og den
subjective Ratio\-%
% BATCH SEAM: p. 96 ends mid-word. "Ratio-" is completed by "nalisme" at the
% head of p. 97, so the discretionary hyphen \-% joins them without printing a
% hyphen. NOTE FOR THE SPLICER: transcription.tex keeps a BLANK LINE between
% consecutive "[text to be added]" markers, and a blank line here would break
% the paragraph in mid-word. The blank line between this fragment and the
% pp. 97--108 fragment must be removed at splice time. (The same seam exists
% between every earlier pair of batches; and note that pp. 1--12 joined broken
% words across page markers with "\-%" while pp. 25--36 onward used a bare
% hyphen, which prints a spurious "- " --- worth normalising in one pass.)

% --- p. 97 ---
\setcounter{footnote}{0}
% JOINT WITH p. 96, verified on PDF 105 at 0.75x: p. 96's last body line ends
% "...Empirismen og den subjective Ratio=" and this page opens "nalisme". The
% footnote block of p. 96 ends on p. 96 (its last note line stops short at
% x1901 of a 2600-wide measure); nothing runs over onto p. 97, whose note
% block is exactly the five lines of its own note 1).
nalisme urgere de verdslige, om man saa tør sige, og reent
menneskelige Momenter, negere de tillige Realiteten af Troens Objecter,
\dsi{} erklære dem for at ligge udenfor deres egen Erkjendelses Grændser. Men
denne Negation er endnu ikke en Paastand om, at den slet ikke er til; thi
ligesom de indenfor deres egen Erkjendelses Grændser paa positiv Maade
fastholde de Gjenstande, der høre til deres Sphære, saaledes kunne de ogsaa
meget vel anerkjende Tilværelsen af en udenfor deres Erkjendelses Grændser
liggende Realitet. --- Men herved forberedes en ny Forsoning mellem Tro og
Viden\footnote{Hvorledes Kant, Jacobi, Schleiermacher, Friis og de eklektiske
Theologer have paa forskjellig Maade modificeret disse den kritske Fornufts
Theoremer, tilkommer det os ikke her nøiere at udvikle, da vor Undersøgelse
kun angaaer Kriticismens Natur i Almindelighed.}. Man kan vel ikke tænke
% PRINTER'S DEFECT in the note, transcribed as printed: "kritske" for
% "kritiske", broken over the line as "kri= | tske". Read at 3x: the second
% syllable is t + the long-s + k ligature + e, with no i at all. NOT in the Rettelser
% list; not corrected. TYPE: the whole note is Fraktur, the four names
% included, so no \textit{} anywhere in it.
Phænomenernes Sum uden at antage en Væsenseenhed af dem, ikke tænke det
Endelige uden at tænke det Uendelige, ikke tænke det Synlige uden at tænke det
Usynlige; men hvad der kan begribes, maa henhøre til den endelige Shære, thi
% PRINTER'S DEFECT, transcribed as printed: "Shære" for "Sphære", the p simply
% absent. Read at 2x: the word is S-h-æ-r-e, and "Sphære" is set correctly
% twice elsewhere on this same page. NOT in the Rettelser list.
hvor der er Begriben, der maa man ogsaa kunne adskille Momenterne, og hvor
Momenterne kunne adskilles, der maae de ogsaa være endelige. Derfor
underkaster Erkjendelsen sig den hele Sphære af endelige Gjenstande og
Momenter, men gjør villigt Afkald paa enhver Begriben af den uendelige og
usynlige Tilværelses Rige, ihvorvel den af sin egen Sphæres Begrændsning og
Indskrænkning med Vished slutter, at denne maa have sin sidste Tilværelses
Grund i hiint. Skal derfor Adgangen til hiint Rige ikke være Mennesket
tillukket, maa dette være i Besiddelse af en anden Modtagelsesevne, end
Erkjendelsen; men saaledes fremkommer der en ny Adskillelse i den menneskelige
Bevidsthed selv, mellem Erkjendelsen paa den ene Side og Villien og Følelsen
paa den anden Side.
% "Gjenstande" here and in l. 6 was read "Gjeustande" by BOTH OCR witnesses,
% which would have made it a fourth turned n. It is not: at 2x both letters
% close at the top, like the n of "endelige" beside it. Set as "Gjenstande".

% --- p. 98 ---
Ved \emph{Villien} er Mennesket istand til med sin hele Subjectivitets Kraft at
hengive sig under det evige og uendelige Rige, og ved denne Hengivelse bevares
baade det endelige og uendelige Moment; thi idet det endelige Subject bliver
underkastet det uendelige Subjects Herredømme, tilfredsstilles dette, medens
paa den anden Side det endelige Menneske tilfredsstiller sig selv derved, at
det med fri Villie hengiver sig selv. --- Ved den rene, fra alle particulaire
Affectioner luttrede \emph{Følelse} bliver den samme Forsoning det religieuse
Gemyt til Deel. I Følelsen er der nemlig en umiddelbar Identitet af det
subjective og objective Moment; man behøver da kun at abstrahere fra alle
Particulariteter og vil da have det almindelige Subject og det uendelige
Object, af hvilket hiint er absolut afhængigt. --- Men heraf følger en anden
særdeles vigtig Adskillelse, en Adskillelse nemlig mellem den religiøse Idee
selv og den Form, hvori den fremtræder; og af saa stor Betydning er denne, at
man fra den kan udlede den store Forandring, der er foregaaet i den nyere
Theologie.
% EMPHASIS on this page, settled on the image at 0.6x: "Villien" and "Følelse"
% are letterspaced; "Idee" (which spacing.py flagged) is NOT --- it is set
% solid, and so is "Ved" before each of the two real ones.

Da nemlig Former og Phænomener efter deres Natur henhøre til det Endeliges
Rige og derfor maae kunne begribes, kan man i Sandhed omvendt slutte, at Gud
ikke kan begribes af os, og at det ikke kan være Religionens Hensigt at gjøre
Gud til Gjenstand for vor Begriben; vi maae da skjelne mellem Gud, som han er
i sig, uendelig og ubegribelig, og de Former og Phænomenener, i hvilke han har
% PRINTER'S DEFECT, transcribed as printed: "Phænomenener" for "Phænomener",
% a syllable set twice. Read at 0.83x; both OCR witnesses agree, and the word
% is spelt correctly four lines above and two lines below. NOT in the
% Rettelser list.
aabenbaret sig til vor Frelse. Det er altsaa i dette Punkt, at den nyere
Theologie afviger fra den ældre. I den ældre Theologie finder Gud sin sande
Bestemmelse igjennem sine Phænomener; han er selv tilstede i dem; de maae
derfor modtages, som de ere os givne i Aabenbaringen; --- i den kritiske
Theologie derimod kan den uendelige Gud ikke indeholdes i eller indskrænkes
ved nogen Form; derfor maae de Phæno\-%
% --- p. 99 ---
\setcounter{footnote}{0}
mener og Former, gjennem hvilke de guddommelige Ting ere os aabenbarede, kun
betragtes fra Menneskets Standpunkt. Det synlige og det usynlige,
Erkjendelsens og Troens Rige kunne altsaa, ihvormeget de end divergere, ikke
holdes fuldkomment adskilte fra hinanden; de have, ligesom Frankrig og
Spanien, Spanien og Portugal, de samme Grændser tilfælleds, og da nu
Rumsbegrebet ikke passende kan anvendes paa dette Forhold, sige vi med Rette,
at der i det endelige Rige findes Symboler og Tegn paa, Mærker og Spor af det
uendelige Rige. --- Først ved Fastsættelsen af disse Betingelser udvirkede
Kriticismen efter saa mange heftige Kampe en Fred imellem \emph{Troen} og
\emph{Fornuften}, ifølge hvilken Alt, hvad der henhører til Erfaringen og
Logiken, skulde være den kritiske Fornuft hjemfalden, dog paa den Betingelse,
at den skulde afholde sig fra Alt, hvad der paa nogen Maade kunde krænke
Troen; paa den anden Side maatte Troen nøies med sit eget Rige, ikke misunde
den frie Fornuft dens jordiske Goder og ikke fordre Andet af den, end at den
skulde være beredt til dens Forsvar og vaage over dens Tarv\footnote{Denne
Pagts Ukrænkelighed skildrer De Wette i følgende stærke Udtryk: Der Verstand
muß wissen, was er zu bestreiten hat, und was nicht; er muß nicht, vom blinden
Wiederspruchsgeist getrieben, seine Waffen gegen dasjenige richten, was er,
als hoch über ihm stehend, achten und schonen soll; und der Glaube, seiner
hohen Abkunft sich bewußt, muß es verschmähen, sich in die niedere Sphære des
Verstandes herabzusenken und diesen freigebornen Sohn der menschlichen
Vernunft in Sklavenfesseln legen zu wollen. Herrschen soll der Glaube im
Menschen, als der der himmlische Genius, gesandt, dessen dunkles Daseyn zu
erhellen und zu verschönern: aber die wahre Herrschaft lässt alle unter ihr
stehende Kräfte frei walten in dem ihnen angewiesenen Gebiet, und findet in
der freien Harmonie des Ganzen ihr schönes Ziel. Ueber Religion und Theologie.
Vorred. P. \textit{VII.}}. Ifølge denne Overeenskomst kan den theologiske
Methodes Beskaffenhed
% THE p. 99 NOTE, read off PDF 108 at 0.83--0.93x over two contact sheets;
% fourteen lines. It is a German quotation set in FRAKTUR, so nothing in it is
% italic except the roman numeral of the reference, which is antiqua.
% Three things in it are as printed and are not corrected:
%   * "Sphære" --- the Danish æ, not the German ä, in a German sentence.
%   * "lässt" for "läßt": at 3x the letters after ä are TWO long s and a t,
%     a doubled long s + t, not the eszett this same note sets cleanly in "muß" (twice) and
%     "bewußt".
%   * PRINTER'S DEFECT: "als der | der himmlische Genius" --- the article
%     doubled across the line break, a dittography. De Wette's own text has it
%     once. NOT in the Rettelser list.
% The note carries no quotation marks of any kind; the German is introduced by
% a colon only.
% --- p. 100 ---
fremstilles paa følgende Maade: Enhver, der har christelig Tro, erkjender, at
han er opdraget til denne Tro i Kirken, eftersom den christelige Religionsidee
skinner frem igjennem Livet rundt omkring ham, gjennem Familielivet og de
hellige Forsamlinger; men er han naaet til en vis Aandsmodenhed, indseer han
tillige, at ikke engang disse Former ere fuldkomne Udtryk af Troens
Gjenstande, da Kirkens Tjenere og Lærere ere Mennesker, der kunne feile som
enhver Anden; han subtraherer derfor den almindelige Idee fra de Former,
Kirken har overleveret, og opsøger sig andre, der ere mere adæqvate; thi vel
kan han ikke selv opstille sig sande Former, men han kan dog, naar
fuldkomnere Former fremstille sig for ham, fryde sig ved dem og acquiescere i
dem, efterdi han selv har Troen. I sin Søgen understøttes han paa den ene
Side af Fornuften og den ydre Empirie, paa den anden Side af Troen, den indre
Empirie. --- Den christelige Religion er en positiv Religion, opstaaet til en
vis Tid og under visse ydre Betingelser; man maa derfor undersøge dens
historiske Documenter. Vi komme saaledes til Prøvelsen af den hellige Skrifts
Authentie, Axiopistie og Integritet, --- og skjøndt vi ved disse Undersøgelser
just ikke komme til nogen mathematisk Vished for, at Skriften er Religionens
rene og uforfalskede Kilde, komme vi dog ved dem til en høi Grad af
Sandsynlighed derfor. Denne Sandsynlighed bliver, naar Troen træder til,
absolut Vished. Man kan nemlig med Rette slutte, at saasandt Gud har meddeelt
os en Aabenbaring, saasandt vil han ogsaa, at denne uforfalsket skal komme til
Alle ad den sikkreste Vei; da nu den menneskelige Tradition, den mundtlige,
som Theologerne kalde den, kan fare vild, holde vi os forvissede om, at
Skriften er Religionens sande Kilde. Men Skriften selv har en doppelt Side,
fra hvilken den bør betragtes, den historiske og den religieuse. Fra sin
historiske Side
% "doppelt" (not "dobbelt") is what is printed --- read at 0.57x, two p's.
% Left as printed. No note and no display matter on this page; the whole page
% is one running paragraph, with no indent anywhere in its 32 lines.
% --- p. 101 ---
\setcounter{footnote}{0}
maa Skriften være Gjenstand for Kritikens Undersøgelse, og man maa her
indrømme, at man kan have forskjellige Meninger og yttre Tvivl om Meget
angaaende de enkelte Skrifters Oprindelse og Uoverensstemmelsen mellem de
enkelte Beretninger, uden derved at skade Religionen; i religieus Henseende
derimod er Alt, hvad der i Skriften læres om Gud og vort Forhold til ham, saa
dybt og klart, at det stedse vil udgjøre Troens sikkreste
Grundlag\footnote{Da den protestantiske Kirke nærmest begrunder de bibelske
Bøgers Anseelse paa den Kraft, hvormed de selv vidne om deres høiere
Beskaffenhed og Oprindelse, berøres den i det Mindste umiddelbart ligesaalidt
af de Tvivl, som den videnskabelige Kritik fremfører mod det Nye Testamentes
Ægthed og Integritet, som den, der er hensjunken i glad Beskuelse og Nydelse
af et Kunstværk, lader sig forstyrre ved de Grunde, hvormed Kunstkritikeren
bestrider dets Anseelse, som et Værk af denne eller hiin berømte Mester.
\textit{Dr.} Scharling. (Theol. Tidsskr. udg. af \textit{Dr.} Scharling og
\textit{Dr.} Engelstoft \textit{4}de Binds \textit{1}ste Hefte
P.~\textit{16}).}. Det religieuse Menneske iagttager derfor ved den første
Gjennemlæsning af Skriften Lærens almindelige Typus, for under dennes
Veiledning at gaae over til en Undersøgelse af de enkelte Dogmer, saa at der
bliver et gjensidigt Forhold mellem det Almindelige og de enkelte Momenter.
Dette maa være den rette Methode for Theologiens Behandling, efterdi den
ligemeget iagttager alle Sandhedens Momenter; den fastholder nemlig det
positive og objective Moment, --- een Idee, der gjennemtrænger Kirken, een og
samme Trang i det religieuse Gemyt, een hellig Skrift endelig under
Aarhundredernes Omvexlinger ---, og bevarer dog midt i denne Troeseenhed den
Enkeltes Frihed; thi indenfor de almindelige Grændser kan Enhver uddanne sig
Læren paa sin Maade, og udvælge sig det, der bedst stemmer med hans
Tilbøielighed\footnote{\textit{Nemini totum est necessarium, alium alia pars
ad salutem ducit. Bengel.}}.
% "Veiledning", not the ABBYY layer's "Anledning" --- read at 0.64x.
% Both reference marks on this page are the plain 1) and 2): the *) that BOTH
% OCR witnesses reported after "Grundlag" is a misreading of the figure 1,
% checked on the image at native scale. spacing.py's RUN on "udvælge sig" is a
% false positive: the words are set solid.
% TYPE: note 1) is Danish Fraktur with the abbreviation "Dr." and the figures
% in antiqua; note 2) is antiqua Latin throughout, hence one \textit{}.
% --- p. 102 ---
\setcounter{footnote}{0}

Ligesom den hellige Skrift har to forskjellige Sider, hvorfra den bør
behandles, saaledes gaaer det nu ogsaa den hellige Historie. Betragter man den
fra dens rationelle Side, om vi saa maae sige, vil man kjende de i
Beretningerne fremsatte Begivenheder, som saadanne, da er den de samme Forhold
underkastet, som den profane Historie; betragtes den derimod fta sin
% PRINTER'S DEFECT, transcribed as printed: "fta" for "fra". Read at 1.5x ---
% the second letter carries a crossbar and a flat head, and is the t of
% "underkastet" on the line above, not the shouldered Fraktur r of "fra" five
% words earlier. Both OCR witnesses read it as t. NOT in the Rettelser list.
religieuse Side, sees det strax, at den staaer uendeligt høit over denne. Thi
skjøndt vi ikke begribe Historien, kunne vi dog ikke undlade, naar den
fremstiller sine ophøiede Exempler for os, igjennem den religieuse Anskuelse
at modtage i vor Bevidsthed et Afbillede af den ophøiede Idee selv, hvorved
Sagen selv bestyrkes og tillige fæster Rod i vor Overbeviisning. Derfor skulle
vi fremfor Alt bestræbe os for at tilegne os Christusbilledet, som det er
afmalet i det N. Test., og vi ville da deraf høste de herligste Frugter til et
religieust Levnet\footnote{Af denne Behandlingsmaade af den hellige Historie
er De Wette i overordentlig Grad fortjent, Noget, vi med Glæde anerkjende, om
vi end ikke billige hans metaphysiske System. Navnlig skylder Moralen ham
megen Tak, idet det var ham, der førte Sædelighedens Princip tilbage til
Christus selv, medens den før ikke indeholdt Andet, end almindelige Love af
Samvittigheden og Fornuften, bekræftede med Christi og Apostlenes Segl. Smlgn.
Lærebog i den christelige Sædelære og sammes Historie, oversat af Mag.
Scharling P.~\textit{3} §~\textit{4} \ldots\ den christelige (Sædelære)
fremstiller de i Aabenbaringen allerede virkelig løste (sædelige) Opgaver i
bestemte Mynsterbilleder \ldots\ eller de bestaaende christelige Sæder ikke
blot i empirisk=historisk (positiv) Form, men ogsaa i eiendommeligt Stof
\ldots\ det Eiendommelige ved den ligger i den samme Sphære, hvori den
christelige Troeslæres reen christelige Læresætninger ligge, i Ahnelsens
Gebeet, hvor Ideen og Virkeligheden gjennemtrænge hinanden. P.~\textit{36}
§~\textit{62} Christus er \ldots\ Viisdom(mens) Princip, men nærmest for
Følelsen, som i ham opfatter Livets Fuldendelse. For at blive os dette Princip
klart bevidst, maae vi opløse det i et forstandigt Begreb.

Det høieste, hvad Christus er \ldots\ følgelig det al Begreb overstigende
Følelsesprincip er Ideen om Guds Søn \ldots\ Guds Søn er tillige
Menneskesønnen \ldots\ og det guddommelige Billede er Menneskehedens rene
Urbillede.}. --- Det ligger i Sagens Natur, at
% *** THIS NOTE CROSSES A PAGE BOUNDARY. IT IS COMPLETE HERE. ***
% Note 1) of p. 102 fills p. 102's whole footnote block (nineteen lines) and
% runs on into the FOUR lines of p. 103's footnote block, below a rule of
% their own, with no mark of their own --- "Det høieste, hvad Christus er
% \ldots\ Urbillede." They were read off PDF 112 and are included above, set
% as a second paragraph because the printing begins them on a fresh line.
% p. 103 therefore carries NO note of its own and gets NO
% \setcounter{footnote}{0}; the next reset is on p. 106. Same shape as the
% p. 72 --> p. 73 note.
% The printed ellipses vary between three, four and five points; all are set
% as \ldots. The figures after P. and § are antiqua, hence \textit{}; the
% German-derived "empirisk=historisk" keeps the Fraktur double hyphen.
% --- p. 103 ---
denne Theologie ikke kan ringeagte, langt mindre udelukke Philosophien, hvis
Maal det jo netop er at bringe Erkjendelsen til større og større Udvikling,
fra sine Enemærker; Menneskets Aand gaaer jo stedse fremad, og hvilken Magt
kunde vel standse den i dens Fremadskriden! Men ogsaa Philosophien er
Menneskeværk, og vi vide Alle, at Philosopherne stedse vakle mellem de meest
modsatte Theoremer. Derfor har denne Methode ogsaa sørget for at udfinde et
Middel imod denne Ustadighed hos Philosopherne. Naar Theologen tager et
philosophisk System for sig, maa han ikke have til Hensigt at forfølge
Systemet selv gjennem dets Tankers utallige Labyrinther, for med Dialektikens
egen Magt at gjendrive dets Indhold og afsløre dets metaphysiske
Vildfarelser; thi det er et uoverkommeligt Arbeide og er saameget
besværligere, som de fleste Systemer ere indviklede og mørke og dertil
anstille saa mange abstracte, for den christelige Fromheds Udvikling aldeles
overflødige Betragtninger over Gud og hans indre Liv, at den sande Theolog ved
at forfølge dem let kunde staae Fare for selv at rives med i Philosophiens
Afgrund og blive Philosoph istedetfor Theolog.

Da det er Kirken og fornemmeligt Skriften, der meddeler os den christelige
Tro, have vi som en Følge heraf i dem visse faste og urokkelige Punkter af
Troeslæren, med hvilke vi kunne sammenligne de Resultater, hvortil de
philosophiske Systemer ad Philosophiens Vei komme. Vi kunne da optage disse,
forsaavidt de stemme med hine Troesartikler, forkaste, hvad der staaer i
Modsigelse til dem, og iøvrigt lade Philo\-%
% --- p. 104 ---
sophen beholde, hvad der tilhører ham som saadan; thi saaledes give vi
Enhver, hvad ham tilkommer. --- Hvor sikker denne Fremgangsmaade er, sees
allertydeligst ved at anføre Exempler paa den Prøve, Theologerne ifølge denne
Methode lade Philosopherne gjennemgaae. \emph{Spinoza} lærte saaledes, at Gud
var Substantsen, og at Verdens Phænomener ikke vare Andet end Former for og
Modificationer af denne Substants; men ved denne Pantheisme ophæves den hele
Lære om Menneskenes Personlighed og Udødelighed, og som en Følge deraf den
hele Sædelære; hans System strider saaledes baade mod Troen og den hellige
Skrift; det bør derfor udelukkes af Theologien.

\emph{Fichte} lærte, at det menneskelige Jeg var alle Tings Ophav, og
forfaldt saaledes til Atheisme og cras Pelagianisme; men en Philosophie, der
% "cras Pelagianisme" as printed --- one s, and set in FRAKTUR like the
% Danish around it, not in the antiqua the book uses for Latin, so no
% \textit{}. Read at 0.57x. NOT in the Rettelser list.
strider saa aabenbart imod al sund Fornuft og kan have en saa fordærvelig
Indflydelse paa Religieusiteten, kan naturligviis ikke fordrages i den
christelige Theologie.

Theologien staaer som en Følge heraf med absolut Myndighed over Philosophien;
og man kan paa denne Maade i een Time prøve ti Philosopher.
% THE PRINTED HEAD AGREES WITH THE INDHOLD HERE. Read on PDF 113 at native
% scale: "§ 19." in bold antiqua, and below it "Undersøgelse af den kritiske
% Theologie." in Fraktur --- word for word the Indhold's entry for p. 104. No
% discrepancy to record, unlike §§ 7, 10, 13 and 15.
\parag{19}{Undersøgelse af den kritiske Theologie.}

Vi indrømme gjerne, at denne Theologie har optaget alle Sandhedens Momenter i
sig; give den Ret i, at den christelige Religion er en positiv Aabenbaring,
hvis Ideer ikke kunne findes eller udvikles ved abstracte, aprioriske
Speculationer; tilstaae, at den hellige Skrift er Rettesnoren for den
% --- p. 105 ---
christelige Lære, og saa fremdeles; --- men fordi den berører alle Sandhedens
vigtigste Momenter, er det endnu ikke afgjort, at Sandheden selv virkeligt er
tilstede i dem, ligesaalidt som man kan sige, at et Menneske lever, fordi alle
de enkelte Dele af hans Legeme ere tilstede.

Det bliver derfor Opgaven at undersøge, hvorledes Momenterne ere forenede, og
hvorledes de skulle forenes, at prøve, om der finder en sand Overgang Sted
imellem de enkelte Momenter. Og det vil da snart vise sig, at det Forbund, der
blev sluttet imellem Troen og Fornuften, ikke er af nogen varig Natur, og at
de ikke have indgaaet denne Fred, fordi den i Virkeligheden tilfredsstillede
begge Parter, men meget mere havde seet sig nødsagede til at indgaae en
Vaabenstilstand for at undgaae den Begge overhængende Fare. --- Betragte vi
% "Vaabenstilstand for" --- the two words are set very tight, with almost no
% space, which is why both OCR witnesses ran them together over an apostrophe.
% At 0.57x there is a space. p. 105 carries no note and no rule.
saaledes først Forholdet mellem Subjectiviteten og Objectiviteten, see vi, at
den subjective Frihed, denne Theologie anpriser, paa eengang er altfor tøiløs
og altfor indskrænket og bundet. Man gaaer ud fra den Tro, der lever i det i
Kirken opdragne Menneskes Bevidsthed, altsaa fra den almindelige Indflydelse,
det christelige Liv udøver over Enhver; man er altsaa enig om ikke videre at
indlade sig paa de særegne Modificationer, under hvilke denne Tro fremtræder
hos de enkelte Individuer, thi at hver Enkelt har og har Ret til at have sin
eiendommelige Subjectivitet, det er en afgjort Sag. Men da nu de samme Forhold
og den samme Opdragelse indvirker forskjelligt paa de forskjellige Gemytter,
og da ingen Lære kan gjøre sig gjældende uden igjennem den subjective Tro,
bliver det ikke den objective Sandhed, der retter og omdanner Gemyttet, men
meget mere den subjective Aand, der vælger og bestemmer, hvad der forekommer
den at være det Gode og Sande. \emph{Jeg vil troe, hvad der meest tiltaler
% --- p. 106 ---
\setcounter{footnote}{0}%
mig}\footnote{Sammenlign hermed \textit{Dr.} H. Martensens Yttring om den
kritiske Theologies Solipsisme: „Da altsaa Gudsideen først faaer Gyldighed ved
den praktiske Fornuft eller ved den absolute Følelse, ikke omvendt, mister den
Charakteren af det Constitutive, og hvis man tillægger den en saadan
Betydning, skeer det kun nominelt; thi hvad der constituerer og
\emph{bestemmer} Alt, er dog med Rette at ansee for den \emph{virkelige} Gud,
den allestedsnærværende Guddom i Bevidstheden, selv om det ikke faaer Navn af
Gud.“ Selvbev. Autonomie, oversat af Petersen.}. Og dog er denne Frihed ikke
% *** QUOTATION MARKS: THIS NOTE DEPARTS FROM THE BOOK'S STANDING FORM. ***
% The note on p. 106 sets the ORDINARY DANISH LOW--HIGH PAIR --- the opening
% mark on the baseline before "Da", the closing mark raised after "Gud." ---
% and not the raised-at-both-ends mark that the body uses throughout. Read at
% 0.98x and 1.0x. It is the same departure already recorded in the notes of
% pp. 48 and 70. Transcribed as printed: this page contributes one opening and
% one closing mark of the low-high system, and no raised mark, and they are
% the ONLY quotation marks anywhere in pp. 97--108, body or notes.
% EMPHASIS inside the note, on the image: "bestemmer" and "virkelige" are
% letterspaced; "constituerer" beside them is not.
% The Martensen quotation is Danish Fraktur, so only the antiqua abbreviation
% "Dr." is italic.
meget langt fra at være Trældom; thi netop fordi den subjective Natur i den
Grad urgeres, at den beholder et Element for sig som sin uomtvistelige
Eiendom, Retten nemlig til at bevare sin umiddelbare Eiendommelighed
ukrænket, bestemmer og befæster paa den anden Side Objectiviteten sine
Grændser saameget desto fastere, saa at de positive Momenter, der paa en
ganske udvortes Maade ere givne, fordre en udvortes Lydighed; men saaledes
blive de objective, positive Elementer, der staae tilbage, ikke Kjød og Blod i
os. Denne Selvmodsigelse vil i det Følgende træde endnu klarere frem. --- Man
% spacing.py's RUN on "Blod os." is a false positive: at 0.64x "ikke Kjød og
% Blod i os" is set solid. The letterspacing on this page is confined to the
% two words inside the note.
udhæver som en stor Fortjeneste ved denne Theologie Adskillelsen imellem
Troens Væsen og dens ydre Form, ifølge hvilken den christelige Tro vel er at
søge i Kirken, men ikke kan acqviescere i dennes Former. De kirkelige Former,
hedder det, ere ikke fuldkomment tilsvarende. Men hvad er en tilsvarende Form?
Mener man dermed en saadan, der umiddelbart udtrykker det guddommelige Væsen
selv, kan der efter denne Theologie slet ikke gives tilsvarende Former; mener
man derimod saadanne, som bevare og angive Troen i dens Sandhed, vil man neppe
kunne undgaae det Dilemma, enten at antage, at Kirken er i Besiddelse af de
sande og tilsvarende Former for Læren, og da maa Kirkelæren have samme
% --- p. 107 ---
Gyldighed, som Skriftens Lære; --- eller at antage Kirkens Former og
Udtryksmaader for i større eller mindre Grad at være feilagtige; men da maa
ogsaa den Tro, Kirken meddeler, være mere eller mindre forvansket, og man vil
da heller ikke kunne komme til nogen sand Skriftfortolkning, efterdi en saadan
udfordrer, at Kirken maa have overleveret os den sande Tro. Et Forsøg paa
under Forkastelse af eller Tvivl paa Kirkelærens enkelte Sætninger at
fastholde den Tro, Kirken har fremkaldt i os, vil bevise Sandheden heraf. Kan
vel saaledes den have Tro, der fjerner alle bestemte Troesforestillinger og
derpaa læser Skriften med Tvivl om Alt det, Troen har lært ham om Gud,
Christus o.~s.~fr.? Maa han da ikke ogsaa tvivle om Guds Tilværelse,
Menneskenes Forløsning og Saliggjørelse ved Christus, --- med andre Ord
begynde sit Skriftstudium med en absolut Negation af Troen, Noget der strider
mod al Fornuft? Overveier man dette, kan man neppe nægte, at den af Kirken
modtagne Tro er knyttet til de almindelige Troesartikler, saa at det er
umuligt for Nogen at tage Skriften for sig uden forudfattede Meninger om de
% NO FOOTNOTE ON THIS PAGE. The Fraktur OCR shows a mark after "Meninger" in
% l. 19; there is none. Checked twice: the line ends clean at 0.64x, and the
% page carries neither a rule (a row-by-row scan for a horizontal run over
% 150 px found nothing between the folio and the foot) nor a note block. Its
% 32 lines are all body. Same for p. 105 and p. 108.
Sandheder, han deri vil finde udtalte. Hovedspørgsmaalet bliver altsaa,
hvorledes Troens Grundartikler kunne afføre sig Kirkens concrete, men ikke
fuldkomment sande Former, og iføre sig Skriftens authentiske.

Kirken lærer, at Gud er til; men hvilke Egenskaber han har, og hvilke Følelser
han nærer imod os, skal Skriften lære. Man fordrer da af den, der vil læse
Skriften paa rette Maade, at han skal abstrahere fra alle Guds Egenskaber,
forsaavidt Kirken har beskrevet dem, og af Skriften selv udfinde Guds concrete
Begreb; men det er umuligt at tænke Gud uden at have en vis Forestilling om
ham, og Kirkens Lære om Guds Væsen vil saaledes stedse, frivilligt eller
ufrivilligt, gjøre sin Indflydelse gjældende paa Læserens Aand.
% --- p. 108 ---
Er han nemlig tilfredsstillet ved den Gudsbeskrivelse, Kirken har givet ham,
vil han i fuldstændig Lighed gjenfinde Kirkens Gud i Skriftens Gud; afviger
han derimod, før han læser Skriften, i eet eller andet Punkt fra Kirkens Lære
om Gud, kan han dog ikke negere Kirkelæren, uden at et positivt Gudsbegreb med
større eller mindre Klarhed maa fremtræde i hans Bevidsthed. (Kan saaledes
Nogen i negativ Henseende ikke erkjende Guds Treenighed, som den læres af
Kirken, maa han dog i positiv Henseende dyrke en anden Gud, hvad enten det nu
er Jødernes eller Muhamedanernes, eller rettere, enten det nu er Gud, som han
skildres af Jøderne, eller af Muhamedanerne; og har han den jødiske eller
muhamedanske Tro, er det en Selvfølge, at han heller ikke kan finde den
treenige Gud i Skriften, da netop den christelige Tro efter denne Theologies
egen Grundsætning er den første Betingelse for den rette Læsning af Skriften).
Men den christelige Læser, hedder det, maa ikke urgere eller med formastelig
Selvtillid holde fast ved sine forudfattede Begreber om de guddommelige Ting,
om han end ikke ganske kan abstrahere fra dem; hans Aand maa have en villig
Bevægelighed til stedse at læmpe sine subjective Meninger efter Skriftens
Udsagn: han maa have den sande Fortolkningsmethode. Hvori bestaaer altsaa den
Methode, efter hvilken vi skulle søge at udfinde Skriftens Mening?
% "udfinde", read at 0.57x. Neither OCR witness got it: the Fraktur model gave
% "udstude" and the ABBYY layer "udstade". There is no footnote mark after
% "Mening?" --- the Fraktur model's stray character there is scan noise.

Allerførst maa man bortrydde de Forhindringer, som Forskjellen i Sproget og i
Tiden lægger Læseren i Veien; og hvor store Fortjenester de lærde Kritikere og
Philologer have i denne Henseende, er bekjendt nok. Men har man nu undersøgt
de almindelige Betingelser for Tale og Handling, som de hellige Forfattere
vare underkastede, tilføiet de nødvendige psychologiske Bemærkninger og
saaledes med temmelig Sikkerhed lært den Mening af de enkelte Ord og Udtryk at
kjende,

% --- p. 109 ---
\setcounter{footnote}{0}
% JOINT WITH p. 108, verified on the image (PDF 117 foot / PDF 118 head, 0.55x):
% p. 108 ends "...at kjende," --- a whole word followed by a comma --- and
% p. 109 opens "som Forfatterne maae antages". NO word is broken at this joint,
% so no \-% is wanted. p. 108 carries no note, and p. 109's footnote block
% opens with its own mark 1) directly under the rule, so nothing runs over
% from p. 108 either.
som Forfatterne maae antages at ville have lagt i dem, --- er man jo dog endnu
meget langtfra at være trængt ind til den evige Betydning af deres subjective
Meninger og Bestemmelser. Her opstaaer nu en doppelt Vanskelighed; thi for det
Første ere Skriftens Ord og Udtryk ikke Andet, end ufuldkomne Betegnelser for
de høieste, guddommelige Gjenstande; og for det Andet gjorde hiin Tidsalders
særegne Forhold det baade af indre og ydre Grunde til en Nødvendighed for
Forfatterne at benytte sig af en hine Tider aldeles ejendommelig Udtryksmaade.
Det hele Udbytte af slige skarpsindige exegetiske Undersøgelser indskrænker sig
da til at udfinde, hvorledes de hellige Forfattere have udtrykt deres
Gjenstande; hvorledes disse Gjenstande selv ere at tænke, kunne de ikke lære
os, thi da Forskjellen mellem Ideen og de Former, hvori den er udtrykt, er saa
stor, lod det sig tænke, at Apostlene havde fremsat Læren paa denne eller hiin
Maade, medens vi ikke kunne tænke den paa samme
Maade\footnote{Et Exempel herpaa finde vi i Læren om Djævelen. Hvad
Adskillelsesprincipet selv angaaer, smlgn. \textit{Dr.} Hagenbach:
Encyclopædie u. Mthdl. der theologischen Wissenschaften P. \textit{174} og
\textit{175}. Sache der Auslegung ist es, in Stellen, wie Rom. \textit{III,
24. 25. IV, 25. V, 9. 10. VIII, 34. 1} Cor. \textit{XV, 3.} Eph. \textit{I, 7.
1} Pet. \textit{I, 18. 19. II, 29. 1} Joh. \textit{I, 7.} Hebr. \textit{VII,
26. IX, 28} u.~s.~w. die Versöhnungs= und Opfertheorie als eine den
Schriftstellern des N.~T. geläufige Idee factisch anzuerkennen. Aber
verschieden kann die weitere Erklärung der Stellen ausfallen, je nachdem man
sich entweder damit begnügt solche Ideen in ihrem historische Zusammenhange
mit dem alttestamentlichen Opferwesen u.~s.~w. aufzupassen, oder je nachdem
man sie auf ein tieferes religiøses Bedürfniß nach Erlösung zurück zu führen
und in ihrem innersten, geheimnißvollen Grunde erfassen zu müssen glaubt.}.
Man maa saaledes, hvad Dogmatiken angaaer, tilstaae, at disse Undersøgelser
vel i negativ Henseende have fjernet mangen en falsk
% p. 109 carries ONE note, marked 1) --- the plain next number.
% *** PRINTER'S DEFECT in the note, transcribed as printed: "religiøses" for
% "religiöses" --- the Danish Ø, an o with a full slash through the bowl, not
% the German umlaut. Read at 2.5x on PDF 118 l. 13 of the note block. The
% compositor's Danish fount reaching into a German quotation; the same
% substitution recurs in the p. 120 note at "sittlich=religiøses" (p. 121).
% Also as printed, uncorrected: "in ihrem historische Zusammenhange" for
% "historischen" --- both OCR witnesses agree.
% DOUBTFUL READING, logged: after "Sache" (l. 4 of the note) there stands a
% small comma-shaped mark BELOW the baseline, and a separate speck before
% "der". At 2.5x it sits far lower than the true commas on the same line (cf.
% "ist es,"), so it is read as ink dirt and no punctuation is set. The
% rejected alternative is "Sache, der Auslegung"; the two OCR witnesses
% disagree with each other here (comma vs. point).
% TYPE: the note is Danish and German FRAKTUR throughout; only "Dr." and the
% figures and roman numerals are antiqua. The scripture references are set as
% runs of antiqua with their internal points and commas, so the \textit{}
% groups follow the runs rather than each figure separately; "Rom.", "Cor.",
% "Eph.", "Pet.", "Joh.", "Hebr." and "P." are Fraktur and stand outside.
% NO letterspacing anywhere on this page: spacing.py's RUN on "lære os," is a
% false positive of justification, checked at 0.55x.
% --- p. 110 ---
Anvendelse af enkelte Skriftsteder, men i positiv Henseende aldeles Intet kunne
udrette; vil man finde de dogmatiske Ideer selv, maa man altsaa slaae ind paa
en ganske anden Vei for sin Undersøgelse. Vi komme saaledes atter her til det
Spørgsmaal, hvorfra vi gik ud, hvorledes vi nemlig af den hellige Skrift kunne
udfinde de positive Former for Troen.

Fortolkningen skal bevæge sig Skridt for Skridt fremad fra de lettere og
fatteligere Udtryk i Skriften til de vanskeligere og mere dunkle. Men da de
umiddelbart indlysende Begreber næsten altid ere hule, tomme og derfor
betydningsløse, bliver Resultatet af en saadan sammenlignende
Behandlingsmaade kun altfor ofte det, at de dybere Ideer aldeles forsvinde.
Spørges der saaledes om Betydningen af Christi Udsagn (Joh. \textit{15, 1, 4}):
Ἐγώ εἰμι ἡ ἄμπελος, ὑμεῖς τὰ κλήματα, og (V. \textit{9. 10}): Ἐὰν τὰς
ἐντολάς μου τηρήσητε, μενεῖτε ἐν τῇ ἀγάπῃ μου, er man i den største Uvished
om, hvorfra Udviklingen skal gaae ud, om man skal antage et saadant Forhold
mellem Christus og hans Disciple, at Kjærligheden ikke kunde have nogen sand
moralsk Kraft, de Christne ikke have nogen sand Kjærlighed, hvis ikke hans Liv
i Ordets egentlige Forstand var deres Livs Kilde, hvis de ikke af hans Marv
sugede deres hele Kraft, --- eller om man skal sige: hvad en reel Forbindelse
mellem Christus og dem, der troe paa ham, betyder, forstaae vi ikke, men hvad
Kjærlighed og Lydighed mod Christi Bud er, det forstaaer Enhver; man kan da
ikke antage denne og lignende Talemaader for Andet, end Billeder, hentede fra
en orientalsk Forestillingskreds. Ligesaa usikker er Fortolkningen, naar
Christus (Joh. \textit{17, 11}) siger, at han er forenet med Faderen (ἵνα ὦσιν
ἓν καθὼς ἡμεῖς); thi man kan ikke vide, om det er Skriftens Mening, at der er
en umiddelbar Identitet mellem Faderen og Sønnen, og at altsaa de Christne,
naar de have Troen, ere indtraadte i det
% p. 110 carries NO note: a row-by-row scan of PDF 119 for a horizontal ink
% run over 150 px found no rule between the folio and the foot, and all 33
% lines are set to the body measure and the body line-pitch (159 px, against
% 116 px in the footnote blocks of pp. 109, 111, 113).
% One paragraph break, at "Fortolkningen skal bevæge sig" (indent +203 px,
% measured, against a body margin of 736 px on this verso).
% GREEK: the fount's script theta in καθὼς is typed as plain θ and its cursive
% kappa as plain κ, per the header --- fount variants, and fatal to the build
% as U+03D1/U+03F0. Accents and breathings are as printed: Ἐγώ, ἵνα, ἓν.
% DOUBTFUL READING: the reference is set "(Joh. 15, 1, 4)"; at 1.5x the mark
% between the 1 and the 4 is read as a comma. The rejected alternative is
% "15, 1. 4" (i.e. verses 1 and 4), which would suit the quotation better ---
% the words quoted are those of v. 5.
% NO letterspacing: spacing.py's RUN on "moralsk Kraft" is justification, not
% Sperrsatz, checked at 0.55x.
% --- p. 111 ---
\setcounter{footnote}{0}
samme reelle Eenhedsforhold med Gud og Christus, --- et Forhold, uforstaaeligt
for Alle, der ikke ved at fødes af Gud blive θείας κοινωνοὶ φύσεως ---, eller
om Meningen er den, at Christus og Faderen ere Eet i samme Forstand, som de
Christne og Christus, nemlig ved et moralsk Eenhedsbaand. Hvad en saadan
moralsk Forening betyder, kan Enhver, der har sund Fornuft, begribe; ved dette
Begreb bør man altsaa blive staaende, thi Alt, hvad der sigter til en
substantiel Eenhed, er kun sagt i figurlig og prægnant Forstand og har ingen
Betydning for Religieusiteten som saadan.

Det gaaer med Undersøgelsen om den hellige Aands Begreb paa samme Maade. Den
beskrives (Joh. \textit{14, 16}) som en Person, og Christus siger om den (Joh.
\textit{14, 26}): ὑμᾶς διδάξει πάντα, og (Joh. \textit{16, 13}): ὁδηγήσει ὑμᾶς
εἰς πᾶσαν τὴν ἀλήθειαν (smlgn \textit{1} Cor. \textit{2, 10. 11}); hos Lucas
derimod kaldes den en δύναμις ἐξ ὕψους; man kan altsaa fortolke dette paa en
doppelt Maade. De, der troe, at den hellige Aand er en Person, nægte dog ikke,
at den tillige er en Kraft, men paastaae kun, at den er mere end en Kraft; de
derimod, der gaae ud fra Tingens umiddelbare Begribelighed, erklære, at de
ikke kunne forstaae, hvorledes Aanden kan være en særegen Person i Gud, og
mene, at det er tilstrækkeligt, at antage den for en guddommelig Kraft, uden
nøiere at bestemme dennes Natur\footnote{Talrige Exempler paa denne Cirkelgang
i Beviisførelsen kunne hentes fra Bretschneiders Dogmatik; smlgn f.~Ex. Om
Triniteten P. \textit{495}. Handbuch der Dogmatik, \textit{1}ster Band,
\textit{2}te Aufl., Leipz. \textit{1822}; Om Christi Person §~\textit{140}, P.
\textit{167}, \textit{2} B.}.

Hertil kommer endnu, hvad Dogmatiken viser, at enhver Lære, som Skriften
fremfører, indeholder Momenter, der ere hinanden indbyrdes modsatte, ja endog
ligefrem ophæve hinanden. Saaledes lærer for Exempel Skriften, at Gud ikke
% p. 111 carries ONE note, marked 1) --- the plain next number. It is Danish
% and German FRAKTUR throughout (including "Handbuch der Dogmatik, 1ster Band,
% 2te Aufl., Leipz. 1822"); only the section sign and the figures are antiqua.
% "§ 140" read at 3x: the middle figure is a 4, with the open triangular
% counter; the rejected alternative is 110, which is what the ABBYY layer
% gives. "smlgn" in the body is printed without a point, as here.
% GREEK: script theta in θείας and ἀλήθειαν, and the cursive kappa of
% κοινωνοὶ, are typed as plain θ and κ, per the header.
% --- p. 112 ---
blot er kjærlig og barmhjertig, men tillige, at han er vred paa dem, der have
syndet. Man kan da af Skriften ikke vide, om man skal urgere den absolute
Barmhjertighed og ikkun antage Guds Vrede for en subjectiv Forestilling hos
Menneskene, eller om Vreden er et objectivt Begreb, der i Virkeligheden
indskrænker den uendelige Barmhjertighed, eller om man endelig skal tænke sig
Guds Natur saaledes, at han vel \textit{ad intra} er den Salige, men
\textit{ad extra} formedelst Verdens Synder indtræder i et saadant
Forsoningsforhold af Vrede og Kjærlighed, at hans uendelige Kjærlighed ved
Incarnationen fjerner den objective Vrede fra dem, der troe paa Christus. Man
maa altsaa her anvende Skriftanalogien og see hen til Skriftens Lære om
Forsoningen ved Christus. Det er imidlertid saa langt fra, at de Steder, der
handle om Christus, løse Knudens Vanskelighed, at de snarere gjøre den endnu
mere indviklet. Thi vel lærer Skriften, at Christi Død er et Beviis paa Guds
Kjærlighed, men den lærer tillige, at Guds og Christi Kjærlighed fornemmeligt
viser sig deri, at Straffen for vore Synder er overført paa den korsfæstede
Christus; forsvarer man nu Vredens Objectivitet, kommer man til Dogmet om
\textit{satisfactio vicaria}; kan man derimod ikke antage, at Gud er vred, maa
man henføre den hele Forsoningslære til en Forandring i Menneskenes
Forestillinger om Gud. Vredens Begreb er saaledes afhængigt af Begrebet om
Christi \textit{satisfactio}, men dette er atter afhængigt af hiint; man seer
da heraf, at den ene Analogie stedse udkræver en anden Analogie, hvorpaa den
kan støtte sig, eller med andre Ord: Analogien selv viser, at man ikke kan
acqviescere ved nogen Analogie.

Man vil maaskee indvende, at denne Skriftens Oeconomie er det bedste Beviis
paa Sandheden af den Yttring, man saa ofte hører, at Gud ikke vil give
Menneskene nogen objec\-%
% --- p. 113 ---
\setcounter{footnote}{0}
tiv Indsigt i de guddommelige Ting, saa at alle de Vanskeligheder, vi før have
omtalt, netop hidrøre derfra, at vi have tillagt den objective Lære en altfor
stor Betydning, --- og at Skriften selv bestyrker denne Mening ved at
fremstille Guds Natur, hans Kjærlighed og Forsoningen i Billeder og
almindelige Begreber, lempede efter Menneskenes Fatteevne, for at de, styrkede
ved Læren om Syndernes Forladelse, kunne være redebonne til enhver god
Gjerning. --- Men ere vi ikke saaledes komne tilbage til det selvsamme Punkt,
hvorfra vi gik ud, til Kirkens almindelige Artikler?\footnote{Det Usikkre i en
saadan Fortolkningsmaade kan vises paa flere Maader. Det hedder f.~Ex. Joh.
\textit{4, 24}, πνεῦμα ὁ θεός, Joh. \textit{1, 5} derimod ὁ θεός φῶς ἐστι;
dersom nu Nogen ikke kunde forstaae, hvad det vil sige, at Gud er en Aand, men
troede lettere at kunne begribe, at Gud er et Lys, kunde man da ikke fristes
til at troe, at de Christnes Gud og Magernes Gud vare den selvsamme Gud?}. Det
hele Resultat af Kriticismens besværlige Arbeide og mangfoldige Undersøgelser
er altsaa den Erkjendelse, at Skriften selv beviser, at den i dogmatisk
Henseende ikke kan lære os Andet, end det, vi alt have lært af Kirken.
Kritikerne have altsaa ved deres store Anstrængelser ikke udrettet Andet, end
at komme i Modsigelse med sig selv, thi de lovede os, at Læsningen af Skriften
skulde forskaffe os en mere concret Form for Troen, naar vi kun paa vor Side
medbragte Troen i dens rene Almindelighed. Vi opfordres saaledes til nøiere at
undersøge, om det ogsaa virkeligt er i Troens sande Interesse, saaledes som de
anbefale os, at renoncere paa enhver Begriben af de objective
Gjenstande.~--- Det er indlysende, at en Tro paa noget aldeles Ubekjendt er
umulig. Jeg modsiger saaledes mig selv, dersom jeg siger, at jeg troer paa
Verdens Skaber, naar jeg paa samme Tid erklærer, at jeg aldeles Intet kan vide
om Maaden, hvorpaa Verden er skabt. Thi undersøger jeg, hvad
% *** THE BODY MARK AND THE NOTE MARK DISAGREE ON THIS PAGE. ***
% The reference mark in the body, after "Artikler?", is a superscript figure
% 1 followed by a bracket --- "1)" --- read at 1.0x on PDF 122 l. 11. The mark
% at the head of the note itself, under the rule, is an ASTERISK, "*)", read
% at 1.0x on PDF 122. This is the same disagreement already recorded at p. 96.
% The note is set here as an ordinary \footnote, which prints 1), i.e. the
% body's mark; the printed "*)" at the foot is recorded only in this comment.
% p. 12's and p. 46's notes are the two earlier pages that print *) at both
% ends, so the asterisk itself is not new; the disagreement is.
% "Gjenstande.~---": the dash is set tight against the point in the printing
% ("Gjenstande.—Det"), unlike the spaced dashes elsewhere on the page.
% GREEK: the script theta of θεός (twice) and the variant phi of φῶς are typed
% as plain θ and φ, per the header. "Joh. 1, 5" is as printed (the place is
% 1 John 1, 5).
% --- p. 114 ---
\setcounter{footnote}{0}
Skabelsen er, kan jeg ikke undgaae at erklære Skabelsen for en Handling af
Guds absolute Frihed; jeg tænker da, saavidt muligt, den absolute Frihed; men
dette udfordrer, at jeg maa have en Forestilling om Friheden, og er nu denne
Forestilling, dette Begreb falsk, da er ogsaa det, hvorpaa jeg troer, falsk,
\dsi{} jeg har en falsk Tro. Og dersom Nogen bekjender, at han paa eengang
troer Christus i Besiddelse af en menneskelig og en guddommelig Natur,
bedrager han sig selv, hvis han paa samme Tid erklærer, at han ikke kan tænke
sig en Forening af dem; thi forsaavidt han har et Begreb om den guddommelige
og menneskelige Natur, forsaavidt kan han ogsaa tænke dem forenede med
hinanden; har han derimod slet ingen Mening om nogen af disse Naturers
Egenskaber, saa at den menneskelige Natur er ham et ubekjendt \textit{x}, den
guddommelige et ubekjendt \textit{y}, --- troer han jo paa en Forening af
\textit{x} og \textit{y}, troer paa Noget, som han ikke selv veed hvad er,
\dsi{} han troer Intet\footnote{Forskjellen mellem den kritiske og speculative
Dogmatik bestaaer ikke deri, at alene denne, ikke hiin erkjender den indre
Erfarings Nødvendighed: de antage begge, at ingen Lære kan udvikles uden at
see hen til det aandelige Livs positive Elementer; de stræbe begge at give
Forestillingerne og Begreberne indre Fasthed. Men Forskjellen ligger i den
forskjellige Maade, hvorpaa enhver af dem benytter de samme positive
Elementer. Den kritiske Theologie behandler alene Forestillingerne og Ideerne
negativt; har den derfor efter at have anført Skriftstederne og fældet en Dom
over de forskjellige kirkelige Meninger angivet, hvorledes Skabelsen ikke er
at tænke, lader den sig nøie dermed, overlader den nærmere Forestilling derom
til Enhvers subjective Anskuelse og gaaer derfra over til at behandle et andet
Dogme paa samme Maade. Og naar nu saaledes Kritikerne have berørt alle enkelte
adspredte Fragmenter af Elementerne uden at lade dem optage og gjennemtrænge
hverandre, undre de sig høiligen over, at speculative Theologer kunne falde
paa at tale om en objectiv Begriben af de guddommelige Ting, som om Elementerne
i en saadan Theorie kun kunne være at gribe i Luften. ”Hvor kan Nogen ville
tænke, hvorledes Skabelsen gik for sig, hedder det; ligger den ikke udenfor
enhver Erfaring! Aarhundreder ligge imellem Verdens Begyndelse og os; --- kan
vel noget Menneske erfare, hvad der ligger Aarhundreder tilbage i Tiden!” ---
Ja, vare saadanne Tankespring nødvendige, da vilde ikke alene metaphysiske
Undersøgelser være saare utilfredsstillende, da vilde Philosophien Dag for Dag
blive mere vanskelig; thi uafbrudt iler Tiden forbi os og vi fjernes mere og
mere fra Begyndelsespunktet! Men Sagen forholder sig ikke saaledes; vel vover
den sande Philosophie aldrig at overskride den menneskelige Erfarings
Grændser; men den gjør en anden Anvendelse af Livets positive Elementer. Hvad
der gjør, at vi kunne have en Forestilling om Skabelsen og Friheden, det Samme
gjør ogsaa, at vi kunne komme til Skabelsens speculative Begreb. Den kritiske
Dogmatik berører kun i Forbigaaende den absolute Friheds Begreb, men den
speculative Dogmatik kræver, netop fordi den er Videnskab, at vi skulle
forfølge det til dets yderste Grændser, og det er ad denne Vei, vi komme til
Skabelsens metaphysiske Begreb. Vil man derfor i Dialektikens Ild, luttre de
positive Momenter, som den kritiske Dogmatik lader staae ubestemte, kommer man
til den speculative Dogmatik.}.
% *** THIS NOTE CROSSES A PAGE BOUNDARY. IT IS COMPLETE HERE. ***
% Note 1) of p. 114 fills p. 114's footnote block (fifteen lines, the rule at
% y 3597 of PDF 123) and runs on into the NINETEEN lines of p. 115's footnote
% block, below a rule of their own (y 3115 of PDF 124) and with no mark of
% their own --- "menterne i en saadan Theorie ... til den speculative
% Dogmatik." They were read off PDF 124 and are included above, as one
% continuous paragraph, because the printing carries the sentence straight
% across the joint ("som om Ele-" / "menterne"). p. 115 therefore carries NO
% note of its own and gets NO \setcounter{footnote}{0}; the next reset is on
% p. 119. Same shape as the p. 102 --> p. 103 and p. 72 --> p. 73 notes.
% QUOTATION MARKS in the note: the RAISED U+201D at BOTH ends of "Hvor kan
% Nogen ... i Tiden!", the book's standing body form. Read at 0.58x on
% PDF 124. This note does NOT depart from the body's system, unlike the notes
% of pp. 48, 70, 79, 84 and 106. Two marks, parity 0.
% "undgaae" (not "undgaas"): read at 2x; the final e is under-inked but its
% bowl and exit stroke are both present. Both OCR witnesses gave "undgaas".
% DOUBTFUL READING: "i Dialektikens Ild, luttre" --- at 1.5x the mark after
% "Ild" sits below the baseline and is read as a comma, though the sense wants
% none. The rejected alternative is a point. Transcribed as printed.
% TYPE: x and y are antiqua against the Fraktur and take \textit{}; \dsi is
% the printed ɔ: sign, twice.
% --- p. 115 ---
Man faaer ofte at høre, at den hellige Skrift ikke udgjør et philosophisk
System, at den er populair; men denne ligesaa trivielle som sande Bemærkning
tages ikke i samme Betydning af Alle. Mene nemlig de, der tale om Skriftens
Eenfoldighed, dermed at bevise, at man kan lade Lærens objective Former staae
uoplyste, naar man blot, alt efter sin Tidsalders større eller mindre
Oplysning, anstiller Reflexioner over fromme Følelser og moralske Forskrifter,
handle de jo tvertimod Skriftens egen hele Oeconomie. Thi Skriften opfordrer
netop selv til at gaae en ganske modsat Vei. Vel er det sandt, at der i
Skriften gives den eenfoldigste Skildring af de guddommelige Ting, men man maa
ligesaafuldt erindre, at dens Udvikling af Fromheden og det sædelige Liv bærer
det selvsamme Eenfoldighedspræg. Det er saa langt fra, at den
% The page opens a NEW PARAGRAPH: its first line is indented +225 px against a
% body margin of 298 px on PDF 124, the standing paragraph indent on this
% recto. Its fifteen body lines are followed by the rule and by the nineteen
% lines that complete p. 114's note; see the comment there. NO reset here.
% The "8*" below the last line of the note block is the printer's signature,
% not text.
% --- p. 116 ---
gaaer ud fra den menneskelige Subjectivitet, at den tvertimod udleder al
Fromhed og Sædelighed derfra, at den objective Gud paa synlig Maade har
aabenbaret sig og ved sin Aand nedlagt et nyt \emph{Viisdoms=} og Livsprincip
i Menneskene. Skriften kjender altsaa aldeles ikke hiin Adskillelse mellem
Erkjendelsens objective og subjective Momenter. Eller kan maaskee Nogen mene,
at f.~Ex. Paulus under Fremstillingen af Prædestinationslæren har tænkt
saaledes ved sig selv: ”Vel veed jeg ikke, hvad Gud er i sig, eller hvilke
Raadslutninger han har fattet, men da Prædestinationens Idee er et fortrinligt
Middel til at indskjærpe Jøderne Beskedenhed og meddele de Christne Trøst, vil
jeg ikke destomindre opstille denne Lære!” --- μὴ γένοιτο! Vil Du derfor være
et Barn i at forstaae den apostoliske Tidsalders objective Gjenstande, maa Du
ogsaa ansee Dig for et Barn i Henseende til subjective, kritiske Reflexioner!
Men opfordrer Din Tidsalder Dig til at anstille Undersøgelser og løse
Problemer, maa Du ikke lige imod Skriftens eget Exempel troe at kunne gjøre
dette uden at see hen til de guddommelige Tings objective Rige!

Den samme Ubestemthed, --- en Følge af Kritikens udvortes Sammenstillen af de
enkelte Momenter ---, maa naturligviis vise sig i den Maade, hvorpaa de
behandle den hellige Historie. Kriticismen kan inddele hele den hellige
Historie i tvende Dele, af hvilke den første indbefatter alle Beretninger om
de hellige Personer, der ere af en saadan Natur, at man endnu den Dag i Dag
troer at kunne paavise Exempler af samme Art (f.~Ex. at Christus gik fra
Judæa til Galilæa, at han tiltalede Folket, at han udviste stor Kjærlighed
mod alle Mennesker). Begivenheder af en saadan Natur findes ligesaavel i den
profane, som i den hellige Historie. Til den anden Klasse maa man derimod
henregne alle de Begivenheder,
% EMPHASIS, settled at 3x on PDF 125: "Viisdoms=" IS letterspaced --- the
% letters stand clear of one another with a full counter of white between
% them. "og" and "Livsprincip" beside it are set SOLID: their letters touch,
% exactly as in "nedlagt et nyt" on the same line. So only the first half of
% the compound is emphasised, and \emph{} closes at the double hyphen. This
% was checked because the reading is surprising; it is what is printed.
% spacing.py flagged the site only as a weak "single?".
% QUOTATION MARKS: the RAISED U+201D at BOTH ends of the imagined Pauline
% speech, "Vel veed jeg ikke ... denne Lære!". The book's standing form.
% Two marks, parity 0. p. 116 carries NO note (no rule, no footnote block).
% The passage "Vil Du derfor være et Barn ... objective Rige!" is NOT a block
% quotation: its lines run to the full body measure and stand at the body
% margin. The indentation the ABBYY layer shows there is its own artefact.
% One paragraph break, at "Den samme Ubestemthed" (indent +237 px).
% --- p. 117 ---
der ere eiendommelige for den hellige Historie, dem nemlig, i hvilke Gud er
umiddelbart tilstede. Med Hensyn til disse har man tvende Fordringer at gjøre:
det maa for det Første godtgjøres, at disse Begivenheder virkeligt have fundet
Sted --- de vilde jo ellers ikke henhøre til Historien ---, og for det Andet
maae de overnaturlige Ideer ikke udenfra paahæftes Begivenheder af ganske
almindelig Natur, men maae kunne udledes af Begivenhederne selv; det
overnaturlige og naturlige Moment maae falde sammen, hiint maa fremtræde
ligesaa reelt og tydeligt som dette. Men dette medfører den største
Vanskelighed for Fortolkeren. Den overnaturlige Form, hvori Begivenhederne,
f.~Ex. Christi Fødsel eller Opstandelse, fremtræde, gjør det nemlig umuligt
for Kritiken at udfinde de reent naturlige Momenter, og om den endog var
istand til at drage disse for Lyset, vilde den dog ikke kunne gjøre dette uden
at krænke den Pagt, den har sluttet med Troen; den maa derfor lade det
naturlige Moment i Fortællingen blive staaende aldeles ubestemt. Og hvad nu
det overnaturlige Moment angaaer, nægter Kritiken vel ikke altid Muligheden af
Mirakler overhovedet; thi ifølge det Princip, hvorfra vi ere gaaede ud, hører
det overnaturlige Rige, og derfor ogsaa Miraklerne overhovedet, hen under
Troens Gebeet; men da det Endelige ikke kan rumme det Uendelige, det Usynlige
ikke være tilstede i det Synlige, vil ikkedestomindre i de enkelte Tilfælde
Miraklets Sandhed stedse være stor Tvivl underkastet. Da man nu ikke kan
acqviescere ved den Maade, hvorpaa den overordentlige Begivenhed er fortalt i
Skriften, indskrænkes ogsaa det overnaturlige Moment til noget aldeles
Ubestemt. Men saaledes forsvinder denne Deel af den hellige Historie af
Historiens eget Rige. De positive Momenter, der blive tilbage, henhøre til den
ydre Indklædning, der ene skyldes den religieuse Phantasie. Menneskets Søns
synlige Skikkelse
% NO FOOTNOTE ON THIS PAGE, and no rule: a row-by-row scan of PDF 126 for a
% horizontal ink run over 150 px found nothing between the folio and the foot,
% and all 32 lines keep the body line-pitch. The stray mark the Fraktur OCR
% reports after "Lyset" is a speck below the L, checked at 0.55x.
% The page is one continuous paragraph: no line on it is indented.
% --- p. 118 ---
kunde ikke i Virkeligheden være det sande Udtryk for Ideen om Guds Søn; hiin
Guds Søn var da kun en Tanke undfanget i Apostlenes religieuse Gemytter.

Forholder Sagen sig saaledes, kan man neppe nægte, at det fremherskende
Særkjende hos Kritikerne er eu udvortes Sammensætning af Positivisme og, om vi
saa maae sige, Liberalisme; de staae altsaa paa den slette Middelveis
Standpunkt, et Standpunkt, der ophæver sig selv. Dette viser sig ogsaa deri,
at man ikke sjelden, naar et System, der, som saadant, kun ad Tænkningens Vei
kan prøves, stiller sig i Modsætning til Kriticismen, hører Kritikerne
apellere til den hellige Skrifts Auctoritet, hører dem paaberaabe sig
Skriftens blotte Navn mod de speculative Dogmatikere, i det de drage en skarp
Adskillelse imellem Guds og Menneskets Ord, ret ligesom Katholikerne fordum i
Kampen mod Luther ideligt apellerede til Fædrene. \emph{Det er sandt, fordi
det staaer i Skriften.} --- Men neppe bliver der Tale om Friheden i den
kritiske Undersøgelse, før alle Positivismens Skrupler strax forsvinde, før de
indrømme, at det kun gjør liden Forskjel, om Skriften er skreven af Apostlene
eller af andre fromme Mænd, før de frimodigt og aabenlyst skjelne imellem
Skriftens særegne Form og Guds Ord. --- \emph{Skriften er hellig, fordi den
taler Sandhed.}

Vi tvivle ingenlunde paa, at der finder et Vexelforhold Sted imellem
Sandhedens Vidnesbyrd om Skriften og Skriftens Vidnesbyrd om Sandheden; men
derom tvivle vi høiligen, at aldeles modsatte Principer kunne lade sig forene
med hinanden. Enten maae vi, som Fædrene, ubetinget lade vor Dom være Skriften
underkastet og, som en Følge deraf, ogsaa antage en Verbalinspiration; eller
vi maae overalt modtage Sandheden i Sandhedens Navn og have da ingen Grund til
at antage, at Gud fra Christi Tid lige indtil Trajan talede
% *** PRINTER'S DEFECT, transcribed as printed: "er eu udvortes" for "er en
% udvortes" --- a TURNED n, which prints as a u. Read at 2x on PDF 127 and
% compared, on the same line, with the n of "Kritikerne" four words earlier:
% the n joins its two stems by a hairline across the TOP, the sort here joins
% them at the BOTTOM and carries the u's top-right flag. Both OCR witnesses
% read u. NOT in the Rettelser. This is the same accident as p. 3's "Tiugen"
% and the turned n in p. 11's § 3 head, "Deu vulgaire Empirie."
% EMPHASIS: TWO letterspaced sentences on this page, both read at 0.55x and
% both flagged as RUNs by spacing.py --- "Det er sandt, fordi det staaer i
% Skriften." and "Skriften er hellig, fordi den taler Sandhed." They are the
% only Sperrsatz in pp. 109--120 apart from p. 116's "Viisdoms=", p. 120's
% Clemens sentence and the De Wette phrase in p. 120's note 3.
% NOT emphasis: a former owner has UNDERLINED IN PENCIL, in the scan, "derom
% tvivle vi høiligen, at aldeles modsatte Principer kunne lade sig forene med
% hinanden." The strokes run under the baseline, break at the margins and vary
% in weight; they are manuscript, not rule. spacing.py's RUN on "Enten vi," is
% that pencil line plus justification, and is a false positive.
% p. 118 carries NO note and no rule; two paragraph breaks, at "Forholder
% Sagen" (+194 px) and at "Vi tvivle ingenlunde" (+237 px).
% --- p. 119 ---
\setcounter{footnote}{0}
umiddelbart gjennem visse fromme Mænd, men fra den Tid af lod Menneskenes Tale
om de guddommelige Ting være uafhængig af enhver guddommelig Indvirkning. ---
Den samme Usikkerhed viser sig, naar Talen er om Forholdet mellem Skriftens
almindelige Idee og de enkelte Skriftsteder. De urgere selv, hvor det
convenerer dem, Skriftens enkelte Udsagn, sige stedse, naar man fremfører
Noget, der strider mod deres Anskuelser: ”Hvor staaer dette skrevet? Derom
findes ikke et Ord i Skriften!” og disputere saaledes i dette Tilfælde κατὰ
ῥητόν; --- anfører man derimod bestemte Skriftsteder imod dem, vide de stedse
at skjelne mellem Sagen selv og dens ydre Indklædning, idet de forsvare deres
Tænkefrihed med den Paastand, at man kun maa slutte fra Skriftens almindelige
Idee og hele Oeconomie; her disputere de altsaa κατὰ
νοητόν\footnote{Exempler paa begge disse Maader at disputere paa findes hos
Bretschneider: \textit{De obedientia activa} §~\textit{157}. (Pag.
\textit{285}. \textit{2} B.).}. Men en saadan Fremgangsmaade er en
”\textit{mystificatio veri}”, der let kan paadrage Forfatterne af denne Skole
Mistanke, som om de sagde Eet, men meente et Andet; thi den overordentlige
Ærbødighed for Skriften seer ofte ud som en Maske, paataget for under dens
Skjul at kunne give den subjective Frihed saa meget større Spillerum.

\parag{20}{Fortsættelse. Overgang til Panlogismen.}

Siig mig, hvilken Metaphysik Du følger, og jeg skal med Undtagelse af Dine
reent private Meninger sige Dig, hvilken dogmatisk Fortolkning af Skriften Du
vil fremføre! Siger Du, at Du slet ingen Metaphysik følger,
% THE § 20 HEAD AGREES WITH THE INDHOLD: the printed head reads "§ 20." in
% centred bold ANTIQUA over the centred letterspaced Fraktur argument
% "Fortsættelse. Overgang til Panlogismen.", word for word and point for point
% the Indhold's entry for p. 119. So there is nothing to record here of the
% kind found at §§ 7, 10, 13 and 15.
% p. 119 carries ONE note, marked 1) --- the plain next number. Its Latin
% title is antiqua and takes \textit{}; "Pag." and "B." are Fraktur and stand
% outside it, as do the section sign's figures' surroundings.
% GREEK: κατὰ ῥητόν and κατὰ νοητόν, read at 1.0x. Both are set with the
% fount's cursive kappa, typed here as plain κ per the header. The rejected
% reading for the second is ῥητόν again (the Fraktur OCR gave YONTOV, which is
% νοητόν); the sense --- the letter against the meaning --- confirms νοητόν.
% QUOTATION MARKS: four on this page, all the RAISED U+201D --- two round
% "Hvor staaer dette skrevet? ... i Skriften!" and two round the antiqua
% "mystificatio veri". Parity 0.
% NO letterspacing in the § 20 opening: "Siig mig, hvilken Metaphysik Du
% følger" looks open at 0.5x, but the same sentence runs on into p. 120's
% first line, "troe vi Dig ikke", which is unmistakably set solid at 3x; the
% two were cropped at the same scale and compared. The openness is the wider
% word-spacing of a short-set first paragraph, not Sperrsatz.
% --- p. 120 ---
\setcounter{footnote}{0}
troe vi Dig ikke; siger Du derimod, at Du ikke selv veed, hvilken Metaphysik
Du følger, ville vi troe Dig\footnote{Es hilft der Theologie nichts, sich
dagegen zu sträuben, zu sagen: sie wolle nichts von Philosophie wissen, es
seyen Philosopheme, also auf der Seite liegen zu lassen. Sie hat es immer mit
Gedanken zu thun, die sie mitbringt; und diese ihre subjectiven Vorstellungen,
Gedanken, ihre Haus= und Privat=Metaphysik sind dann die Reflexionen,
Meinungen u.~s.~f. der Zeit \ldots\ diese allgemeinen Vorstellungen sind
nichts, als was sich auf der Heerstraße findet, was auf der Oberfläche der
Zeit umherschwimmt. Hegel: Geschichte der Phil. (W. \textit{15} B. Pag.
\textit{268}).}. Det er da ikke ved fornyede historiske eller philologiske
Undersøgelser, at Kritiken kan haabe at gjøre noget Fremskridt, men ene og
alene ved en Forandring af sin Metaphysik. Kan altsaa Kritiken forsvare sin
Metaphysik gjennem alle dens enkelte Momenter imod Dialektikens Indvendinger,
da vil ogsaa den kritiske Theologie staae fast og hævde sit Navn, som den
høieste Videnskabsform, hvortil Menneskene i dette Liv kunne komme; kan den
ikke dette, ere vi berettigede til at opsøge en anden Form for den sande
Videnskab. I ethvert Tilfælde staaer dette fast, at den ikke maa søge sit
Forsvar i den hellige Skrift, men at den maa erindre, hvad Clemens siger ---
en Yttring, hvortil der i vor Tid ofte appelleres\footnote{Sammenlign Fr.
Baader, Vorles. üb. Dogm. \textit{1}. H. Pag. \textit{17}. J. Martensen. S.
A. §~\textit{1}.} ---: ”\emph{Philosophere maa man, selv da, naar man vil
bekjæmpe Philosophien!}”\footnote{Sandheden heraf have ogsaa de aandrigeste og
berømteste kritiske Theologer selv erkjendt: f.~Ex. De Wette: \ldots\ ich
schmeichele mich mit der Hoffnung, daß ich wenigstens die allgemeine Ansicht,
welche für die Theologie von entscheidender Wichtigkeit ist, klar genug
gemacht habe. Und diese ist keine andere, als \emph{die Unterscheidung der
Verständigen, idealen und ästhetischen Ueberzeugung}, welche ich für den
Schlüssel der ganzen Theologie halte. Relig. und Theol. Pag. \textit{IV}.
sammenlign ogsaa Pag. \textit{162---172}.

Saaledes søger ogsaa en Kritiker at bekjæmpe Weise med metaphysiske Vaaben i
sin Bedømmelse af dennes Bog: ”Die ewangelische Geschichte” (Kritische
Prediger=Bibliothek v. D. J. F. Röhr, \textit{20}. Band, \textit{29}. Heft,
Pag. \textit{200---219}). Vel synes han, idet han dadler Weises Lære om en
substantiel Eenhed af Christus, Menneskene og Gud, at appellere til Skriftens
Auctoritet: ”Diese Einheit aber der Menschen mit Gott ist nicht eine
substantiale, ein Hineinbilden unserer creatürlichen Selbstheit in die göttl
che Substantz, es ist eine moralische nach der ausdrücklichen Lehre Jesu und
aller seiner Apostel” \ldots\ men betragte vi hans Indvendinger nærmere, viser
det sig let, at de henhøre til Metaphysiken; han forkaster nemlig først hiint
substantielle Forhold, fordi han ikke kan begribe det: ”Was sollen wir uns
aber auch unter jenem Hineinbilden in die göttliche Substantz Vernünftiges
denken?” Dernæst er han overbeviist om, at det substantielle Forhold strider
imod Guds sande Idee og derfor ikke kan tjene til Moralitetens Fremme: ”und
indem so der Rationalismus die pantheistischen Gedanken von einer
Menschwerdung Gottes, einem sich Hineinbilden in die göttliche Substanz
verwirft, weil sie einerseits unerweisbar und der reinen Gottes=Idee
wiedersprechend, desshalb aber auch andererseits für Beförderung der Moralität
gleichgiltig erscheinen \ldots” Og for at vi ikke skulle indvende, at hans
Paastande, af hvilken den første kan forekomme Læseren rimeligere, end den
sidste, ere uden sand Beviiskraft, erklærer Forf., at han ønsker enhver
Undersøgelse philosophisk gjennemført, ja han forsikkrer endog, at alle
saadanne Spørgsmaal alt for længe siden vilde have været tilfredsstillende
besvarede, hvis ikke tilfældigviis den speculative Philosophie for en Tid
havde fortrængt Kriticismen: ”Hätte man hier die Bahn verfolgt, welche der
unsterbliche Kant und die Anhänger seiner Schule betraten, die völlige
Aussöhnung der philosophischen Speculation mit dem reinen und einfachen
Christenthume würde zum Heile der ewangelischen Kirche bald erfolgt, aller
unnützen philosophischen Grübelei über Dinge, die wir nun einmal nicht zu
erkennen vermögen, und die darum auch für unser wahres sittlich=religiøses
Leben keinen Werth haben, würde vielleicht für immer vorgebeugt, und somit ein
Höhepunkt deutscher Philosophie erreicht worden seyn, den diese Wissenschaft
in keinen unserer cultivirten Nachbaarstaaten so bald erreichen dürfte(!!!).
Da trat die Speculation des absoluten Philosophie dazwischen \ldots\ Trotz dem
aber werden diese Extreme philosophischer Speculation ihre heilsamen Folgen
haben; sie werden darauf zurückführen, daß wir den Grund der Erkenntniß
göttlicher Dinge allein in den Anforderungen unserer sittlichen Natur suchen
\ldots” Efterat have anført alle disse vægtfulde Grunde henter han sin sidste
Indvending fra Skriften, idet han, for at give sine Bemærkninger større
Eftertryk og med Skriftens Taushed bringe Weise til Taushed, disputerer κατὰ
ῥητόν: ”Wo lesen wir in den heiligen Urkunden ein Wort von einem Hineinbilden
des creatürlichen Selbst in die göttlichen Substanz?”}.
% p. 120 carries THREE notes, marked 1) 2) 3) --- the plain next numbers, and
% the body marks and the note marks agree.
% *** NOTE 3) CROSSES TWO PAGE BOUNDARIES. IT IS COMPLETE HERE. ***
% It fills p. 120's footnote block (seven lines under the rule at y 3277 of
% PDF 129), then ALL of p. 121's footnote block (forty lines under a rule of
% its own, below only TWO lines of body), and then ALL of p. 122's footnote
% block (sixteen lines), and there it ends. NEITHER p. 121 NOR p. 122 has a
% mark anywhere in its footnote block: every line of both blocks stands either
% at the notes' text column or at the notes' paragraph indent, and no line
% stands out at the mark position (about 126 px left of the text column, as
% measured on p. 109 and p. 120). The joints are verbal as well as
% geometrical: "Schlüssel der gan-" / "zen Theologie halte", and "würde
% vielleicht" / "für immer vorgebeugt". The whole note is therefore set here,
% at its mark, in two paragraphs --- the printing begins "Saaledes søger" on a
% fresh indented line, exactly as at the p. 102 note.
% CONSEQUENCE FOR THE NEXT BATCH (pp. 121--132): p. 121 and p. 122 carry NO
% note of their own, take NO \setcounter{footnote}{0}, and their footnote
% blocks must NOT be transcribed again --- they are already above.
% EMPHASIS: the Clemens sentence "Philosophere maa man ... Philosophien!" is
% letterspaced in the body (read at 0.55x, and flagged RUN by spacing.py), and
% inside note 3 "die Unterscheidung der Verständigen, idealen und ästhetischen
% Ueberzeugung" is letterspaced, beginning at the "die" that ends its line.
% QUOTATION MARKS: fourteen on this page and its two continuation blocks, all
% the RAISED U+201D at both ends --- the Clemens sentence, and in note 3 "Die
% ewangelische Geschichte", "Diese Einheit ... Apostel", "Was sollen ...
% denken?", "und indem ... erscheinen", "Hätte man ... suchen" and "Wo lesen
% ... Substanz?". NOT ONE note in pp. 109--120 uses the low-high pair „…“;
% the departure recorded at pp. 48, 70, 79, 84 and 106 does not recur here.
% Parity 0.
% AS PRINTED, uncorrected, in note 3 (none of it in the Rettelser):
%   * "ewangelische" / "ewangelischen" for "evangelische(n)", twice, with a
%     Fraktur w both times.
%   * "Weise" for Weisse (C. H. Weisse), three times --- once "Weises".
%   * "die göttl che Substantz" --- the i of "göttliche" has failed to print;
%     there is a clean word-width gap where it stood. Read at 0.58x.
%   * "Substantz" with the tz ligature twice (p. 121, read at 2x: the sort is
%     t joined to z, not the ß of "Bedürfniß") but "Substanz" with a plain z
%     twice (pp. 121, 122); "Trotz dem", also the tz ligature.
%   * "desshalb" set with TWO long s, not ß. Read at 3x and compared with the
%     certain ß of "Bedürfniß" in the p. 109 note: that ß is one sort, a long
%     s carrying a z-tail; here there are two separate stems, each with its
%     own top flag. Likewise "wiedersprechend" for "widersprechend";
%     "in keinen unserer cultivirten Nachbaarstaaten"; "des absoluten
%     Philosophie"; "dürfte(!!!)" with no space before the bracket.
%   * "sittlich=religiøses" --- the Danish ø again in a German word, as in the
%     p. 109 note. Read at 0.58x.
% DOUBTFUL READING: "J. Martensen" (note 2). The sort is the Fraktur I/J,
% which this fount does not distinguish --- the identical sort stands for
% Röhr's certain "J." in note 3, and the Fraktur H of "Heft" three lines above
% it is a wholly different letter with no descender, so the printed initial is
% not H. H. L. Martensen is meant (S.~A. = Den menneskelige Selvbevidstheds
% Autonomie, cited as "Dr. H. Martensen" on p. 106), so the initial is wrong
% however it is typed. Set as "J." to match the treatment of the same sort in
% "J.~G. Walch" at p. 84; the rejected alternatives are "I." and "H.".
% The printed ellipses run to three and four points; all are set as \ldots.

% --- p. 121 ---
% NO NOTE OF ITS OWN, so no \setcounter{footnote}{0} here. The whole footnote
% block on this page (and on p. 122) is the RUNOVER of the note that begins on
% p. 120 --- there is no mark at its head, and its first paragraph completes a
% German word broken at the foot of p. 120 ("...zen Theologie halte."). The
% runover text is transcribed in the comment block below and MUST BE APPENDED
% to p. 120's \footnote{} when the pp. 109--120 fragment is spliced. Checked at
% 0.6x: the second paragraph ("Saaledes søger ogsaa en Kritiker...") carries a
% plain paragraph indent and NO footnote mark, so it is not a note of its own.
Vi opfordres saaledes til at gjøre Kritikernes Metaphysik til Gjenstand for en
metaphysisk Undersøgelse. --- Det følger
%
% ============================================================================
% RUNOVER of the p. 120 note, set on pp. 121--122 below the rule.
% CALLER: append everything between the two ==== lines to the footnote that
% opens on p. 120. It is NOT a separate note and must not be given a mark.
% The word broken at the p. 120/121 joint needs `\-%` at the foot of p. 120.
% ----------------------------------------------------------------------------
% zen Theologie halte. Relig. und Theol. Pag. \textit{IV}. sammenlign ogsaa
% Pag. \textit{162---172}.
%
% Saaledes søger ogsaa en Kritiker at bekjæmpe Weise med metaphysiske Vaaben i
% sin Bedømmelse af dennes Bog: ”Die ewangelische Geschichte” (Kritische
% Prediger=Bibliothek v.\ D.\ J.\ F.\ Röhr, \textit{20}. Band, \textit{29}.
% Heft, Pag. \textit{200---219}). Vel synes han, idet han dadler Weises Lære om
% en substantiel Eenhed af Christus, Menneskene og Gud, at appellere til
% Skriftens Auctoritet: ”Diese Einheit aber der Menschen mit Gott ist nicht eine
% substantiale, ein Hineinbilden unserer creatürlichen Selbstheit in die göttl
% che Substanz, es ist eine moralische nach der ausdrücklichen Lehre Jesu und
% aller seiner Apostel”\ldots\ men betragte vi hans Indvendinger nærmere, viser
% det sig let, at de henhøre til Metaphysiken; han forkaster nemlig først hiint
% substantielle Forhold, fordi han ikke kan begribe det: ”Was sollen wir uns
% aber auch unter jenem Hineinbilden in die göttliche Substanz Vernünftiges
% denken?” Dernæst er han overbeviist om, at det substantielle Forhold strider
% imod Guds sande Idee og derfor ikke kan tjene til Moralitetens Fremme: ”und
% indem so der Rationalismus die pantheistischen Gedanken von einer
% Menschwerdung Gottes, einem sich Hineinbilden in die göttliche Substanz
% verwirft, weil sie einerseits unerweisbar und der reinen Gottes=Idee
% wiedersprechend, desshalb aber auch andererseits für Beförderung der Moralität
% gleichgiltig erscheinen \ldots” Og for at vi ikke skulle indvende, at hans
% Paastande, af hvilken den første kan forekomme Læseren rimeligere, end den
% sidste, ere uden sand Beviiskraft, erklærer Forf., at han ønsker enhver
% Undersøgelse philosophisk gjennemført, ja han forsikkrer endog, at alle
% saadanne Spørgsmaal alt for længe siden vilde have været tilfredsstillende
% besvarede, hvis ikke tilfældigviis den speculative Philosophie for en Tid
% havde fortrængt Kriticismen: ”Hätte man hier die Bahn verfolgt, welche der
% unsterbliche Kant und die Anhänger seiner Schule betraten, die völlige
% Aussöhnung der philosophischen Speculation mit dem reinen und einfachen
% Christenthume würde zum Heile der ewangelischen Kirche bald erfolgt, aller
% unnützen philosophischen Grübelei über Dinge, die wir nun einmal nicht zu
% erkennen vermögen, und die darum auch für unser wahres sittlich=religiöses
% Leben keinen Werth haben, würde vielleicht
%       [p. 122 begins here, still inside the same note]
% für immer vorgebeugt, und somit ein Höhepunkt deutscher Philosophie erreicht
% worden seyn, den diese Wissenschaft in keinen unserer cultivirten
% Nachbaarstaaten so bald erreichen dürfte(!!!). Da trat die Speculation des
% absoluten Philosophie dazwischen \ldots\ Trotz dem aber werden diese Extreme
% philosophischer Speculation ihre heilsamen Folgen haben; sie werden darauf
% zurückführen, daß wir den Grund der Erkenntniß göttlicher Dinge allein in den
% Anforderungen unserer sittlichen Natur suchen \ldots” Efterat have anført alle
% disse vægtfulde Grunde henter han sin sidste Indvending fra Skriften, idet
% han, for at give sine Bemærkninger større Eftertryk og med Skriftens Taushed
% bringe Weise til Taushed, disputerer κατὰ ῥητὸν: ”Wo lesen wir in den heiligen
% Urkunden ein Wort von einem Hineinbilden des creatürlichen Selbst in die
% göttlichen Substanz?”
% ----------------------------------------------------------------------------
% READINGS SETTLED ON THE IMAGE IN THIS RUNOVER (all at 0.6x or better):
%   * "Pag. IV." --- bold ANTIQUA roman numeral, four; NOT "LIV" as the Fraktur
%     OCR gave. Read at 0.95x.
%   * "göttl che Substanz" (l. 8 above) --- the "i" of "göttliche" failed to
%     print; transcribed as the printer left it.
%   * "desshalb" --- read at 3.2x: TWO long s, not the ß ligature (compare
%     "daß"/"Erkenntniß" on p. 122, where the ß is unmistakably long-s + ezh).
%     Rejected reading: "deßhalb".
%   * "Trotz dem" --- read at 2.2x: the ligature after "Tro" is t + ezh (short first
%     element), not ß. Rejected reading: "Troß dem".
%   * "wiedersprechend", "Nachbaarstaaten", "ewangelische", "Weise" (for
%     Weisse) --- all as printed.
%   * κατὰ ῥητὸν: printed with a GRAVE on the omicron, where the grammar has
%     ῥητόν; kept as printed. The fount's cursive kappa (U+03F0) is normalised to plain kappa.
%   * The printed spaced points are rendered \ldots per the house convention.
%   * Quotation marks: the raised ” at BOTH ends throughout this note, exactly
%     as in the body. Not one low-high Danish pair anywhere in it.
%     12 marks, 6 pairs.
% ============================================================================
% --- p. 122 ---
af Kriticismens egen Natur, at den adskiller de positive Momenter fra hverandre
(thi κρίσις er nerop den Virksomhed, hvorved man adskiller, afsondrer
% PRINTER'S DEFECT, transcribed as printed: "nerop" for "netop" (r for t). Read
% at 0.95x --- the letter has no ascender and no cross-stroke, so it cannot be
% a t. Both OCR witnesses agree on "nerop".
Gjenstandene). Erkjendelsens Rige berører Troens Sphære, men ikke saaledes, at
der nogensinde fremkommer nogen Identitet af dem; thi hvad der gjør, at de
berøre hinanden, det Samme gjør tillige, at de bortstøde hinanden. Enhver af
dem synes at acquiescere i sig indenfor sine egne Grændser og ingen fremmede
Elementer at optage i sig: de forholde sig altsaa ligegyldige til hinanden. Var
nu denne Ligegyldighed absolut, vilde Theologien være aldeles overflødig, ja
fornuftstridig; thi har den theologiske Videnskab ingen anden Gjenstand, end de
endelige Tings Rige, og maa paa den anden Side den over al Erkjendelse ophøiede
Tro selv varetage sine Gjenstandes Interesse, saa er ogsaa Troen sikker
indenfor sit eget Riges Grændser, saa kan kan og maa den ogsaa paa egen Haand
% PRINTER'S DEFECT, transcribed as printed: "saa kan" ends one line and "kan og
% maa den" opens the next --- the word "kan" is set twice (dittography across
% the line break). Read at 0.5x; both OCR witnesses show the doubling.
opfatte de guddommelige Ideer. Men denne gjensidige Udelukkelse ophæver sig
selv ved en nærmere Betragtning. Thi skal den sig selv overladte Tro fremstille
sine Gjenstande i deres Orden og rette Forhold, maae ogsaa de religieuse
Følelser staae i Mod\-%
% --- p. 123 ---
sætning til og indskrænke, begrændse hverandre, og saaledes vil nu
Erkjendelsen, der er alle endelige, begrændsede Gjenstandes Ordner og Styrer,
gjennembryde alle Troens Grændser. Paa den anden Side er heller ikke
Erkjendelsens Rige indenfor sine egne Grændser sig selv nok; thi skulle dettes
endelige Elementer bestemme hverandre indbyrdes, ophæves Grændserne imellem
dem, eftersom det Endelige som saadant er ligt det Endelige og ikke kan
begrændses af Andet end af sin egen Modsætning, \dsi{} af det Uendelige. Det er
saaledes aldeles vist, at Følelsens uendelige Rige paa ethvert Punkt berører
Erkjendelsens endelige Rige, og at de som en Følge deraf ikke kunne vedblive at
forholde sig ligegyldige til hinanden.

Og dog kan man paa den anden Side ligesaalidt haabe, at de skulle afgive nogen
Begrændsning for hinanden, thi de tilbagestøde hinanden, ligesaaofte de berøre
hinanden; man kunde saaledes med Rette kalde den Theologie, der fremgaaer af
denne Methode, fortryllet, efterdi den paa eengang undviger og elsker Tanken.
Det er Erkjendelsen, der bestemmer Former og Skikkelser, saa at Enhver, der vil
fastholde Troens Gjenstande, er nødt til at apellere til denne, men da den
theologiske Selvmodsigelse ikke blot gjælder i Almindelighed, men ogsaa gjælder
i ethvert enkelt Tilfælde, og da den Sætning, at det Endelige ikke kan rumme
det Uendelige, paa mangfoldig Maade vender tilbage, griber Erkjendelsen stedse
\emph{efter} Troens Gjenstande, men \emph{griber} dem i Virkeligheden aldrig.
% Both letterspaced, read at 0.5x and confirmed on the line: "e f t e r" and
% "g r i b e r". spacing.py flagged both as "single?"; they are real.
Det er forgjeves, at den prøver Begivenhederne, thi den naturlige Pragmatisme
er Rettesnoren for alle dens Undersøgelser: kunde den nu ved at følge denne Vei
bevise, at Begivenhederne have fundet Sted, vilde jo alt Overnaturligt i dem
forsvinde; --- finder den derimod, at de staae i Modsigelse til den sædvanlige
Tingenes Orden, maa den jo erklære dem for
% "Tingenes", not "Tiugenes": checked at 0.95x against the n of "den" on the
% same line. The p. 3 turned-n defect does not recur here.
% --- p. 124 ---
usande. --- Erkjendelsen kræver altsaa, at Begivenheder, der ere overnaturlige,
skulle sees at have tildraget sig paa naturlig Maade, medens den dog paa samme
Tid erklærer, at disse Riger ere hinanden indbyrdes modsatte. Ønske vi derfor
at anskue Troens Gjenstande, maae vi gjøre Afkald paa Tænkningen; kun under den
Betingelse nedstige disse fra Himlen i begrændsede, bestemte Former. Da
% PRINTER'S DEFECT: the "e" of "bestemte" failed to print --- at 1.55x only two
% specks of ink remain between the b and the st-ligature, so the page reads
% "b-stemte" with a bare speck where the e should be.
% Broken or under-inked type, not a compositor's error; the intended letter is
% supplied here and the defect logged at its site.
saaledes Anerkjendelsen af de guddommelige Ting forudsætter en Opgivelse af
Erkjendelsen, staaer som en Følge deraf Theologien udenfor Tænkningens Sphære;
ja jo oftere man tænker over dette Forhold, destomere forviklet seer det ud, og
Forviklingen gaaer i det Uendelige. Christus deles i denne Theologie imellem
begge Riger; til Følelsens og Troens Rige henhører Guds Søn, til Erkjendelsens
Rige derimod Menneskets Søn. Man skulde altsaa vente, at disse Christi Naturer
paa Grund af deres gjensidige Ligegyldighedsforhold maatte isolere sig fra
hinanden, saa at det var Troen og Følelsen uvedkommende, paa hvad Maade
Mennesket Christus blev opfattet, og Troen ikke behøvedes for at begribe hans
menneskelige Natur. --- Men betragte vi Sagen nøiere, see vi let, at den
menneskelige Natur ikkun har den Hensigt at bestemme den guddommelige, og at
den ikke paa nogetsomhelst Punkt kan holdes adskilt fra samme, da jo den
guddommelige Herlighed overalt fremskinner igjennem den, saa at i Mennesket
Christus den hele Christus er aabenbaret for den æsthetiske Følelse: vi synes
altsaa umiddelbart at skue, tænke, begribe Christus. Men saa ofte vi nu ville
tænke denne vor Tanke, vender atter det Ligegyldige og Modsatte i Forholdet
mellem Christus som Menneske og Christus som Gud tilbage; Mennesket Christus er
da bundet til den lavere Sphære, medens den guddommelige Christus opfarer til
Himlen? Erkjendelsen urgerer atter det fra Gudsnaturen isolerede Menneske,
% The question mark after "Himlen" is as printed (read at 0.5x); the sense
% would take a full stop. Recorded, not corrected.
medens Følelsen griber efter den fra Men\-%
% --- p. 125 ---
neskenaturen isolerede Gud: Christus bliver et almindeligt Menneske, en
unævnelig Gud: han forsvinder som et skuffende Taagebillede, saa ofte vi gribe
efter ham; og vi staae saaledes trufne af Apostelens Ord: Ὠ ἀνόητοι, τίς ὑμᾶς
ἐβάσκανε; οἱς κατ' ὀφθαλμοὺς Ἰησοῦς Χριστὸς προεγράφη.
% GREEK (Gal. 3,1), read at 2.4---2.6x. Normalisations, per the header:
%   * the fount's cursive kappa (U+03F0) is normalised to plain kappa and its
%     script theta (U+03D1) to plain theta, per the header.
%   * Ὠ: only a smooth breathing is printed before the capital omega --- no
%     circumflex, where the grammar has psili AND perispomeni (U+1F6E).
%     Kept as printed.
%   * οἱς: the rough breathing is printed over the OMICRON and there is no
%     circumflex; the placement is normalised to the standard one, as the
%     header prescribes for this fount. The grammar carries a circumflex on
%     the iota as well (U+1F37); none is printed.
%   * Ἰησοῦς: the circumflex of the -ου- diphthong sits over the omicron in the
%     printing; placement normalised.
%   * ἐβάσκανε (not ἐβάσκανεν) as printed; the mark after it is the Greek
%     question mark.
%   * κατ' --- the elision mark is set as a raised apostrophe.

Denne Vanskelighed har fornemmelig sin Grund deri, at den profane Erkjendelse,
der ikke kan begribe den guddommelige Natur, men alene fatter de Ting, som
hører til denne Verden, staaer urørt og kold tilbage, medens Følelsen og
Villien indvies Christus for ved ham at helliges og besjeles af Guddomsfylde.
Saalænge derfor Modsætningen bestaaer mellem Christus og Verden, saalænge vil
og hiin Splid bestaae; skal Troen bestemmes, er den nødsaget til at tye hen til
den hedenske Erkjendelse. Da altsaa Følelsen er uendelig og derfor uden bestemt
Skikkelse, er det Erkjendelsen, der fastsætter Grændserne og bestemmer, hvad
der ligger indenfor dens eget, og hvad der ligger indenfor Troens Gebeet. Den
kan da ogsaa let, naar den finder for godt, forvise den christelige Tro af sit
Rige og bemægtige sig det hele Herredømme. Er Fornuften bleven sig denne sin
Myndighed bevidst, vil den snart constituere sig selv som absolut Eneherre og
uden Sky erklære, at alle discrete, himmelske og jordiske, overnaturlige og
naturlige Momenter ere at søge indenfor dens egen Tænknings Grændser; men
saaledes fremgaaer \emph{den objective Rationalisme}. --- Tænkningen er
Fornuft, Fornuften er Gud; men Tænkningen er en Negation af Phænomenerne;
saaledes forsvinder det sidste Spor af Positivismen. Ere Phænomenerne negerede,
da er man kommen til den speculative Logik, det er, den rene Værens Begreb,
Gud, som han var før Verdens Skabelse. Men ligesom den rene Væren ikke kan være
uden gjennem et virkeligt Phænomen, saaledes er heller ikke Gud uden igjennem
Verden, og herved har man befriet sig fra den
% --- p. 126 ---
sidste Rest af den ældre Supranaturalisme. Verden er dernæst et Complex af
endelige Phænomener; betragtes den umiddelbart for sig, er den i Modsætning til
Gud. Men da denne Modsætning ikkun angaaer Tingene i deres umiddelbare
Tilværelse, vil der indtræde en fuldkommen Overeensstemmelse mellem Gud og
Verden, naar man opgiver enhver, om man saa tør sige, raa Umiddelbarhed.
Tilværelsens enkelte Momenter ere da, i og for sig betragtede, Intet, skulle
kun betragtes og bedømmes i deres dialektiske Overgang i andre Momenter; der er
saaledes kun \emph{en logisk Progres}. Vel sættes den absolute Aands
% The letterspacing takes in "en": checked at 0.95x, "e n  l o g i s k
% P r o g r e s" against the solid "kun" that precedes it.
Aabenbaring som sammes endelige Maal; men dennes absolute Kriterium er ene og
alene Immanentsen af Idealitet og Realitet, Væsen og Phænomener. Man kan
saaledes med Rette kalde dette System \emph{Panlogisme}, efterdi dets Udvikling
er en Bevægelse fra Logiken igjennem logiske Momenter til Logiken.

Den sande Opfattelse af Tingene i denne deres nye, kosmiske Mediation fordrer,
at man skjelner mellem Principets potentielle Væren og dets Aabenbaring. I
Principets dunkle Skjød er den umiddelbare Identitet af Væsen og Phænomen
skjult tilstede, men først gjennem dennes Udtrædelse i sine mangfoldige,
forskjellige Momenter, fremkommer der en sand, concret Immanents. Væsenet maa
derfor fra sin rene Værens Sphære stige ned i Phænomenernes ydre Rige; der er
altsaa noget Mere at søge i Phænomenerne end de blotte Phænomener, og heri
ligger Sporen til videre Fremadskriden. Thi forsaavidt Væsenet trænger til
Phænomenet, forsaavidt sætter og bevarer det samme; men naar det igjennem dette
har udviklet sin Tilværelse, forsaavidt det indenfor dettes Grændser var
muligt, nedbryder det Grændserne og tilintetgjør det forældede Phænomen for i
dettes Sted at lade et nyt og fortrinligere
% --- p. 127 ---
fremgaae. Og spørger man, hvorfor det fortrinligere nu først fremtræder efter
hiint og ikke strax i Begyndelsen er tilstede, maa hertil svares, at det andet
Phænomens Fortrin netop er medieret igjennem det førstes Tilværelse og
Undergang; thi netop ved at sætte og atter destruere det første har Væsenet
erholdt ny Kraft. Under Phænomenernes Alteration holder Væsenet stedse fast ved
sig selv, og bevarer, ligesom Substantsen, i ethvert nyt Phænomen, det skaber,
Alt hvad der i ethvert foregaaende Phænomen er væsentligt, medens det
destruerer Alt, hvad der er forkrænkeligt og uadæqvat.

Den Maade, hvorpaa denne Udvikling foregaaer, er forskjellig i de forskjellige
Sphærer. I det lavere Naturrige er Væsenets Kraft til at vedligeholde sig selv
svagere; derfra hidrører den tautologiske Gjentagelse og eensformige
Cirkelgang, der bemærkes hos dette Riges Phænomener. Progressen i sin sande
Udvikling maa søges i det Rige i Tilværelsen, hvor der midt i
Uovereensstemmelsen imellem Væsenet og Phænomenet tillige viser sig det fasteste
Baand imellem begge: hvor Væsenet vender tilbage til sig selv gjennem en
Negation og Restauration af Phænomenerne: hvor Alterationen af de ydre Momenter
tillige er en Assimilation af de indre: hvor Phænomenerne kunne bevares i en
Erindring \dsi{} i \emph{Menneskeslægtens Historie}.

Gud har altsaa efter at være udtraadt af sin rene Idealitets Rige sat synlige
Phænomener i Naturens Rige for gjennem Negationen af disse at komme til den
subjective Aand, \dsi{} til den menneskelige Selvbevidsthed. Vor Selvbevidsthed
henhører til to Sphærer. Bundne til de endelige og sandselige Phænomeners Rige
ere vi Syndere, hjemfaldne til Døden; men opløftede til det absolute Væsens
Sphære gjennem Negationen af Phænomenerne have vi Gud nærværende i os og ere
saaledes forsonede med ham. \emph{Historiens absolute}
% The letterspacing runs on across the page boundary: "H i s t o r i e n s
% a b s o l u t e" ends p. 127 and "G r u n d l a g  e r  d e r f o r
% R e l i g i o n e n , d e n s  a b s o l u t e  M a a l  G u d s
% A a b e n b a r i n g ." fills the first two lines of p. 128. It is set as
% two \emph{} groups so that neither straddles the page marker; the rendered
% text is identical.
% --- p. 128 ---
\emph{Grundlag er derfor Religionen, dens absolute Maal Guds Aabenbaring.}

Da saaledes hele den kosmiske Mediation maa betragtes i Religionens Lys,
gjælder det Samme, vi her have udviklet om Forholdet mellem Væsenet og
Phænomenet i Almindelighed, tillige om det Baand, hvorved Religionen forbinder
Menneskene med Gud. Gud er først umiddelbart tilstede for den menneskelige
Selvbevidsthed, saa at han kan paapeges synlig i et bestemt, enkelt Phænomen,
\dsi{} Gud aabenbares. Men i enhver Aabenbaring er et doppelt Moment at
iagttage, et phænomenalt og et aandeligt. Thi Organet for Aabenbaringen,
Phænomenet, kan hverken anerkjendes eller dyrkes som saadant, hvis ikke det
samme Religionsprincip, der er tilstede i Phænomenet, ogsaa er tilstede i den,
der betragter Phænomenet. Religionens Idee uddannes derfor i enhver Betragter
fra begge Sider; paa den ene Side betragter han det særegne Phænomen og knytter
sig fast til dette; --- paa den anden Side søger han at udfinde og forklare sig
Phænomenets Væsen og negerer saaledes det særegne, tilfældige Phænomen for at
gjøre det Almindelige, Væsentlige, der fremskinner i det, til sin sande
Eiendom. Og her gjælder den Sætning, som \emph{Gud er i Phænomenet, saaledes er
han ogsaa i Bevidstheden.} Ligesom derfor ingen Religion og ingen
Gudsaabenbaring kan gives uden igjennem et sandseligt Phænomen, saaledes er paa
den anden Side intet Phænomen et umiddelbart eller absolut Udtryk af Gudsideen.
Her maa derfor finde en Mediation Sted, for at det empiriske Phænomen kan blive
negeret og et nyt, fuldkomnere, Religionen mere tilfredsstillende Phænomen ved
Ideens egen Kraft træde i dettes Sted. Denne Negation og Restauration foretager
den menneskelige Selvbevidsthed sig efter et vist Instinct, saa at
% --- p. 129 ---
\setcounter{footnote}{0}
det Samme, der bevirker, at Phænomenet udsmykkes i Ideens Interesse, tillige
skaffer det Tiltro som objectivt og reelt.

Herfra hidrøre alle den hellige Histories Beretninger om Guds Aabenbaringer og
Miraklerne; og man vil nu let forstaae Overeensstemmelsen og
Uovereensstemmelsen mellem Panlogismen og Kriticismen. Begge urgere nemlig paa
det Strængeste en pragmatisk Naturudvikling, som umuligt kan afbrydes ved en
høiere Magts Mellemkomst\footnote{Diese Ueberzeugung ist so sehr Bewußtseyn der
neuen Welt geworden, daß im wirklichen Leben die Meinung oder Behauptung, eine
übernatürliche Ursache, eine göttliche Wirksamkeit habe irgendwo unmittelbar
eingegriffen, geradezu als Unwissenheit oder Betrug betrachtet wird. Strauß.
L.~J. § \textit{14} Pag. \textit{91}.}; Begge drage Miraklet i Tvivl. Thi det
% One note, marked 1) --- the plain next number. The mark stands after
% "Mellemkomst", with a space before the semicolon in the printing. The note is
% German Fraktur throughout; only the figures are antiqua, hence \textit{}.
sande Mirakelbegreb udfordrer, at der maa finde en virkelig Modsætning Sted
imellem Guds Villie og Naturens Love, saa at Naturens Momenter fremgaae af
hverandre efter en vis Nødvendighedslov; vare nemlig Naturgjenstandene i deres
almindelige Udvikling umiddelbare Udtryk af Guds Villie og Raadslutninger,
vilde de jo alle være Mirakler, og der vilde da slet intet Mirakel finde Sted;
men i Modsætningen maa der tillige være en Identitet tilstede, thi havde Gud
ikke selv fastsat Naturlovene, og kjendte han dem ikke, kunde han heller ikke
forandre dem. Troen paa Mirakler forudsætter altsaa en Identitet af Guds Villie
og Naturlovene, men antager tillige, at denne Identitet i den historiske
Udvikling gaaer over til en Modsætning, saa at der kunne indtræde Collisioner,
--- antager, at der i det guddommelige Forsyn i Almindelighed viser sig en
Omsorg for Menneskeslægten i Særdeleshed, ifølge hvilken Gud, uagtet vor
phænomenale Verden ikke svarer til Guds eget Væsen, hvor han finder det
fornødent, i særegne Øiemed benytter og for\-%
% --- p. 130 ---
herliger de naturlige Phænomener: saa at ikke blot Menneskene troe i et eller
andet Phænomen at see en særegen Aabenbaring af Gud, men ogsaa Gud selv ifølge
en særegen Raadslutning fremstiller sig i dette Phænomen for Menneskene.

Kriticismen tvivler paa Muligheden af Mirakler, fordi Gud, uagtet han er
fuldkommen i sig, er saa ophøiet, at det vilde være under hans Værdighed at
nedstige i sandselige Phænomener, --- saa viis og uforanderlig i sine
Raadslutninger, at han aldrig vilde forandre de Love, til hvilke han fra
Begyndelsen har knyttet Naturgjenstandenes Udvikling. Den skarpe Modsætning
imellem Guds abstracte Væsen og Phænomenernes Rige gjør det altsaa umuligt for
den til Deisme heldende Kritiker at troe paa Miraklerne.

Panlogisterne derimod statuere en Immanents af Gud og Verden; dog gjælder
Immanentsen kun den almindelige Verdensorden, --- tages Naturgjenstandene i
deres Enkelthed, gjør Modsætningen sig gjældende. At Gud nu selv i
Virkeligheden skulde være paa særegen Maade tilstede i dette eller hiint
Phænomen, er af tvende Grunde nmuligt. For det Første maatte Gud da have en
% PRINTER'S DEFECT, transcribed as printed: "nmuligt" for "umuligt" --- a
% turned u. Measured at 1.2x against "Grunde" three words earlier on the same
% line: the u of "Grunde" carries the fount's breve, the first letter of
% "nmuligt" carries none, while the u of "-mu-" does. Cf. "absolnt" p. 2.
oververdslig, i sig reflecteret Selvbevidsthed, for at kunne tænke paa en
særegen Aabenbaring --- men dette strider aabenbart imod den kosmiske
Mediation, vi før omtalte ---, og for det Andet maatte den uendelige Gud være
tilstede i et endeligt Phænomen; men dette er ligesaa umuligt. Phænomenernes
Forklarelse, Tanken om specielle Aabenbaringer kan da ikke henføres til nogen
objectiv Virksomhed af Gud, men kan kun hidrøre fra subjective, menneskelige
Forestillinger. Uagtet nu saaledes de hellige Beretninger ene og alene maae
antages at hidrøre fra subjectiv Religieusitet, kunde de dog ikke opstaae uden
en eller anden tilstrækkelig Grund, da de gjælde og finde Tiltro hos Flere paa
eengang. Grunden og Anledningen er altsaa
% --- p. 131 ---
\setcounter{footnote}{0}
det objective Moment i Mirakelberetningerne, Formerne derimod, hvori de ere
indklædte, skyldes Tilskuerne og Tilhørerne: man kan derfor med Rette kalde
saadanne Beretninger \emph{Myther}.
% "M y t h e r" is letterspaced --- read at 0.5x. spacing.py reported no
% letterspacing at all on this page; it is a false negative on a short word,
% exactly as BATCH-AGENT.md warns.

Der ere altsaa i Mytherne tvende Momenter at fastholde, et objectivt og et
subjectivt; paa mangfoldig Maade gaae disse over i hinanden, og det er i
Udviklingen og Bedømmelsen af disse Momenters gjensidige Overgang i hinanden,
at Kriticismen afviger fra Panlogismen. Kritikerne anprise nemlig
Hverdagserfaringen og betragte Gjenstandene udvortesfra; de ville derfor have
en haandgribelig Virkelighed; thi da Mediationens Begreb er dem dunkelt, kunne
de ikke forstaae, hvorfor man ikke strax fra Begyndelsen af sætter denne
Adskillelse mellem den objective Virkelighed og de subjective Forestillinger om
den. De erkjende vel, at det har været nødvendigt paa Grund af hiin Tids
aandelige Umyndighed eller idetmindste overeensstemmende med dens særegne
aandelige Standpunkt at udbrede saadanne Traditioner; men Deismen driver dem
ofte saavidt, at de hellere ville erklære hine Beretninger for Producter af
List og Bedrageri, end antage Gud for deres Ophav\footnote{Smlgn. Strauß. Leb.
Jesu (\textit{1838}) Einleit. §§ \textit{5}. \textit{6}.}.
% One note, marked 1) --- the plain next number.

Paa en ganske anden Maade bedømmes de i Panlogismen. Den Adskillelse,
Kriticismen urgerer mellem Gudsideen, som den foresvæver det fromme Gemyt, og
Gud selv, strider fuldkomment mod Panlogismens Princip. Indrømmer derfor
Kriticismen, at Traditionerne om Mirakler og Aabenbaringer ere udgaaede fra en
Gudsidee, beviser Panlogismen, at de ere fremkomne ved en i den menneskelige
Bevidsthed immanent Virksomhed af Gud:
% The colon after "Gud" is as printed --- two dots, read at 2x, the upper one
% lightly inked. A new paragraph follows, where the sense wants a full stop.
% Recorded, not corrected. The Fraktur OCR read a semicolon here; rejected.

Men da saaledes disse Gudsaabenbaringer, af hvad Be\-%
% --- p. 132 ---
\setcounter{footnote}{0}
skaffenhed de nu end ere, udgjøre et nødvendigt Stadium i
Aabenbaringshistorien, henhøre de ogsaa til den kosmiske Mediation og ere ikke
at ringeagte. Det er saa langtfra, at de maae betragtes som opdigtede Fabler,
at de meget mere overlevere os Phænomener restaurerede og udsmykkede i den
religieuse Sandheds Interesse. Mythen indeholder altsaa ifølge sit Begreb ikke
alene et Aabenbarings=, men ogsaa et Illusionsmoment\footnote{Wenn man die
Religion im Verhältniß zur Philosophie bestimmt als das Bewußtseyu desselben
% PRINTER'S DEFECT in note 1, transcribed as printed: "Bewußtseyu" for
% "Bewußtseyn" --- a turned n. Read at 2.4x: the final letter has the bowl of a
% u but no breve, where "Bewußt-" two syllables earlier shows a breved u.
% Both OCR witnesses read a u here and read "Bewußtseyn" correctly in p. 129's
% note. Rejected reading: "Bewußtseyn".
absoluten Inhalts, aber nicht in Form des Begriffs, sondern der Vorstellung: so
ist leicht zu sehen, daß nur unter und über dem eigentlichen Standpunkte der
Religion das Mythische fehlen kann, innerhalb der eigentlich religiøsen Sphäre
% ERRATA (Trykfeil, p. 132, Note l. 1): printed "das Mystische", corrected here
% to "das Mythische". The uncorrected reading was verified on the image at
% 1.25x --- the page sets M-y-(long s)-t-i-(long s)-c-h-e, with the st
% ligature and no "th" at all.
% PRINTER'S DEFECT at the same site, transcribed as printed: "religiøsen" for
% "religiösen" --- a Danish ø sort set for the German ö. Read at 2x: the o
% carries a clean diagonal stroke, not the fount's ö, which is set correctly in
% "religiöses" in the p. 121 runover. The i of "reli-" is under-inked to a bare
% dot at the same place.
aber dasselbe wesentlich und nothwendig vorhanden ist. S.~S. Pag. \textit{98}.
§ \textit{14}.}. De, der have overleveret os Mytherne, troede, at
Beretningerne vare objective, \dsi{} at Gud selv udenfor den menneskelige
Bevidsthed havde paataget sig hine endelige Former. Men denne Vildfarelse
aabenbares og fjernes i Tidernes Løb; thi ethvert efterfølgende historisk
Stadium overgaaer stedse det foregaaende, og den guddommelige Idees
Aabenbaring bliver altid mere og mere fuldkommen\footnote{Wie aber Bildung
überhaupt Vermittlung ist: so wird die fortschreitende Bilder der Völker auch
der Vermittlungen immer deutlicher bewußt, welche die Idee zu ihrer
Verwirckligung bedarf; und so erscheint jene Differenz der neuen Bildung und
der alten Religionsurkunden \ldots\ Das Göttliche kann nicht so (theils
überhaupt unmittelbar, theils noch dazu roh) geschehen seyn; oder das so
Geschehene kann nichts Göttliches gewesen seyn --- so wird die Differentz sich
aussprechen. Strauß L.~J. § \textit{1}: Pag. \textit{2}.}. Man indseer nu, at
% Two notes on this page, marked 1) and 2) --- the plain next numbers.
% In note 2, transcribed as printed and read at 1.25x: "die fortschreitende
% Bilder der Völker" (Strauß has "Bildung"); "Verwirckligung" (for
% "Verwirklichung"); "Differenz" at its first occurrence but "Differentz" at
% its second, four lines later --- the two spellings are visibly different in
% the printing and are both kept. The mark after "§ 1" is a colon, read at 1.5x.
de Ting, hvorpaa man hidtil troede, vare falske, og dette bliver en Kilde til
Sorg og Bedrøvelse; men man seer tillige, naar Erkjendelsen har naaet sin fulde
Udvikling, at det er den samme Gudsvirksomhed, der før var Kilden til

% --- p. 133 ---
Mytherne, som nu atter har bevirket Erkjendelsen af disses Usandhed. Det er da
saa langtfra, at Gud forlader os, naar vi have opdaget Vildfarelsen, at han
meget mere nu træder os nærmere, og dette er Kilden til Trøst og Glæde. Thi
ogsaa heri afviger Panlogismen fra Kriticismen, at hiin har en altid fortsat
Gudsaabenbaring, medens denne, efter at have frakjendt de hellige Efterretninger
deres Gyldighed, tillige har tabt de sidste Spor af en virkelig guddommelig
Nærvæværelse. --- Gud existerer nemlig ikke udenfor Menneskeslægten: Guds Aand
% PRINTER'S DEFECT, p. 133: "Nærvæværelse" for "Nærværelse". Read at 0.46x: the
% line ends "guddommelig Nærvæ=" and the next line begins "værelse.", so the
% syllable "væ" is set twice. Transcribed as printed. Both OCR witnesses show
% the same doubling. p. 133 carries NO footnote --- its foot is blank, so
% p. 132's note does not run over onto this page.
og Menneskenes Aand bevare stedse deres Identitet, Gud træder ud i Menneskene og
Menneskene vende tilbage til Gud.

Hvad der saaledes gjælder om de hellige Phænomener i Almindelighed, kan
naturligviis ogsaa anvendes paa den absolute Religions fortrinligste Phænomen.

Ved den rette Forstaaelse af den mythiske Christus, \dsi{} af Christus, saaledes
som han er skildret i den hellige Skrift, komme følgende trende Momenter i
Betragtning.

Først maa man betragte Phænomenet, saaledes som det kunde opfattes gjennem
Hverdagserfaringen, --- tænke sig Christus, saaledes som han vilde have viist
sig for vor Tids oplyste Mennesker, dersom han var fremstaaet iblandt os. Fra
denne Side betragtet maa han antages at have været et Menneske, der var sig
Religionens Idee fuldkomment klart bevidst; han vilde ellers ikke have bevirket
en saa stor Forandring i den religieuse Sphære. --- Men at han opfattede
Religionens absolute Idee, medfører ingenlunde, at han var en evig Personlighed,
udrustet med Magt til at gjøre Mirakler; ja en saadan Antagelse strider endog
imod Religionsideen selv. Thi forsaavidt Christus var et Menneske, kunde han
hverken, som enkelt Individ, rumme Guds hele Fylde i sig, eller aabenbare denne
for sine Samtidige. Skriftens Fremstilling kan
% --- p. 134 ---
\setcounter{footnote}{0}
saaledes ikke give Hverdagserfaringen Mere, end visse Fragmenter og Grundrids af
det sande Messiasbillede\footnote{Smlgn. Strauß: L.~J. Schlußabhandlung.}. ---
% p. 134, note: printed mark is the plain 1). The final mark after
% "Schlußabhandlung" is read at 2x as a baseline point with a short stroke above
% it; set as a full stop. The alternative rejected is a colon --- the colon in
% "Strauß:" earlier in the same line has two clearly ROUND dots, and this mark's
% upper element is a horizontal stroke, i.e. an ink defect over a point.
% "L. J." = Leben Jesu; the Fraktur I/J glyph is set as J throughout this book.
Han fremtraadte i Tidens Fylde, da Tidsalderen allerede havde et rigt Apparat af
Forbilleder og billedlige Betegnelser; da nu hiint religieuse Princip, der stod
klart for Christi Bevidsthed, kom til og ligesom gav disse Elementer et nyt Liv,
og da det særegne Phænomen i Ideens Tjeneste blev Dødens Bytte, havde den deraf
undfangede religieuse Idee saa stor en Magt i deres Bevidstheder, der havde
været den store Mesters Disciple, at de paany restaurerede hans Liv og hans
Personlighed og saaledes overleverede de Troende Ideens Fylde anskueliggjort i
et Phænomen, der, som saadant, kunde blive Gjenstand for Tilbedelse. Vi maae
derfor, uagtet hiin Christus, Apostlene have beskrevet os, er sammensat af saa
modsatte Momenter, indrømme, at vor Tids Mennesker ikke vilde have kunnet
begribe den religieuse Idee paa indre Maade, hvis den ikke først paa denne
udvortes Maade var bleven anskueliggjort for os. Man maa altsaa ikke ringeagte
hiin Christus, thi vel døde han som enkelt Menneske, som enestaaende Phænomen
(κατὰ σάρκα), men han opstod forherliget som Symbol paa Menneskeslægten
% Greek on this page read at 0.46x. The fount's cursive kappa is normalised to
% plain κ, per the standing convention; accents and breathings are as printed.
(κατὰ πνεῦμα), og vi skylde saaledes ham, at vi have lært Sandhedens absolute
Begreb at kjende. --- Enhver, der nøies med Panlogismens Gud, maa altsaa ogsaa
bifalde den mythiske Behandlingsmaade af den hellige Historie. Uagtet der
saaledes ikke bliver nogen Realitet tilbage uden med en aldeles momentan
Tilværelse, og uagtet det subjective og phænomenale Liv stedse opsluges af det
objective og væsentlige: bliver dog Gud i Ideernes Rige den samme, vindicerer
dog Menneskeslægten uagtet de enkelte Menneskers Dødelighed sit Evighedspræg,
tager dog den sande
% --- p. 135 ---
Messias Sæde paa den panlogistiske Lyksaligheds Bjerg og lærer sit Folk:

”Salige ere de, som troe paa Himmelens og Jordens speculative Immanents, thi
% QUOTATION MARKS, p. 135: the raised ” (U+201D) at BOTH ends, as everywhere in
% the body of this book. Verified at 0.46x on the opening mark before "Salige"
% and the closing mark after "Rige!".
Himmeriges Rige er deres; salige ere de, som rense sig fra alle Fordomme om en
over Verden værende Gud, thi i Ideens rene Speil skulle de see Gud; salige ere
de, der forsmaae egen personlig Udødelighed, thi deres Løn skal være stor i
Erindringens Rige!”

\parag{21}{Panlogismens metaphysiske Prøvelse.}
% § 21 head: the PRINTED head reads "Panlogismens metaphysiske Prøvelse." ---
% metaphysiske, without the h after "met". The Indhold (PDF 8) gives
% "methaphysiske". BOTH READINGS ARE RECORDED; the printed head is what is set.
% This is the fifth head/Indhold disagreement in the book (cf. §§ 7, 10, 13, 15).
% Read at 0.46x on the argument line, PDF 144.

Vil Nogen indvende, at ene og alene Forsømmelse af Skriftstudiet kan have
fremkaldt denne Udvikling af den christelige Religion, svare vi, at det meget
mere er skarpsindige, lærde, i Skriften velbevandrede Mænd, der anbefale samme.
% "anbefale samme" and "forsømt nogen": both OCR witnesses put a point between
% the words. At 0.46x each mark sits high, at mid x-height, not on the baseline
% --- scan speckle, not punctuation. No point is set.
De have ikke forsømt nogen af Kritikernes indledende Undersøgelser; udrustede
med det hele Apparat af hellig Philologie skjelne de med den største Fiinhed
imellem alle Skriftens ζεύγματα og øvrige Figurer; ja de kjende endog alle
Forsøgene paa at forsone Christi Genealogier til Punkt og Prikke, --- men Alt
dette omvender dem ikke. Der finder nemlig følgende gjensidige Forhold Sted
mellem Livet og Philosophien: Livet sætter Momenterne som positive, adskilte,
men bevarer dem tillige alle paa eengang, lader dem alle bestaae ved Siden af
hverandre, forbinder de enkelte Momenter med hinanden uden at ophæve deres
Enkelthed; Tænkningen derimod negerer og restaurerer Livets Phænomener, begriber
gjennem en Anskuelse, anskuer gjennem en Begriben, og urgerer som en Følge deraf
snart Elementernes Adskillelse, snart deres Identitet. Og heri er netop hele den
kritiske Deismes og den speculative Pantheismes Retmæssighed og Uretmæssighed
begrundet.
% --- p. 136 ---
Hos Orientalerne, hvor den menneskelige Sjels sløve Øie ikke kan opdage nogen
anden Modsætning, end Modsætningen mellem Phænomenernes Mangfoldighed og
Væsenets abstracte Eenhed, opsluges alt Endeligt af det Uendelige; \emph{der}
% EMPHASIS, p. 136: "der" is letterspaced --- the only letterspaced word in
% pp. 133--144. Read at 0.46x: "d e r" with clear inter-letter space, against
% the ordinary "derfor" two words later. The ABBYY layer independently renders
% it "d e r". spacing.py did NOT flag it (a false negative on a short word).
foregaaer derfor ogsaa Forsoningen paa den meest umiddelbare Maade. ---
Bevidstheden, der tænker, er den sandselige Natur underlagt; det naturlige
Moment er saaledes overveiende i Bevidstheden, det er derfor ved et vist
umiddelbart Instinct, at denne fatter sin Tilværelses Grund. Den kan da kun hæve
sig over Naturen ved frivilligt at give Naturen, hvad den tilhører; den vender
derfor igjennem Askesen tilbage til sit oprindelige Intet (Væsenet i dets rene
Abstraction). Bevidstheden om Naturen er altsaa her indskrænket ved
Bevidsthedens Natur. --- Med det christelige Liv forholder det sig ganske
anderledes; det tillader ikke nogen Nedsynken i Naturdrømme; først gjennem
mangfoldig Position og Negation af de givne Kjendsgjerninger, først efterat
Aanden havde overvundet de mangfoldige Besværligheder, der stode den i Veien,
først da kunde Panlogismen fremtræde som det endelige Resultat.

Vi have alt forhen (smlgn §~\textit{7}) bemærket, at den Restauration, hvortil
% "smlgn" is printed WITHOUT a point here, and the § without one either, against
% the "smlgn.\ §.~\textit{10}" of p. 18. Read at 0.46x. Set as printed.
den dialektiske Stimulus fører Bevidstheden, er forskjellig efter de
Phænomeners forskjellige Beskaffenhed, der negeres. Dog er Restaurationen ikke
altid en saadan Gjendindsætten af de negerede Phænomener, at disse komme til at
indtage den samme Værdighed, de før havde; thi den dialektiske Restauration er
tillige en Prøvelse. Vi maae derfor undersøge, om Panlogismen vil være istand
til at hindre de Momenter af den christelige Religions ældgamle Væsen, som den
har bundet med den dialektiske Negations Lænker og nedstødt i de subjective
Forestillingers Skyggeverden, fra paany at træde frem, --- eller om det er
Sandhedens evige Ideer, paa hvis Ruiner den speculative Philosophie har opbygget
sit
% --- p. 137 ---
\setcounter{footnote}{0}
Pallads, og som derfor atter efter længere eller kortere Tids Forløb ville
sprænge deres Fængsel og hævne den dem tilføiede Overlast. Frygter ikke, ville
de sige, vi ere gode Christne! Den Gud vi prædike, er det ikke den absolute
Aand, der negerer Naturen i al Evighed, netop for at forherlige den? Er ikke den
Gud, der igjennem den metaphysiske Trilogie uafbrudt træder ud af sig selv og
vender tilbage til sig selv, den Treenige? Her kan der ikke være nogen Tale om
en Vanhelligelse af Aabenbaringen; vi vanhellige ikke Gud, som om han nogensinde
af Misundelse tillukkede Himlen, nei den kjærlige Gud gjør bestandigt
Menneskene deelagtige i sit Liv! Hvad Andet betyder vel den personlige Eenhed,
end netop denne inderlige Forening? Ja Christus er Religionen selv; han har
været en historisk Person, i hvis Bevidsthed Identiteten af begge Momenter lod
sig tilsyne; men hvor han nu befinder sig, hvad vedkommer det os? Det er
Menneskeslægten, der er hiin θεάνθρωπος, denne er reelt tilstede i ethvert
% θεάνθρωπος: the fount's script theta is normalised to plain θ (a hard build
% requirement --- U+03D1 is fatal under textalpha). Read at 0.46x.
Menneske, --- kun at den kommer til en Aabenbaring! Vi sige da gjerne det
forsvindende Individ, Christus, vort Farvel, for at komme i Besiddelse af den
evige, nærværende og saaledes i Sandhed historiske Christus\footnote{\ldots{}
% p. 137, note: mark 1). The note opens with an ellipsis printed as four spaced
% points; set as \ldots. Every word of it is FRAKTUR --- the German AND the
% German book-title "der historische Christus und die Philosophie" --- so
% nothing here takes \textit{} except the antiqua figures 64, 159 and 768.
% Read off the image at 0.92x in two x-halves; "S. 64." (not 614) and "Diess"
% (long-s + s, distinct from the ß of "daß" in the same note) both confirmed.
durch die Belebung der Idee der Menschheit in sich, namentlich nach dem Momente,
daß die Negation der Natürlichkeit und Sinnlichkeit, welches selbst Negation des
Geistes ist, also die Negation der Negation, der einzige Weg zum wahren
geistigen Leben für den Menschen sei, wird auch der Einzelne des
gottmenschlichen Lebens der Gattung theilhaftig. (Med disse Ord gjendriver
Strauß den Indvending, som Schaller i sit Skrift: der historische Christus und
die Philosophie. S.~\textit{64}. fremsætter imod ham). Diess allein ist der
absolute Inhalt der Christologie: daß dieser an die Person und Geschichte eines
Einzelnen geknüpft erscheint, gehört nur zur geschichtlichen Form derselben.
Strauß, Schlußabhandlung L.~J. §~\textit{159} Pag.~\textit{768}.}. Thi hvad er
vel den sande
% --- p. 138 ---
Historie Andet end Virkelighedens evige Nærværelse? Hvorfor skulde man da
attraae en subjectiv Udødelighed? Hvo som søger sit Liv, han skal miste det! En
saadan Egoisme fandtes ikke hos den Christus, hvis Priis Du selv saa ofte har
paa Dine Læber. Han opoffrede sin subjective Personlighed for at opstaae i
Andres Bevidstheder og fare op til en taknemlig Ihukommelses Himmel. --- I
Christus altsaa skulle vi leve, døe og opstaae!

Vi ville nu noget nøiere prøve hiin Immanents af Idealitet og Realitet, hvorpaa
hele denne Religionsanskuelses Fortolkningsmethode er bygget. En abstract
Idealitet udfordrer, at alle de Momenter, der ved Negationens Kraft gjensidigt
udelukke hinanden, tillige maae have denne Negation saa gjennemreflecteret inden
for deres egne Grændser, at Position og Negation overalt falde sammen. Men den
uendelige Bestemmelses og Begrændsnings Sandhed er afhængig af Negationens
Virkelighed. Nu indeholder den abstracte Idealitet kun en abstract Identitet med
sig selv; sin sande Modsætning holder den udenfor sig. Derfor savne dens
Momenter indenfor dens Grændser overalt en sand Modsætning; vel synes det Ideale
at begrændses af det Ideale, Negationen af Negationen, men denne Begrændsning
finder ikke virkeligt Sted. Idet den abstracte Idealitet saaledes har sit sande
Negative udenfor sig, hvilket strider imod Idealitetens Natur, forsvinder den.
Kun den ydre Negations, den empiriske Virkeligheds, Rige bliver tilbage,
opbygget paa Idealitetens Ruiner. --- Det er her den første Grundsætning, at de
positive Ting udelukke hverandre; denne gjensidige Udelukkelse tænkes vel først
som en saadan, at hver enkelt Ting hviler uforstyrret indenfor sine retmæssige
Grændser; men da Tingens Grændse selv er en Ting, og Tingen saaledes bliver sin
egen Modsætning: bliver der en uophørlig Modsætning uden Hvile, gjennemtrænger
en uop\-%
% --- p. 139 ---
hørlig Drift til Alteration og Tilintetgjørelse Tingens Marv, tilintetgjør den
abstracte Realitets Rige sig selv. Men i det Realiteten ved denne Destrueren
gjennemfører Negationen af sig selv, vender Idealitetsforholdet atter tilbage.
Da altsaa hverken den rene Idealitet eller den abstracte Realitet ere Noget i
sig, følger naturligviis heraf, at denne rastløse Oscillation mellem Idealitet
og Realitet ikke blot er Begyndelsen til, men ogsaa Principet for vor Tilværelse
heri Livet. Begyndelse og Princip falde sammen. Da der nu Intet kan fremgaae i
Slutningen, som ikke alt har været i Begyndelsen, er det ikke at undre over, at
Fluctuationen gaaer i det Uendelige. Vel søger Idealiteten altid at finde sin
Tilværelse i Realitetens Rige; men Empiriens enkelte Momenter ere forgjængelige;
derfor ere alle dens Forsøg forgjeves; --- var Idealiteten sand og fuldkommen i
sig, vilde den ogsaa have sin Bestaaen i sig selv; den vilde da ogsaa kunne
sammenfatte Realitetens adspittede Momenter og aabenbare sig selv i den
% PRINTER'S DEFECT, p. 139: "adspittede" for "adsplittede" --- the l is not
% there. Read at 0.46x, and both OCR witnesses agree. The same word IS set
% correctly as "adsplittede" twice within three pages (p. 141 "adsplittede
% Fragmenter", p. 142 "adsplittede Liv"), which is what makes this a slip and
% not an orthographic variant. Transcribed as printed.
empiriske Sphære; men nu vedbliver Realiteten at være forkrænkelig formedelst
Idealitetens Afmagt, Idealiteten at være afmægtig formedelst Realitetens
Forkrænkelighed.

Det er denne Selvopløsning, hvoraf hiin θεάνθρωπος, som skal indtage Christi
Plads, lider. I den Gud, der boer i Idealitetens Sphære, seer Menneskeslægten
sin Fader, men denne Gud er ikke et i sig befæstet, urokkeligt Subject, ikke
Jehova, der i sin urokkelige Fastholden af sig selv har alle sine Modsætninger
uden for sig, men det er Faderen, der søger sin Tilværelse deri evigt at avle
Sønnen \dsi{} Menneskeslægten. Men hvilken Tilværelse har nu denne Guds Søn? Han
er ikke den rene Idealitet --- thi denne tilhører alene Faderen ---, heller ikke
den abstracte Realitet, der udfolder sig gjennem de enkelte Ting og de enkelte
Mennesker; --- han er selve Overgangen fra Faderens Idealitet til Verdens
Realitet,
% --- p. 140 ---
Midleren (\textit{mediator}) mellem Gud og Verden. Vi maae da nøiere prøve denne
% "mediator" is Latin set in ANTIQUA against the Fraktur --- verified at 0.26x
% on the whole-page sheet, where the roman letterforms are unmistakable.
Mediation.

Man maa skjelne mellem Menneskeslægtens almindelige Bevidsthed og de enkelte
Menneskers enkelte Bevidstheder. Christus er denne almindelige Bevidsthed ---
ikke hiint evige individuelle Jeg, der eengang aabenbarede sig paa Jorden og nu
regjerer i Himmelen (da maatte man jo atter tye til den i Bethlehem fødte, paa
Golgatha korsfæstede Jesus, hvem den speculative Immanents alt længe har
antiqveret!) ---, men den almindelige Bevidsthed, der bereder sig Tilværelse
igjennem de enkelte Bevidstheder, saa at den af samme Grund, fordi den nemlig
ikke vil være Noget i og for sig, men kun have Alt igjennem de Enkelte, fordrer,
at de Enkelte ikke maae ville være Noget i sig selv, men at de skulde opgive
deres Subjectiviteter og vende tilbage i den. Denne almindelige Bevidstheds
% spacing.py flagged "Bevidstheds" here at 2.00; at 0.46x it is ordinary
% setting in a loosely justified line --- NOT letterspaced. No \emph.
Udvikling igjennem de enkelte Bevidstheder, --- den historiske Operation, ved
hvilken Spliden mellem Menneskeslægten og de enkelte Mennesker mere og mere
ophæves, er Aanden, --- dog ikke den Aand, der som det tredie Jeg forener
Faderens og Sønnens Liv, fører Guds Negativitet tilbage til at være en immanent
Negativitet indenfor Gudslivets Grændser og derved tillige kan tillade de
enkelte Bevidstheder at have deres eget Negativum, den sande Idealitet, det
evige Liv, i sig, --- men den Aand, der meddeler Øieblikkets Virkelighed alle de
forbigangne Tiders ideale Skatte og oprigtigt erklærer ikkun i denne Immanents
af Idealitet og Realitet at have sin Tilværelse. Men den øieblikkelige
Virkelighed, der har udfoldet sig igjennem de enkelte Bevidstheder, viser sig
allerede at være bleven en anden (altereret); den objective Aand begrændses
altsaa ved Summen af de endelige Bevidstheder. Thi betragtet efter sin reelle
Existents har den menneskelige Bevidsthed andre Bevidst\-%
% --- p. 141 ---
heder satte udenfor sig, saa at der ikke er nogen Idealitet tilstede i deres
Sum; ideelt betragtet optager derimod den ene Bevidsthed den andens Liv i sig.
Den enkelte reelle Selvbevidsthed optager altsaa i sig de øvrige Bevidstheders
aandelige Fylde. Det er derfor i de Enkelte, man maa søge Immanentsen af
Idealitet og Realitet. Men denne Immanents opløser sig stedse i adsplittede
Momenter; forsaavidt nemlig den reelt existerende Bevidsthed tænker, har den
ogsaa Immanentsen i sig, men forsaavidt den reelt existerer, indtræder
Adsplittelsen. Thi den Idealitet, der er tilstede i den reelle Bevidsthed, er
kun betroet den for en Tid; vel erkjender den tænkende Bevidsthed, at denne
Idealitets Fylde er uendelig, men den vil aldrig fuldkomment kunne udtømme den.
Da nemlig den reelle Selvbevidsthed, der, som det hedder, har sin Negation
udenfor sig, men ikke i sig, skal døe, opsøger Idealiteten sig stedse en anden
reel Selvbevidsthed for igjennem denne at bevare Identiteten med sig selv. Men
da enhver reel Selvbevidsthed er samme Skjebne underkastet, trænger Modsigelsen
og Modsætningen igjennem lige til denne Aands inderste Kjerne. Thi den reelle
Selvbevidsthed vedbliver at være det nødvendige Organ, uden hvilket Aanden ikke
kan existere, men udsat, som den er, for uendelig Alteration vil den aldrig
kunne optage Aanden i sig; dette incommensurable Forhold vil da vedvare i al
Evighed.

Da Sagen forholder sig saaledes, kan man ikke forvexle den panlogistiske
% DOUBTFUL GLYPH, p. 141: "forvexle". At 12x the sixth letter is barely
% separable from the Fraktur r, and BOTH OCR witnesses read one ("forverle",
% "forverte"), which is the documented x->r failure of this scan. Compared
% directly, at the same magnification, with the x of "Afvexlinger" on p. 143;
% the forms agree as far as the 1-bit bitmap allows. Rejected: "forverle",
% which is not a word. "forvexl-" already occurs three times in pp. 1--108.
Treenighed med den kirkelige. Panlogismens Gud er ikke vor Fader,
ὁ μόνος ἔχων ἀθανασίαν, men en beruset Gud, der ikke kan samle sine adspredte
Tanker og drukken raver om mellem Tilværelsens adsplittede Fragmenter. Dens
Gudsøn er ikke den evige Søn, der erklærer om sig selv, at han havde Herlighed
hos Faderen, før Verden blev til, men den Søn, hvem Faderen avlede, ansporet af
Idealitetens
% --- p. 142 ---
og Realitetens Modsigelse, --- der gjør den efter Tilværelsen gribende Faders
Haab til Intet, der spotter hans abstracte Idealitets Nøgenhed og derfor følges
af Faderens Forbandelse. Vi kjende ogsaa denne Aand, der udsendes af en saadan
Fader og en saadan Søn; det er denne Verdens Aand, der overalt forraader
Faderens og Sønnens σκάνδαλον. Dens Immanents er en Modsigelse, dens Forsoning
en Rædsel. Jeg, siger den, aftørrer alle Taarer og stiller alle Smerter. Thi
for at kunne gribes af Smerte, maa man have Livets positive Fylde i sig, en
Livsfylde, der lider under en til det Inderste gjennemtrængende Modsigelse; men
var Livet fuldkomment negeret og havde det ikkun sit Negativum udenfor sig,
vilde der indtræde en pludselig Død, og der vilde da ingen Smerte være; det
smertefulde, i sit Inderste opløste og adsplittede Liv forbliver altsaa sig selv
ligt, og saaledes bevares Identiteten. Smerten er den menneskelige Tilværelses
inderste Rod; thi saasnart Mennesket samler alle sine Kræfter for at anskue sin
egen Natur i fuldstændig Klarhed, falde Idealititeten og Realiteten sammen: der
% PRINTER'S DEFECT, p. 142: "Idealititeten" for "Idealiteten" --- an extra "it".
% Read at 0.46x; both OCR witnesses agree. Transcribed as printed.
er en sand Identitet af dem, thi netop denne udgjør Livet selv, men der er
tillige, som vi have seet, en Grundmodsigelse imellem dem, og hermed er Livets
Opløsning og Smerten givet. Deri bestaaer altsaa Synden, at Mennesket hendrages
af Kjærlighed til Livet; men Kjærlighed til Livet avler Frygt for Døden, og
denne er netop Straffen for hiin; gribes vi da af begge, kunne vi ikke nære
Tvivl om, at Synd og Straf gjennemtrænge hinanden, og at Forsoningen netop
bestaaer i denne Smerte. Smerten er den fuldkomneste Immanents af Modsigelse og
Identitet, den er netop den evige Spore til Stræben efter det evige Liv, den er
Guds helligste Mysterium; Smerten er Salighed; --- men det er en Salighed, der
minder om Djævelens Skoggerlatter! Den menneskelige Selvbevidst\-%
% --- p. 143 ---
\setcounter{footnote}{0}
hed hører det og gyser; thi jo opmærksommere den betragter dette Riges
Elementer, desto skarpere og klarere seer den Smertens Modsigelse under tusinde
forskjellige Skikkelser gjennemstrømme alle dets enkelte Dele. Men naar den i
dette Mørkets Helvede hører Kirken lovsynge sin Christus, der blev korsfæstet,
men opstod af Graven og sidder ved Faderens høire Haand, see! da reiser han sig
atter, som det Menneske, der vaagner af en rædselsfuld Drøm og glad hilser
Dagens Lys; hvad der drager ham er ikke en beregnende Egoisme, der kun veed at
søge sit Eget, nei det er den objective Magt, til hvilken han hengiver sig af
sit ganske Hjerte, den sande Længsel efter den sande Gud, og han henrives ikke
mere af en blind Lidenskab, naar han saaledes har erfaret den dialektiske
Modsigelses bittre Smerte. Viig bort, Satan, du som fornægter det Hellige! Til
os komme Negationens Negation!

\parag{22}{Tilbagevenden til Kirken. Den speculative Theologie.}
% § 22 head agrees verbatim with the Indhold. Unlike § 21's, this argument is
% set in a LARGER Fraktur and is not letterspaced; \parag renders it the
% standard way, which is what every other head in this edition gets.

\begin{verse}
% VERSE, p. 143: eight lines, indented as a block and set about three quarters
% the body's size (line pitch measured at 116 px against the body's 157 px on
% the same page). The plain `verse` environment is used, following
% texts/hoeffding/danske-filosoffer; no size change is applied, and the smaller
% fount is recorded here instead. Quotation marks are the body's raised ” at
% both ends. "Gylden=Sol" carries the Fraktur double-hyphen, kept as =.
”Lad hvisle kun i Ormegaard,\\
At nu jeg er lagt øde;\\
Jeg kroner ligefuldt mit Aar\\
Med Frugtbarhed og Grøde;\\
Ja! op jeg staaer som Ax i Vang,\\
Som Mai i Bøgeskove,\\
Og prægtig under Fuglesang,\\
Som Gylden=Sol af Vove!”\footnote{Grundtvig: Den christelige Tro.}
\end{verse}

Disse Ord, hvormed Digteren besynger den christelige Troes Seier under Tidernes
forskjellige Afvexlinger, gjenlyde ikke blot i det kirkelige Liv, men ogsaa i
Videnskaben. Den menneskelige Fornuft bevæger sig i to modsatte Retninger, fra
% No paragraph break at "Den menneskelige Fornuft": the preceding line runs the
% full measure (2852 px, the page's own right edge), measured row by row.
Gud til Verden og sig selv, fra sig selv og Verden
% --- p. 144 ---
\setcounter{footnote}{0}
til Gud. Vel var det ikke at vente, at denne Overgang i vor syndige og endelige
Verden skulde foregaae roligt og uden Afbrydelse; men vi see dog stedse, at
Fornuften, uagtet den ofte synes aldeles at vende sig bort fra Christi Sandhed,
atter fortryder dette Frafald, naar den tilfulde har smagt de sørgelige Følger
deraf. --- Man maa unægteligt erkjende Panlogismen for det høieste
Selverkjendelsestrin, hvortil den abstracte Fornuft kan komme, og man kan ikke
med Rette bebreide den Uvillie til at vende tilbage til det positive Liv; men er
det ogsaa den rette Vei, hvorpaa den vender tilbage til dette? Hvis der ikke var
en Verdensorden, som kunde negeres, vilde det menneskelige Jeg heller aldrig
kunne komme til en Udvikling af de logiske Begreber. Men heraf følger, at det
Begrebssystem, Panlogismen opstiller, kun er et formelt, at det overalt fattes
den positive Virkelighed og paa ethvert Punkt er opfyldt af en Trang, en Længsel
efter denne. Og dette erkjender Panlogismen selv tilfulde. Men det er netop dens
πρῶτον ψεῦδος, at den tillægger denne Idealitet guddommelig Auctoritet og
% πρῶτον ψεῦδος: the fount sets the circumflex of the -ευ- diphthong over the
% epsilon; the placement is normalised to the standard one, as the header
% prescribes for -ου-. Read at 0.46x.
antager dens Afmagt for at være Almagt. Denne Idealitet er nemlig intet Andet,
end den menneskelige Fornufts egen Skikkelse. Dog ligger Feilen ikke saa meget
deri, at den antager en Identitet af den guddommelige og den menneskelige
Fornuft, som deri, at den erklærer Immanentsen af begge for at være umiddelbart
given i denne. Thi en Identitet, der ikke støtter sig paa en Modsætning, er
efter den speculative Logiks Vidnesbyrd ligesaavist holdningsløs og slet, som en
Sammenblanding af den guddommelige og menneskelige Fornuft efter Kirkens
Vidnesbyrd\footnote{Hvor kraftigt og skarpsindigt J.~H. Fichte har bekjæmpet
% p. 144, note: mark 1). The initials are the Fraktur I/J glyph; set as J,
% following "J.~G. Walch" at p. 71 and the man's own usage (Immanuel Hermann
% Fichte wrote himself J. H. Fichte). The ABBYY layer reads I. H.
denne Rationalisme, er bekjendt nok.} er det sande Kjendemærke paa Rationalisme.

% [text to be added: pp. 145--156]

% --- p. 157 ---
% JOINT p. 156/157 verified on the image: p. 156 ends "Den sig fra Evighed" with
% NO word broken across the leaf, and p. 156 carries no note, so nothing runs
% over onto this page. p. 157's first line is NOT indented (x464 against the
% page's own 456 left edge) --- it continues p. 156's paragraph.
% NO FOOTNOTE on this page: 32 body lines run to the foot, no rule, no small
% type. Verified by a row-by-row ink scan of the last 1100 rows, which finds
% nothing below the last text line but scan-edge noise at x < 260.
af sin absolute Frihed bevidste Gud fattede den Raadslutning at udvikle en
menneskelig Frihed og Personlighed; derfor maae ogsaa alle saadanne Phænomener,
som Menneskene ikke kunne Andet end antage for bestemte Yttringer af den
guddommelige Bistand, sigtende til Virkeliggjørelsen af hiin Raadslutning,
nødvendigt antages for i Virkeligheden at have hørt med til de guddommelige
Raadslutninger. Betragtes derfor de hellige Phænomener i en saadan
Accommodations Lys, maa man ogsaa tillægge dem Virkelighed i dette Ords dybeste
Forstand. Betragtes nemlig saadanne Phænomener i og for sig, ere de ifølge deres
Natur hverken tilsvarende Udtryk for de Ting, der høre til det timelige, eller
for dem, der høre til det evige Liv. Thi forsaavidt de overnaturlige Ting
aabenbares, forsaavidt maae de ogsaa give sig Udtryk i Phænomener, der høre til
vor Empirie, altsaa iføre sig Former, der ere uovereensstemmende med deres egen
Natur; og uagtet de nedstige i vor Sphære, maae dog de Phænomener, hvori de
aabenbares, paa særegen Maade forherliges for ikke at forvexles med andre,
% "forvexles", NOT the OCR's "forverles" --- settled by direct bitmap comparison,
% not by eye. The word was cut out and set beside "forvexles" in the BODY of
% p. 36 (PDF 45, line 9, x1392--1771), which the pp. 25--36 batch already
% transcribed with an x and which the Fraktur model itself reads with an x: the
% two words are glyph for glyph identical, and the suspect letter is 41 px wide
% against 42 px for the confirmed x and 30--34 px for the r of "for" in the same
% word. Both OCR witnesses read r here (ABBYY's known Fraktur x -> r failure).
% Same letter, same call, at "Omvexlinger" on p. 161 below.
almindelige Begivenheder, eller --- Væsenets Nedstigen i den lavere Sphære maa
tillige være Phænomenets Opløftelse i den høiere Sphære. Derfor er ogsaa
% "Opløftelse" (ft), not the Fraktur OCR's "Opløstelse" --- read at 8x; the
% sense (Nedstigen / Opløftelse) requires it and the glyph is an f, not a long s.
Illusionen stedse et nødvendigt Moment i ethvert timeligt
Aabenbaringsphænomen. --- De hellige Phænomeners Virkelighed forudsætter derfor
altid en særegen Styrelse af Gud. Naar Menneskene see et særegent og
forunderligt Phænomen (τὸ τέρας), slutte de deraf, at Gud vil tilkjendegive dem
% GREEK, read at 5--6x. τὸ carries a grave and τέρας an acute, both as printed.
% Fount variant normalised per the standing convention: the tailed rho (U+03F1)
% of τέρας is typed as plain ρ. No \textit{} --- Greek in the body of this book is
% set raw (cf. πρῶτον ψεῦδος, p. 144).
sin umiddelbare Nærværelse (τὸ σημεῖον); er nu denne Slutning falsk, bliver
% σημεῖον: the perispomeni is over the IOTA here, i.e. in its standard place ---
% no normalisation needed, unlike the -ου- diphthongs the header describes.
selv det meest ophøiede Phænomen uden Virkning; bestaaer den derimod som sand,
lægger Phænomenets svævende, doppelte Natur ingen Hindring iveien for
% "doppelte" is AS PRINTED (pp, German fashion), not "dobbelte" --- read at 8x
% and confirmed by both OCR witnesses. Not in the Rettelser; not corrected.
Aabenbaringens Virksomhed. Det er nemlig Guds overnaturlige Kraft, der giver
Phænomenet dets særegne Præg; og at det ikke tilfulde udtrykker Væse\-%
% --- p. 158 ---
% NO FOOTNOTE on this page. spacing.py reports "a footnote block appears to
% start at y=1103" and flags "til [note]" --- both are false. The page is 33
% solid body lines with no rule and no small type; the flag comes from the heavy
% scan noise down both margins of PDF 167, which also wrecks its Fraktur OCR.
net selv, hidrører derfra, at han accommoderer sig efter Menneskets Fatteevne.
Det er saaledes uden Betydning for Miraklets Sandhed, om Engelen Gabriel ligesom
med eet Slag aabenbarede sig for Maria, eller om hun ved Guds Aands Medvirkning
først senere saae Ahnelser, der fra Begyndelsen af ikkun havde rørt sig i den
jomfruelige Sjels Inderste, anskueliggjorte i fuldstændig Klarhed. Den efter sin
Natur usynlige Gabriel bliver nemlig paa Guds Befaling synlig, for at give os en
Forestilling om Christi overnaturlige Undfangelse. Gabriel er kun et tjenende
Redskab i Guds Haand. --- Det Samme gjælder ogsaa om den Sky, der efter
Skriftens Beretning, omgav Christus, da han opfoer til Himmelen. Skyen dannede
% "efter Skriftens Beretning, omgav" --- the comma after "Beretning" with none
% before "efter" is AS PRINTED, read at 4x. Both OCR witnesses give it.
nemlig en synlig Overgang fra det jordiske til det himmelske Rige. Man kan da
med ligesaamegen Ret sige, at den var til, som at den ikke var til; det er
derfor ligegyldigt med Hensyn til Miraklets Sandhed, om det bestod deri, at
Forvandlingen foregik synligt i den udvortes Verden, eller om Guds Aand ikkun
for Apostlenes Tanker fremstillede Phænomenet i denne Form, naar kun det staaer
fast, at Gud har været en sanddru og ikke en vildledende Aand. Man kan saaledes
fuldtvel have forskjellige Meninger om Phænomenernes Oprindelse og Uddannelse og
desuagtet have den urokkeligste Tro paa deres fulde Virkelighed. Men denne
Meningsforskjellighed er derfor ikke afhængig af vort subjective Skjøn; thi vi
bør ligesaavel tage Hensyn til den vulgaire Empirie \textit{ad extra}, som til
% "ad extra" and "ad intra" are set in ANTIQUA against the Fraktur --- verified
% on the image at 2x, the only two antiqua words in this batch.
Guds virkelige og sande Accommodation \textit{ad intra}.

Den ældre Kirke urgerede eensidigt Guds objective Virksomhed, men gav saaledes
ikke den menneskelige Frihed dens fulde Ret. Inspirationsmiraklet udvikledes
nemlig paa sensualistisk Maade, om vi saa maae sige, saa at den hellige Aands
Virksomhed ikke saameget sattes i Fortolkningen af de
% --- p. 159 ---
% NO FOOTNOTE on this page.
stedfundne Begivenheder, som i den tro Ihukommelse af dem; som derfor Aanden
havde undertrykket Apostlenes Selvvirksomhed, saaledes søgte nu ogsaa den
kirkelige Tro at bringe Fornuften i Trældom. Kriticismen har lært os, at
Aabenbaringen ikke foregik paa en saa slavisk Maade. Ligesom nemlig det
religieuse Forhold imellem Gud og Menneskene først bliver sandt ved disses
sædelige Kraftanstrængelse, saaledes har ogsaa Inspirationen frigjort Apostlenes
Bevidsthed paa en saadan Maade, at de kunde fælde en friere Dom over de
stedfundne Begivenheder. Og det er saa langt fra, at Aabenbaringens Sandhed
svækkes ved Uovereensstemmelsen imellem den hellige Histories enkelte Momenter,
at den meget mere bestyrkes derved. Thi var Christus ikke fra sin Herlighed hos
Faderen nedsteget i Hverdagserfaringens Rige, det er, i vor Tilværelses
Skrøbelighed, vilde han ikke have været et sandt Menneske (υἱός τοῦ ἀνθρώπου);
% GREEK, read at 4--6x. Fount variants normalised per the standing convention:
% script theta -> θ and tailed rho -> ρ in ἀνθρώπου. The dasia over the iota of
% υἱός and the acutes are as printed; τοῦ carries its perispomeni over the
% UPSILON here, in the standard place.
% *** ONE PRINTED ACCENT DROPPED, and recorded. *** The fount sets a further
% perispomeni over the OMICRON of the final -ου of ἀνθρώπου, so the page reads
% in effect ἀνθρώπ + omicron-with-perispomeni + υ. The grammar gives that
% syllable no accent at all (the word is accented on the ω), and Unicode cannot
% carry a perispomeni on omicron in any case, so the mark is dropped and logged
% here. This is the same habit the header records for οὐ / ποιοῦσι (p. 35),
% ἐμοῦ (p. 36) and τοῦ αἰῶνος τούτου (p. 47) --- the compositor evidently has a
% marked-omicron sort for -ου- and used it for both diphthongs of "τοῦ
% ἀνθρώπου" though only the first needs it. REPORTED to the caller.
han vilde da være bleven staaende i Skyerne, ligesom Menneskets Søn i Propheten
Daniels Vision (\emph{Bar enash}), og vi maatte da kalde hans jordiske Liv en
% *** HEBREW IN THE PRINTING --- SUBSTITUTION, NOT A TRANSCRIPTION. ***
% The page sets the pointed Hebrew phrase itself, read at 6x: bet + dagesh +
% patach, resh --- aleph + hataf-segol, nun + qamats, shin + shin-dot, i.e.
% בַּר אֱנָשׁ, "bar enash", "son of man", the Aramaic of Daniel 7:13 (there
% כְּבַר אֱנָשׁ; the page quotes it without the prefixed kaf). Code points
% U+05D1 U+05BC U+05B7 U+05E8 / U+05D0 U+05B1 U+05E0 U+05B8 U+05E9 U+05C1.
% This edition is built with pdfLaTeX + Libertinus and has NO Hebrew font, so
% the phrase is rendered by its name, exactly as אֱלֹהִים at p. 33 and הִנֵּנִי
% at p. 39 are. If cjhebrew is installed the glyphs should be restored: add
% \usepackage{cjhebrew} to the preamble and set "\cjRL{b.ar 'E na:s}" in place
% of "\emph{Bar enash}" --- that input string is reconstructed from the pointing
% and has NOT been compile-tested. THIS IS THE THIRD HEBREW WORD IN THE BOOK;
% see the OPEN ITEM in RESUME-NOTES.
Skintilværelse (δοκετίσμος). Men ved saaledes at fremtræde i Phænomenernes Rige
% δοκετίσμος is AS PRINTED, read at 6x: epsilon where the word wants eta and the
% acute on the iota, not the ultima (standard δοκητισμός). The cursive kappa is
% typed as plain κ per the standing convention. Not in the Rettelser.
har han givet Aanden Frihed til at samle Erfaringens adsplittede Momenter i
% A STRAY MARK stands between "Frihed" and "til" --- a single small comma-shaped
% tick at ascender height, read at 2.6x. It is NOT the book's quotation mark,
% which is the two-stroke raised ” (compare "Kommer og seer!" on p. 161 below),
% and there is no quotation anywhere on this page. Both OCR witnesses turn it
% into a quote mark; it is transcribed as nothing and recorded here. PDF 168
% carries scattered specks of the same size.
Ideens egen klare, anskuelige Totalitet. Her gjælder altsaa ogsaa den gamle
Sætning, at Uovereensstemmelsen i det Uvæsentlige bestyrker Overeensstemmelsen i
det Væsentlige, saa at vi ikke have en sindrig Digtning, men en virkelig
Aabenbaring for os. Og ligesom Aanden hævdede sin Frihed til at skabe Formerne
--- saa at den hellige Aand tilkommer samme guddommelige Ære som Sønnen ---,
saaledes indrømmede den ogsaa Apostlene og den apostoliske Tidsalder Frihed til
at fremstille den samme Sandhed med forskjellige Ord. Den gjenfødte vel
Apostlene, men helliggjorde dem ikke med eet Slag, saa at enhver Yttring af den
menneskelige Villies
% --- p. 160 ---
% NO FOOTNOTE on this page.
Skrøbelighed pludseligt blev dem umulig; den oplyste vel deres Aand, men ikke
saaledes, at den pludseligt fjernede enhver subjectiv Eensidighed eller
Vildfarelse fra dem. --- Skriftens Skildring af de hellige Phænomener indeholder
derfor paa eengang den guddommelige Sandhed i dens høieste Virkelighed og de
frieste, reent subjective Bevægelser i de hellige Forfatteres Bevidsthed; de
hellige Phænomener ere ikke enten hiint eller dette, men ere netop kun i den
inderlige gjensidige Overgang af begge disse Momenter i hinanden, og kunne ikke
bestemmes ved nogensomhelst udvortes Grændse.

Men en Aabenbaring og Overlevering, for hvis hele Oeconomie Friheden er
Principet, kan naturligviis ikke tilegnes, før Frihedens sande Methode er
funden. Da nu den ældre Kirke og Kriticismen søgte at tilegne sig den
overnaturlige Idee ved at fastholde Phænomenerne for sig, var det en Selvfølge,
at de kun kunde slutte ved Analogie og omhyggelig Analyse af de enkelte
Momenter. Deraf fulgte, at paa den ene Side den ældre Tro maatte urgere
Beretningernes ordrette Overeensstemmelse, paa den anden Side Kriticismen mere
og mere maatte udstrække sin Tvivl til Miraklerne overhovedet, saasnart den
havde opdaget Uovereensstemmelsen i det Enkelte; \emph{thi kunne de ringere
% THE ONLY LETTERSPACED PASSAGE IN THIS BATCH, and it is unmistakable: three
% and a half whole lines set open, from "thi" through "det." spacing.py returns
% four consecutive RUNs over it and the ABBYY layer spaces the letters out.
% The closing point is inside the emphasis, as it is set.
Mirakler ikke bestaae i Undersøgelsens Prøve, hvormeget mindre maae da de større
kunne det.} Men denne Positivisme og Probalisme ophæves fra Grunden af, naar vi
% "Probalisme" is AS PRINTED, not "Probabilisme" --- read at 4x and given the
% same way by both OCR witnesses. Not in the Rettelser; not corrected.
gjennem Negationen af Phænomenerne trænge ind til Ideen selv: vi holde da ogsaa
urokkeligt fast ved Grundmiraklerne, for i deres Lys at kunne betragte de
speciellere Mirakler. Christusphænomenets Restauration bør da foregaae paa en
saadan Maade, at de enkelte Beretninger om hans Liv beholde hver sin Plads, for
saaledes hver paa sin Maade at aabenbare Incarnationens Idee.

% --- p. 161 ---
% NO FOOTNOTE on this page. The rule below is the Slutning divider, not a
% footnote rule --- what follows it is display and body type, not small type.
% p. 161's first line IS indented (the profile's x303 is a ghost speck; the
% indent is plain at 2x), so a new paragraph opens the page.
Vi komme saaledes atter tilbage til det kirkelige Christusbillede, men paa en
saadan Maade, at vi, efter at have underkastet dets enkelte Træk en særskilt
Analyse, indsee, at den samme Christus, som Kirken troede at kunne anskue i
Hverdagserfaringens udvortes Rige, ikke maa betragtes udenfor Aandens Speil.

Vel ophørte Inspirationen som saadan, men den hellige Aand vedbliver sit Liv i
Kirken: under alle Tidens Omvexlinger kjæmper den med Vælde mod denne Verdens
% "Omvexlinger" --- the same Fraktur x that the OCR reads as r at "forvexles",
% p. 157 above; see the comparison logged there.
Aand; hvergang den hellige Historie forvanskes, er det den, der paa ny
opfrisker, hævder og ved sit guddommelige Lys forklarer Kirkens guddommelige
Christusbillede; og den vil vedblive at belyse det, saalænge indtil Begrebets og
Anskuelsens Momenter fuldkomment gjennemtrænge hinanden; --- saa længe indtil
% A small raised speck stands between "hinanden;" and the dash. PDF 170 carries
% heavy show-through from p. 162 across its whole lower half (the mirrored
% "Slutning." head is legible in the blank foot of PDF 171), and this is part of
% it: transcribed as nothing, recorded here.
Christus aabenbares i sin Herlighed og hans Røst lyder til Alle: ”Kommer og
% BOTH quotation marks are the RAISED ” (U+201D), verified at 3x --- the book's
% standing form. The only pair in this batch, so the parity test is even.
seer!”

% *** A PLAIN CENTRED RULE stands between the last line of § 23 and the Slutning
% head. *** Measured, not eyeballed, on PDF 170: it runs x1324--2027 on a text
% block of x290--3110, i.e. 0.249 of the measure, centred on 1675 against the
% measure's own centre of 1700; it is 7--8 px thick at 300 ppi and UNIFORM ---
% no taper, unlike the swelled rule before the Anden Deel head (p. 53), and
% rather heavier than the hairline after the § 1 argument (p. 1, 347 px wide and
% 6 px thick). \centerline rather than a center environment, per the p. 1 and
% p. 53 precedents, so that check.py's hand-rolled-head test is not tripped.
% NOTE FOR THE CALLER: the printed gaps are about 1.4 lines from the last text
% line to the rule and 1.6 from the rule to the head; \division adds
% 2.0\baselineskip of its own above the head, so the space below the rule will
% come out a little generous. Adjust there if it matters.
\par\vspace{1.0\baselineskip}\centerline{\rule{0.25\textwidth}{0.5pt}}

% *** THE SLUTNING. *** The head is printed "Slutning." --- WITH a point, in
% LARGE BOLD FRAKTUR, centred, and NOT letterspaced (measured: 575 px for nine
% letters, set solid; spacing.py's "single? Slutning. (1.35)" is noise). It is
% larger than the body and about the size of the "Første Deel." line of p. 1,
% but it stands alone: there is no argument line under it, and no number. That
% is exactly \division, which the preamble sets as \large\bfseries centred. It
% agrees with the Indhold, which lists "Slutning ... 161" without a point; the
% printed head has the point and is what is set here. Last head in the book.
\division{Slutning.}

Dersom man indrømmer Sandheden af den Udvikling, vi saaledes have givet, vil man
ogsaa give os Ret i at antage den speculative Behandlingsmaade af den hellige
Historie for den eneste sande.

Den bevarer paa den ene Side den positive Charakteer, thi i Erkjendelsen af, at
den objective Fornuft ikke kan begribes uden gjennem en omhyggelig Iagttagelse
af de givne Gjenstande, modtager den med Glæde de Skatte, som de Lærde have
samlet, og lader ingen af Kritikens Undersøgelser upaaagtet. Men paa den anden
Side veed den speculative Theologie vel at skjelne imellem Videnskabens Apparat
og Videnskaben selv; derfor forfølger den med Omhyggelig\-%
% --- p. 162 ---
% *** THE LAST PAGE OF THE BOOK. *** Fourteen body lines and then nothing: no
% tail ornament, no rule, no colophon, no printer's imprint and no signature
% mark. Verified by a row-by-row ink scan of every row below the last text line
% (y2993 of 6645) as well as on the image. What that scan turns up in the blank
% two-thirds of the leaf is SHOW-THROUGH of p. 161 on the other side of the
% same sheet, mirror-reversed and legible as such --- the Slutning divider rule
% at y3430 and the mirrored "Slutning." head at y3783 --- plus binding shadow at
% the extreme foot. PDF 172--175 are blanks and binding. NO FOOTNOTE.
hed alle philosophiske Undersøgelser; ikke i blind Efterfølgelse af dette eller
% "blind" is NOT letterspaced; spacing.py's "single? blind (3.36)" is a false
% positive, checked at 1.5x. Nothing on this page is emphasised.
hiint System, men for, som den abstracte Philosophies Lærerinde, efter omhyggelig
Prøvelse af ethvert af dens Systemer at føre den skabte Fornuft tilbage til dens
sande Afhængighed af den uskabte Fornuft og for igjennem denne Forsoning af
begge at fortolke Religionens Phænomener. --- Den er langt fra al Nyhedssyge;
thi dens Bestræbelse gaaer netop ud paa at drage alle de sande Momenter, der
findes i de ældre Systemer, for Lyset. Men fremfor Alt fordrer den en sand
% "ældre", not the OCR's "ælvre". The page is heavily inked here and the d's
% ascender is nearly closed up, so the call was measured rather than eyeballed:
% the glyph group carrying it tops out 8 px above the line's x-height, exactly
% like the d of "findes" earlier in the same line, where a Fraktur v would stop
% at x-height. Both OCR witnesses read v.
Overgang fra den ene Erkjendelsessphære til den anden og en indre dialektisk
Forbindelse imellem de enkelte Dele; den forkaster derfor enhver Sammenblanden
af modsatte Elementer, forsmaaer alle Eklekticismens Kunstgreb og flyer den
slette Middelveis Middelmaadighed.

\end{document}
